\documentclass[12pt]{article}
\usepackage[T1]{fontenc}
\usepackage[margin=1in]{geometry}
\usepackage{amsmath,amssymb}
\usepackage{graphicx}
\usepackage{float}
\usepackage{array}
\usepackage{longtable}
\usepackage{booktabs}
\usepackage{threeparttable}
\usepackage{caption}
\usepackage{authblk}
\usepackage[authoryear,round]{natbib}
\usepackage{xurl}
\usepackage[hidelinks]{hyperref}
\usepackage{microtype}
\emergencystretch=1.5em
\renewcommand{\thesection}{\Roman{section}}
\renewcommand{\thesubsection}{\thesection.\Alph{subsection}}
% Possessive citation: one macro so a bibliography change cannot desynchronize
% hand-rolled \citeauthor{X}'s (\citeyear{X}) constructions.
\newcommand{\citepos}[1]{\citeauthor{#1}'s (\citeyear{#1})}

\title{Mortgage Lock-In and the Composition of the Federal Reserve's Agency-MBS Runoff}
\author{Eugene Ong}
\affil{Carnegie Mellon University}
\date{}
\begin{document}
\maketitle
\begin{abstract}
\noindent U.S. mortgages are fixed-rate, repaid at face value, and not portable; when rates rise, that combination locks households into old, cheap loans. I measure what it cost the Federal Reserve: repayments on its mortgage bonds ran \$764.7 billion below the monthly ceilings set for June 2022 through November 2025. Under the production floor form, most of that gap was never about behavior: switch the lock-in response off and the model still accounts for 85.7\% of it (35.6\% under the additive form), because scheduled amortization and baseline turnover (itself read from realized, partly behavioral turnover) fall short of a \$35 billion monthly ceiling however the rate-responsive margin behaves. A loan-by-loan model run on Freddie Mac data accounts for 91.3\%. \par What lock-in itself adds is a range rather than a number. The design bounds it between $+2.9$ and $+8.7$ points under its production floor form (the floor read's sampling error at my central elasticity, the widest uncertainty source with a coverage property though not the widest read of that source; a wild-cluster interval on 31 clusters; the percentile read under-covers; the unestimated baseline seasoning ramp spans $+0.9$ to $+15.6$ at the same calibration), with a form-conditional hull of $+3.5$ to $+13.1$ points, and it identifies a marginal, not a level and not monthly timing. Inside that range, $+5.6$ points, or \$42.6 billion, is the value at the calibration I headline, and the corrections I can measure to the baseline turnover floor or to the accounting basis no longer all run one way: calendar-standardizing the floor read moves it to $+4.7$, while matching that read's housing-activity level to the window's moves it to $+10.7$. That is the range's coordinate at one anchor---the middle of the three depth cuts fixed before the reads---and not its centre; I attach no posture to where it sits. The cost to households who could not move is real. The institutional cash-flow cost is small, on an accounting that counts repaid principal at face value, and under the production floor form; the additive form roughly doubles the margin.
\end{abstract}

\vspace{0.5\baselineskip}
{\footnotesize\noindent\textbf{Keywords:} mortgage prepayment; lock-in effect; quantitative tightening; mortgage-backed securities; agent-based models; hazard models.\\
\textbf{JEL codes:} E52, E58, G21, R31.\par}
\newpage
\section{Introduction}\label{sec:intro}

For most of its history, the Federal Reserve operated under a ``Treasuries-only'' framework, using sovereign debt with fixed, predictable maturities. During the 2008 Global Financial Crisis, and again during the 2020 COVID-19 pandemic, the central bank departed from this framework, launching aggressive Quantitative Easing (QE) programs that purchased trillions of dollars in Agency MBS to stimulate the consumer economy directly.

This intervention suppressed 30-year mortgage rates to historic lows of roughly 2.6\%, triggering a large refinance boom. As inflation surged in 2022, the Federal Reserve attempted to reverse this stimulus through Quantitative Tightening (QT), establishing a passive roll-off target intended to drain up to \$35 billion per month from its SOMA portfolio as homeowners relocated naturally, a cap phased in at \$17.5 billion per month and stepped up to \$35 billion three months later.

That the lock-in effect slowed the runoff of the Federal Reserve's mortgage portfolio is not a new claim. The Federal Reserve's own staff and officials have said as much (\citealp{na2024}; \citealp{perli2024}; \citealp{hammack2025}). This paper's contribution is not that mechanism but its decomposition. That decomposition is an accounting exercise within a fully disclosed model rather than a causal identification: switching the lock-in elasticity off partitions the realized shortfall between an elastic channel and a no-elasticity baseline (scheduled amortization plus baseline turnover held at its observed level; the floor that sets that level is itself read off realized, partly behavioral turnover, so the baseline is not a non-behavioral object), and the elasticity's causal content is inherited from the external estimates that calibrate it, not established here (Section~\ref{sec:identification}). I ask how much of the runoff shortfall the lock-in elasticity accounts for, as opposed to the scheduled amortization and baseline turnover that would have undershot a \$35 billion monthly cap regardless of any rate response above that baseline. Most of the shortfall sits in the no-elasticity baseline. A loan-level model with the lock-in elasticity switched off still recovers 85.7\% of the \$764.7 billion shortfall-against-cap benchmark constructed in Section~\ref{sec:method-benchmark} at the calibration this paper headlines, and 88.7\% at its in-sample calibration point, measured on the paper's benchmark-consistent accounting basis (Section~\ref{sec:hazard-interp}). The Federal Reserve's own ex-ante projection implies the same composition, though not independently: the projection may already have embedded part of the lock-in channel (Section~\ref{sec:method-benchmark}). The New York Fed's May 2022 staff baseline already anticipated the large majority of the realized cap-shortfall: 75.6--88.5\%, depending on a disclosed intra-2022 allocation choice (Section~\ref{sec:method-benchmark}), a range wide enough that its overlap with the mechanical null's recovery carries no precision.

The elasticity's identified footprint is therefore a bounded margin, not a pinned magnitude and not a sign. The floor that dominates the level is calibrated in-window on the same locked-in cohorts the margin decomposes, so I re-measure it outside that window, on pre-episode 2017--2019 discount cohorts, and re-run both legs at the result: the margin runs $+4.3$ to $+6.8$ points across the off-window floor's defensible range, and $+2.9$ to $+8.7$ points once that floor's own sampling error is propagated with a wild-cluster small-sample correction, which is the binding layer (Section~\ref{sec:robustness-floor}). The value at the mid-grid anchor inside that interval is $+5.6$ points ($+\$42.6$ billion), and it is the range and not that value that the design delivers: the corrections I can measure to the floor or to the accounting basis no longer run one way either---calendar-standardizing the floor read moves the value to $+4.7$ points, matching that read's housing-activity level to the window's moves it to $+10.7$---and the specification variants that are not floor corrections run both ways as well, so the range and not any interior member carries the identified content (Section~\ref{sec:identification}). The in-window production calibration returns $+9.2$ points (within-stratum sampling interval $[+9.17, +9.23]$ points from the 200-replicate paired loan-level bootstrap, a scheme the cluster construction supersedes: resampling the 130 strata puts the paired marginal at $[+8.27, +10.19]$ points, a 30.8$\times$ wider read (run \texttt{bootstrap\_pathb\_cluster} at this floor); the operative uncertainty remains calibration, Table~\ref{tab:uncertainty}), reported here as the in-sample calibration point rather than the headline. Neither figure is the entire shortfall (Section~\ref{sec:identification}).

A pre-committed sweep of the two calibration inputs that dominate the level, the baseline turnover floor and the elasticity band, bounds the margin between $+2.1$ and $+13.2$ points, positive at every point of the calibration box and at every off-window anchor tested, though that positivity checks the specification rather than identifying the mechanism (Section~\ref{sec:identification}). A functional-form test that replaces the floor's hard maximum with a competing-risks combination leaves the margin inside that box and nearly floor-invariant (Section~\ref{sec:robustness-floor}). (``Pre-committed,'' here and throughout, means the grid, parity gates, and interpretive thresholds were committed to the repository before any run executed; I reserve ``pre-registered'' for third-party registries, which none of this is; Appendix~\ref{app:verdicts} tabulates every such rule together with how its outcome was adjudicated after the run.)

I reach this decomposition through two structurally distinct estimators, a 10,000-household agent-based model and a loan-level hazard framework, that disagree sharply about how much household choice explains. The decomposition the hazard framework supports is the paper's central empirical result; the disagreement between the two estimators I read as a supporting diagnostic rather than a clean paradigm contrast, since a real-covariate variant of the ABM recovers far more (Section~\ref{sec:robustness-crossdesign}).

This result sits inside a two-sided exchange: the pool homogeneity that generates To-Be-Announced (TBA) market liquidity is purchased at the price of household mobility, because the standardized contracts that TBA delivery presupposes carry the ``par-payoff'' obligation, under which a homeowner must extinguish an existing mortgage at full face value to relocate. The obligation itself is a property of the U.S. fixed-rate note enforced by the due-on-sale clause (made broadly enforceable by the Garn--St.~Germain Act of 1982) rather than of TBA eligibility, and it binds on portfolio and non-agency loans alike; the FHA/VA loans backing Ginnie Mae pools (20.4\% of the SOMA book) are assumable by statute, a carve-out from the book-wide mobility cost that Section~\ref{sec:pathb} scores by observed speeds. Realized take-up of that statutory right is not measured anywhere in this design, so the 20.4\% share bounds the carve-out from above rather than sizing it. The observed speeds bound it from above a second way, and in the same direction: an assumed loan does not prepay, so materially high take-up would show up as \emph{slower} Ginnie voluntary speeds, and over the window Ginnie's voluntary component instead runs $1.35$ points per year \emph{above} the conventional one. Both bounds are ceilings, not sizes; what neither prices is the frictions---second-lien gaps, equity buy-outs, qualification---that stand between the statutory right and its exercise. Extension-protected institutional cash flow is often named as a third design objective; Section~\ref{sec:conclusion} states what my evidence says about it and under what conditions. The exchange this paper measures is the two-sided one. The decomposition above disciplines this frame rather than dissolving it: household mobility is constrained, as the household-level lock-in literature documents, while the cash-flow shortfall is governed mostly by loan-level and institutional mechanics, so the trade-off's binding cost is denominated in mobility rather than in institutional cash flow (Section~\ref{sec:conclusion}). Section~\ref{sec:abm-danish}'s Danish counterfactual initially implied a \$925.5 billion institutional gap. Recalibrating against separately estimated Danish elasticities (Section~\ref{sec:abm-danish}) reduces this by an order of magnitude, to $+\$61.2$ billion (8\% of the benchmark), the zero-refinance edge of a band swept to $+\$256.8$ billion at 3\% refinance-in-place (Section~\ref{sec:abm-danish}), under the shared-accounting hybrid pipeline when the transplant changes only the payoff rule (the Danish moving channel's estimated flatness anchored at the U.S. zero-gap intercept, baseline turnover floor respected). What that construction can and cannot demonstrate is set by the anchor: a Danish leg pinned at the zero-gap U.S. hazard must prepay faster than a U.S. book that carries a negative rate gap almost everywhere, so a positive gap---and its survival of the refinance-in-place sweep---is forced rather than found, and the informative content is the magnitude. The bracketing anchor that also imports Denmark's baseline mobility flips the sign, to $-\$99.9$ billion---but that leg runs its Danish households at 3.39\% annual prepayment, below the 4\% turnover floor both frameworks impose on the U.S. side, which is why I bracket with it rather than headline it. I treat the corrected magnitude as informative about where the trade-off actually binds, and the sign as anchor-determined---and, on the cash-flow side, incidence-determined as well: pricing the market-value buyback credit under a cash-haircut reading reverses the gap to $-\$89.4$ to $-\$117.7$ billion at the committed proxy discounts, so the institutional relief is established under face accounting only, which is the benchmark's own denomination (Section~\ref{sec:pathb}). The mobility-versus-cash-flow ranking above rests on external evidence, not on a within-paper comparison: I measure the institutional leg and import the mobility leg at an elasticity I never re-measure, so the two are nowhere commensurable. The one household-side magnitude I carry is \citepos{batzer2024} roughly \$2.4 trillion of foregone capital gains, computed on universal relocation and therefore an upper bound, and a stock of foregone gains is not a cash flow over forty-two months. The ranking's direction is what that supports; its ratio is not.

Unlike standard U.S. Treasuries, U.S. mortgages carry asymmetric, state-dependent payoff profiles. As the Federal Reserve's policy tightening pushed the 30-year fixed mortgage rate past 7.0\% (a rise decomposable into the 10-year Treasury yield's climb and a widening of the primary-mortgage-to-10-year spread from roughly 170 to over 300 bp at its 2022--2023 peak, FRED series MORTGAGE30US less DGS10; the decomposition of that spread into MBS supply, prepayment-risk repricing, and convexity-hedging components follows \citet{boyarchenko2019}, whose framework predates and does not document the 2022--2023 magnitudes), the par-payoff rule imposed a substantial financial penalty on relocation, producing a ``lock-in effect'' that suppressed prepayment activity across the U.S. housing market. Because mortgages did not prepay, principal cash flows from the Federal Reserve's MBS portfolio ran far below the redemption caps, extending the portfolio's effective maturity. Cumulative runoff fell \$764.7 billion short of the phased caps over the tightening cycle (methodology and interpretation in Section~\ref{sec:method-benchmark}). How much of that shortfall the lock-in channel itself accounts for is the question Sections~\ref{sec:abm} and~\ref{sec:hazard} decompose.

Standard macroeconomic models often underestimate this extension risk because they assume frictionless, rational household behavior. In practice, the lock-in effect was compounded by historic housing inventory shortages and depressed consumer sentiment.

This paper quantifies the mechanics of this shortfall using two structurally distinct estimators. The first is a 10,000-household Agent-Based Model parameterized with dynamic macroeconomic frictions, used to estimate stranded cash flows and to decompose that estimate against the empirical benchmark. The second, introduced after the ABM's residual proved persistent and unresponsive to additional behavioral realism, is a reduced-form hazard framework estimated on loan-level survival data that abandons the household utility function altogether. As Sections~\ref{sec:abm} and~\ref{sec:hazard} show, synthetic-population household choice explains only a small share of the shortfall (13.6\% on the fifty-seed mean, 10--17\% across seed noise, 11.9\% on the frozen production draw), while the loan-level hazard framework recovers the large majority of it at the aggregate level: 91.3\% in its literature-calibrated microsimulation form on the benchmark-consistent shared accounting basis at the off-window floor this paper headlines, 88.2--93.7\% across that floor's defensible range, and 97.9\% at the in-sample calibration point (92.2\%--99.2\% across covariate calibrations at that in-sample point), and 100.0\% (approximately 94.3\%--101.3\% across those calibrations; 98.1\% when sample and book are first put on one coupon convention) once the sample's coupon mix is reweighted to the SOMA book under the ex-ante priors. Its estimated cohort-hazard form (Path A) lands at 112.4\% but is excluded from headline ranges and does not corroborate the result: its forward simulation carries no useful propagated precision, and no inferential object has been recomputed under its production specification (Section~\ref{sec:patha}). Two qualifiers travel with the two-estimator contrast wherever I state it: the null that switches off the lock-in elasticity (the $\beta_1 = 0$ null) already recovers 85.7\% of the benchmark at the off-window floor and 88.7\% at the in-window production calibration, so the hazard framework's margin over a mechanical model is $+5.6$ points at the off-window floor ($+9.2$ at the in-window production calibration), not the level (Section~\ref{sec:identification}). A cross-design variant of the ABM recalibrated on real loan covariates recovers 59.3\% on the frozen seed, 60.2\% across fifty seeds with a 95\% band of 55.1--66.0\%, and 76.3\% at the book's composition (effective sample 1{,}332 of 10{,}000 loans; a re-draw returns 70.6\%; Section~\ref{sec:robustness-crossdesign}), so the contrast is not clean evidence about modeling paradigm alone. The hazard framework is the instrument; the decomposition it yields is the paper's central empirical result. (Section~\ref{sec:robustness} shows, on the production specification, that population size does not drive this gap.) A counterfactual simulation based on the Danish market-value repurchase regime further shows that the institutional cash-flow component of systemic duration risk is a consequence of U.S. institutional design: under the rule-only transplant it matches the lock-in marginal in origin, and in size at the in-sample calibration point on which both were run.

This paper's affirmative contributions are three. First, a decomposition of the QT mortgage-runoff shortfall into mechanical and elastic components, with the elastic part delivered as a bounded margin rather than a point (Sections~\ref{sec:abm}, \ref{sec:hazard} and~\ref{sec:identification}). Second, an expectations-based complement that measures the shortfall against the Federal Reserve's own ex-ante projection rather than the never-binding caps (Section~\ref{sec:method-benchmark}). Third, the cap-design arithmetic. That the caps would sit above achievable runoff was projected by New York Fed staff from the month the \$35 billion ceiling took effect \citep{nyfed2022}; the observation is theirs, and what I add is its quantification and its ex-ante construction: the ceiling was set roughly 1.7 to 1.9 times as high as the Federal Reserve's own contemporaneous projection of achievable runoff, and both components of the mechanical path a binding cap would have to sit below --- scheduled amortization from the book's coupon, age, and term composition, and baseline turnover at any assumed floor --- are computable in advance of the cycle, as a band rather than as a point (Sections~\ref{sec:identification} and~\ref{sec:discussion-capdesign}). Two further elements of the design are not contributions, and I label them here as what they are. The paired cross-design test --- the ABM's decision rule on real loan covariates, and the same survival structure on a fully synthetic population, under interpretive thresholds fixed ex ante --- is a validation device, and it ran against this paper's own ABM reading rather than for it: on all fifty seeds the real-covariate leg crossed the threshold, fixed before that run, at which that reading is undercut (Section~\ref{sec:robustness-crossdesign}). The order-of-magnitude reduction of the Danish market-value counterfactual, with the incidence bracket that prices it, is a robustness result on a counterfactual rather than a measurement of the U.S. book: the \$925.5 billion figure it corrects is this paper's own earlier extrapolation, and the corrected gap's sign is anchor- and incidence-determined (Sections~\ref{sec:abm-danish} and~\ref{sec:pathb}).

Table~\ref{tab:headline} collects the core results and their operative uncertainties in one place; every value repeats a committed, gated figure.

{\footnotesize
\begin{longtable}{@{}p{4.6cm}p{4.6cm}p{5.3cm}@{}}
\caption{Core results and operative uncertainties. Each row repeats committed values derived in the referenced exhibit; recovery percentages are quoted on the benchmark-consistent shared basis. The uncertainty hierarchy is established in Table~\ref{tab:uncertainty}: simulation-seed noise is negligible at reporting precision and sampling is quoted under the stratum-cluster scheme; the calibration box is the operative uncertainty --- except the headline marginal, whose binding layer is the floor read's own propagated sampling error (Table~\ref{tab:uncertainty}). Table~\ref{tab:crosswalk} (Appendix~\ref{app:ledger}) maps each headline quantity to its accounting basis and source run. The Path B recovery is estimated on the Freddie 2017--2021 conventional universe; Section~\ref{sec:method-benchmark} derives that universe's coverage of SOMA book face on the book's own agency $\times$ vintage joint cells, and states the external-speed overlay that bounds each exclusion (Section~\ref{sec:pathb}).}
\label{tab:headline}\\
\toprule
Object & Headline value & Uncertainty / bounds \\
\midrule
\endfirsthead
\toprule
Object & Headline value & Uncertainty / bounds \\
\midrule
\endhead
Cap-relative benchmark, June 2022--November 2025 (Section~\ref{sec:method-benchmark}) & \$764.7 billion & -- \\
Expectations-based complement (Section~\ref{sec:method-benchmark}) & \$87.8 billion, 11.5\% of the cap-relative benchmark & Fed's ex-ante projection implied 75.6--88.5\% of the cap-shortfall, by intra-2022 allocation \\
No-elasticity baseline, $\beta_1 = 0$ (Section~\ref{sec:identification}; Table~\ref{tab:bases}) & 85.7\% of benchmark at the headline off-window floor; 88.7\% in-sample, both under the production floor form (35.6\% under the additive form) & 83.9--86.9\% across the off-window floor range under that form; moves one-for-one with the floor level \\
Lock-in marginal, off-window floor (Sections~\ref{sec:identification} and~\ref{sec:robustness-floor}; Table~\ref{tab:oosfloor}) & $+2.9$ to $+8.7$ points, basis-invariant --- the floor read's own sampling error propagated at the central elasticity (the 6.5\% mid-band value), not an interval on the elasticity, whose dimension Table~\ref{tab:lowband} sweeps; $+5.6$ points $= +\$42.6$ billion at the mid-grid anchor inside it, an anchor convention rather than a central tendency; the range carries the identified content (Section~\ref{sec:identification}) & floor-read wild-cluster bootstrap-$t$ $[+2.9, +8.7]$, the binding layer (31 clusters; percentile read $[+3.0, +8.0]$, under-covering; the corrected ladder and its degrees of freedom are Table~\ref{tab:ladder})\tnote{a} \\
\addlinespace
Lock-in marginal, in-sample calibration point (Section~\ref{sec:identification}; Table~\ref{tab:uncertainty}) & $+9.2$ points $= +\$70.3$ billion, basis-invariant & sampling: within-stratum $[+9.17, +9.23]$, stratum-cluster $[+8.27, +10.19]$ (the operative scheme); calibration box $+2.1$ to $+13.2$; Fannie replication $+8.68$ (Section~\ref{sec:pathb-fannie}) \\
Path B central level (Table~\ref{tab:estimators}) & 91.3\% at the headline off-window floor; 97.9\% in-sample (100.0\% composed as committed, 98.1\% on a single coupon convention) & 88.2--93.7\% across the off-window floor range; covariate priors 92.2--99.2\% (in-sample); Ginnie-speed overlay at the headline floor 84.0\% central against a 79.5\% null, and 89.3\% at the in-sample point (raw; of the \$66.3B raw adjustment the Ginnie-specific piece is $+\$38.5$B, within the \$20--47B bound) \\
Production ABM (Table~\ref{tab:estimators}) & 13.6\% seed mean & 10--17\% seed band; frozen draw 11.9\%; the operative uncertainty is calibration, not seed noise. The swing that establishes this is measured on a different design: on the real-covariate cross-design, switching one calibration choice moves recovery across 20.9--59.3\%, against a fifty-seed cross-design sweep whose own band is 10.9 points (Section~\ref{sec:robustness-crossdesign}). That is a span across two calibration choices on that design, not an interval around this row's 13.6\%, and does not contain it \\
Cross-design ABM, real covariates (Table~\ref{tab:crossdesign}) & 59.3\% recalibrated; fifty-seed mean 60.2\%, 95\% band 55.1--66.0\% (76.3\% at book composition; effective sample 1{,}332 of 10{,}000, re-draw 70.6\%) & 20.9\% under the frozen scale; fifty-seed mean 22.0\% (12.6\% at book composition) \\
Danish rule-only counterfactual (Table~\ref{tab:danish}) & $+\$61.2$ billion & positivity is forced by the zero-gap anchor, so the refinancing sweep cannot flip it; the Danish-level bracketing anchor does, at $-\$99.9$ billion. The gap is incidence-conditional: $+\$61.2$ billion under face accounting (the benchmark's own denomination), reversing to $-\$89.4$ to $-\$117.7$ billion under cash accounting at the committed proxy discounts (run \texttt{buyback\_credit\_bracket}; Section~\ref{sec:pathb}) \\
\bottomrule
\end{longtable}
}

\noindent{\footnotesize \emph{Notes to Table~\ref{tab:headline}.} a Sensitivity catalogue for the headline marginal, each entry a committed figure derived where it is cited and none of them a sampling interval: off-window floor range $+4.3$ to $+6.8$ (max form; $+3.9$ to $+7.5$ joint with the band); additive form $+11.2$, floor-invariant; form-conditional hull $+3.5$ to $+13.1$, upper end from a 61.2\%-recovery cell and not read as an equally credentialed member (the concave transform and the additive form do not compose additively: jointly $+9.2$, interaction $-1.47$); mixture curve in the floor's strictly-involuntary share: $+9.7$ at $s{=}0.25$, $+10.7$ at $s{=}0.4$ (run \texttt{floor\_form\_mixture}); Ginnie overlay $+4.4$ (conventional-share scaling $0.797\times$); floor-stability check $+\$45.1$ billion.\par}

\paragraph{Definitions used throughout.} Four accounting bases attach to every recovery percentage (Section~\ref{sec:method-benchmark}): the \emph{standalone-scorer} basis (no adjustment), the \emph{benchmark-consistent shared} basis (nets a common \$69.6 billion curtailment flow identically from every U.S. leg; the main-text convention), the \emph{full-book} basis (reweights to the SOMA book's coupon mix), and the \emph{composed} basis (both). Two floor calibrations attach to every marginal and level: the \emph{in-sample calibration point} (the 4\% baseline turnover floor measured in-window on the locked-in 2023--24 cohorts; returns $+9.2$ points) and the \emph{off-window calibration} (the 4.70--5.33\% floor measured on pre-episode 2018 discount cohorts; returns $+5.6$ points and is the headline). The floor itself is a measured object: realized-turnover reads put it at 3.8--3.9\% in-window (flat across seasoning cuts), 4.70--5.33\% on the off-window 2018 leg (all of it on cohort-months aged twelve to twenty-four, so that leg's mature-seasoning test is not computable rather than passed), 5.51\% under an imputation-heavy age standardization (84\% imputed weight, an indicative bound derived from the same 2018 leg), and 5.52\% on an out-of-agency Fannie Mae read of the pooled off-window period; the last two lines sit above the clean band, as does every mature-versus-young turnover contrast the data make measurable, which is why the band's lower edge is the soft one (Section~\ref{sec:robustness-floor}). The ``baseline turnover floor'' (sometimes called an involuntary-turnover floor) names that measured object without decomposing it: what these reads measure is total deep-discount-cohort turnover, which mixes discretionary life-cycle moves and cash-out refinancings with strictly involuntary events, and the strictly-involuntary share is plausibly well under half (Section~\ref{sec:pathb}), a share the mixture curve of Section~\ref{sec:robustness-floor} prices directly. ``Trapped liquidity'' and ``shortfall'' denote the net shortfall of realized roll-off against the phased redemption caps (\$764.7 billion over June 2022--November 2025), and the \emph{lock-in marginal} is the difference between the central simulation and its $\beta_1 = 0$ null, the paper's identified object, whose operative uncertainty is the calibration-and-form envelope of Table~\ref{tab:uncertainty} rather than any sampling interval. That envelope is a grid, not a confidence region, and the distinction is load-bearing: the elasticity band it sweeps is \citeauthor{liebersohn2024}'s \emph{specification} range with my adopted midpoint at its centre, not a published point estimate with a published standard error, and no sampling error of theirs is propagated anywhere in this paper. So the envelope has no coverage property. The $\beta_1 = 0$ leg those readings are differenced against is the \emph{no-elasticity baseline}: it carries scheduled amortization from the book's own coupon, age, and term composition together with baseline turnover held at its observed level, and the marginal is the imported elasticity's increment above it. What the partition separates is that baseline from that increment, not mechanics from behavior---the floor that sets the turnover level is measured on realized turnover which is itself partly behavioral, so the baseline is not a non-behavioral object and I do not call it one. Where a genuine sampling interval exists---the floor reads' cluster bootstrap, and the loan- and stratum-level bootstraps of Section~\ref{sec:pathb}---I quote it as such and rank it against the envelope explicitly rather than merging the two. One global property of the design should be stated here as well as in the limitations, because every number below inherits it: \emph{no outcome-holdout months exist anywhere in this paper}, without exception. Its unavoidability is scoped to the cash-flow margin. The identified object has nothing to hold out: the marginal is a difference between two simulations, so no realized outcome corresponds to it. A holdout on the level's fit is conceivable but not attempted, since it would mean splitting the benchmark itself---a single cumulative path against one cap schedule. The household margin does admit outcome moments this design does not use, since it imports the mobility elasticity rather than estimating it (Section~\ref{sec:limitations}). The turnover floor that sets Path~B's level is calibrated on in-window turnover, Path~A trains on nineteen of the forty-two QT months, and Path~B's level is evaluated on the window that sets it. Each of these is disclosed again where it arises, and the nineteen-month floor variant reported later is an input-stability check rather than an outcome holdout: what it holds out is calibration months, not the benchmark. Every recovery percentage in this paper is accordingly an in-sample or calibrated quantity in the time dimension, and should be read as a description of fit rather than as a prediction (Section~\ref{sec:limitations}).

\paragraph{How to read this paper.} The short path for a referee: Sections~\ref{sec:lit}--\ref{sec:method-benchmark} for the claim and the design; Section~\ref{sec:pathb} for the operative estimator; Section~\ref{sec:identification} for what is and is not identified; Section~\ref{sec:robustness-floor} for the calibration sensitivity the headline rests on; Section~\ref{sec:conclusion} for the bottom line. Everything after the online-appendix banner (Appendix~\ref{app:ledger} through Appendix~\ref{app:verdicts}) is apparatus: the run ledger, the specification blocks, the sweep mechanics, and the verdict-adjudication audit. The seven qualifications that scope the headline live in Section~\ref{sec:identification}; floor and form sensitivity in Section~\ref{sec:robustness-floor}; every non-mechanical adjudication in Appendix~\ref{app:verdicts}.

\section{Literature Review}\label{sec:lit}

The theoretical foundation of mortgage prepayment is the assumption of ``ruthless exercise,'' wherein borrowers refinance instantaneously as market conditions dictate \citep{dunn1981}. Subsequent literature recognized behavioral frictions: \citet{schwartz1989} estimate an empirical prepayment function from aggregate GNMA experience with maximum-likelihood hazard techniques and embed it in a valuation framework (the econometric ancestor of the loan-level survival approach I adopt in Section~\ref{sec:hazard}), while \citet{richard1989} formalize the industry practice of modeling prepayment as an empirical function of the refinancing incentive, seasoning, seasonality, and a pool's own prepayment history. \citet{stanton1995} introduced heterogeneous transaction costs, smoothing the instantaneous refinancing cliff into a predictable S-curve, and \citet{agarwal2013} give the rational benchmark its sharpest form, a closed-form optimal refinancing threshold against which observed behavior can be scored. \citet{chernov2018} recover the prepayment function implied by MBS market prices and find it driven by a macroeconomic turnover component alongside the rate response---the same turnover-versus-rate-response split this paper's mechanical null operationalizes on the cash-flow side. \citet{quigley1987,quigley2002} established the lock-in effect as a persistent friction limiting mobility during rising-rate environments. \citet{ferreira2010} provided early large-sample evidence for it, finding that rising mortgage rates and negative home equity each depress owner mobility---the latter an equity-lock-in channel distinct from, and complementary to, the rate channel this paper models (the equity dimension is left unmodeled here, a boundary disclosed in Section~\ref{sec:pathb}).

The 2022--2025 rate shock generated substantial empirical literature documenting the breakdown of standard prepayment assumptions. \citet{fonseca2024} documented the steep decline of U.S. mobility and refinancing activity, while \citet{batzer2024} estimated that U.S. households would have faced roughly \$2.4 trillion in foregone capital gains had they chosen to relocate. \citet{liebersohn2024} quantify the same lock-in mechanism in hazard form, estimating that each percentage point of prevailing rate above a borrower's locked-in coupon reduces quarterly mobility probability by 5.5\% to 7.7\% depending on specification: the elasticity band my Path B microsimulation adopts as its calibration; \citet{graybill2026}, in a Federal Reserve Bank of Philadelphia working paper, independently estimate the same margin from transaction records in a housing-market-equilibrium setting, with deep rate gaps cutting the sale hazard by roughly 7--11\% per percentage point, a corroboration of that band from data that never touch credit records, and one that sits at or above it rather than at its centre: only the 7\% low end of their range falls inside the 5.5\% to 7.7\% band, and its midpoint sits above the high edge (Section~\ref{sec:pathb}).

A related strand of literature decomposes how much of the observed mobility decline a rational, rate-gap-driven framework can explain, leaving the remainder as an unexplained residual. \citet{aladangady2024} construct a counterfactual rate-gap distribution holding the 2021 rate environment fixed and find that rate-gap-driven lock-in explains only 44\% of the actual decline in mortgage borrower mobility between 2021 and 2022, after correcting for selection bias from forgone refinancing opportunities. Section~\ref{sec:identification} returns to that 44\% share and maps it into this design's own units as a calibration target for the scaled-null variant (run \texttt{scaled\_null\_housing\_activity}). This decomposition approach is methodologically analogous to the one I take in Sections~\ref{sec:abm} and~\ref{sec:hazard}, though I apply it to dynamically modeled CPR and cash-flow shortfalls rather than to mobility rates directly, and use two structurally distinct estimators rather than one. \citet{gerardi2024} similarly compare a structural search-and-matching model of lock-in against empirical lifecycle migration patterns, finding that their model accounts for most, but not all, of the observed disruption to housing market liquidity.

A separate behavioral finance literature provides theoretical grounding for what such residuals likely represent. \citet{keys2016} document that more than 20\% of U.S. borrowers fail to refinance even when doing so would be unambiguously in their financial interest, attributing this to procrastination and variation in financial sophistication. \citet{defusco2020} document a structural counterpart to these behavioral frictions, showing that employment-documentation requirements and out-of-pocket closing costs bind most severely in recessions and inhibit refinancing pass-through. \citet{andersen2020} reach a similar conclusion using Danish administrative data, finding that inattention and inertia, rather than transaction costs, account for a substantial share of household failure to act on refinancing opportunities. \citet{andersen2022} further show, using Danish housing data, that households exhibit reference dependence and loss aversion relative to their original purchase price, weighting losses roughly 2.5 times more heavily than equivalent gains. \citet{johnson2019} examine sluggish refinancing using administrative, survey, and field-experimental data, finding that among the behavioral factors they test, suspicion of lenders' motives (rather than standard rational frictions such as time preference or risk aversion) is most consistently related to households' failure to accept profitable refinancing offers, further supporting a behavioral interpretation of the residual. \citet{bissiri2014} formalize this distinction directly, defining ``behavioral risk'' as the share of prepayment decisions not driven by a fully rational financial criterion, an operationalization more recently applied by \citet{perotti2024} to model behavioral uncertainty as a stochastic overlay on an otherwise rational prepayment framework and to derive robust hedging strategies that bound, rather than eliminate, the exposure such uncertainty creates.

A further strand treats prepayment as a loan-level survival process shaped by pool composition and servicer behavior rather than by borrower optimization. Prepayment ``burnout'' (the tendency for a pool's aggregate refinancing sensitivity to decline over time as its most rate-sensitive borrowers prepay earlier and self-select out of the remaining pool) is a standard feature of industry MBS valuation models (one of the four factors of \citepos{richard1989} canonical prepayment function, alongside the Public Securities Association seasoning convention), but is difficult to reconcile with a fixed, representative-agent population. Recent theory formalizes why: in a heterogeneous pool the observed pool-level hazard is the survival-weighted mean of the individual hazards, and its drift carries a negative selection term proportional to the cross-sectional dispersion of those hazards \citep[arXiv preprint]{lesniewski2026}, so pool speeds decline with seasoning even when no individual hazard changes. My hazard framework (Section~\ref{sec:hazard}) is designed to accommodate this loan-level survival structure.

Fixed-income literature separately highlights the systemic risks generated by the lock-in effect. \citet{jiang2023} documented that duration extension in bank asset portfolios, prominently including low-coupon MBS, generated roughly \$2 trillion in unrecognized mark-to-market losses across the U.S. banking sector, contributing to self-fulfilling solvency runs including the collapse of Silicon Valley Bank. \citet{smith2023} demonstrated that QT relies heavily on predictable MBS runoff schedules, which are structurally broken by these duration extensions.

Closest to this paper, the Federal Reserve's own staff and officials have already described the mechanism and measured its first-order footprint. \citet{na2024} report that ``a vast majority of the agency MBS held in the SOMA are far out-of-the-money to refinance and may be subject to the `lock-in' effect,'' and that because ``monthly principal payments have been below the \$35 billion-per-month redemption cap \ldots the size of the SOMA agency MBS portfolio has decreased at a much slower pace than in the case of a binding cap''; they put June 2022 to June 2024 runoff at roughly \$450 billion against the \$820 billion a binding cap would have retired, an implied shortfall near \$370 billion. The same note reports the composition of those below-cap payments: scheduled principal accounts for roughly half, with prepayments mostly due to turnover making up the remainder, an external, staff-measured anchor for the scheduled-plus-turnover decomposition my $\beta_1 = 0$ null quantifies. That anchor and my null have different denominators: Na et al.'s composition describes realized principal \emph{payments} (roughly half scheduled, half turnover-driven prepayment, and my null is built from the same two ingredients: scheduled amortization plus a turnover floor calibrated from realized, partly behavioral turnover), while my 88.7\% (in-sample; 85.7\% at the off-window floor) describes the null's share of the shortfall \emph{against the caps}; a payments mix of scheduled-plus-turnover running below \$35 billion in essentially every month is the state in which the shortfall accrues mechanically, so their composition and my null recovery are the same fact read against two denominators, on the composition the null is built to reproduce. \citet{perli2024}, then Manager of the System Open Market Account, stated that the cap ``has not been binding \ldots resulting in very slow prepayments,'' and \citet{hammack2025} observed that some of the low-coupon holdings ``could be very long lived.'' New York Fed staff had projected the below-cap regime from its first month: a September 2022 staff post put monthly agency MBS principal payments at \$20 to \$30 billion over the following months, below the \$35 billion cap that took effect that month \citep{nyfed2022}. I treat the qualitative mechanism as established by these sources rather than as my own finding. Four features separate my exercise from theirs. First, this characterization attributes the below-cap runoff to lock-in without isolating the elasticity's contribution from the scheduled amortization and baseline turnover that would have undershot the cap under any rate-driven prepayment response above that baseline; Sections~\ref{sec:abm} and~\ref{sec:hazard} supply that decomposition and find the identified lock-in margin is bounded between $+2.9$ and $+8.7$ points of the benchmark at a turnover floor measured off the locked-in population (a wild-cluster interval on the floor read's 31 clusters; the narrower percentile read under-covers), $+5.6$ percentage points at the mid-grid anchor inside that interval ($+9.2$ at the in-window production calibration), and qualified in Section~\ref{sec:identification}. Second, I extend the same shortfall-against-cap construction from the \$450 billion-through-June-2024 quantity to the full June 2022 to November 2025 window, where it reaches \$764.7 billion (Section~\ref{sec:method-benchmark}). Third, I measure the shortfall against the Federal Reserve's own ex-ante projection as well as against the caps, separating the component the Federal Reserve itself anticipated from the genuine surprise (Section~\ref{sec:method-benchmark}). Fourth, the decomposition turns the staff observation into a statement about cap design: scheduled amortization and a turnover floor read from realized, partly behavioral turnover undershoot a \$35 billion monthly ceiling in essentially every month, so a cap at that level could not have bound under any rate response above that floor (Section~\ref{sec:identification}). A more recent \citet{eyal2026} FEDS Note also decomposes the Federal Reserve's balance-sheet reduction, but into macroeconomic drivers (inflation, real GDP growth, and active SOMA reductions relative to nominal GDP) rather than into the loan-level sources of the MBS runoff shortfall itself, so it does not speak to the lock-in-versus-mechanical question I take up here.

The Danish mortgage securitization system offers a tested structural alternative to negative convexity. \citet{berg2018} outlined the Danish ``Balance Principle,'' which requires strict one-to-one match funding. \citet[\S3.2.3]{berger2026}, in a January 2026 SSRN working draft, estimate that under these rules household mobility is close to flat in the coupon gap, because borrowers hold the right to repurchase their outstanding mortgages at market value, neutralizing extension risk for the bondholder. That flatness is a slope result on Danish data, and Section~\ref{sec:abm-danish} is careful about what it licenses: this paper imports the slope and anchors the level at the U.S. hazard, since Berger et al.\ estimate no moving level for U.S. households under a Danish payoff rule. What this strand supplies is the rule and the slope; what it lacks for the question here is a priced U.S. counterfactual. Section~\ref{sec:abm-danish} scores the rule-only transplant on the SOMA book directly, at an order of magnitude below a mechanism-extrapolated read, and Section~\ref{sec:pathb} prices the buyback credit's incidence, which decides whether the institutional gap survives the accounting.

A final strand situates the lock-in effect within the broader theory of monetary policy transmission. \citet{berger2021} demonstrate that the Federal Reserve's ability to stimulate the economy through rate cuts is path-dependent, because prepayment behavior depends on the outstanding distribution of in-the-money refinancing incentives. This paper extends that logic to the tightening direction: the same path dependence that weakens easing after a low-rate period also weakens tightening after a low-rate period, since the resulting stock of deeply in-the-money, rate-locked mortgages cannot prepay regardless of how aggressively the central bank subsequently raises rates. \citet{eichenbaum2022} formalize this state dependence for the easing direction (the potency of monetary policy depends on the distribution of outstanding rate gaps). \citet{beraja2019} provide complementary regional evidence that the strength of the refinancing channel varies systematically with the local distribution of in-the-money mortgages, and \citet{dimaggio2020} supply direct evidence for the easing-direction channel this paper mirrors, showing that the Federal Reserve's QE1 agency-MBS purchases transmitted to the real economy specifically by inducing refinancing. This QT setting also connects the paper to the broader literature on how central-bank balance-sheet operations transmit through the supply of long-duration assets: \citet{krishnamurthy2011} document that quantitative easing operates through several distinct channels, including one specific to agency MBS that runs through prepayment, and \citet{vayanos2021} formalize, under preferred-habitat demand, why shifts in the outstanding maturity and supply of held bonds move the term structure at all. Lock-in is thus an input to the balance-sheet-supply channel these papers describe, not a phenomenon separate from it. \citet{du2024} assemble the cross-country QT record---seven central banks, active and passive runoff---and find announcement effects concentrated in yields; the MBS runoff shortfall this paper decomposes is the passive-runoff margin of that record, measured at the portfolio cash-flow level. \citet{vickery2013} document that the liquidity premium of the TBA market derives specifically from the homogeneity of deliverable pools, finding that TBA-eligible securities trade at significantly lower cost than otherwise similar specified pools; this homogeneity requirement is the institutional constraint that, in Section~\ref{sec:conclusion}, I argue makes the Danish market-value repurchase structure difficult to reconcile with the current U.S. secondary market absent a broader redesign of agency securitization rules.

My two estimators sit at the transparent end of two active methodological literatures. On the survival side, a growing body of work fits prepayment and default hazards with nonparametric and machine-learning methods, including random survival forests \citep{ishwaran2008}, competing-risks Bayesian additive regression trees \citep{sparapani2020}, and deep neural-network mortgage hazard models \citep{sadhwani2021}. The last of these is the benchmark for what flexibility buys at scale: fitted to origination and monthly performance records for over 120 million U.S. mortgages, the deep network of \citet{sadhwani2021} learns highly nonlinear, interaction-rich prepayment and default responses, including the burnout-type history dependence that Path~A carries as an explicit stratum-level regressor and the ABM probes by pre-committed falsification (Sections~\ref{sec:patha} and~\ref{sec:abm-results}), without hand-specified functional form. I instead estimate a fixed-effects Poisson hazard and a literature-calibrated microsimulation whose functional forms and calibration inputs are written out in full, because the object I need is not maximal predictive fit but an auditable decomposition: the null of Section~\ref{sec:identification} that switches the lock-in elasticity off is interpretable only if every term it holds fixed is visible, which a flexible black-box hazard would obscure. Those machine-learning survival methods are a natural robustness avenue for the loan-level fit, and I probe its tabular end directly: Section~\ref{sec:patha} reports a gradient-boosted Poisson comparator and an isotonic-recalibration comparator, each given Path A's information set on the identical pre-committed holdout, and each recalibrating levels without gaining cross-cell predictive power; the deep loan-level variants remain future work on data-availability grounds. A complementary strand attacks nonstationarity rather than nonlinearity: landmarking re-anchors a simple hazard among the survivors at successive prediction times \citep{vanhouwelingen2007,putter2022}, and recent credit-risk work pairs landmark targets with explicit probability recalibration to keep survival predictions stable under covariate and base-rate drift, demonstrated on the same Freddie Mac performance data my hazard paths use \citep{peng2026}. My design meets the same concern with fixed, visible machinery (the stratum fixed effects, calendar-month dummies, and dynamic friction index of Section~\ref{sec:patha}) and prices the residual time-variation it cannot remove through the temporal moving-block bootstrap reported there. On the agent-based side, my calibration is fixed by the pre-committed parameterization of Appendix~\ref{sec:method-abm} rather than by an automated estimator, but it sits alongside a maturing literature on data-driven agent-based-model calibration and estimation, including simulated-minimum-distance and Bayesian methods \citep{grazzini2017,platt2020} and the empirical-validation taxonomy of \citet{fagiolo2007}, within which the pre-committed stress-test design of Section~\ref{sec:abm} is one validation strategy among several. The paired design-crossing that Section~\ref{sec:robustness-crossdesign} builds on top of it --- the ABM's decision rule re-run on real loan covariates, and the same survival structure run on a fully synthetic population, with interpretive thresholds fixed before the cross-design runs --- is offered as a validation design in its own right. Its nearest relative in that literature is indirect inference, which estimates an auxiliary model on real and on simulated data and minimizes the distance between the two sets of parameters \citep{platt2020}; the pair here crosses designs not to estimate parameters but to attribute a result gap to modeling paradigm versus data source, and it is not, to my knowledge, a named member of the validation taxonomy of \citet{fagiolo2007} or of the estimation approaches surveyed by \citet{grazzini2017} and \citet{platt2020}.

\section{Empirical Methodology}\label{sec:method}

\subsection{Two Structurally Distinct Estimators}\label{sec:method-estimators}

The first estimator is a 10,000-household Agent-Based Model (ABM), specified in Appendix~\ref{sec:method-abm}, which draws a synthetic population calibrated to aggregate FRED medians and tests whether a utility-maximizing household decision rule, rational or behaviorally extended, can account for the shortfall; its results are reported in Section~\ref{sec:abm}. The second is a hazard framework, specified in Section~\ref{sec:hazard}, which abandons the household utility function entirely and models loan-level survival directly on real Freddie Mac loan-level data: a stratum-month discrete-time hazard model estimated by Poisson pseudo-maximum likelihood on the full 2017--2021 origination universe aggregated to stratum-months (Path A), and a loan-level microsimulation with competing prepayment and default risks on a stratified 75,000-loan sample drawn from that universe (Path B).

Both estimators are evaluated against the same \$764.7 billion empirical benchmark (Section~\ref{sec:method-benchmark}) and the same June 2022--November 2025 QT window, so their results are directly comparable even though the estimators draw on different data and different behavioral assumptions. The two estimators differ in modeling paradigm (household choice versus loan-level survival), and, on the hazard side, in the use of real rather than synthetic loan-level data; Section~\ref{sec:robustness} finds that the difference in population size between them (10,000 households versus Path B's 75,000 simulated loans, with Path A estimated on the full origination universe) does not drive the gap in their explanatory power, though the data-source difference remains an unresolved confound that qualifies the paradigm interpretation throughout. Figure~\ref{fig:architecture} maps the pipeline.

\begin{figure}[H]
\centering
\includegraphics[width=\textwidth]{fig7_architecture}
\caption{Pipeline architecture. Data sources (FRED macro series, NY Fed SOMA holdings, Freddie Mac loan-level performance) feed three estimators (the synthetic-population ABM and hazard Paths A and B), whose simulated roll-off paths are scored through one shared macro-accounting layer against the \$764.7 billion empirical benchmark of Section~\ref{sec:method-benchmark}. The two robustness cross-links are the cross-design test (real Freddie covariates into the ABM, Section~\ref{sec:robustness-crossdesign}) and the symmetric companion test (synthetic population into the hazard framework, Section~\ref{sec:robustness-synthesis}).}
\label{fig:architecture}
\end{figure}

The agent-based model's specification, and the structural extensions whose grounding Section~\ref{sec:abm}'s results reference, are given in full in Appendix~\ref{sec:method-abm} and Appendix~\ref{sec:method-frictions}.

\subsection{Empirical Validation Benchmark}\label{sec:method-benchmark}

I construct the empirical shortfall benchmark from the New York Fed's SOMA current face value series, netting realized roll-off against the actual phased redemption cap schedule (\$17.5 billion per month from June through August 2022, stepping up to \$35 billion per month from September 2022 onward) strictly within the active QT window of June 2022 through November 2025. Net accumulation, rather than one-sided accumulation of shortfall months only, is essential for comparability with the Danish counterfactual of Section~\ref{sec:abm-danish}, where simulated roll-off routinely exceeds the cap. Figure~\ref{fig:benchmark} traces the construction: monthly roll-off against the phased cap, and the cumulative paths whose terminal gaps are the cap-relative benchmark and the expectations-based complement introduced below. The monthly series carries one artifact that no result in this paper reads as behavior: four months come in at exactly zero because the back-out's non-negativity clip pins any month whose reported roll-off falls below scheduled amortization and posts the displaced mass into the following month's reading, a month-boundary allocation artifact of the reporting series rather than a behavioral observation (run \texttt{h1\_zero\_months\_diagnosis}; Appendix~\ref{sec:robustness-seasonalfloor} traces the mechanism and shows that repairing the affected pairs flips this paper's frozen monthly-timing rule from pass to concede), which is one reason every claim made from this benchmark is a claim about levels.

\begin{figure}[H]
\centering
\includegraphics[width=0.94\textwidth]{fig9_benchmark_construction}
\caption{The empirical benchmark. Left: realized monthly SOMA agency MBS roll-off against the phased redemption cap; the shaded area is the monthly shortfall, which accrues on a net basis (the two circled opening months are net settlement inflows). Right: the same construction cumulatively. The terminal gaps are the \$764.7 billion cap-relative benchmark and, against the New York Fed's May 2022 ex-ante projection, the \$87.8 billion expectations-based complement. Monthly series from the frozen production run's metrics; the projection's monthly path is derived from the published year-end levels under the pre-committed uniform-spread rule and asserted, in the generator, to reproduce the frozen expectations-run totals.}
\label{fig:benchmark}
\end{figure}

This methodology yields an empirical shortfall of approximately \$764.7 billion (unrounded: \$764.748 billion; every percentage-of-benchmark figure in this paper divides by this value), which I adopt as the benchmark for the ABM's and the hazard framework's simulated trapped liquidity throughout Sections~\ref{sec:abm} and~\ref{sec:hazard}. The SOMA face underlying this benchmark spans the Federal Reserve's full agency MBS book (including Ginnie Mae collateral, 20.4\% of face), and 66.2\% of that face lies in the 2017--2021 vintages covered by the Freddie Mac loan-level sample of Section~\ref{sec:hazard}; the hazard framework's loan-level foundation therefore directly samples neither Ginnie Mae collateral nor roughly one-third of the book by vintage (the Ginnie dimension is externally bounded and scored by a published-speed overlay, and the vintage dimension is bounded by observed out-of-sample-vintage speeds from the replication ingest, both in Section~\ref{sec:pathb}; what each bound leaves open is stated in Section~\ref{sec:limitations}). The two exclusions must be netted on the book's own joint cells rather than counted separately: 33.7\% of face lies in out-of-window vintages (the 2022 cohort 23.1\%, pre-2017 10.6\%), and the 20.4\% Ginnie face share cuts across vintage rather than nesting inside it, so the conventional-and-in-window intersection is 51.0\% of face on the book's own agency $\times$ vintage joint cells (bounded above by 51.8\% for cell truncation; 52.7\% under an independence approximation), and the Freddie 2017--2021 conventional universe on which Path B is estimated directly represents about \emph{half} of SOMA book face. Two-thirds is the vintage-only marginal and overstates coverage. Each exclusion is bounded by an external-speed overlay (Section~\ref{sec:pathb}): the Ginnie Mae share at \$20--\$47 billion by published observed speeds, the out-of-window vintages at \$11.7 billion by observed cohort speeds, each signed toward overstating trapped liquidity. The empirical CPR back-out nets cohort-weighted scheduled amortization; an income-scaled curtailment assumption averaging approximately 0.84\% annualized, consistent with GSE daily prepayment report estimates, enters on the simulated side (added to each simulated roll-off and netted as a common dollar item in the shared accounting layer of Section~\ref{sec:robustness-hybrid}), so the back-out's CPR series retains the curtailment component. Two implementations of this back-out exist in the pipeline (the ABM side nets cohort-weighted, term-aware scheduled amortization, while the hazard side uses a single-pool amortization approximation), and both amortize at the pass-through WAC, while scheduled principal physically follows the borrowers' note rate, which is that WAC plus a guarantee-fee and base-servicing wedge of $+0.80$ points: recomputing the hazard-side back-out on the note-rate basis raises the CPR level by $+0.27$ points as a deterministic schedule difference while every trapped-dollar figure is bit-identical across the correction --- the benchmark's dollars never pass through the CPR scale --- and every peak lag and $U_2$ is unchanged at reporting precision (run \texttt{coupon\_convention\_amortization}; Table~\ref{tab:theil} is scored on the corrected basis), and any correlation computed against them differs only at the third decimal; Table~\ref{tab:estimators}'s note states which rows use which. 

\paragraph{Accounting bases.} Every recovery percentage in this paper is quoted on one of four accounting bases, which differ only in two orthogonal choices, each a pure accounting layer with no modeling content: whether the common curtailment flow is netted out, and whether the estimation sample is reweighted to the SOMA book's coupon mix. The \emph{standalone-scorer basis} makes neither adjustment. The \emph{benchmark-consistent shared basis} nets the common curtailment flow just described (\$69.56 billion over the window) from every U.S. leg identically, a flat 9.1 percentage points of the benchmark for every U.S. estimator. The \emph{full-book basis} instead reweights the standalone figure from the estimation sample's coupon mix to the SOMA book's weighted-average coupon (WAC), with no curtailment netting. The \emph{composed basis} applies both adjustments: the same 9.1-point netting to the full-book (book-coupon-reweighted) figure. Because the netting is bit-identical across the central and null legs of every comparison, it cancels exactly from the lock-in marginal and from every cross-estimator gap, an invariance that holds at every cell of the calibration box and under a $\pm 25\%$ stress on the assumed curtailment rate (derivation and stress in Section~\ref{sec:robustness-hybrid}; the four bases are collected for every hazard leg in Table~\ref{tab:bases}). The cancellation is term-by-term monthly, not merely aggregate: the netted flow is the product of a shared monthly rate and the shared actual-holdings path, so any time profile for curtailment (calendar-seasonal, or income elasticities that differ across sub-periods) cancels identically from the marginal and from every cross-estimator gap, moving only the common wedge. The functional form of the income-to-rate mapping is immaterial to the marginal for the same reason: any alternative map (concave, threshold, or piecewise in income) produces some monthly rate series, which is again a shared profile, so it can move only the common wedge. Main-text headline figures are quoted on the shared basis throughout.

I construct a second, expectations-based benchmark alongside the cap-relative one, because the cap-relative construction carries a tension the complement addresses: the caps were widely expected to be non-binding for agency MBS, so a shortfall measured against them mixes an anticipated mechanical component with genuine surprise. Measured instead against the Federal Reserve's own ex-ante projection, the baseline path published in the New York Fed's May 2022 Open Market Operations annual report \citep{omo2021}, whose projections start from the February 2022 balance sheet, incorporate the FOMC's May 2022 Plans, and model runoff from June 2022, with the published year-end figures spread uniformly within each printed year and clipped to the window under a rule fixed before the numbers were seen (the report prints year-end levels only, and no staff projection at monthly granularity is public for any vintage covering this window; the spread rule was pre-committed precisely because the source is coarser than monthly), projected agency MBS runoff over the window is \$740.6 billion against realized runoff of \$652.8 billion: an expectations-based shortfall of \$87.8 billion, 11.5\% of the cap-relative \$764.7 billion. Equivalently, the Federal Reserve's own ex-ante projection already implied 75.6--88.5\% of the realized cap-shortfall, the range spanning defensible intra-2022 allocations of the projection's printed annual net: \$676.9 billion, or 88.5\%, under the pre-committed uniform-spread central construction. That a source predating every estimate in this paper already anticipated the large majority of the shortfall is the substantive point, and it stands on its own. What does not stand is a sharper reading it invites. Under the uniform-spread construction the 88.5\% sits 0.2 points from the 88.7\% the $\beta_1 = 0$ mechanical null recovers on the benchmark-consistent basis, and I presented that near-equality as convergent evidence. I withdraw the presentation: the 0.2 points require two conditions at once. They require the in-sample 4.0\% turnover floor, because at the off-window floor the paper headlines the null recovers 85.7\% and falls 2.8 points \emph{below} the projection rather than 0.2 above it; and they require the uniform-spread allocation, because the settlement-aware allocation puts the anticipated share at 75.6\%, which the null clears by 13.2 points in-sample and 10.1 points at the headline floor. A gap that ranges from $-2.8$ to $+13.2$ points across two disclosed choices is not a convergence, and I now report the projection range and the null's recovery as two separately stated quantities. Both recoveries are production-floor-form reads: under the additive form the null recovers 35.6\% on the benchmark-consistent basis at the headline floor (44.7\% standalone; Section~\ref{sec:robustness-floor}), so the projection is being compared against one form's null rather than against a form-invariant object. The agreement that survives is coarse and is the one I rely on: under the production floor form both objects put the anticipated mechanical component in the large majority of the shortfall. That majority is a fact about where the caps were set rather than about the size of the rate response --- the caps sat above any prepayment path this book could have delivered --- so leading with it understates lock-in, and the reading that speaks to mortgage markets is the marginal's share of the expectations-based shortfall, reported below. On the allocation choice itself, the central 88.5\% rides on the uniform-spread rule's treatment of the printed 2022 annual net, which nets first-half settlement inflows against June--December runoff, and the settlement-aware allocation that assigns those inflows to the first half (the realized pattern: holdings rose \$92.3 billion over January--May 2022 in the benchmark's own SOMA series) is what lowers the anticipated share toward roughly three-quarters (75.6\%), leaving the large-majority reading, and every gate and threshold of the run, unchanged. Parametric within-year ramps bound the remaining allocation freedom: linear 2022 ramps move the anticipated share only to 88.0--89.0\%, the December-2025 window clip spans 86.1--90.6\% between its all-front and linear-back extremes, and the extreme bracket that pushes all 2025 runoff to the clipped December raises the share to 115.3\%; the mechanical-majority threshold survives at every allocation. Every percentage-of-benchmark figure in this paper remains a cap-basis quantity; dollar objects and estimator rankings are basis-invariant, the same dollar decomposition rescales by the $8.7\times$ benchmark ratio on the expectations basis, and the identified lock-in marginal amounts to 49\% of the expectations-based shortfall under the central uniform-spread construction at the headline off-window floor ($+\$42.6$ billion), or 80\% at the in-sample calibration point ($+\$70.3$ billion); because that shortfall is itself allocation-dependent, the ratio falls as the anticipated share falls (to roughly 23\% and 38\% respectively at the 75.6\% allocation) and is undefined at the extreme 115.3\% bracket, where the expectations-based shortfall turns negative. Every one of these ratios rests on one condition I cannot test: it reads the expectations-based shortfall as a residual free of lock-in, when the projection's own prepayment model may already have embedded part of the channel. I cannot recover that model from the published path of holdings, so whatever lock-in the projection already anticipated sits inside the \$740.6 billion rather than in the \$87.8 billion residual, and the ratios are upper bounds on lock-in's share of the genuine surprise, exact only if the projection embedded none of the channel; a second condition points the same way and is likewise untestable from the published record: the projection conditions on the spring-2022 market-implied rate path, well below what was realized, so part of the residual is rate-path surprise rather than lock-in, and the ratios bound the lock-in share of a mixed surprise. Both conditions are properties of the projection rather than of any one quotation of it, so they govern every site at which these ratios appear --- the introduction, this section, and the conclusion --- and each ratio is to be read as an allocation-conditional upper bound on lock-in's share of a mixed lock-in-and-rate-path surprise, not as a partition of that surprise. One question of standing precedes this accounting, and I state it rather than leave it implicit: a shortfall against the MBS cap is not by itself evidence that a stated policy objective was missed. The Committee's operating object was the aggregate securities portfolio's decline, Treasury runoff ran against its own caps throughout the window, and nothing in this paper shows that reserves or the aggregate path came in off the intended course because MBS ran slow. What the \$764.7 billion measures is the \emph{composition} of the decline --- MBS principal that did not arrive on the book's own schedule --- and the duration consequences that composition carries (Section~\ref{sec:discussion-fallout}); it is a statement about how the decline was delivered and what that delivery cost in portfolio structure, not a claim that the decline itself fell short.

The benchmark measures shortfall relative to the redemption cap, and the cap is a ceiling on runoff rather than a forecast of it: New York Fed staff projected the below-cap regime from the month the \$35 billion cap took effect \citep{nyfed2022}, so the shortfall relative to the cap is an upper bound on the shortfall relative to any reasonable ex ante expectation, an upper bound the complement above now quantifies: the expectations-based shortfall is \$87.8 billion, 11.5\% of the cap-relative figure. I retain the cap as the headline counterfactual for two reasons. First, the cap is the only policy-committed quantity available: the staff projection underlying the complement is a baseline forecast, not an FOMC-stated target, and adopting any forecast, the Federal Reserve's or a model's, as the headline benchmark would render it circular with respect to models of expected prepayment, the object under study. Second, the cap-relative construction is a policy-relevant quantity: it measures the gap between the maximum balance-sheet reduction the FOMC authorized and the reduction the mortgage structure delivered. Throughout the paper, ``trapped liquidity'' and ``shortfall'' accordingly denote the shortfall relative to the cap.

\section{Agent-Based Model Results: A Synthetic-Population Stress Test}\label{sec:abm}

This section is a stress test of the household-choice hypothesis, not the paper's main estimate, and the first thing to report about it is what its own cross-design leg establishes about its scope, because that result governs how everything below should be read. Re-running the same decision rule on real Freddie Mac structural covariates, under the recalibration rule the model already applies elsewhere, recovers 59.3\% of the benchmark on the committed draw and 60.2\% averaged over fifty seeds---above the 50\% threshold fixed before that run at which this section's paradigm reading would be undercut, and undercutting on all fifty of them---rising to 76.3\% once the sample is reweighted to the SOMA book's own coupon composition. Three things bound that reading and travel with it wherever it appears: the same reweight moves the frozen-calibration leg the other way, down to 12.6\%; the covariates swapped in are structural only, since the loan records contain no analogue for mobility desire or loss aversion, so no leg here replaces a synthetic behavioral population with a real one; and the verdict the run supports is that the contrast cannot, on that test alone, be cleanly attributed to modeling paradigm rather than to calibration and data source (Section~\ref{sec:robustness-crossdesign}). What this section falsifies is therefore a \emph{synthetic-population} household-choice specification: no specification tested here explains the Federal Reserve's trapped liquidity, and why it does not is not cleanly separable from the population it was tested on. A second qualifier travels with the verdict and is developed where it binds: the multi-vintage stage moved against its own predicted direction and the inversion is undiagnosed (Appendix~\ref{sec:robustness-extensions}).

\subsection{The Synthetic-Population Estimate: Baseline, Production, Extensions, and the Burnout Falsification}\label{sec:abm-results}

Section~\ref{sec:identification} \emph{declines} that same cross-design recovery as evidence that the imported lock-in elasticity is too small, on the ground that a recovery level is not what that design identifies; read as a scope limit on this section's verdict, which is how it is used above, it survives that objection, and neither use licenses the other. The production estimate is the synthetic-population special case, and it is the case in which the test bites: every specification tested here --- rational baseline, behavioral extensions, and a burnout variant --- moves further from, or fails to close, the \$764.7 billion benchmark, so I report them once and move to the hazard framework of Section~\ref{sec:hazard}. A purely rational baseline projects a structural cash flow deficit of \$544.0 billion (71.1\% of benchmark). The production specification adds scheduled amortization, a live CUSIP-level parse of SOMA holdings, a corrected empirical CPR back-out, a settlement-lag kernel, and the previously excluded 15-year MBS book in cohort weighting and scheduled amortization (structural terms only); its headline quantity is the seed distribution, not any single draw, at a fifty-seed mean of \$103.7 billion, or 13.6\% of the benchmark, with the frozen seed-42 draw (\$91.0 billion, 11.9\%, near the 34th percentile) retained as the reproducibility anchor the run manifests freeze. The kernel is retained in production and both frozen run manifests record it as applied, but it is a documented timing null rather than a fit improvement: a pre-fold-in ablation left the peak cross-correlation unchanged at $r = +0.192$ at lag $-3$ with the kernel on or off, and the kernel-applied production run peaks at $+0.195$ at the same lag. The production specification is not kernel-free, and every figure reported here is the kernel-applied run's. Seed dispersion dominates the third digit: the seed-level standard deviation is \$24.5 billion, a $\pm$1~SD share band of 10.4--16.8\%; the 95\% CI of the mean, \$96.9--\$110.5 billion, prices only simulation noise in the mean itself; and the 34th percentile is pinned only to within roughly $\pm$7 percentage points from 50 draws, so the ABM's share is better read as on the order of 10--17\% than as a three-digit point estimate (Figure~\ref{fig:mc}, Appendix~\ref{sec:method-abm}). The path fit is worse than the level fit and no extension repairs it (Figure~\ref{fig:abmpaths}, Appendix~\ref{sec:method-abm}): simulated mean CPR runs at 11.68\% against an empirical 5.14\%, the raw path correlation is negative at lag 0 ($r = -0.318$) with no value inside a $\pm$3-month window exceeding the +0.19 at the three-month lead, and the path $R^2$ is $-6.984$ (throughout, $R^2$ is computed as $1-\mathrm{SSE}/\mathrm{SST}$ against a mean-only null, so a negative value indicates fit worse than a constant at the empirical mean). Layering in three literature-grounded behavioral extensions --- a debt-to-income constraint with asymmetric loss aversion and a rate-shock freeze, income-scaled curtailment, and multi-vintage coupon composition --- moves the model to 54.9\%, then 45.1\%, then 33.7\% of the benchmark (Figure~\ref{fig:waterfall}); the first move is a recalibration artifact rather than evidence for loss aversion, since holding a real-world turnover floor fixed under a steeper loss-aversion penalty mechanically inflates rate sensitivity, and none of the three touches the CPR path's central failure, each operating on the level of simulated mobility rather than its timing. A fourth variant fails more informatively: path-dependent prepayment ``burnout,'' implemented as closed-population survivor selection, collapsed simulated CPR to approximately zero and raised trapped liquidity to 147\% of the benchmark, because a closed population of 10,000 agents depletes its mobile mass during the low-rate 2020--2021 window before QT even begins --- an absorbing state incompatible with a real MBS pool that continuously receives new originations, and \citepos{lesniewski2026} selection identity in its degenerate limit. It is a negative finding, not used in production, and it is why burnout is treated at the cohort rather than the per-agent level in Section~\ref{sec:hazard}. Figure~\ref{fig:recovery} places these levels beside the other estimators. The decision rule, every behavioral parameter and its source, the macroeconomic-friction and structural extensions, the mechanism-level detail behind all three behavioral stages, and the 15-year fold-in are specified in Appendices~\ref{sec:method-abm}, \ref{sec:method-frictions}, \ref{app:abmspec}, \ref{sec:robustness-extensions} and~\ref{sec:robustness-pipeline} and are reproduced end to end by the replication package (Appendix~\ref{app:ledger}); this section reports results only.

\begin{figure}[H]
\centering
\includegraphics[width=\textwidth]{fig1_recovery_by_estimator}
\caption{Benchmark recovery by estimator, on the shared accounting basis (standalone-scorer values in parentheses, both bases collected in Table~\ref{tab:bases}; ABM-family rows run through that layer natively). The production ABM row shows the fifty-seed mean (open diamond) with its $\pm 1$~SD band alongside the frozen seed-42 draw (solid dot); Path A's dot is an aggregate-context, non-corroborating exhibit excluded from headline figures (Section~\ref{sec:patha}). Raw lag-0 correlations are reference-only (figure footer; Section~\ref{sec:pathb}); level recovery alone does not rank models.}
\label{fig:recovery}
\end{figure}

\begin{figure}[H]
\centering
\includegraphics[width=\textwidth]{fig3_abm_waterfall}
\caption{The ABM stress-test bridge. Each stage applies the next correction or extension to the household-choice model, and every pre-committed behavioral extension moves the recovery away from, not toward, the benchmark (the multi-vintage stage against its own expected direction, Appendix~\ref{sec:robustness-extensions}): rational baseline 71.1\% $\to$ behavioral gates 54.9\% (a recalibration artifact, Appendix~\ref{sec:robustness-extensions}) $\to$ income-scaled curtailment 45.1\% $\to$ multi-vintage composition 33.7\% $\to$ production corrections 13.2\% (a data-and-accounting stage, not a behavioral one) $\to$ 15-year fold-in 11.9\% (the production headline of Section~\ref{sec:abm-results}; the fold-in itself is worth $-1.3$ points) $\to$ native 15-year gate 11.1\% (the Berger-run calibration-stage variant of Section~\ref{sec:abm-danish}). Stage levels from the frozen run manifests catalogued in the run ledger (Appendix~\ref{app:ledger}).}
\label{fig:waterfall}
\end{figure}

\subsection{Interpretation}\label{sec:abm-interp}

Every ABM specification tested here either widens the gap to the benchmark or leaves it unchanged, its CPR path stays uncorrelated or negatively correlated with the empirical path throughout, and Section~\ref{sec:robustness}'s scale-sensitivity replication finds no size effect on the production specification, so what remains is either that household-level decision-making does not govern the majority of the shortfall, or that this specification is missing a mechanism. The cross-design leg reported at the head of this section bears directly on that either-or and does not settle it: handing the same rule real structural covariates moves the recovery from 13.6\% to 60.2\% averaged over fifty seeds, which is what a missing-population problem looks like rather than a missing-mechanism one, but the behavioral layer stays synthetic in both legs, so the alternative is narrowed rather than eliminated. The clean falsification is the one that bypasses that free parameter entirely: the simulated-minimum-distance two-moment fit (Section~\ref{sec:robustness-crossdesign}) returns joint infeasibility under a verdict vocabulary fixed before execution, its nearest admissible fit overshooting the benchmark at 107.1\%, and it does not run through the recalibration scale that qualifies the cross-design reading. Section~\ref{sec:hazard} tests the first branch directly on a different data source and modeling paradigm; whichever explanation holds, these results weigh against household mobility choices, rational or behavioral, as the primary price the Federal Reserve paid for secondary-market liquidity, and the institutional counterpart --- whether the par-payoff rule itself, rather than household behavior, drives the trap --- is examined as a securitization-policy alternative in Section~\ref{sec:abm-danish}. One asymmetry in that comparison should be stated rather than left for a reader to notice. Section~\ref{sec:identification} insists that this design identifies a differential and not a level, and the ABM is nonetheless compared with Path B on levels, because the ABM admits no counterpart to the $\beta_1 = 0$ null. The reason is structural. In \eqref{eq:pathB} the elasticity is a separable term sitting over a floor that is measured independently and survives the elasticity's removal, so the null is still anchored to observed baseline turnover and lands at 85.7\% of the benchmark. The ABM has no such term: its 4--5\% turnover floor is not a component of \eqref{eq:abmmove} but a \emph{product} of it, obtained by searching the mobility-desire scale $\theta$ until the penalty-bearing rule returns 4--5\% CPR at an 8\% market rate, so zeroing the penalty removes the floor along with the channel. With the production $\theta = 43{,}882.8$ held fixed, the resulting leg returns 40.7\% CPR at the calibration point against the production rule's 4.8\%, over-retires the book, and scores $-259.4$\% of the benchmark; re-deriving $\theta$ under the null restores the anchor to 4.7\% but scores $+116.3$\%, because the re-search fits the null to the very turnover level the differential is meant to measure. The two readings imply marginals of $+270.4$ and $-105.3$ points. They do not bracket a magnitude; they disagree in sign, and neither is interpretable. I therefore report no ABM differential, and the estimator contrast of Section~\ref{sec:hazard} remains a contrast of levels, with the qualification that entails (run \texttt{abm\_null\_feasibility}; a feasibility probe, not a pre-committed estimate, and nothing in this paper's ranges depends on it).

\section{The Hazard Framework: A Structurally Distinct Estimator}\label{sec:hazard}

This is the estimator on which the paper's central decomposition rests. The ABM of Section~\ref{sec:abm} does not approach the hazard framework's recovery at the aggregate dollar level (91.3\% in the microsimulation form on the shared accounting basis of Section~\ref{sec:robustness-hybrid} at the headline off-window floor, 97.9\% at the in-sample calibration point; the estimated cohort-hazard form corroborates at 112.4\% but is excluded from headline ranges for the reasons Section~\ref{sec:patha} states): the ABM's production specification recovers 11.9\% on the frozen draw (fifty-seed mean 13.6\%), and Section~\ref{sec:robustness-crossdesign}'s cross-design leg leaves open whether that contrast belongs to modeling paradigm or to data source. Even its highest-recovery behavioral extension (the 54.9\% figure I attribute in Section~\ref{sec:abm-results} to a mechanical recalibration artifact rather than to loss aversion as a mechanism) falls well short. (The uncorrected rational baseline reaches 71.1\%, but that figure predates the scheduled-amortization, CUSIP-parse, and CPR-back-out corrections described in Section~\ref{sec:abm-results} and is not part of the production estimate family.) The hazard framework's recovery range, by contrast, rests on a different data source and modeling paradigm. Both hazard paths are evaluated against the same aggregate benchmark as the ABM, but only Path B's recovery is accurate on the benchmark-consistent accounting basis, and no estimator's monthly co-movement with the empirical path, Path B's included, survives removal of the window's shared trend (Section~\ref{sec:pathb}). Path A's monthly path fit is no better than the ABM's, a limitation Sections~\ref{sec:patha} and~\ref{sec:hazard-interp} quantify and assess, because it bears directly on how much evidential weight each recovery figure can carry.

\subsection{Motivation}\label{sec:hazard-motivation}

Section~\ref{sec:abm}'s pattern is difficult to reconcile with the hypothesis that the residual reflects a missing rational or behavioral household mechanism that a sufficiently rich utility function would eventually capture. It is more consistent with the hypothesis that the binding constraints on the empirical prepayment path lie at least partly outside the household decision function: in loan-level heterogeneity, servicer pipeline dynamics, and cohort-level path dependence that a fixed, representative-household population struggles to represent.

The hazard framework introduced in Section~\ref{sec:method-estimators} is built to test this hypothesis. It abandons the household utility function: no mobility-desire parameter, no loss-aversion coefficient, no utility-maximizing decision rule; it instead models loan-level survival directly, using proportional hazards (proportional above a baseline turnover floor; Section~\ref{sec:pathb}) estimated on real loan-level data, with the closed-population survivor-selection failure of Section~\ref{sec:abm-results} motivating a cohort-level rather than per-loan treatment of burnout. The framework retains the same \$764.7 billion empirical benchmark and the same QT window definition as the ABM for comparability, and comprises two paths: a stratum-month discrete-time hazard model estimated on aggregated loan-level data, and a loan-level microsimulation calibrated to literature elasticities with explicit competing risks.

\begin{table}[H]
\centering
\caption{Compact specification of the two hazard paths. Every entry restates
content specified in Sections~\ref{sec:patha} and~\ref{sec:pathb} or
Appendix~\ref{app:params}.}
\label{tab:specbox}
\footnotesize
\begin{threeparttable}
\begin{tabular}{@{}p{2.2cm}p{5.9cm}p{5.9cm}@{}}
\toprule
 & Path A (estimated cohort hazard) & Path B (literature-calibrated microsimulation) \\
\midrule
Unit; sample & stratum-month; full Freddie 2017--2021 origination universe (296 four-way cells) & loan-month; stratified 75{,}000-loan draw from the same universe \\
Outcome; estimation & prepaid UPB with log-exposure offset; Poisson pseudo-maximum likelihood, log link \eqref{eq:pathA} & monthly prepayment hazard \eqref{eq:pathB}; nothing estimated on the benchmark; every parameter fixed a priori, though the dominant level-setter (the turnover floor) is calibrated in-window (Sections~\ref{sec:pathb}, \ref{sec:robustness-floor}; Appendix~\ref{app:params}) \\
Age profile & seasoning spline, knots $\{12,24,36,60,84,120\}$ months (seven columns) & 100~PSA ramp to 6\% CPR at month 30 \\
Rate-gap term & estimated: $+0.65$ standardized ($+0.29$ per 100 bp), spec v4 & imported from \citet{liebersohn2024} via \eqref{eq:beta1}; 5.5--7.7\% band \\
Stratification & 295 stratum fixed effects (vintage $\times$ coupon $\times$ FICO $\times$ LTV) & stratum-level burnout state $B_{s(i),t}$; standardized FICO/LTV ex-ante priors \\
Time effects & eleven calendar-month dummies (spec v4) and the dynamic friction index & none; competing default hazard \eqref{eq:default} and the Markov delinquency pipeline; the omitted housing-activity level term is priced rather than modelled, by the scaled-null variant of Section~\ref{sec:identification} (run \texttt{scaled\_null\_housing\_activity}) \\
Turnover floor & none (stratum fixed effects set baseline levels) & 4\% CPR floor, max form; competing-risks form tested (Section~\ref{sec:robustness-floor}) \\
Inference & stratum-cluster and temporal moving-block bootstraps (Table~\ref{tab:bootstrap}); no claim survives them---the rate-gap coefficient's sign is not distinguishable from zero under any bias-respecting construction (Section~\ref{sec:patha}) & stratified loan-level bootstrap (200 replicates) for sampling; calibration box for parameters (Section~\ref{sec:robustness-floor}) \\
Validation & temporal holdout, reported as a diagnostic only (Section~\ref{sec:patha}) & Fannie Mae external replication (marginal $+8.68$ points against Freddie's $+9.20$; the pre-committed envelope it clears is 11.06 points wide) and the synthetic companion at 106.0\% (Sections~\ref{sec:pathb-fannie}, \ref{sec:robustness-synthesis}) \\
Role & aggregate context only; does not corroborate the result; excluded from headline figures & headline estimator: at the off-window floor central 91.3\%, $\beta_1 = 0$ null 85.7\%, marginal $+5.6$ points; at the in-sample calibration point 97.9\%, 88.7\%, $+9.2$ points (all shared basis) \\
\bottomrule
\end{tabular}
\end{threeparttable}
\end{table}

\subsection{Path B: Literature-Calibrated Microsimulation}\label{sec:pathb}

The first hazard path brings competing risks, literature-calibrated elasticities, and a cleaner isolation of the Danish counterfactual, without reintroducing the utility-maximizing machinery of the ABM. It simulates a stratified sample of 75,000 loans (40,234 of which are still active when the QT window opens; Appendix~\ref{app:params}), each carrying static covariates (FICO, region, original LTV, coupon, and a stratum identifier matching the Path A specification) and dynamic state (loan age, balance, and a regime-specific rate gap). Equity enters only through the static original-LTV covariate: no updated CLTV and no local house-price feed exists anywhere in the pipeline, so equity dynamics over the window are unmodeled---a disclosed design boundary for this turnover-dominated, deeply out-of-the-money window, with the FICO/LTV covariate ablation reported below bracketing what the static block contributes. Each loan's monthly hazard of prepayment is specified as
\begin{equation}\label{eq:pathB}
h^{\mathrm{prep}}_{i,t} = \max\!\left\{\, \underline{h},\; h_0\!\left(a_{i,t}\right) \exp\!\left( \beta_1\, g_{i,t} + \beta_F\, z^{\mathrm{FICO}}_{i} + \beta_L\, z^{\mathrm{LTV}}_{i} + \beta_B\, B_{s(i),t} \right) \right\},
\end{equation}
with a competing hazard of default
\begin{equation}\label{eq:default}
h^{\mathrm{def}}_{i,t} = h^{\mathrm{def}}_0 \exp\!\big( \gamma_r\, \sigma_{i,t} + \gamma_F\, z^{\mathrm{FICO}}_{i} + \gamma_L\, z^{\mathrm{LTV}}_{i} \big), \qquad \sigma_{i,t} = \max\{0,\ -\tilde g_{i,t}\},
\end{equation}
where $\sigma_{i,t}$ is rate stress ($\tilde g_{i,t} = g_{i,t}/100$ is the rate gap of \eqref{eq:pathB} expressed in decimal units, so $\sigma$ is the decimal magnitude of an out-of-the-money gap), and $h^{\mathrm{def}}_0 = 3\times10^{-4}$ monthly, a prior implying $0.36\%$ per year against a realized Freddie Mac default rate of $0.31\%$ over the window on this paper's own GMAR extraction---corroboration the calibration did not use, since the prior was fixed before that series was assembled. Every Path B parameter is fixed a priori in the configuration and none is estimated on the benchmark; Appendix~\ref{app:params} collects the full parameter vector with provenance, together with the delinquency transition matrix. Table~\ref{tab:hazard-notation} maps every symbol here to its Path A counterpart. Here $g_{i,t}$ is the rate gap (the loan's coupon minus the prevailing market rate, in 100 bp units, negative when the borrower is locked in, so that the positive $\beta_1$ of \eqref{eq:beta1} suppresses prepayment exactly when the par-payoff penalty binds); $B_{s(i),t} \in [0,1]$ is the cohort-level burnout state of loan $i$'s stratum $s(i)$ (the stratum's cumulative prepaid UPB as a share of its original UPB, capped at one, initialised at zero for every stratum at the window open so that it accumulates only within-window prepayment against an original-UPB denominator that still includes pre-window-terminated loans---window-cumulative, not life-cumulative---so $\beta_B = -0.5$ scales the hazard by $e^{-0.5 B} \in [0.61, 1]$); and $\underline{h}$ is the baseline turnover floor: the monthly single-month-mortality (SMM) equivalent of 4\% annual CPR, $\underline{h} = 1-(1-0.04)^{1/12} \approx 0.0034$. It floors total turnover, not a strictly involuntary component of it: deep-discount-cohort turnover at this level includes discretionary life-cycle moves and cash-out refinancings alongside strictly involuntary events, no component decomposition exists in this design, and the strictly-involuntary share---death, divorce, and forced relocation---is plausibly well under half. Where the floor's semantics do institutional work---selecting the U.S.-intercept Danish anchor (Section~\ref{sec:abm-danish}) and retaining the max floor form (Section~\ref{sec:robustness-floor})---the operative premise is only that baseline turnover cannot fall below its observed deep-discount level when the payoff rule changes, which the total-turnover reading supports; those choices are conventions consistent with that premise, not consequences of a strictly-involuntary decomposition. The floor's level is anchored to the same empirical observation as the ABM's mobility calibration (Appendix~\ref{sec:method-frictions}): turnover on deeply out-of-the-money discount cohorts during 2023--2024, which lies \emph{inside} the June 2022--November 2025 evaluation window. While no Path B parameter touches the benchmark dollar figure, the specification's dominant level-setter is measured from in-window prepayment behavior. The verification content below is the marginal over the $\beta_1 = 0$ null rather than the aggregate level itself. The pre-committed floor sweep of Section~\ref{sec:robustness-floor} shows that even this marginal is conditional on the floor calibration: it spans $+10.8$ to $+5.5$ points at the central elasticity across the empirically plausible 3--5\% floor range ($+9.2$ at the production 4\% floor), with that interval, not any point inside it, the identified object; its positivity at every floor and band point swept is a consistency check the window's rate configuration cannot let fail (Section~\ref{sec:identification}). The floor's functional form matters too: replacing the hard maximum of \eqref{eq:pathB} with a competing-risks combination, under which the elasticity is never censored, moves the production marginal to $+11.25$ points and makes it nearly floor-invariant ($+11.29$ to $+11.21$ across the same 3--5\% range; Section~\ref{sec:robustness-floor}). The max-form marginal's floor-dependence is censoring-driven rather than evidence about the elasticity. The floor binds in 36\% of the 1,683,124 evaluated loan-months in the production U.S. run. A substantial share of the aggregate level rides on the floor rather than the elasticity, consistent with the no-lock-in null of Section~\ref{sec:hazard-interp}, which shows the mechanical model recovers 97.8\% of the benchmark under the standalone scorer (88.7\% on the shared accounting basis). Those are production-floor-form levels: under the additive form the same null recovers 44.7\% at the off-window floor against the production form's 94.8\%, both on the standalone scorer (Section~\ref{sec:robustness-floor}). Proportionality in \eqref{eq:pathB} accordingly holds above the floor, not through it. The baseline hazard $h_0$ follows the 100~PSA seasoning ramp (CPR rising 0.2 percentage points per month of loan age to 6\% at month 30) applied at monthly frequency via $\mathrm{SMM} = 1-(1-\mathrm{CPR})^{1/12}$. The rate-gap coefficient is calibrated from \citet{liebersohn2024}, who estimate that a 100 bp increase in the prevailing rate above the locked-in coupon reduces quarterly mobility probability by 5.5\% to 7.7\% depending on specification; I adopt the 6.5\% midpoint of that band as the ex-ante central scenario, with the band edges as sensitivity bounds\footnote{The central 6.5\% is the midpoint of that band, adopted as the ex-ante central scenario rather than a published point estimate, with 5.5\% and 7.7\% as sensitivity bounds. Quarterly probability shocks are mapped to the proportional-hazard coefficient $\beta_1$ via the survival-function transform in \eqref{eq:beta1}, evaluated at a baseline turnover probability of 6\% in the transform's quarterly slot---a turnover-proxy calibration auxiliary rather than a measured quarterly mobility rate, and well above \citepos{liebersohn2024} own 1.5\% zero-gap quarterly moving level---that affects only the conversion level, not the elasticity band itself, and only weakly (derivation and range in Appendix~\ref{app:params}). \citet{fonseca2024} report a related but distinct estimand (a roughly 9\% reduction in moving rates per percentage point of mortgage-rate delta), which I cite as independent corroboration of the lock-in magnitude, not as a source of the band.}, converted to a monthly hazard via the survival-function transform
\begin{equation}\label{eq:beta1}
\beta_1 = -\ln\!\left[\frac{1-\big(1-(1-\delta)\,P_q\big)^{1/3}}{1-\big(1-P_q\big)^{1/3}}\right],
\end{equation}
where $\delta$ is the proportional quarterly mobility decline per 100 bp of rate gap, $P_q = 0.06$ a reference baseline turnover probability occupying the transform's quarterly slot---a conversion auxiliary rather than a measured quarterly mobility rate, well above \citepos{liebersohn2024} own 1.5\% zero-gap quarterly moving level (Section~\ref{sec:robustness-crossdesign}), and retained because the conversion is nearly insensitive to it over $P_q \in (0, 0.12]$ (Appendix~\ref{app:params})---and the inner terms are the monthly single-month mortalities implied by constant-hazard compounding within the quarter, $1 - P_q = (1 - h_m)^3$. Equation~\eqref{eq:beta1} is the transform implemented in the codebase: at the central $\delta = 0.065$ it evaluates to $\beta_1 = 0.0686$ (unrounded $0.068571$), the value carried in the production configuration and rounded to $0.069$ in Table~\ref{tab:params}. The superseded $0.0693$ is a distinct object, the continuous-hazard variant $-\ln\!\big[\ln(1-(1-\delta)P_q)/\ln(1-P_q)\big] = 0.0693$ (unrounded $0.069258$), which the simulation does not use; the two forms differ by $0.0007$ and both round to $0.069$, but only \eqref{eq:beta1} reproduces the reported results. The transform yields the dimensionless log-hazard-ratio implied by a quarterly probability decline (it replaces a divide-by-three approximation, not a hazard itself), and \eqref{eq:pathB} then applies $\beta_1$ linearly per 100 bp of gap. At QT-era gaps of roughly $-380$ bp this compounds the local quarterly estimate to a hazard ratio of about $0.77$ ($\approx e^{3.8 \times -0.069}$), a log-linear extrapolation beyond the range in which the source elasticity was estimated. I flag it as a functional-form assumption and test it directly: a piecewise-concave transform (full marginal effect over the first 200 bp of lock-in, half beyond) moves the recovery from 107.0\% to 105.7\% (\$808.4 billion) and slightly improves the contemporaneous path correlation ($+0.281$ versus $+0.190$, peak lag unchanged at $-3$). The log-linear extrapolation is not load-bearing for the level, and is for the marginal---but the asymmetry is one of proportion, not of magnitude, and the same absolute $\pm 1$-point convention cannot support both readings. A pre-committed rerun of the paired legs under the concave transform (run \texttt{concave\_marginal}) puts the concave marginal at $+7.87$ points against the production $+9.20$ at the in-sample floor, and at $+5.06$ against $+5.57$ at the off-window point floor. At the in-sample floor the concave null is bit-identical to the production null (\$748.185 billion on both legs, because $\beta_1 = 0$ leaves nothing for the transform to act on), so the level move and the marginal move are numerically the \emph{same} 1.32 points. That 1.32 points is 1.2\% of a 107.0\% recovery and 14\% of a $+9.20$-point marginal, and it is that relative reading, not the absolute one, on which the transform is immaterial to the level and material to the marginal; on the absolute $\pm 1$-point materiality threshold the floor-form test of Section~\ref{sec:robustness-floor} applies, both moves clear it, and the transform dimension travels with the marginal's operative uncertainty in Table~\ref{tab:uncertainty}.

One more assumption travels with the import, and nothing here tests it. The decline \citet{liebersohn2024} estimate is a proportional fall in a quarterly ZIP-code moving probability, while \eqref{eq:pathB} applies it to a loan's total prepayment hazard, which contains refinancing as well as moving and which is floored at a turnover level that, as noted above, itself mixes discretionary life-cycle moves and cash-out refinancings. No split of realized prepayment into moving and refinancing components exists in this design, so treating the two objects as one is a specification assumption rather than a result. Two further transport assumptions travel with the import and are named here rather than left implicit: the source estimand is a ZIP-code-level moving hazard applied as a loan-level log-hazard coefficient, an aggregation transport that assumes response homogeneity; and it is estimated on a population not conditioned on mortgage survival, while \eqref{eq:pathB} applies it to a pool from which 34,734 of the 75,000 sampled loans had already prepaid before the window opened --- a survival-selected pool depleted of the rate-responsive, so this transport's direction is signed toward attenuation rather than amplification, and priced: under the gamma-frailty selection identity the attenuation is $a = S^{\theta}$ at the measured pre-window survival share, and the off-window marginal reads $+5.1$/$+4.6$/$+3.7$ points at $\theta = 0.25/0.5/1$ (run \texttt{attenuation\_sensitivity}; $\theta$ is not estimable in this design and is displayed as the free parameter --- Path A's burnout coefficient, the one object that speaks to it, is not distinguishable from zero under either bootstrap scheme).

Each month, both hazards are evaluated for every active loan. In the production configuration prepayment is fractional: each active loan's balance is drained deterministically by $h^{\mathrm{prep}}$ (a per-loan uniform draw for prepayment is a non-production discrete branch), while a uniform draw governs the default hazard, sending hits into the delinquency pipeline (current loans to D30, already-delinquent loans to default), where they transition through the Markov delinquency-resolution matrix of Appendix~\ref{app:params}. A proportional rescaling guards the case $h^{\mathrm{prep}} + h^{\mathrm{def}} > 1$. It is a safety clamp, not a competing-risks construction. It never binds in production: across all 1,683,124 evaluated loan-months per regime the maximum hazard sum is 0.016, two orders of magnitude below the threshold, at which level the exact exponential cause-allocation construction $1-\exp(-(\lambda_p+\lambda_d))$ with shares $\lambda_p/(\lambda_p+\lambda_d)$ differs from mine only at second order in the hazards. In the taxonomy of standard competing-risks formulations \citep{putter2007}, Path B is a discrete-time multi-state model parameterized entirely through cause-specific hazards: Equations~\eqref{eq:pathB} and \eqref{eq:default} are monthly event probabilities conditional on the loan remaining active (the estimand of discrete-time cause-specific regression, as implemented in open-source packages such as PyDTS \citep{meir2025}), with the delinquency pipeline of Appendix~\ref{app:params} extending the two-cause structure to a multi-state process with cures rather than an immediately absorbing default. No subdistribution (Fine--Gray) hazard \citep{fine1999} appears in the specification, and none is needed: that device places covariate effects directly on cumulative incidence functions, whereas the simulation obtains every cumulative-incidence functional (prepaid UPB, the delinquency stock, trapped liquidity) by forward composition of the cause-specific hazards. One caveat qualifies the mapping: in production, prepayment is a fractional balance-flow intensity rather than a per-loan absorbing transition, so the cause-specific reading is exact for default and holds at the balance-aggregate level for prepayment. Prepaid balances settle to simulated SOMA cash flow in the month they occur; no settlement delay is modeled on this path.

\begin{table}[H]
\centering
\footnotesize
\setlength{\tabcolsep}{4pt}
\begin{threeparttable}
\caption{Path B parameters.}
\label{tab:params}
\begin{tabular}{@{}llll@{}}
\toprule
Parameter & Value & Units & Provenance \\
\midrule
$\beta_1$ (central) & $0.069$ & per 100 bp gap & \eqref{eq:beta1} at $\delta = 0.065$ \citep{liebersohn2024} \\
$\delta$ band & $0.055$--$0.077$ & quarterly decline & \citet{liebersohn2024} \\
$P_q$ & $0.06$ & quarterly slot & turnover proxy (conversion auxiliary; not a measured quarterly mobility rate) \\
$\beta_F$ (FICO, z) & $-0.15$ & per SD & calibration prior; panel estimate $-0.39$ (see text) \\
$\beta_L$ (LTV, z) & $+0.10$ & per SD & calibration prior; panel estimate $+0.21$ (see text) \\
$\beta_B$ (burnout) & $-0.5$ & per unit of $B \in [0,1]$ & calibration prior \\
$\underline{h}$ (floor) & $0.0034$ & monthly SMM & 4\% annual CPR (baseline turnover) \\
PSA speed & 100 & -- & industry seasoning convention \\
$h^{\mathrm{def}}_0$ & $3\times10^{-4}$ & monthly & calibration prior \\
$\gamma_r$ (rate stress) & $2.0$ & per decimal unit & calibration prior ($\approx$2\% per 100 bp) \\
$\gamma_F$ (FICO, z) & $-0.20$ & per SD & calibration prior \\
$\gamma_L$ (LTV, z) & $+0.25$ & per SD & calibration prior \\
\bottomrule
\end{tabular}
\end{threeparttable}
\end{table}

The transform above prices how the elasticity extrapolates in the gap; a second mapping assumption is what the elasticity applies to. \citepos{liebersohn2024} estimand is a moving probability, and \eqref{eq:pathB} applies the implied $\beta_1$ to the entire voluntary hazard, moving and refinancing alike. A pre-committed bracket prices that convention (run \texttt{moving\_share\_bracket}; the gap response enters through a mixture, with $1-s$ of the voluntary hazard held at its no-response level and $s$ of it carrying the full response, the turnover floor unscaled, and parity at $s = 1$ against the committed central leg plus an algebraic identity probe fixed in the script header before any run): at the production calibration the marginal is $+\$37.7$ billion ($+4.9$ points) at $s = 0.5$ and $+\$19.2$ billion ($+2.5$ points) at $s = 0.25$, against the production $+\$70.3$ billion at $s = 1$, close to proportional in $s$. Re-run at the 4.991\% off-window floor the paper headlines, the bracket is $+\$25.7$ billion ($+3.36$ points) at $s = 0.5$ and $+\$13.9$ billion ($+1.82$ points) at $s = 0.25$ against $+\$42.6$ billion at $s = 1$ (run \texttt{moving\_share\_bracket\_offwindow}), mildly super-proportional because the higher floor censors the elasticity in more loan-months and censored contributions do not scale with $s$. For this book's rate configuration the convention is therefore signed: applying the imported elasticity to less of the hazard can only shrink the marginal, and by roughly its share. No loan-purpose field exists in the ingested data, so $s$ is a bracketing parameter rather than an estimate, and the production convention $s = 1$ remains the reported form. \citet{graybill2026} corroborate the imported band on exactly this margin: their transaction-record estimate of the sale-hazard response --- roughly 7--11\% per percentage point of rate gap, from records that touch no credit file --- is an estimate of the moving component that $s$ varies.

The floor enters \eqref{eq:pathB} as a single constant $\underline{h}$ broadcast to every loan-month, carrying no rate-gap, calendar-time, or covariate argument, so the specification treats baseline turnover as acyclical and assigns every rate-responsive movement in prepayment to the elastic channel. The marginal-over-null design cancels the floor's level, since both the central and the $\beta_1 = 0$ legs carry the identical constant, but it does not cancel a floor that itself co-varies with the rate cycle: to the extent that involuntary moves are themselves suppressed when a household is locked in, that co-movement is absorbed into the estimated elastic marginal. The identified marginal at the in-window calibration is therefore best read as total rate-attributable prepayment suppression, an upper bound on the strictly-voluntary elasticity's contribution rather than a clean voluntary-versus-involuntary decomposition. At the off-window calibration the ordering reverses: the demotion's window component (roughly \$20.9 billion; Section~\ref{sec:robustness-floor}) is in-window baseline turnover that the off-window floor classifies as rate-suppressed and the marginal no longer counts, so the off-window $+\$42.6$ billion bounds total rate-attributable suppression from below (the matched-depth $+\$63.5$ billion restores that component), and the two calibrations bracket two estimands rather than conservative and liberal readings of one. The level sweep of Section~\ref{sec:robustness-floor} already brackets any plausible constant-floor excursion (the 3--5\% range keeps the marginal within $+10.8$ to $+5.5$ points at the central elasticity), and the within-run cyclical variant is now tested directly: replacing the constant with $\underline{h}_t = 4\% \times (1 + \kappa z_t)$, where $z_t$ is the window-standardized 30-year mortgage rate and $\kappa = 0$ nests (and in the run exactly reproduces) the production specification, moves the marginal between $+7.1$ and $+9.4$ points across $\kappa \in [-0.5, +0.5]$, every cell inside the level sweep's bracket.

That range is mostly not a measure of cyclicality, and I no longer read it as one. A separate pre-committed run (\texttt{floor\_cyclical\_permutation}) reproduces all seven committed $\kappa$ cells bit-exactly and then decomposes them. Scrambling the floor path's time order within each $\kappa$ cell, preserving the multiset of monthly floor values and their mean exactly, leaves 85.0\% of the range span standing: 72 draws, twelve per non-zero $\kappa$, span $+7.20$ to $+9.12$ points against the committed $+7.09$ to $+9.36$. Switching only the combination rule removes it: under the competing-risks (additive) form the identical $\kappa$ grid spans $+11.22$ to $+11.27$ points, a span of $0.045$ points, so changing the rule alone eliminates 98.0\% of the max-form span, and the additive form moves less across the whole $\kappa$ grid than it moves across the 3--5\% constant-floor level sweep. The mechanism is the one-sided Jensen effect that Appendix~\ref{sec:robustness-seasonalfloor} diagnoses at length for the seasonal floor: under a hard maximum, dispersion in $\underline{h}_t$ lifts the effective hazard whenever the high cells bind and is truncated away where they do not, regardless of \emph{when} it occurs, so $\kappa$ is a dispersion dial before it is a cyclicality dial. Decomposing the low endpoint confirms it. At $\kappa = +0.5$, which departs $-\$16.09$ billion from production, level equivalence accounts for $-\$2.14$ billion (13.3\%), order-free dispersion for $-\$11.42$ billion (71.0\%), and genuine time-order for $-\$2.53$ billion (15.7\%): 84.3\% of that endpoint is order-free. The high endpoint is not an endpoint at all, departing from production by $+\$1.23$ billion, $+0.16$ points, at $\kappa = -0.1$.

A cyclicality signal does survive, and it is real but small. The true marginal falls outside the entire twelve-draw permutation range at every one of the six non-zero $\kappa$ (0 of 72 draws bracket it), and the order components are monotone in $|\kappa|$ and sign-consistent with the co-movement story the concession above states: $+\$5.29$, $+\$4.77$, and $+\$2.69$ billion at $\kappa = -0.5$, $-0.25$, $-0.1$, and $-\$1.42$, $-\$1.75$, and $-\$2.53$ billion at $+0.1$, $+0.25$, $+0.5$. A floor rising with rates does lower the marginal and one falling with rates does raise it. But scrambling the time order removes most of the grid's asymmetry while leaving its level drop intact: the permuted means differ across the sign of $\kappa$ by only $\$0.22$, $\$1.80$, and $\$6.01$ billion at $|\kappa| = 0.1$, $0.25$, and $0.5$, against a true spread of $\$13.82$ billion at $|\kappa| = 0.5$ ($+\$68.08$ billion at $\kappa = -0.5$ against $+\$54.26$ billion at $+0.5$). A two-sided response that survives losing the time order is the signature of dispersion rather than of co-movement: the symmetric part of the grid's movement is large and is dispersion, the asymmetric part is small and is cyclicality. What the grid bounds is therefore the marginal's sensitivity to within-run floor variation of this amplitude, of which at most \$5.29 billion in any cell is attributable to the direction of co-movement. That the concession's residual is bounded by the level sweep already reported holds a fortiori. One correction to the specification as earlier described: the window mean is pinned at the production 4\% only where the floor's non-negativity clip does not bind, and it binds in the two $|\kappa| = 0.5$ stress cells, whose realized window means are 4.0005\% and 4.0828\% rather than exactly 4\% (the artifact flags the clip per cell). Of those two, $\kappa = +0.5$ is the cell that sets the bracket's low end at $+7.09$ points, while $\kappa = -0.5$ sits mid-grid at $+8.90$; the invariant therefore fails in the cell that matters most for the reported range, and the convexity-matched flat floor equivalent to the $\kappa = +0.5$ path is 4.0977\%, not 4\%.

Aggregation from the 75,000-loan sample to the SOMA book applies balance weights: simulated prepaid UPB is scaled by the ratio of the Fed's MBS book to the sample balance, with no reweighting toward Ginnie Mae collateral (20.4\% of book face) or vintages outside 2017--2021; the Freddie-estimated hazards are extrapolated to those segments as-is. This extrapolation is not behaviorally one-to-one: Ginnie Mae pools are FHA/VA collateral with a distinct borrower composition and a servicer buyout channel absent from the conventional book, so Freddie voluntary speeds cannot be mapped onto the Ginnie share as a like-for-like behavioral prediction. The framework therefore does not model Ginnie behavior; it bounds the composition error that treating the Ginnie share as conventional induces. Two bounds constrain this extrapolation. Along the coupon dimension, Section~\ref{sec:robustness-hybrid}'s full-book reweighting (sample WAC 3.9\% mapped to the book's 2.49\%) moves Path B from 107.0\% to 109.1\%, a composition effect of about two percentage points --- a mixed-basis figure, since 3.9\% is the unweighted mean of borrower note rates while 2.49\% is the face-weighted mean of security pass-through coupons; converted to one convention the effect is $+0.2$ points (Section~\ref{sec:robustness-hybrid}, run \texttt{coupon\_convention\_reweight}). Along the agency dimension, an external bound now applies: Ginnie Mae's fixed-rate speeds ran 1.2--2.8 percentage points above Freddie Mac's across the window (5.2\% versus 4.0\% in February 2023; 10.3\% versus 7.5\% in May 2025; Ginnie Mae Global Markets Analysis Report, Recursion data), with the involuntary buyout channel supplying a large minority of that difference rather than most of it. Over the forty-two window months the mean Ginnie--Freddie CPR gap is $+2.14$ points, of which the voluntary component---the conditional repayment rate (CRR), prepayment net of the default-driven channel---supplies $+1.35$ points, or 63\%, against the buyout channel's (CDR) $+0.85$ points, or 40\%; the two shares sum past one hundred because the CPR $\approx$ CRR $+$ CDR extraction carries a small residual. The snapshot that most favours the buyout reading, CDR 2.1\% against 0.4\% in May 2025, is the buyout-dominated month rather than the representative one. Treating the 20.4\% Ginnie face share as conventional therefore mis-states book-wide CPR by at most $0.204 \times (1.2\text{--}2.8) \approx 0.24$--$0.57$ percentage points. At the null-to-central sensitivity of roughly \$83 billion per CPR point---the average book face over the window, times one percentage point of annual prepayment, times the window's three and a half years--- the composition error is bounded at \$20--\$47 billion (2.6--6.2\% of benchmark), signed toward overstating trapped liquidity. The bound is one-sided because the observed differential is: the published speeds run Ginnie above Freddie across the window, so the symmetric lower branch, a Ginnie-slower error signed toward understating trapped liquidity, does not arise in this data.

The vintage dimension carries no bound of either kind: the 2022 vintage alone is 23.1\% of book face (and pre-2017 vintages another 10.6\%; Appendix~\ref{sec:robustness-wal}), lies outside the 2017--2021 Freddie sample. Its lock-in state is not qualitatively different from the sampled vintages': 72.2\% of the 2022 vintage's book face carries a pass-through coupon below 3.0\% (SOMA CUSIP joint coupon $\times$ vintage cells, run \texttt{composition\_shift}), so the vintage is predominantly late-2021 and H1-2022 production sitting four to five points below the window's market rates, with only 13.2\% of it at 4.0\% or above. What the extrapolation lacks is not a different lock-in state but a different \emph{seasoning} state, and neither the coupon reweighting nor the Ginnie speed comparison constrains it. An observed-speed bound in the same template does: the per-quarter cohort cells retained from the Fannie Mae replication ingest carry observed QT-window prepayment for the out-of-sample vintage groups. Over the window the 2022 vintage prepaid at 4.48\% against the sampled universe's 4.30\% (an unrounded differential of $+0.17$ points) and the pre-2017 vintages at 5.26\% ($+0.96$ points; the 2015--2016 cohorts alone are 99.99\% of that tail's exposure and reproduce its speed unchanged), a share-weighted book-CPR error of $+0.14$ points. At the same \$83-billion-per-CPR-point sensitivity, a bound of \$11.7 billion, 1.5\% of benchmark, below the static Ginnie bound's own lower edge and more than a quarter of the headline lock-in marginal (27.6\% of the off-window \$42.6 billion; a sixth, 16.7\%, of the in-sample \$70.3 billion), signed the same conservative direction: out-of-sample vintages prepaid \emph{faster}, so trapped liquidity is if anything overstated. The same template applied as a marginal overlay --- scoring the 33.7\% out-of-window vintage share by its observed Fannie speeds in \emph{both} legs, so that share contributes zero marginal --- confines the identified marginal to the in-window-vintage share: $+3.7$ points, pure share scaling ($0.66\times$ against a nominal retained share of $0.663$, with the placebo within $0.0002$), the vintage counterpart of the Ginnie overlay's $+4.4$ (run \texttt{vintage\_overlay}); the two overlays cannot be composed without the agency-by-vintage cross-tabulation, since 2.4\% of book face is the 2022--Ginnie intersection. Three caveats travel with the bound: Fannie observed speeds proxy the book's, the pre-2017 leg is a seasoned acquisition tail rather than a random draw of old originations, and the 2022 leg is a full-year speed applied to an H1-tilted face.

Burnout is deliberately cohort-level rather than per-loan: a fractional burnout measure is not well-defined for a one-time, absorbing transition (Section~\ref{sec:abm-results}). Instead, each stratum tracks its cumulative prepaid UPB as a share of its original UPB (the $B \in [0,1]$ state defined under \eqref{eq:pathB}) after all prepayment transitions in a given month, broadcast to all of that stratum's remaining active loans before the next month's hazard is evaluated, and entering \eqref{eq:pathB} with the calibrated prior coefficient $\beta_B = -0.5$ (Appendix~\ref{app:params}); Path A's separately estimated burnout coefficient, statistically indistinguishable from zero (Section~\ref{sec:patha}), plays no role here. Because the simulated burnout state is endogenous, I test whether the channel is load-bearing with a $\beta_B = 0$ rerun on the production seed (run \texttt{burnout\_ablation}): trapped liquidity moves by \$6.58 billion ($-0.86$pp of benchmark, to 106.17\%), with the peak lag unchanged at $-3$. Two qualifications travel with that figure. It is measured at the in-window 4.0\% floor; at the 4.991\% off-window floor the paper headlines the same ablation is worth only \$3.48 billion ($-0.45$pp), because the floor's bind share rises from 36.3\% to 68.8\% of evaluated loan-months and the hard maximum makes burnout inert wherever it binds. And because the burnout state is endogenous to realized prepayment, the ablation moves the $\beta_1 = 0$ null as well as the central leg, so the identified marginal is not burnout-invariant: switching the coefficient off raises it from $+9.20$ to $+9.94$ points at the production floor and from $+5.57$ to $+6.51$ at the off-window floor. Both moves sit inside the $\pm 1$-point materiality convention, the second only barely, so I report them rather than resting anything on them. The FICO and LTV covariates receive the same treatment: a $\beta_F = \beta_L = 0$ rerun moves the recovery to 108.3\% ($+1.3$pp), and a rerun under coefficients estimated Path-A-style from the panel moves it to 101.3\% ($-5.7$pp), both at the in-window 4.0\% floor; at the 4.991\% off-window floor the same endpoints span more, not less --- 102.5\% against 94.3\% standalone, an 8.2-point swing --- while the marginal they bracket moves by only $+0.1$ points against $+1.9$ in-window, because both legs shift together where the floor censors more (run \texttt{covariate\_ablation\_offwindow}); Appendix~\ref{app:params} reports the estimation, whose imprecise point estimates are consistent with the priors' turnover-regime signs without rejecting alternatives. Production retains the ex-ante priors (the estimates are too imprecise to license replacing a pre-specified calibration mid-stream), and I flag the covariate block as the specification's largest single sensitivity: the two endpoints correspond to 99.2\% and 92.2\% on the shared accounting basis (the common netting makes the conversion exact), a swing comparable in size to the lock-in marginal itself.

The static Ginnie bound is now extended to a time-varying overlay built from the published series the bound's two snapshots came from. Extracting the full monthly agency CPR series (June 2017 through November 2025) from the December 2025 Global Markets Analysis Report (whose aggregation was revised in that issue to weighted-average unpaid principal balance, one consistent dollar-weighted methodology across the whole window) and scoring the 20.4\% Ginnie face share by the observed Ginnie speed in \emph{both} Path B legs moves the shared-basis central from 97.9\% to 89.3\% and the null from 88.7\% to 81.9\% in the raw construction, both legs at the in-sample calibration point, and confines the identified marginal to the conventional share: $+7.3$ points, positive under every variant tested (total-CPR, voluntary-only, and placebo), scaling with the conventional face share almost exactly ($0.797$ against $0.796$). Composed with the off-window headline calibration (run \texttt{ginnie\_overlay\_offwindow}), the overlay pair is 84.0\%/79.5\% central/null on the shared basis with a $+4.4$-point marginal, and the overlay marginal is pure conventional-share scaling ($0.797\times$) --- a change of estimand, a re-scoring onto the conventional sub-book, rather than a correction to the headline, identical across the primary, voluntary-only, and placebo series; the headline marginal therefore carries the conventional-extrapolation convention wherever quoted, with $+4.4$ points its overlay-scored member. The raw level shift overstates the Ginnie-specific effect, though: a placebo that scores the same 20.4\% share by the observed GSE-mean series moves the level by $+\$27.8$ billion on its own (the model's level calibration plus the full-universe-versus-SOMA composition difference, common to any observed-series overlay), so the Ginnie-specific differential contribution is $+\$38.5$ billion, 5.0\% of benchmark (the placebo and Ginnie-specific components, $\$27.8 + \$38.5 = \$66.3$ billion, reconcile with the raw 8.6-point shared-basis shift on the \$764.7 billion benchmark, $\approx\$65.8$ billion, to within rounding of the 0.1-point recovery shares): inside the static \$20--\$47 billion bound above, and consistent with the window-mean Ginnie--Freddie differential of 2.1 points the series itself shows against the bound's 1.2--2.8-point snapshots. The full-universe series if anything overstates the correction for SOMA's seasoned, low-coupon Ginnie holdings, the same signed direction as the static bound; and structural Ginnie inclusion remains out of scope for the modeling reason above: the overlay treats the buyout-inclusive Ginnie speed as data, which is what runoff accounting requires of a passthrough holder. This overlay and the age-standardized floor read of Section~\ref{sec:robustness-floor} cut the same way, and the joint cell has now been run under a pre-committed landing rule (run \texttt{b5\_joint\_cell}): composed, they return $+2.9$ points ($+\$22.4$ billion), within $0.1$ points of the proportional projection, and that measured point is the assembly's lower member.

Two Danish counterfactual mechanisms exist in this framework's pipeline, and only one is live in production. The first is a \emph{rate-gap swap}: holding the baseline hazard, rate-gap coefficient, covariates, and random draw logic fixed, the Danish leg expresses the market-value buyback in rate-equivalent units, so that the buyback discount offsets the locked-in spread; $g_{\mathrm{DK}} = g_{\mathrm{US}} + \Delta y_{\mathrm{buyback}} \approx 0$ whenever the loan is out of the money, and $g_{\mathrm{DK}} = g_{\mathrm{US}}$ when rates have fallen and the buyback price is capped at par (at the representative 3.0\% coupon / 6.8\% market state the U.S. gap is $-380$ bp while the Danish effective gap is zero, so the U.S. hazard is suppressed by the full locked-in spread while the Danish hazard resets toward its baseline seasoning curve). This swap defines the \emph{mechanism-substitution} variant (the specification behind the superseded 47.1\% Danish CPR of Section~\ref{sec:abm-danish}), and it is \emph{not} the production mechanism. The production Danish leg is the Berger-recalibrated specification adopted in Section~\ref{sec:abm-danish}: it replaces the rate-gap swap with the imported flat moving elasticity, so the production Danish hazard does not respond to the buyback discount at all. The swap machinery remains in the codebase and its PV-versus-par classification is still computed each month, but no production quantity consumes its output, which is why the discount-rate diagnostic of Appendix~\ref{sec:method-abm} moves the Danish leg by exactly \$0.00 billion. Table~\ref{tab:discount} reports that inertness: shifting the buyback discount rate down by 25 to 150 bp (through the primary--secondary spread range) flips thousands of monthly present-value-versus-par classifications while moving the Danish CPR path, roll-off path, and trapped-liquidity dollar by exactly zero on every scoring basis. Two limits apply to that result. I state both because together they retract the immunity I previously claimed from it. First, the diagnostic executed on the Danish-level bracketing leg, not on the production U.S.-intercept leg: its own artifact records a mean Danish CPR of 3.39\%, the bracketing anchor's speed. Under the production anchor the same independence does hold, but by construction rather than by measurement, because that branch replaces the rate gap with zeros before the hazard is evaluated. Second, and more important, the result is about the hazard, not about the accounting. The Danish leg credits every retired balance at par: its dynamic-balance loop books roll-off as balance times drain, with no price term anywhere in it. Netting the roughly \$199 billion scheduled component from the leg's \$669.3 billion of Danish roll-off leaves on the order of \$470 billion of buyback-triggered retirement credited at par, against a proxy discount of 32 to 34\% at the representative 3.0\% coupon / 6.8\% market state (36 to 38\% at the book's 2.49\% weighted-average coupon). A market-value credit would reduce the Danish leg's cash receipts by an amount larger than the $+\$61.2$ billion gap itself, and in the direction that shrinks it. The pipeline does not distinguish a cash haircut from a balance adjustment, and the incidence depends on which the buyback is, so a pre-committed bracket prices both readings from the committed pair rather than choosing one (run \texttt{buyback\_credit\_bracket}; gates and the interpretive rule fixed in the script header before execution): under the balance-adjustment reading face retirement is unchanged and the gap stands at $+\$61.2$ billion, while under the cash-haircut reading receipts fall by the proxy discount on the leg's \$470.7 billion of early face and the gap reverses, to $-\$89.4$ to $-\$117.7$ billion across the committed 32--38\% discount grid. The two readings are not symmetric answers to one question. The benchmark is a face-value object---Section~\ref{sec:method-benchmark} constructs it from the SOMA current face value series against the par-denominated caps---so the balance-adjustment reading is the benchmark-consistent one, and on the benchmark's own denomination the gap stands. The cash-haircut reading answers a different question, reserves actually drained, and its reversal has a name: the $1-P$ wedge on bought-back face is exactly the par windfall a moving household pays the bondholder under the U.S. rule and keeps under the Danish one, so the cash-flow cost is the disappearance of a household-to-bondholder transfer, not an added resource cost. The gap's sign is therefore incidence-conditional, the channel is priced rather than un-priced, and the discount-rate diagnostic still should not be read as immunity to the pricing critique.

Simulated on the 40,234 window-start survivors of the 75,000-loan stratified sample drawn from the same Freddie Mac 2017--2021 origination universe (Appendix~\ref{app:params} reports the full attrition accounting), this specification recovers \$818.5 billion, or 107.0\%, of the empirical benchmark, with a raw CPR correlation against the empirical path of $r = +0.190$ at zero lag and a peak cross-correlation of $r = +0.404$ at lag $-3$. Both are descriptors of the published series, not credentials: the monthly benchmark cannot adjudicate timing, since four clip-induced zero months dominate its differenced signal (Appendix~\ref{sec:robustness-seasonalfloor}). Throughout, the cross-correlation at lag $k$ pairs the simulated series at month $t$ with the empirical series at month $t+k$ (equivalently, the empirical series at month $t$ with the simulated series at month $t-k$; both series demeaned over the overlap, boundary months dropped), so at the peak $k=-3$ the simulated value at month $t$ is matched with the empirical value from three months earlier. That definition was inverted in print---the two series interchanged, which would have implied the opposite direction---but the direction claims themselves were always computed under the implemented convention (Appendix~\ref{app:ledger}). A direction check with synthetic data confirms that under this implementation a peak at $k=-3$ means the simulated path \emph{trails} the empirical path by three months. At this peak, the empirical SOMA-derived CPR path moves approximately three months ahead of the simulated path (a peak selected across the seven examined lags, so its nominal significance is overstated; no inference rests on this single coefficient). The level and band reported here follow the rate-gap coefficient's unit convention documented in Appendix~\ref{sec:robustness-pipeline} (the superseded figure is catalogued, with provenance, in Appendix~\ref{app:ledger}).

Importing Path B's elasticity unmodified, rather than fitting it to the benchmark, has two consequences. First, unlike Path A, whose rate-gap and burnout coefficients are estimated directly from the Freddie Mac sample, Path B's rate-gap elasticity is imported unmodified from \citet{liebersohn2024} and never touches the \$764.7 billion benchmark during calibration. The two are mutually reinforcing rather than redundant: Path A's fitted rate-gap coefficient is an in-sample estimate of the same suppression channel on the same loans ($+0.29$ per 100 bp; Section~\ref{sec:patha}), and the Fannie Mae replication reproduces its sign on a second agency's loan book at roughly $3.6\times$ the per-basis-point magnitude, so the two in-sample fits are directionally consistent with the imported band---but no more than that: a pre-committed test of $H_0\!: \beta_g \le 0$ that respects the diagnosed bootstrap bias (run \texttt{patha\_sign\_test}: BCa with a 296-refit jackknife acceleration, $p = 0.093$ stratum scheme and $0.412$ temporal; a 200-draw month-permutation of the rate-gap series, $p = 0.241$) does not reject at any conventional level, so the in-sample estimates are not statistically distinguishable from zero under bias-respecting constructions and the elasticity's evidential basis is the external literature alone. Within that basis, \citet{graybill2026} corroborate the adopted band's location independently: their transaction-record estimate puts the sale-hazard response at roughly 7--11\% per percentage point of rate gap, in a housing-market-equilibrium setting whose records touch no credit file. Like the Fonseca--Liu anchor below, it corroborates from at or above the adopted band rather than from its centre: only the 7\% low end falls inside the 5.5\%--7.7\% range and the midpoint sits above the high edge, so the adopted central calibration is the conservative choice against this anchor as well. A prepayment-hazard coefficient is also not unit-comparable to \citepos{liebersohn2024} quarterly-mobility decline, and Path A measures prepayment rather than pure mobility, so what remains is directional consistency, not triangulation of sign or direction, not magnitude. A second literature anchor corroborates from outside the band: propagating \citepos{fonseca2024} roughly 9\% moving-rate reduction per 100 bp through the same transform ($\beta_1 \approx 0.096$) and the same band machinery yields a marginal of $+\$87.9$ billion ($+11.5$ points) at the production floor, above the Liebersohn--Rothstein high edge ($+\$79.5$ billion, $+10.4$ points) yet inside the pre-committed floor-and-band box, so the adopted central calibration is the conservative choice among the available literature estimates; because that estimand is a moving-rate reduction rather than a quarterly-mobility decline, I read it as directional corroboration, not a competing point estimate.

Path B is accordingly not a second attempt to fit my benchmark but a structural verification exercise: it asks whether an elasticity estimated in a different empirical setting, dropped into a PSA-baseline hazard with competing risks and no tuning against my target, reproduces the Federal Reserve's actual QT runoff shortfall. Because the aggregate level is floor-dominated, landing within 8 percentage points of the benchmark is not, by itself, evidence for the imported elasticity, since the residual over the $\beta_1 = 0$ null that the elasticity must explain is itself a function of the floor calibration (the marginal spans $+2.1$ to $+13.2$ points across the calibration box of Section~\ref{sec:robustness-floor}). The defensible verification claim is the bounded interval the calibration box implies; the marginal is also positive at every floor, band, and floor-form point tested, but under this window's rate configuration that positivity verifies the specification rather than identifying the mechanism (Section~\ref{sec:identification}), and I do not claim the level validates the elasticity's magnitude. But the timing is not part of that evidence: as Section~\ref{sec:hazard-interp} reports, a rerun with the elasticity set to zero reproduces the same three-month lead, so the lead cannot discriminate the elasticity from the mechanical model.

Second, propagating the full Liebersohn--Rothstein 5.5\%--7.7\% sensitivity band through the simulation yields a recovery range of \$809.9 billion to \$827.7 billion (105.9\%--108.2\% of benchmark under the standalone scorer; 96.8\%--99.1\% on the shared accounting basis of Section~\ref{sec:robustness-hybrid}), with the peak cross-correlation lag stable at $-3$ across the entire band (peak $r$ between $+0.371$ and $+0.423$; given the moving-block intervals above and the seven-lag search, these differences carry no evidential weight). Stronger assumed mobility suppression therefore maps monotonically to more trapped liquidity, as expected, and the central 107.0\% estimate should be read as a point within this interval; Figure~\ref{fig:band} plots the band. That band is a calibration interval, not sampling error, and the sampling dimension is now quantified separately: a stratified loan-level bootstrap specified ex ante in its script header (200 replicates resampling the 75,000-loan draw with replacement within strata, central and null legs paired per replicate, engine seed held at production) puts the central recovery's 95\% interval at $[107.0, 107.1]\%$ under the standalone scorer (replicate standard deviation 0.03 points), the null's at $[97.8, 97.9]\%$, and the paired lock-in marginal's at $[+9.17, +9.23]$ points (standard deviation 0.016); re-run under the cluster scheme at this same floor, the paired marginal reads $[+8.27, +10.19]$ points (sd 0.52; run \texttt{bootstrap\_pathb\_cluster}), an understatement of 30.8$\times$ in-window against 37$\times$ off-window. Those intervals are narrow because the scheme holds the design fixed: it resamples loans within each stratum and preserves every stratum's count, so between-stratum composition variance is never drawn. Drawing the 130 strata themselves with replacement, at the off-window floor the paper headlines (run \texttt{bootstrap\_pathb\_cluster}, 200 replicates), widens the marginal's interval to $[+4.63, +6.92]$ points with a standard deviation of $0.60$ points---a width of $2.29$ points against the within-stratum scheme's $0.062$, so the committed scheme understated loan-sample uncertainty by roughly a factor of $37$. Two properties of that interval are not corrected for and should be stated: drawing strata with replacement makes the loan count itself random---it runs from 54,134 to 101,404 across replicates against a fixed-design 75,000---so every composition weight is a ratio of random totals rather than of fixed ones, and the percentile interval carries no correction for the resulting ratio bias. A fixed-$n$ variant that rescales weights to the design total would settle how much of the width is that bias; it is not run here. The committed point marginal falls inside that interval. Two limits on the wider figure: the strata are very unequal, so the effective cluster count is $25.8$ rather than $130$ (the largest single stratum is $9.77\%$ of sample balance) and a percentile interval on $26$ effective clusters is itself noisy; and the resampled loan count varies from roughly $54{,}000$ to $101{,}000$, which is harmless here only because the engine renormalises the pool to Fed holdings each month, so the level is set exogenously and only composition varies. Even at the wider figure the sampling layer is less than half the width of the floor-read interval of Section~\ref{sec:robustness-floor}, so the operative uncertainty statements for Path B's level and marginal remain the calibration box and floor-form results there rather than sampling intervals; the shared-basis equivalents follow by the common 9.1-point netting, which cancels from the marginal entirely.

\begin{figure}[H]
\centering
\includegraphics[width=0.82\textwidth]{fig6_elasticity_band}
\caption{Path B trapped liquidity across the Liebersohn--Rothstein elasticity band. The quarterly mobility-decline parameter $\delta$ is swept across its published 5.5\%--7.7\% range (central 6.5\% marked), mapping monotonically to recovered trapped liquidity of \$809.9--\$827.7 billion on the standalone-scorer basis (shared-basis equivalents in Table~\ref{tab:bases}) against the \$764.7 billion benchmark line: stronger assumed mobility suppression yields more trapped liquidity, with the peak cross-correlation lag stable at $-3$ across the entire band. Table~\ref{tab:floorband} extends this band across the baseline-turnover-floor grid.}
\label{fig:band}
\end{figure}

\begin{table}[!t]
\centering
\caption{The calibration box: lock-in marginal over the $\beta_1 = 0$ null, in percentage points of the \$764.7 billion benchmark, across the baseline turnover floor and the Liebersohn--Rothstein band $\delta$. The 6.0\% row's level is reached only by the refi-contaminated pooled and 2019 anchors that Section~\ref{sec:robustness-floor} disqualifies; it is retained as a grid point, not a defensible anchor, and the box's lower edge should be read from the defensible off-window rows of Table~\ref{tab:oosfloor}. The bind column is the share of evaluated U.S. loan-months at which the floor lifts the hazard (central elasticity). Every marginal is basis-invariant (the underlying levels are standalone-scorer figures; both bases in Table~\ref{tab:bases}). Pre-committed runs; artifacts catalogued in Appendix~\ref{app:ledger}.}
\label{tab:floorband}
\begin{tabular}{lcccc}
\toprule
Floor (annual CPR) & $\delta = 5.5\%$ & $\delta = 6.5\%$ & $\delta = 7.7\%$ & Floor binds \\
\midrule
2.0\% & $+9.7$ & $+11.3$ & $+13.2$ & 3.9\% \\
3.0\% & $+9.3$ & $+10.8$ & $+12.6$ & 12.1\% \\
3.5\% & $+8.9$ & $+10.3$ & $+11.8$ & 21.6\% \\
4.0\% (production) & $+8.1$ & $+9.2$ & $+10.4$ & 36.3\% \\
4.5\% & $+6.7$ & $+7.5$ & $+8.4$ & 53.5\% \\
5.0\% & $+5.0$ & $+5.5$ & $+6.0$ & 69.1\% \\
6.0\% & $+2.1$ & $+2.3$ & $+2.4$ & 88.7\% \\
\bottomrule
\end{tabular}
\end{table}

The peak alignment does \emph{not} show the simulated path leading SOMA cash flow. The empirical path moves first: a settlement-delay mechanism, in which loan-level prepayment decisions show up in SOMA cash two to three months later, predicts the \emph{opposite} direction of offset. I therefore attach no settlement interpretation to the timing alignment. The direct input-timing diagnosis behind that reading is in Appendix~\ref{app:patha-detail}.

\subsection{External Replication on Fannie Mae Loan-Level Data}\label{sec:pathb-fannie}

Beyond the Freddie Mac sample on which Path B is built, the same construction admits a direct external replication on a second agency's loan book. I re-ran the full Path B construction on Fannie Mae's single-family loan performance data (24 acquisition quarters, 2017Q1--2022Q4; 17.6 million loans and 826 million loan-months, staged into the identical origination/performance layout so that the panel builder, the estimator, and the stratified 75{,}000-loan sampling design run unchanged on the same 2017--2021 origination universe) under a spec fixed and committed before the run: zero-balance code 01 maps to voluntary prepayment as for Freddie; codes 06 and 16 (non-credit removals with no Freddie counterpart) are censored, with a worst-case bracket of 0.38\% of removals if every such removal were instead counted as a prepayment; and the pre-committed acceptance gate required the Fannie lock-in marginal to be positive and inside the Freddie floor/band envelope of Section~\ref{sec:robustness-floor}. The result replicates: central \$821.0 billion (107.4\% of benchmark), $\beta_1 = 0$ null \$754.6 billion (98.7\%), lock-in marginal $+\$66.4$ billion ($+8.68$ points of benchmark) against Freddie's $+\$70.3$ billion ($+9.20$ points): inside the envelope, within half a point of the Freddie estimate. The envelope was the calibration box of Section~\ref{sec:robustness-floor}, which spans 2.11 to 13.17 points---an 11.06-point window around a 9.20-point Freddie estimate, admitting a replication 77\% below it or 43\% above. A gate that wide cannot fail against anything but a sign error, and I do not treat its passing as evidence. The agreement is a different matter: 17.6 million staged Fannie loans, run through the same pipeline end to end, land 0.52 points from the Freddie point, which is a good deal more than the gate demanded and is the part of this exercise that is reassuring. What that agreement certifies needs the same care: every behavioral parameter is fixed a priori and identical across the two runs, and the two stratified samples are compositionally indistinguishable (population-stability indices at or below 0.018, Appendix~\ref{app:composition}), so the 0.52-point agreement certifies pipeline determinism and cross-agency compositional stability; the external-validity content lives in the sign reproduction and the staging inputs, not in the closeness itself. The Path A fixed-effects refit on the Fannie cohort panel reproduces the rate-gap coefficient's sign (its per-basis-point magnitude is about $3.6\times$ the Freddie fit's, and its holdout diagnostics carry the same no-forward-precision character as Section~\ref{sec:patha}'s, so it retains aggregate-corroboration status only). Two scope notes are part of the claim, both fixed ex ante: whether Freddie's code 01 absorbs exits that Fannie books under 06/16 is not settled by public documentation, so the replication claim is sign and envelope-conditional marginal consistency rather than a hazard-level match; and Fannie's terms prohibit redistribution, so the replication package ships the Fannie ingestion and replication code only, with no Fannie data and, conservatively, no derived cells; the committed artifact carries headline statistics alone.

\subsection{Path A: Stratum-Month Discrete-Time Hazard Model}\label{sec:patha}

Path A is the aggregate corroboration leg: a stratum-month discrete-time hazard fit to Freddie Mac loan-level performance data. Its aggregate recovery is \$928.9 billion, 121.5\% standalone --- reported as corroboration only and excluded from every headline figure --- and its estimation detail, diagnostics, and cross-correlation exhibits are collected in Appendix~\ref{app:patha-detail}.

The same caution applies to the paths themselves. Path B's $+0.190$ lag-0 levels correlation is carried by the window's shared trend: in first differences it is $-0.23$ and after linear detrending $-0.20$. The same constructions for the other estimators give $-0.20$ and $-0.30$ for the ABM (from $-0.32$ in levels) and $-0.21$ and $-0.29$ for Path A (from $-0.44$ in levels): detrending collapses the three estimators' apparent differences into a cluster between $-0.20$ and $-0.30$. Against the i.i.d.\ $1.96/\sqrt{n}$ reference ($\pm 0.30$ at $n = 42$ for the detrended series; $\pm 0.31$ at $n = 41$ for first differences), Path A ($-0.29$) and Path B ($-0.20$) are inside the band under both empirical-series conventions, while the ABM's detrended value straddles the band's edge across them ($-0.301$ common-series, $-0.303$ against its own back-out). Indistinguishability verdicts therefore rest on the wider moving-block bands of Table~\ref{tab:estimators}, under which none of the three is distinguishable from zero. Once the shared trend is removed, no estimator's contemporaneous co-movement is positive; Path B's distinction is the accuracy of its aggregate level, not its timing. (In levels, which are trend-contaminated over the window and reported only for completeness, the simulated path co-moves most strongly with the rate from four months earlier while the empirical path is contemporaneous.) No single channel carries the offset: the $\beta_B = 0$ and $\beta_1 = 0$ ablations both leave the peak at $-3$, and re-running the simulation with the rate input shifted by one to three months in either direction never moves the peak to lag 0 (it reaches $-1$ under a two-to-three-month input lead) while moving recovered trapped liquidity by less than \$1.5 billion across the entire scan. The offset is therefore a property of the fitted path's smooth aggregate dynamics set against a seasonal empirical series (mechanism not isolated), and Path A's seasonal re-estimation moving its own peak from $-3$ to $-2$ (Section~\ref{sec:patha}) is consistent with seasonality carrying part of it. Figure~\ref{fig:ccf} plots the full cross-correlation functions for all three estimators. Extending this input-timing diagnosis to the cohort level confirms the residual is structural rather than compositional: across the six 30-year coupon cohorts spanning the SOMA low-coupon mass (2.0--4.5\%), every simulated cohort CPR tracks the contemporaneous rate change with near-mechanical strength (first-difference correlation between $-0.87$ and $-0.94$; five of six cohorts cross-correlation phase-locked at lag~0, the 4.5\% coupon peaking at lag $-6$ with a weaker $r=-0.29$, and reproducing the aggregate $-0.318$), whereas no realized cohort speed, measured independently in the Freddie 2017--2021 performance panel and in the Fannie replication cohorts, shows a rate response distinguishable from zero at any lag within $\pm 6$ months (Freddie lag-0 first-difference correlations of $-0.05$ to $-0.16$, Fannie $-0.03$ to $-0.14$, all inside the $\pm 2/\sqrt{n}$ band, and a maximum magnitude of $0.34$ over all cohorts and all lags, with no coherent sign or lag). The aggregate empirical non-response is not a composition artifact, and the simulated lead is not a datable phase shift at the cohort level. One honest scope limit travels with this exhibit: realized cohort-level monthly speeds exist only in the loan-level agency panels that anchor Path B, not in the SOMA holdings from which the aggregate empirical CPR is backed out, so the cohort split is run in the Path B universe rather than on the SOMA book itself, which is why the estimator-level timing falsification above stays at the SOMA aggregate. A cumulative counterpart completes this exhibit (run \texttt{episode\_confrontation}). Because monthly differences are seasonality-dominated, the cross-section is also read in cumulated levels: over the window, the age-matched Freddie panel's realized prepayment runs from 9.44\% annualized CPR (28.1\% of balance cumulatively) in the shallowest surviving gap bucket ($[-1,0)$ points, mean coupon 6.02\%) down to 5.24\% (16.4\%) at gaps of $-3$ points and deeper (mean coupon 3.37\%): a gradient of $+4.20$ CPR points, with a stratum-cluster 95\% interval of $[+3.59, +4.66]$ and a one-sided permutation $p$ of 0.005---the one realized-data exhibit in this paper that moves in the lock-in direction. It is not corroboration of the imported elasticity's magnitude: the production hazard's implied gradient across the same buckets is $+0.94$ points at the central $\beta_1$, so the realized gradient is roughly $4.5\times$ the model's, and the excess carries everything else that travels with coupon across cohorts---credit composition, within-window refinancing on the high-coupon shallow buckets, residual seasoning past the age cut. Detection power for the implied gradient is 0.68: the design could have resolved $\beta_1$'s implication, and what it found instead was larger and confounded. The composition limb of that confound has now been tested rather than asserted: standardizing the deep bucket to the shallow bucket's vintage $\times$ FICO $\times$ LTV composition leaves the gradient at $+3.44$ points (95\% CI $[+1.17, +5.14]$, imputed weight share $0.3\%$), and an exact decomposition attributes $+0.77$ of the raw $+4.20$ to composition and $+3.43$ to differences within common cells (run \texttt{episode\_confrontation\_within}) --- against the production hazard's implied $+0.94$, the excess is not composition in the dimensions the panel can hold fixed. I therefore read the cross-section as a signed, significant level gradient whose composition component is measured and small, consistent with lock-in operating but still unusable as a measurement of its size: what remains unresolved is within-window refinancing on the high-coupon shallow buckets, which no field in this design separates from moving, and a cross-sectional level gradient is not a dollar marginal. The monthly-timing nulls above stand unchanged.

\subsection{Identification: the Marginal over the Null}\label{sec:identification}

The design does not identify the aggregate recovery level: the level is floor-dominated, and the $\beta_1 = 0$ null (the elasticity switched off, baseline turnover held at its observed level) already recovers 85.7\% of the benchmark on the shared accounting basis at the off-window floor this paper headlines, and 88.7\% at the in-sample calibration point (Section~\ref{sec:hazard-interp}; standalone figures in Table~\ref{tab:bases}). What it identifies is the \emph{marginal} of the central Path B run over that null, writing $R$ for recovered trapped liquidity,
\begin{equation}\label{eq:marginal}
\text{lock-in marginal} \;=\; R_{\text{central}} - R_{\beta_1 = 0},
\end{equation}
which is basis-invariant because the shared-accounting netting is common to both legs (Section~\ref{sec:robustness-hybrid}). Its value moves with the baseline turnover floor, and because that floor is the design's dominant level-setter and is calibrated in-window on the very cohorts the marginal decomposes, the headline value is the one obtained with the floor measured off that population: $+\$42.6$ billion, $+5.6$ points of the benchmark, at the defensible off-window floor of Section~\ref{sec:robustness-floor}, and $+4.3$ to $+6.8$ points across that floor's range. At the in-window production calibration the same difference is $+\$70.3$ billion, or $+9.2$ points. The recovery \emph{level} moves with the floor for the same reason the marginal does, and I now report it at the same two calibrations rather than quoting a level from one floor beside a marginal from another: the central leg recovers 91.3\% of the benchmark on the shared basis at the off-window floor (88.2--93.7\% across its range) against 97.9\% at the in-window calibration, and the null 85.7\% against 88.7\%. Only the central and null legs were re-run off window, so those are the two levels quoted at both calibrations; the concave-gap variant was also re-run there, at 99.8\% central against a 94.8\% null on the standalone scorer; every remaining estimator level in this paper---Path~A, the full-book variant, the composed estimator, the covariate-prior and elasticity bands, the Ginnie overlay and the Danish legs---was run only in window, and any level appearing without a floor label should be read at the in-window 4.0\% calibration. I call the in-window configuration the \emph{in-sample calibration point} throughout. Differencing does not, however, immunize the marginal against non-stationary baselines: because the baseline multiplies the very rate-gap channel the null switches off, a structural break in $h_0$ scales the marginal rather than cancelling from it. Re-running both legs under a $\pm 20\%$ multiplicative break from January 2023 onward, imposed identically on both legs and fitting nothing to the benchmark, carries the marginal from $+5.0$ to $+12.1$ points around the committed $+9.2$ (reproduced bit-exactly at the unit factor), while a $\pm 20\%$ break in the turnover floor alone moves it the other way, to $+6.6$ through $+10.4$ points. Every perturbation leaves the marginal inside the calibration box and positive. What survives a baseline break is the box membership, not the point value.

Three qualifications fix its evidential content. First, the identified content is the marginal's \emph{bounded range}, and neither its point magnitude nor its sign: the box spans $+2.1$ to $+13.2$ points, and across the empirically plausible floor range the marginal spans $+5.0$ to $+12.6$ points over the full band, or $+10.8$ to $+5.5$ points at the central elasticity, moving to $+11.25$ points under the competing-risks floor form. The share of the sample the floor censors moves with its level: at the 4.991\% off-window floor the paper headlines, the central leg's hazard is pinned to the turnover floor in 68.8\% of the 1,683,124 evaluated loan-months, against 36.3\% at the in-sample calibration point and 60.0\% to 77.3\% across that floor's defensible range. In most evaluated loan-months at the headline calibration the elasticity therefore no longer sets the hazard, and its contribution to the marginal there is truncated at the floor rather than carried in full; these are unweighted loan-month counts, not balance-weighted shares, so they do not say what fraction of the dollar marginal the uncensored months carry. The $\beta_1 = 0$ leg is pinned less often---35.8\% of loan-months at the headline floor and 14.3\% at the in-sample calibration point---and wherever it is pinned, switching the elasticity off cannot raise the hazard at all. The marginal is thus floor-level sensitive under the production max form, while the form-robust statement is roughly $+9$ to $+11$ points at the production floor (Section~\ref{sec:robustness-floor}); no form-robust point exists at the off-window anchors, where the two forms return $+5.6$ and $+11.2$, so the defensible statement there is the form-conditional range $+3.5$ to $+13.1$ points, whose endpoints are now the endpoints of a measured mixture curve in the floor's strictly-involuntary share: the curve rises to $+9.7$ points by $s = 0.25$ and plateaus near $+11.2$ (run \texttt{floor\_form\_mixture}), so the paper's own floor semantics point at the hull's upper region, and the $s = 0$ headline is the curve's conservative endpoint. Its upper end comes from the additive form at the 4.695\% floor and the band's 7.7\% high edge, whose central leg recovers 61.2\% of the benchmark, and Section~\ref{sec:robustness-floor} states that it should not be read as an equally credentialed member of the range. That hull is conditional on each floor read being taken at its point estimate, and it is therefore not the paper's widest disclosed interval: propagating the floor reads' own cluster-bootstrap sampling error puts the max-form member at $+2.9$ to $+8.7$ points after wild-cluster correction (the percentile read is $+3.0$ to $+8.0$; Section~\ref{sec:robustness-floor}) and leaves it open below $+4.3$ on an age-standardized read (Section~\ref{sec:robustness-floor}). The floor read and the loan/stratum resample are reported side by side rather than maximized over: convolved as independent draws the pair spans $+2.8$ to $+9.0$ points (run \texttt{layer\_convolution}), which is the interval a reader who wants one number for sampling error should use, and which I do not substitute for the floor read because independence between a 2018 cohort-month ratio and a window-period composition resample is an assumption rather than a measurement. The hull's lower edge is thus narrower than the layer the paper calls binding, and the two must be quoted together rather than the hull alone.

The floor's provenance now has an out-of-agency read as well as an out-of-window one, and it does not narrow the range. Applying the committed off-window selection rule unchanged to the Fannie Mae cohort-month panel---the same fourteen columns, the same 2017 start, staged for the external replication of Section~\ref{sec:pathb-fannie}---returns 5.52\% on the headline cell against Freddie's 5.19\%, on 421 cohort-months against 438 (run \texttt{fannie\_floor\_read}; the Freddie panel replays every committed cell of that grid through the identical code path first). That is above the clean band's 5.334\% top, but it is not a second read of the clean leg: the committed rule applied off window covers the pooled 2017--2019 period, whose Freddie counterpart is the 5.19\% just quoted rather than the 2018 leg's 4.991\%, so the agency difference is 0.337 points of CPR (5.522\% against 5.185\%) and the rest of the distance from the clean read is period. The age-standardized read of 5.51\% that the same table already flags lands 0.01 points away, but it perturbs something else again---an age correction to the 2018 leg itself, on 84.0\% imputed weight---so the near-coincidence of the two is not a second measurement of either. What they share is direction: both sit above the clean band's top, and both say the clean read understates the mature book's turnover. That direction is what the clean leg cannot test on its own. The 2018 leg holds no out-of-the-money cohort-month past the twenty-four-month age cut, so its mature-seasoning test is not computable rather than passed (Appendix~\ref{app:floormech}), and the one age contrast it does support is confounded by the PSA ramp. Where a mature contrast is measurable at all---two agencies, two rate regimes, three depth cuts---it runs one way in twelve cells of twelve, the age-$\geq$-24 read the higher one in each: in window, where the rate configuration suppresses refinancing as it did on the 2018 leg, by 0.08 and 0.19 points of CPR on Freddie and Fannie; on the refi-contaminated pooled 2017--2019 period, by 3.74 and 4.61 points, which bounds the correction from above rather than measuring it (run \texttt{fannie\_floor\_read}). A floor read taken entirely on cohort-months aged twelve to twenty-four therefore sits below the mature book's baseline turnover, and the marginal it implies sits above the one a mature book would give: the error is signed against the headline, and its size is not priced anywhere in this paper. Read against the committed floor-to-marginal grid, a 5.52\% floor brackets the marginal between the 5.334\% row's $+4.27$ points and the 6.0\% row's $+2.28$, that is, below the $+4.3$ lower edge this paper quotes. I do not restate the headline on the Fannie floor, because the production calibration is the Freddie one and the two books differ in composition as well as in this read; what the out-of-agency read establishes is that the band's \emph{lower} edge is the soft one, which is the direction the age-standardized caveat already pointed.

Because that sensitivity is what makes the floor's provenance decisive, Section~\ref{sec:robustness-floor} re-measures the floor outside the evaluation window and re-runs both legs there: the marginal lands at $+4.3$ to $+6.8$ points across the defensible off-window range, inside the box and well clear of either edge, and stays positive at all twenty-seven floor-by-band cells run. The covariate-prior swing of Section~\ref{sec:pathb} is the design's largest single level sensitivity at both calibrations, and at the headline floor it exceeds the marginal itself: re-run at the 4.991\% floor, the two covariate endpoints span 8.2 points of recovery against the $+5.6$-point marginal (run \texttt{covariate\_ablation\_offwindow}). The marginal is the counter-fact: switching the covariate block between its endpoints moves the off-window marginal by only $+0.1$ points, against $+1.9$ in-window, because both legs shift together where the floor censors more.

That the marginal is positive in all twenty-seven cells needs a word of its own, because it reads as the robust identified content, and is not. Over the forty window months the cohort panel covers (June 2022 through September 2025; the panel ends two months before the window does) the exposure-weighted mean coupon on the surviving book is 3.134\%, with 97.40\% of exposure at or below 4.79\% and 99.53\% at or below 5.09\%, against a window-minimum 30-year rate of 5.2311\% (the August 2022 monthly mean --- the window's rate dip --- not a June 2022 figure; \texttt{sign\_forcing\_stats}). The two uncovered months cannot overturn this: the book is closed at 2017--2021 originations, so no higher-coupon loan can enter it, and the out-of-the-money share can only hold or rise as the low-coupon cohorts season. Essentially the entire book is out of the money for essentially the entire window, so the rate gap is negative in every cell of every sweep reported here, and a specification in which a more negative gap suppresses prepayment cannot return a negative marginal at any floor or band in the box. The sign is fixed by the rate configuration rather than by the estimated elasticity. The twenty-seven-cell check accordingly verifies the specification---it would catch a wiring error that flipped a coefficient or crossed the central and null legs, and that is worth having---rather than identifying the mechanism, and I retain it on that footing. This is why the reported range, and not the sign, carries the paper's identified content. The band's monotone response is forced the same way: \eqref{eq:beta1} is monotone in $\delta$ and the hazard is monotone in $\beta_1$ at that gap, so the band's ordering is a composition of monotone maps. I retain that ordering as a wiring check on the same footing, not as identified content; the sweep's endpoints stay the convention range they are.

Second, timing carries no credential: the $\beta_1 = 0$ null also peaks at lag $-3$ (Section~\ref{sec:hazard-interp}), so the lead cannot discriminate the elasticity from the mechanical model. Third, the marginal's sampling precision is not the operative uncertainty, and the premise is now the cluster scheme's rather than the retracted one's: the stratum-cluster bootstrap puts the paired marginal at $[+8.27, +10.19]$ points at this calibration (the within-stratum $[+9.17, +9.23]$ understates by 30.8$\times$; Section~\ref{sec:pathb}), and the calibration box, at 11.1 points of width against the cluster read's 1.9, still governs the interval.

A fourth qualification concerns composition. This subsection's decomposition exhibits are computed at the in-sample calibration point, so the total they partition is the $+\$70.3$ billion figure; the composition findings are shares and intensities, which the floor's level shifts far less than it shifts the total. A group-ablation decomposition of the marginal itself (zeroing $\beta_1$ cell by cell over a vintage $\times$ coupon grid, with everything else held at production) shows the $+\$70.3$ billion is broad-based rather than cohort-specific: all fifteen cells contribute positively, the cells sum to within 1.03\% of the paired-run total, no single cell contributes more than 21\%, and the 2020--2021 vintages carry 72\% of the marginal on 73\% of sample balance, with per-balance intensity ranging from $0.82\times$ to $1.75\times$. Extending the same group-ablation engine to the dimensions that grid does not cover, the marginal is broad-based on all three: every cell of the FICO ($<$680 / 680--740 / 740$+$), LTV ($\leq$80 / $>$80), and Census-region partitions contributes positively, the per-grid cell sums reproduce the $+\$70.3$ billion total within interaction residuals of 0.8, 0.4, and 0.5 percent, and every cell's per-balance intensity lies in 0.88--1.26, tighter than the vintage $\times$ coupon grid's span. That verdict does not, on its own, discriminate, and I say so rather than leave the reader to assume otherwise: running the identical engine and the identical pre-committed decision rule over twelve random three-way partitions of the same 75,000 loans returns \emph{broad-based} in all twelve, on per-balance intensities of 0.995--1.007 and additivity residuals of 0.49 percent (\texttt{subgroup\_marginals\_placebo}). A partition with no economic content earns the label, so the label alone carries no evidential weight. What the comparison does establish runs in the opposite direction. The random partitions' intensities concentrate on unity --- the widest of the twelve spans 0.011 --- while the economic groupings span 0.381, more than thirty times wider. The informative statement is therefore not that these cells are homogeneous, which chance alone delivers, but that they disperse considerably more than chance: the tilts below are structure, not sampling noise, and the generality claim rests on the sign homogeneity and on the bounded magnitude of that dispersion rather than on the verdict label. The mild tilts run toward below-740-FICO ($1.23\times$--$1.26\times$) and above-80-LTV ($1.23\times$) borrowers, with regional intensities spanning only 0.93--1.06. The marginal is a book-wide property of the rate-gap distribution interacting with seasoning and the floor; Figure~\ref{fig:marginalcells} maps the cells. The benchmark-basis question is likewise answered by direct computation: Section~\ref{sec:method-benchmark}'s expectations-based complement shows the Federal Reserve's own ex-ante projection already implied 75.6--88.5\% of the realized cap-shortfall. The 0.2-point distance between the 88.5\% upper end and the null's 88.7\% benchmark-consistent recovery is not corroboration: the coincidence is doubly conditional, holding only at the in-sample 4.0\% floor \emph{and} under the uniform-spread allocation, which Section~\ref{sec:method-benchmark} itself judges the less faithful of the two constructions. At the headline off-window floor the null recovers 85.7\%, 2.8 points \emph{below} the uniform-spread projection rather than 0.2 above it; against the settlement-aware allocation's 75.6\% the null sits 13.2 points above in-sample and 10.1 points above at the headline floor. Two quantities that move by several points under either condition cannot be reported as converging, so I report the projection range and the null separately and draw no agreement claim from their proximity. Every percentage-point magnitude quoted for this decomposition is a cap-basis quantity, with dollars and rankings basis-invariant.

The monthly-path and coupon-bucket timing exhibits are collected in Appendix~\ref{app:patha-detail}. Because neither simulated leg carries calendar-month seasonality, seasonal structure enters these exhibits only through the empirical comparison, and per the pre-committed posture of Section~\ref{sec:hazard-interp} they are fit diagnostics, not timing credentials.

A fifth qualification states what kind of claim ``identifies'' makes. The marginal is a within-model counterfactual decomposition: the difference between two runs of a fully disclosed hazard structure that differ only in the elasticity term. It is not causal identification in the econometric sense: no quasi-experimental variation in rate gaps across otherwise-comparable loans is exploited anywhere in this paper, and none is claimed. The causal content of $\beta_1$ itself is inherited rather than asserted: the elasticity is imported unmodified from \citepos{liebersohn2024} estimated band, with \citet{fonseca2024} and the transaction-record band corroboration of \citet{graybill2026} as independent support, and no Path B parameter is estimated on the benchmark (Section~\ref{sec:pathb}). What the switch-off adds is accounting, at a stated price: the marginal inherits any misspecification of the held-fixed structure, which is why the channels I can name (floor level and functional form, the rate-gap transform, the covariate priors) are each swept directly (Sections~\ref{sec:pathb} and \ref{sec:robustness-floor}), and why Table~\ref{tab:uncertainty} reports calibration, not sampling, as the operative uncertainty.

A sixth qualification concerns three readings I decline in three separate places. Path A's fitted rate-gap coefficient is $+0.29$ per 100 bp against the imported $\beta_1$'s $0.069$, though the two are on different scales and I compare them on sign rather than magnitude (Section~\ref{sec:patha}); the realized cross-cohort gap gradient runs at roughly $4.5\times$ the gradient the central $\beta_1$ implies (run \texttt{episode\_confrontation}); and the real-covariate cross-design variant recovers 59.3\% of the benchmark, 76.3\% once reweighted to the book's own composition, against the production ABM's 13.6\% seed mean (Section~\ref{sec:robustness-crossdesign}). I stand by each reason for declining: the first is not unit-comparable and does not survive a bias-respecting sign test; the second survives composition standardization ($+3.44$ against the implied $+0.94$, run \texttt{episode\_confrontation\_within}) and is declined on the narrower ground that remains --- within-window refinancing on the shallow buckets is inseparable from moving, and a level gradient is not a dollar marginal; the third is a recovery level, which this design does not identify, and the same reweight moves the frozen cross-design variant down to 12.6\%. Stated together rather than three sections apart: none of the three points toward a smaller lock-in channel. I keep the imported value for the reason Section~\ref{sec:pathb} gives---Path B is a structural verification exercise whose elasticity is estimated in a different setting and never touches this benchmark, so re-fitting it to any in-sample reading would end that exercise---and among the literature estimates available the adopted band is already the conservative one, since \citepos{fonseca2024} anchor puts the marginal at $+11.5$ points against the central calibration's $+9.2$ at the production floor. Other disclosed variants move the marginal the other way, the Ginnie-speed overlay to $+4.4$ points and the age-standardized floor read to $+3.8$ (Table~\ref{tab:uncertainty}), so I do not read the three as a case for a larger headline.

A seventh qualification assembles those readings against the ones that run the other way, because this paper states each correction where it arises and never lines them up, and a reader who lines them up sees something the sections do not show one at a time. Holding the production hard-maximum form and the imported elasticity fixed, every correction to the floor or to the accounting basis moves the headline down. The in-sample calibration returns $+9.2$ points; re-anchoring the floor outside the window returns the headline $+5.6$; the Ginnie-speed overlay at that same floor returns $+4.4$; the age-standardized floor read of 5.51\% implies a marginal near $+3.8$; and the Fannie Mae read of the same selection rule on the pooled off-window period, 5.52\%, brackets the marginal below the $+4.3$ edge of the clean band (above), its like-for-like Freddie counterpart 5.19\% rather than the 2018 leg's 4.991\%. The anchor convention does not flatter the headline: extending the depth ladder to the five reads the 10\% support rule admits puts the middle of them at 4.869\%, below the adopted 4.991\%, and a lower floor implies a marginal above the headline rather than below it (run \texttt{matched\_depth\_reconciliation}). Carrying the overlay's $0.797$ scaling up to the age-standardized floor projected a marginal near $+3.0$; the joint cell has now been run under a pre-committed landing rule (run \texttt{b5\_joint\_cell}) and measures $+2.9$ points, an interaction within $0.1$ of proportional --- where the first joint cell this paper ran for two adjustments to the marginal returned an interaction of $-1.47$ points (Section~\ref{sec:robustness-floor}): corrections here do not compose additively, and a composed figure would carry a precision nothing in the design supports. Nor are they independent of each other. The floor-level corrections and the censoring share are one mechanism read twice, since raising the floor is what crowds the elasticity out: at the 4.991\% headline floor the central leg is already pinned to the floor in 68.8\% of evaluated loan-months, against 36.3\% at the in-sample calibration point, so a further increase in the floor has progressively less of the hazard left to act on. That share is not a statement that the estimate rests on a third of the data, and it should not be read as one: a pinned month's contribution to the marginal is truncated at the floor rather than erased, the $\beta_1 = 0$ leg is pinned in only 35.8\% of loan-months at the same floor so the two legs still separate wherever the null clears it, and these are unweighted loan-month counts rather than balance-weighted shares, so nothing follows from them about what fraction of the dollar marginal the uncensored months carry. What the assembly settles is where in the interval the mass sits, not whether the finding holds. Within the production form the widest disclosed layer is the baseline level: sweeping the seasoning ramp over 75 to 150 PSA on both legs carries the marginal from $+0.9$ to $+15.6$ points at the off-window floor and from $+3.3$ to $+16.5$ at the in-sample calibration (run \texttt{psa\_level\_sweep}), a span with no coverage property --- it is a range over a convention I did not estimate, and its endpoints are not draws from anything. Aggregate fit is not available as a defense of the 100~PSA convention: at the off-window floor the $\beta_1 = 0$ null recovers 101.9\% of the benchmark under the standalone scorer at 75~PSA against 94.8\% at the production 100, while the central leg's recovery is closer to the benchmark at 100 (100.4\%) than at 75, so the two legs rank the ramps opposite ways. A conditional estimate of the ramp's \emph{shape} now exists, and it does not favour the headline. Transporting the age profile of Path~A's own estimated seasoning spline onto $h_0$---level held fixed by construction, because the spline's level is absorbed by its fitted constant and the stratum fixed effects and is not separately identified---returns a median marginal of $+9.44$ points, with a 5th-to-95th-percentile spread of $+8.99$ to $+9.81$ across the 107 of 200 bootstrap replicates whose normalised profile stays inside the usable range (run \texttt{h0\_reanchor}). Every one of those replicates lies above the $+5.6$ this paper headlines, and all but one lies above the $+8.7$ upper edge of the binding interval quoted below, though all of them lie inside the $75$--$150$~PSA span just given. With 46.5\% of replicates discarded it is a conditional estimate and not the ramp's sampling interval, so the convention span above stays labelled unestimated. I had pre-committed to the opposite sign: the prediction recorded before the run was that re-anchoring would push more loan-months onto the floor and so \emph{lower} the marginal. The bind share did rise as predicted, from 35.8\% to between 50.4\% and 58.2\% on the null leg, and the marginal rose anyway---a failed prediction, reported as one. The floor read itself transports in two directions, and both were fixed before either ran (run \texttt{floor\_transport\_two\_sided}). Standardizing its calendar raises it: the 2018 leg is read in August--December, on the declining side of the committed seasonal profile, so re-weighting to the window's own calendar mix carries 4.991\% to 5.2156\% and the marginal to $+4.70$ points. Matching its housing-activity level lowers it: active listings averaged 1.27 million over those 2018 months against 0.80 million over the window, so the same multiplicative transport carries the floor to 3.1533\% and the marginal to $+10.69$. I report the two at their own floors and do not net them---two floor-side adjustments composed with an interaction of $-1.47$ points the one time this paper ran them jointly, so a cancellation would be a finding requiring its own cell rather than an arithmetic consequence of these two. The activity leg is a listings-based proxy for an argument about existing-home \emph{sales}, which is not constructible here: the sales series carries no observations over the 2018 leg months at all. Its magnitude also inherits a one-for-one mapping from activity level to turnover level that nothing here estimates; the sign is what it establishes. Post-stratifying the draw onto the book's composition moves the marginal, and by less than the paper's own share arithmetic implies. Raking the pool's balance weights to the SOMA coupon marginal and to the book's vintage mix within the 2017--2021 support the draw has---both margins hit to better than $10^{-10}$ (run \texttt{book\_composition\_marginal})---returns $+5.11$ points against the unweighted $+5.57$, a move of $-0.47$. The vintage margin alone returns $+4.67$, so the coupon margin offsets about half of the vintage correction rather than compounding it, and the vintage overlay's $+3.7$, which is pure share scaling, overstates the composition correction by roughly a factor of four. Two limits travel with the number. The book's coupon-by-vintage \emph{joint} distribution is not observable---the SOMA cohort tabulation collapses vintage to a single value-weighted origination date per coupon bucket---so this matches two margins and not a cell grid. And the pre-2017 and 2022 vintages, 33.7\% of book face, have no support in the draw at all and are not re-weighted here; they remain covered by the vintage overlay. I pre-committed both a direction and a band: the direction held and the marginal fell, while $+5.11$ lies just outside the $[+2.5, +5.0]$ band I fixed before the run, which I report rather than re-draw. The widest layer that does have a coverage property is the floor reads' own sampling error, $+2.9$ to $+8.7$ points after wild-cluster correction of the percentile read's under-coverage (run \texttt{floor\_inference\_correction}; the percentile read is $+3.0$ to $+8.0$), and that is the interval this paper quotes as binding; the headline sits just below that interval's midpoint, while every correction listed above falls in its lower half and the uncomposed pair falls near its bottom. The convention span and the sampling interval are nested rather than alternative: the seasoning-ramp sweep is the outer statement, and $+2.9$ to $+8.7$ is a sampling interval on the floor read alone, nested inside it at the production 100~PSA ramp. Which rung of that inner layer to quote is a further choice the rule does not settle: the Webb wild-$t$ read quoted here is the third-narrowest of the nine undemoted rungs, and the widest rung with both endpoints interior to the floor grid---CR2 at Bell--McCaffrey degrees of freedom, $+2.4$ to $+9.1$ points---stands beside it rather than behind it. I attach no posture to where $+5.6$ sits among these readings: with the downward specification variants lined up below beside the upward counterweights, no interior member is privileged, and the range rather than any point is what the design delivers---with the form-conditional hull of $+3.5$ to $+13.1$ points remaining the wider statement, since the form dimension is not a floor correction and is not disciplined by this list. Three things stop the list from becoming a retraction. The sign is fixed by the window's rate configuration rather than by any of these calibrations, as the sign-forcing argument above establishes, so nothing here bears on the direction of the result; the interval's exclusion of zero is deliberately not among the three, because a forced sign makes zero-exclusion a restatement of that configuration rather than independent evidence for it. Specification choices that are not floor corrections run both ways, and this assembly now lines both directions up. Upward: the additive form returns $+11.2$ points at the same off-window anchors, nearly floor-invariantly, and it is what sets the hull's upper end, while \citepos{fonseca2024} anchor would put the marginal at $+11.5$ points at the production floor against the adopted band's $+9.2$. Downward: the concave transform alone returns $+5.06$ at the off-window point floor against the production $+5.57$ (run \texttt{concave\_marginal}), the sub-band elasticity at $\delta = 3.25\%$ returns $+3.47$ (Table~\ref{tab:lowband}), and the moving-share bracket re-run at the headline floor returns $+3.36$ points at $s = 0.5$ and $+1.82$ at $s = 0.25$ against $+5.57$ at $s = 1$ (run \texttt{moving\_share\_bracket\_offwindow}) --- a convention bracket rather than an estimate, and mildly super-proportional because censored contributions do not scale with $s$. The survival-selection attenuation of the same imported coefficient is a bracket of the same kind and belongs on this side: $a = S^{\theta}$ at the measured pre-window survival share carries the marginal from $+5.1$ at $\theta = 0.25$ to $+3.7$ at $\theta = 1$ (run \texttt{attenuation\_sensitivity}), with $\theta$ free (Section~\ref{sec:pathb}) --- a different axis from the elasticity band, not a widening of it. The assembly is not exhaustive over runs that score a marginal at a labelled floor: the baseline-level PSA sweep, the covariate ablation and the floor-form mixture do so as well and are reported at their own sites rather than as rows here, and I state no rule that would admit these rows while excluding those. A calibrated companion prices the omitted housing-activity level term (Table~\ref{tab:specbox}) from the paper's own catalogued anchor: \citet{aladangady2024} attribute 44\% of the 2021--2022 mobility decline to rate-gap lock-in, and scaling the seasoning baseline in both legs until the non-rate channel traps $0.56/0.44$ times the rate channel's trapped liquidity puts the root at $\phi^{*} = 0.754$ --- effectively 75.4 PSA --- and the marginal at $+0.9$ points at the headline floor and $+3.5$ at the in-sample calibration (run \texttt{scaled\_null\_housing\_activity}; the $\phi = 1$ cells reproduce the committed values bit-exactly). The variant documents an upward bias; the corrected member lies below the headline, at the bottom of the seasoning-ramp span just swept, and at the headline floor it falls outside the quoted $+2.9$ to $+8.7$ interval---an externally anchored calibration the quoted interval does not cover, not a challenge to the sign; at the in-sample calibration it falls inside. That the floor side runs one way is also an absence: the floor is a level of total turnover rather than of a strictly involuntary component of it (Section~\ref{sec:pathb}), so it carries the activity of the period it is read on, and re-reading it at the window's own level enters at the floor rather than at the seasoning ramp just scaled, so it cuts the other way --- a lower anchor raises the marginal. It is the one floor-side correction I can name that would move the headline up, and it turns on whether the window's turnover level is the lower of the two, which I do not measure here. And the three readings of the preceding paragraph point to a larger lock-in channel, not a smaller one. I decline those three for the reasons given there, and I decline to compose these for the reason given here: the identified content is the range; the floor and basis corrections concentrate in its lower half, and the form-side counterweights sit above it. Table~\ref{tab:assembly} tabulates this assembly.

\begin{table}[H]
\centering
\caption{The assembly at a glance. Every value repeats a committed figure from the text of this subsection; the binding interval and the form-conditional hull are deliberately not restated here (see the text and Table~\ref{tab:uncertainty}).}
\label{tab:assembly}
\footnotesize
\begin{tabular}{@{}p{3.1cm}p{2.6cm}p{2.4cm}p{7.7cm}@{}}
\toprule
Reading & Marginal (pp) & vs.\ $+5.6$ & Status \\
\midrule
In-sample calibration point & $+9.2$ & above & in-window point, not the headline \\
Off-window re-anchor & $+5.6$ & headline & mid-grid anchor (the middle of the three depth cuts fixed before the reads) inside the binding interval \\
Ginnie overlay at the corrected floor & $+4.4$ & below & change of estimand (re-scoping onto the conventional sub-book), not a correction \\ Vintage overlay at the corrected floor & $+3.7$ & below & change of estimand (re-scoping onto the in-window-vintage share), not a correction (run \texttt{vintage\_overlay}) \\
Age-standardized floor (5.51\%) & $\sim +3.8$ & below & indicative bound, 84\% imputed weight \\
Fannie floor read (5.52\%) & below $+4.3$ & below & bracketing; second agency, pooled period; like-for-like Freddie 5.19\% \\
Overlay $\times$ age-standardized & $+2.9$ & below & measured composed lower member (run \texttt{b5\_joint\_cell}); within $0.1$ of proportional \\ Moving-share bracket, $s = 0.5$, off-window & $+3.4$ & below & convention bracket, not an estimate (run \texttt{moving\_share\_bracket\_offwindow}) \\ Moving-share bracket, $s = 0.25$, off-window & $+1.8$ & below the interval & the bracket's deep stress; edge by construction \\ Book-composition post-stratification & $+5.11$ & below & two margins raked (SOMA coupon $\times$ book vintage) on the 2017--2021 support; the joint is unobserved and 33.7\% of face is out of support (run \texttt{book\_composition\_marginal}) \\ Survival-selection attenuation, $\theta = 0.25$ to $1$ & $+5.1$ to $+3.7$ & below & bracket in $\theta$ on the imported coefficient ($a = S^{\theta}$, $\theta$ not estimable here); a different axis from the elasticity band (run \texttt{attenuation\_sensitivity}) \\ Housing-activity-scaled baseline (Aladangady-anchored) & $+0.9$ & below the interval & scaled-null variant; upward-bias entry, root $\phi^{*} = 0.754$ (run \texttt{scaled\_null\_housing\_activity}) \\ Sub-band elasticity, $\delta = 3.25\%$, off-window & $+3.5$ & below & band-extension member (Table~\ref{tab:lowband}) \\ Concave transform alone, off-window & $+5.1$ & below & transform member (run \texttt{concave\_marginal}) \\
Concave $\times$ additive joint cell & interaction $-1.47$ & --- & the one joint cell run \\
Additive form & $+11.2$ & above & counterweight, functional-form choice \\
Fonseca--Liu anchor & $+11.5$ & above & counterweight, conservative-band evidence \\ Cross-sectional gradient, composition-standardized & gradient $+3.4$ vs implied $+0.9$ & above (directional) & directional; not a marginal re-estimate (run \texttt{episode\_confrontation\_within}) \\
\bottomrule
\end{tabular}
\end{table}

\begin{table}[!t]
\centering
\footnotesize
\begin{threeparttable}
\caption{Out-of-sample floor anchors and the lock-in marginal each implies, in percentage points of the \$764.7 billion benchmark. Floors are exposure-weighted annualized turnover CPR on out-of-the-money cohort-months measured \emph{outside} the June 2022--November 2025 evaluation window, except the early-window calibration row, whose floor is calibrated on the window's own first nineteen months. The pooled and 2019-leg rows read 2017--2019 performance; the three 2018-leg rows that set the headline read 2018 reporting months on 2017-vintage cohorts, and every cohort-month they contain is aged twelve to twenty-four months, so no mature cell enters them. Every re-run row, the early-window one included, reports its floor's \emph{full-window} marginal on a common basis, and the held-out-months evaluation is given in the note; the age-standardized and Fannie rows are grid reads rather than paired re-runs, so they carry the central elasticity alone. Contaminated anchors are reported for completeness and are not this paper's off-window anchor; the pooled anchor's 6.07\% measurement is evaluated at the box's adjacent 6.0\% grid row. Every marginal is basis-invariant and reproduces the committed production artifacts exactly at the 4\% row. Pre-committed run \texttt{oos\_identification}; artifact catalogued in Appendix~\ref{app:ledger}.}
\label{tab:oosfloor}
\begin{tabular}{llcccc}
\toprule
Floor anchor & Status & Floor & $\delta = 5.5\%$ & $\delta = 6.5\%$ & $\delta = 7.7\%$ \\
\midrule
2018 leg, gap $\leq -0.005$ & defensible & 4.70\% & $+6.1$ & $+6.8$ & $+7.5$ \\
2018 leg, gap $\leq -0.0025$ & defensible & 4.99\% & $+5.0$ & $+5.6$ & $+6.1$ \\
2018 leg, gap $\leq 0$ & defensible & 5.33\% & $+3.9$ & $+4.3$ & $+4.6$ \\
\midrule
Pooled 2017--2019 & refi-contaminated & 6.07\% & $+2.1$ & $+2.3$ & $+2.4$ \\
2019 leg, gap $\leq 0$ & refi-contaminated & 6.91\% & $+0.7$ & $+0.7$ & $+0.8$ \\
\midrule
2018 leg, age-standardized & indicative bound\tnote{b} & 5.51\% & --- & $+3.8$ & --- \\
Fannie Mae, pooled off-window & cross-agency read\tnote{c} & 5.52\% & --- & $+2.3$ to $+4.3$ & --- \\
\midrule
In-window (production) & in-sample point & 4.00\% & $+8.1$ & $+9.2$ & $+10.4$ \\
Early-window calibration & full window\tnote{a} & 3.97\% & $+8.1$ & $+9.3$ & $+10.5$ \\
\bottomrule
\end{tabular}
\end{threeparttable}
\end{table}

\noindent{\footnotesize \emph{Notes to Table~\ref{tab:oosfloor}.} a Temporal holdout. The 3.97\% floor is calibrated on deeply out-of-the-money turnover in the window's first nineteen months (June 2022--December 2023). Summed over the strictly held-out remaining twenty-three months (January 2024--November 2025), which enter that calibration nowhere, the central-minus-null marginal is $+\$45.1$ billion at the central elasticity, against $+\$44.8$ billion at the production floor over the identical months. The table row above reports the same floor's full-window marginal, so that every row is comparable; the two quantities are different objects and should not be read off one another. \quad b The 2018 leg's age-bucket reads reweighted to the in-window deep-discount population's age mix, its empty mature buckets imputed from the in-window age gradient: 84.0\% of the weight is imputed, so this is an indicative bound and not a measurement (run \texttt{floor\_uncertainty}). \quad c The committed off-window selection rule applied to the Fannie Mae panel at gap $\leq -0.0025$ over the pooled 2017--2019 period, not to the 2018 leg: the like-for-like Freddie read of that cell is 5.185\%, so the agency difference is 0.337 points of CPR and the rest of the distance from the 2018 leg's 4.991\% is period (run \texttt{fannie\_floor\_read}). Both correction rows are read off the committed floor-to-marginal grid rather than re-run.\par}

\begin{table}[!t]
\centering
\caption{The lock-in marginal as a curve in the imported elasticity. Paired central-minus-null runs at each quarterly mobility decline $\delta$ (the elasticity import of Section~\ref{sec:pathb}, mapped to $\beta_1$ via \eqref{eq:beta1}), with the floor and every other production convention held fixed, at the in-sample 4.0\% floor and the off-window 4.991\% headline floor (run \texttt{band\_low\_extension}; the band-cell rows reproduce the committed production and off-window cells within the sweep tolerance of $\pm\$0.01$ billion). The $\delta = 0$ rows are the $\beta_1 = 0$ nulls, so their marginal is zero by construction; $\delta = 3.25\%$ is half the adopted central; $\{5.5, 6.5, 7.7\}\%$ is the adopted Liebersohn--Rothstein band.}
\label{tab:lowband}
\footnotesize
\begin{threeparttable}
\begin{tabular}{@{}rrrrrr@{}}
\toprule
& & \multicolumn{2}{c}{In-sample floor (4.0\%)} & \multicolumn{2}{c}{Off-window floor (4.991\%)} \\
\cmidrule(lr){3-4}\cmidrule(lr){5-6}
$\delta$ (\%) & $\beta_1$ & marginal (\$B) & (pp) & marginal (\$B) & (pp) \\
\midrule
0.00 & $0.0000$ & $0.0$ & $0.00$ & $0.0$ & $0.00$ \\
1.00 & $0.0103$ & $+12.8$ & $+1.67$ & $+9.4$ & $+1.23$ \\
2.00 & $0.0206$ & $+24.9$ & $+3.26$ & $+17.7$ & $+2.31$ \\
3.00 & $0.0311$ & $+36.4$ & $+4.76$ & $+24.9$ & $+3.25$ \\
3.25 & $0.0337$ & $+39.1$ & $+5.12$ & $+26.5$ & $+3.47$ \\
4.00 & $0.0417$ & $+47.1$ & $+6.16$ & $+31.0$ & $+4.06$ \\
5.00 & $0.0523$ & $+57.1$ & $+7.46$ & $+36.3$ & $+4.74$ \\
5.50 & $0.0577$ & $+61.7$ & $+8.07$ & $+38.6$ & $+5.05$ \\
6.50 & $0.0686$ & $+70.3$ & $+9.20$ & $+42.6$ & $+5.57$ \\
7.70 & $0.0817$ & $+79.5$ & $+10.39$ & $+46.6$ & $+6.09$ \\
\bottomrule
\end{tabular}
\begin{tablenotes}[flushleft]\footnotesize
\item The $\beta_1$ column follows \eqref{eq:beta1}'s sign convention, under which $\beta_1$ is positive and suppression enters \eqref{eq:pathB} through the signed gap $g_{i,t}$, negative for a locked-in borrower. The production helper \texttt{rothstein\_beta1} returns the same quantity without that leading minus, so a replicator reading the code will see these nine values negated; the marginal columns are unaffected either way.
\end{tablenotes}
\end{threeparttable}
\end{table}

\begin{figure}[tbp]
\centering
\includegraphics[width=0.78\textwidth]{fig11_marginal_decomposition}
\caption{Group-ablation decomposition of the lock-in marginal over the vintage $\times$ coupon grid (production configuration; each cell's $\beta_1$ zeroed with everything else held fixed). All fifteen cells contribute positively; the 2020--2021 rows carry 72\% of the cell sum on 73\% of sample balance, and no cell exceeds 21\%. Cell values and shares are read from the committed decomposition artifact, whose additivity gate bounds the cells' sum within 1.03\% of the paired-run \$70.3 billion in-sample total.}
\label{fig:marginalcells}
\end{figure}

Table~\ref{tab:uncertainty} consolidates the three uncertainty layers that attach to the headline level and marginal (sampling, simulation seed, and calibration) on both accounting bases. The hierarchy it displays is the one argued above: sampling error under the operative cluster schemes sits an order below the calibration layers, the draw-seed contribution is smaller still (under \$0.01 billion across ten paired reruns), and the operative uncertainty on every headline quantity is calibration --- except the headline marginal, where the binding layer is the floor read's own propagated sampling error (Section~\ref{sec:robustness-floor}), the layer Table~\ref{tab:uncertainty}'s headline row quotes --- binding among the layers with a coverage property; the unestimated baseline-level convention spans wider (run \texttt{psa\_level\_sweep}, Section~\ref{sec:identification}).

{\footnotesize
\begin{longtable}{@{}p{2.4cm}p{3.3cm}p{4.0cm}p{4.8cm}@{}}
\caption{Consolidated uncertainty for the headline Path B quantities. Sampling intervals are the stratified loan-level bootstrap's 95\% percentile intervals (200 replicates, Section~\ref{sec:pathb}); that scheme resamples within strata and holds the 130 strata fixed, so the headline row also carries the stratum-cluster interval, which is the wider of the two loan-sampling reads; shared-basis intervals follow from the standalone intervals by the exact common netting of Section~\ref{sec:robustness-hybrid} (9.10 points of benchmark), which cancels from the marginal entirely. Calibration entries restate Sections~\ref{sec:pathb} and~\ref{sec:robustness-floor}.}
\label{tab:uncertainty}\\
\toprule
Quantity & Point (shared; standalone) & Sampling 95\% CI & Calibration range \\
\midrule
\endfirsthead
\toprule
Quantity & Point (shared; standalone) & Sampling 95\% CI & Calibration range \\
\midrule
\endhead
Path B central level (in-sample calibration) & 97.9\%; 107.0\% & $[97.9, 98.0]$\%; $[107.0, 107.1]$\% & covariate priors 92.2--99.2\%; elasticity band 96.8--99.1\% (shared basis); off-window floor 91.3\% shared (88.2--93.7\% across its range); covariate endpoints at the off-window floor 94.3--102.5\% standalone, an 8.2-point swing exceeding the marginal there (run \texttt{covariate\_ablation\_offwindow}) \\
$\beta_1 = 0$ null level (in-sample calibration) & 88.7\%; 97.8\% & $[88.7, 88.8]$\%; $[97.8, 97.9]$\% & moves one-for-one with the floor level: 85.7\% shared at the off-window floor, 83.9--86.9\% across its range (Section~\ref{sec:robustness-floor}) \\
Lock-in marginal, off-window floor (headline) & $+5.6$pp (basis-invariant) & floor-read wild-cluster bootstrap-$t$ $[+2.9, +8.7]$pp at the central elasticity $\delta = 6.5\%$ (the binding layer among those with a coverage property---the unestimated baseline-level convention spans wider---and a sampling interval on the floor read alone at the production hard-maximum form and the 100~PSA ramp, pricing neither the elasticity nor the read-to-read spread across the floor reads in the next column; percentile read $[+3.0, +8.0]$pp; the elasticity dimension is Table~\ref{tab:lowband}); loan/stratum cluster bootstrap $[+4.63, +6.92]$pp (width $0.39\times$ the corrected floor read, meaning the Webb wild-$t$ interval, against $0.455\times$ the demoted percentile read; 25.8 effective clusters); the two layers convolved as independent, $[+2.80, +8.99]$pp (run \texttt{layer\_convolution}; independence is assumed, not measured --- under maximal positive dependence the width is $8.12$pp) & $+4.3$ to $+6.8$ across the defensible off-window floor range (Table~\ref{tab:oosfloor}); age-standardized floor 5.51\% $\to$ $+3.8$ (84\% imputed weight; band open below $+4.3$); Fannie Mae read of the same selection rule on the pooled off-window period 5.52\% (run \texttt{fannie\_floor\_read}), above the clean band and bracketing the marginal below $+4.3$, its like-for-like Freddie counterpart 5.185\% rather than the 2018 leg's 4.991\%, so 0.337 points of that distance is agency and the rest is period; $+3.9$ to $+7.5$ joint with the band; additive form $+11.2$ (form-conditional hull $+3.5$ to $+13.1$); concave transform $+5.1$; the two \emph{jointly} $+9.2$, an interaction of $-1.47$ points, so these two entries do not add (run \texttt{concave\_additive\_marginal}); Ginnie overlay $+4.4$ (cuts the same way as the age-standardized read; now run jointly with it in run \texttt{b5\_joint\_cell}: composed $+2.9$, within $0.1$ of proportional); floor-stability check $+\$45.1$B; baseline-level sweep 75--150 PSA $+0.9$ to $+15.6$ (a convention range, no coverage property; run \texttt{psa\_level\_sweep}); covariate endpoints move this marginal by only $+0.1$ while spanning 8.2 points of level (run \texttt{covariate\_ablation\_offwindow}); survival-selection attenuation $+5.1$/$+4.6$/$+3.7$ at $a = 0.86/0.73/0.54$, the frailty-identity mapping $a = S^{\theta}$ at the measured $S = 40{,}234/75{,}000$ and $\theta = 0.25/0.5/1$ (run \texttt{attenuation\_sensitivity}; a transport sensitivity on the imported coefficient, listed in Table~\ref{tab:assembly} as a bracket in $\theta$ rather than as a correction) \\
\addlinespace
Lock-in marginal, in-sample point & $+9.2$pp (basis-invariant) & within-stratum $[+9.17, +9.23]$pp (sd $0.016$); stratum-cluster $[+8.27, +10.19]$pp (sd $0.52$; run \texttt{bootstrap\_pathb\_cluster}) & box $+2.1$ to $+13.2$; floor form $+11.2$; within-run floor dispersion $+7.1$ to $+9.4$; regime-split baseline break $+5.0$ to $+12.1$; PSA level sweep $+3.3$ to $+16.5$ (run \texttt{psa\_level\_sweep}); covariate endpoints move this marginal $+1.9$ \\
\bottomrule
\end{longtable}
}

\noindent{\footnotesize \emph{Notes to Table~\ref{tab:uncertainty}.} Seed noise: the marginal varies by less than \$0.01 billion across ten paired draw seeds (Section~\ref{sec:hazard-interp}). The Fannie Mae external replication reproduces the marginal at $+8.68$ points against Freddie's $+9.20$; it clears the pre-committed envelope, but that envelope is 11.06 points wide and the agreement, not the gate, is what carries the information (Section~\ref{sec:pathb-fannie}). \quad Both cluster bootstraps above rest on few clusters: 31 carry the floor read, and the loan/stratum scheme's 130 strata are unequal enough to be worth 25.8 effective ones. Percentile cluster bootstraps are known to under-cover at cluster counts in that range, and this rung's printed lower edge additionally turns on an unlabelled convention: 26 of its 1{,}000 draws fall outside the committed floor grid, all 26 above the 6.0\% edge and therefore all on the low-marginal side, and the row excludes them---truncating them to the edge instead would print $+2.28$ to $+8.01$ in place of $+3.0$ to $+8.0$; run \texttt{floor\_inference\_correction} and its Webb-weight successor \texttt{floor\_inference\_correction\_v2} correct the floor read, Table~\ref{tab:ladder} sets the corrected ladder out layer by layer with its degrees of freedom, and the headline row quotes the wild-$t$ interval as the binding layer among those with a coverage property. That interval prices one layer: the floor read's own sampling error around the 4.991\% point read, propagated through the frozen floor-to-marginal grid at the production hard-maximum form, the 100~PSA ramp and the central elasticity $\delta = 6.5\%$, so none of its width comes from the elasticity, which is swept separately in the calibration column and in Table~\ref{tab:lowband} --- swept over the published 5.5\%--7.7\% range as a convention, not estimated, so no layer of this paper's accounting prices sampling error in the elasticity itself. No such correction is run for the loan/stratum scheme, whose interval is still better read as a lower bound on sampling uncertainty than as calibrated 95\% coverage; one input to the convolved line being a lower bound makes the convolved line a lower bound as well. The ladder's reads agree in location and none crosses zero, and the closest agreement among them carries no information: the Rademacher and Webb constructions land with both endpoints within a tenth of a point, a re-printing rather than a substantive move, because both are sign-symmetric reweightings of the same residuals, so each bootstrap distribution turns on the draw's sign for the one cluster carrying a third of the total leverage---nearly twice the next largest (Table~\ref{tab:ladder}). Their coincidence re-prints that cluster's influence instead of testing it. The small-sample degrees of freedom and the restricted inversion are wider. The bootstrap imposes the point estimate as the data-generating truth over unrestricted residuals (B $= 9{,}999$), studentizes on the floor's own scale, and maps endpoints through the frozen floor-to-marginal grid. \quad The within-run floor dispersion range is the $\kappa$ grid of Section~\ref{sec:pathb}. It bounds the marginal's sensitivity to floor variation within the window; it does not bracket floor cyclicality, and should not be read as doing so. Scrambling the floor path's time order leaves 85.0\% of its span standing, and switching the combination rule from the hard maximum to the additive form eliminates 98.0\% of it (the $\kappa$ grid then spans $+11.22$ to $+11.27$ points). Appendix~\ref{sec:robustness-seasonalfloor} diagnoses the Jensen mechanism responsible. The co-movement component that does survive is monotone in $|\kappa|$ and sign-consistent, but reaches at most \$5.29 billion in any cell.\par}

\begin{table}[!t]
\centering
\caption{The inference ladder on the binding layer---binding among the layers with a coverage property; the unestimated baseline-level convention spans wider. Every row is one rung on that layer --- the floor read's own sampling error, propagated through the frozen floor-to-marginal grid at the central elasticity $\delta = 6.5\%$ --- so no row's width comes from the elasticity, which is swept separately in Table~\ref{tab:lowband}. All reads are on the 2018 leg's mid-grid anchor, which 31 stratum clusters carry. The Webb-weight wild-cluster bootstrap-$t$ is the primary construction and the interval this paper quotes; the percentile read is reported demoted rather than as a competing interval. Runs \texttt{floor\_uncertainty} (percentile), \texttt{floor\_inference\_correction} (Rademacher wild-$t$; CR1--CR3 at $t(30)$), and \texttt{floor\_inference\_correction\_v2} (Webb primary, Bell--McCaffrey degrees of freedom, restricted inversion).}
\label{tab:ladder}
\footnotesize
\begin{threeparttable}
\begin{tabular}{@{}llll@{}}
\toprule
Rung & Marginal (pp) & Degrees of freedom & Status \\
\midrule
Percentile cluster bootstrap & $+3.0$ to $+8.0$ & --- (1000-replicate) & demoted; under-covers at 31 clusters \\
CR1 $t$ & $+3.2$ to $+8.3$ & $t(30)$, $G-1$ & no leverage adjustment \\
CR2 $t$ & $+3.0$ to $+8.6$ & $t(30)$, $G-1$ & conventional df \\
CR3 $t$ & $+2.8$ to $+8.8$ & $t(30)$, $G-1$ & conventional df \\
Wild-$t$, Rademacher & $+2.8$ to $+8.7$ & bootstrap, $B = 9{,}999$ & the committed construction, replayed \\
Wild-$t$, Webb & $+2.9$ to $+8.7$ & bootstrap, $B = 9{,}999$ & primary; the binding layer \\
CR1 $t$, Bell--McCaffrey & $+2.8$ to $+8.8$ & $t(6.2)$, data-driven & small-sample df \\
CR2 $t$, Bell--McCaffrey & $+2.4$ to $+9.1$ & $t(5.1)$, data-driven & small-sample df \\
CR3 $t$, Bell--McCaffrey & $+2.3$ to $+9.6$ & data-driven, CR3's own & small-sample df \\
Restricted wild-cluster inversion & $+2.3$ to $+9.1$ & 401-point grid inversion & agrees in location \\
\bottomrule
\end{tabular}
\end{threeparttable}
\end{table}

\noindent{\footnotesize \emph{Notes to Table~\ref{tab:ladder}.} The Bell--McCaffrey degrees of freedom are data-driven and estimator-specific, so one value cannot serve all three cluster-robust variants: the 6.2 and 5.1 shown are CR1's and CR2's, and CR3's own is smaller still. CR1 applies no leverage adjustment --- it is the narrowest of the three sandwiches --- and is the standard error the wild rows studentize with, so its two rows are the ladder's baseline rather than competing reads; its Bell--McCaffrey row printing CR3's conventional-df interval is an accident of this leverage profile, not an identity. The effective cluster count is 5.9 by the same Herfindahl-inverse convention as the loan/stratum scheme's 25.8, against 31 nominal clusters, and the run's recorded maximum single-cluster leverage is 0.33, ten times an equal 31-cluster share. That concentration is why the wild-cluster correction is applied and the percentile read demoted. It does not rank the wild construction against the Bell--McCaffrey degrees of freedom, and I do not claim it does: the wild rows re-draw cluster-level weights at that same leverage profile, the Bell--McCaffrey rows keep the sandwich and correct the reference distribution's degrees of freedom for it, and at 5.9 effective clusters neither correction dominates. So the Webb row is the primary read and the interval this paper quotes, but it is the third-narrowest of the nine undemoted rungs, and the widest rung with both endpoints interior to the floor grid---CR2 at Bell--McCaffrey degrees of freedom, $+2.4$ to $+9.1$ points---belongs beside it rather than behind it. The two rungs printed wider still understate their own width: CR3 at Bell--McCaffrey degrees of freedom and the restricted inversion have \emph{lower} endpoints at floors above the committed grid's 6.0\% end and are censored to it, and the marginal falls as the floor rises, so both run below the $+2.3$ printed; their upper endpoints are interior.\par}

\subsection{Interpretation}\label{sec:hazard-interp}

Table~\ref{tab:estimators} summarizes the three estimators evaluated in this paper against the common \$764.7 billion empirical benchmark.

\begin{table}[!tp]
\centering
\footnotesize
\begin{threeparttable}
\caption{Estimator comparison against the \$764.7 billion empirical trapped-liquidity benchmark.}
\label{tab:estimators}
\renewcommand{\arraystretch}{0.92}
\setlength{\itemsep}{0pt}
\begin{tabular}{@{}>{\raggedright\arraybackslash}p{3.2cm}>{\raggedleft\arraybackslash}p{2.4cm}>{\raggedleft\arraybackslash}p{2.4cm}>{\raggedright\arraybackslash}p{2.3cm}>{\raggedright\arraybackslash}p{4.0cm}@{}}
\toprule
Estimator & Trapped liquidity & \% of benchmark & CPR path corr. (lag) & Path diagnostics \\
\midrule
Empirical benchmark (SOMA, phased cap, 30+15yr) & \$764.7B & 100.0\% & -- & mean CPR 5.14\% (ABM-basis back-out; hazard-basis 5.52\%, 5.79\% at the note-rate WAC) \\
Production ABM (Section~\ref{sec:abm-results}) & \$103.7B seed mean (frozen draw \$91.0B) & 13.6\% (seed band 10--17\%; frozen draw 11.9\%) & $r = -0.318$ (0) & mean CPR 11.68\%; raw $R^2$ = $-6.984$ \\
Hazard Path B: literature microsimulation & \$818.5B & 107.0\% (band 105.9--108.2\%) & $r = +0.404$ ($-3$) & mean CPR 4.76\%; $r = +0.190$ at lag 0; 3-month offset (empirical leads; reference only, Section~\ref{sec:pathb}) \\
Hazard Path B: concave-gap variant (Section~\ref{sec:pathb}) & \$808.4B & 105.7\% & $r = +0.281$ (0) & piecewise transform, kink 200 bp; robustness on the log-linear extrapolation \\
Hazard Path A: stratum-month Poisson PML (spec v4, calendar-month) & \$928.9B & 121.5\% & $r = -0.378$ ($-2$) & mean CPR 3.34\%; holdout diagnostics under spec v3 \\
Cross-design ABM, real covariates, recalibrated (Section~\ref{sec:robustness-crossdesign}) & \$453.5B & 59.3\% (fifty-seed mean 60.2\%) & $r = -0.336$ (0) & mean CPR 8.14\%; no positive alignment at any lag \\
Cross-design ABM, real covariates, frozen calibration (Section~\ref{sec:robustness-crossdesign}) & \$159.5B & 20.9\% & $r = -0.311$ (0) & mean CPR 11.63\%; no positive alignment at any lag \\
Hazard Path B: Fannie Mae external replication (Section~\ref{sec:pathb}) & \$821.0B & 107.4\% & -- & null 98.7\%; marginal $+8.68$pp \\
\bottomrule
\end{tabular}
\end{threeparttable}
\end{table}

\noindent{\footnotesize \emph{Notes to Table~\ref{tab:estimators}.} The table reports the three primary estimators, the two Section~\ref{sec:robustness-crossdesign} cross-design variants, and the Fannie Mae external replication of Path B, whose path-correlation cell is dashed because the committed replication artifact carries headline statistics only (Section~\ref{sec:pathb}); the Danish counterfactual specifications are collected in Table~\ref{tab:danish}. Dashes denote quantities that are not applicable to the row. The benchmark row's ``30+15yr'' scope is not symmetric across estimators: 15-year holdings enter the empirical back-out and the ABM through real weights and real scheduled amortization, with voluntary CPR for those cohorts drawn from the same-coupon 30-year surface (the structural-only fold-in of Appendix~\ref{sec:robustness-pipeline}), while the hazard rows are 30-year-collateral estimates from which 15-year face is excluded (Table~\ref{tab:composition}, Term row). Peak-lag entries select the maximum-magnitude cross-correlation (max $|r|$) within a $\pm$3-month window; under a most-positive-$r$ convention the production ABM's best alignment is $+0.19$ at lag $-3$ (Section~\ref{sec:abm-results}); the two conventions summarize the same lag profile, and Path B's lag $-3$ peak satisfies both. \quad Path B is the only estimator whose aggregate recovery is close to the benchmark-consistent basis (97.9\%, and 100.0\% composed as committed --- 98.1\% converted to a single coupon convention --- at the in-sample calibration point; 91.3\% at the off-window floor the paper headlines). Its raw lag-0 correlation ($+0.190$) is carried by the window's shared trend: in first differences ($-0.23$) and after linear detrending ($-0.20$) the simulated--empirical co-movement is indistinguishable from zero ($\pm 0.31$ at $n = 41$), as it is for every estimator. The two cross-design rows substitute real Freddie Mac structural covariates into the ABM under two calibration choices; neither reproduces the empirical path's timing (no estimator does; see above), but the recalibrated variant's dollar recovery exceeds my own ex-ante threshold for undercutting the paradigm interpretation (Sections~\ref{sec:robustness-crossdesign} and~\ref{sec:conclusion}). Moving-block bootstrap 95\% intervals (block length 6): ABM lag-0 $[-0.60, +0.04]$; Path A lag-0 $[-0.68, -0.02]$; Path B lag-0 $[-0.13, +0.44]$, peak lag $[-0.21, +0.49]$ (no positive lag-0 coefficient is distinguishable from zero; block-length sensitivity at 4/6/8 changes no zero-inclusion verdict in this table except Path A's marginally negative levels endpoint, detailed in Section~\ref{sec:pathb}). Path A's contemporaneous levels correlation is negative on the levels construction, which strengthens the offsetting-errors reading of its aggregate fit (detrended, all three estimators cluster at $-0.20$ to $-0.30$; Section~\ref{sec:pathb}), and no timing evidence is claimed for any estimator. Table~\ref{tab:theil} decomposes each estimator's monthly error into bias, variance, and covariance shares and reports cumulative dollar-error profiles; it is a fit diagnostic and changes no verdict in this table. \quad Correlation-column convention: the ABM row's correlations are computed against the ABM pipeline's own empirical back-out (cohort-weighted, term-aware scheduled amortization), while the hazard rows use the hazard-side back-out; against the common hazard-side series the ABM lag-0 value is $-0.316$, a third-decimal, no-verdict difference. Section~\ref{sec:method-benchmark} notes the two constructions. \quad Shared-accounting-basis equivalents, together with the full-book and composed bases, are collected in Table~\ref{tab:bases} (basis defined in Section~\ref{sec:method-benchmark}, derived in Section~\ref{sec:robustness-hybrid}; the common curtailment netting is a flat 9.1pp for every U.S. hazard leg). Path A's point estimate carries no useful forward precision and is excluded from headline figures; its refit-and-resimulate interval (since rerun jointly at both specifications; the joint interval is floor-censored and reported beside, not in place of, the committed one), and the spec v3 provenance of every Path A inferential object, are reported in Section~\ref{sec:patha}.\par}

\begin{table}[H]
\centering
\footnotesize
\begin{threeparttable}
\caption{Recovery of the \$764.748 billion benchmark by accounting basis. Every U.S. hazard leg is quoted on one of these bases: the shared basis nets the common \$69.56 billion curtailment flow from the standalone-scorer figure (a flat 9.1 percentage points), and the composed basis applies the same netting to the book-coupon-reweighted (full-book) figure. The lock-in marginal (central minus $\beta_1 = 0$ null) is basis-invariant because the netting cancels.}
\label{tab:bases}
\begin{tabular}{@{}lrrrrr@{}}
\toprule
Hazard leg & Standalone & Shared & Full-book & Composed & Cum.\ runoff \\
 & (scorer) & ($-9.1$pp) & (book WAC) & ($-9.1$pp) & error (\$B) \\
\midrule
Path B, central (in-sample calibration) & 107.0\% & 97.9\% & 109.1\% & 100.0\% & $+53.8$ \\
Path B, $\beta_1 = 0$ null (in-sample calibration) & 97.8\% & 88.7\% & -- & -- & $-16.6$ \\
\addlinespace
Path B, central (off-window floor) & 100.4\% & 91.3\% & -- & -- & $+2.8$ \\
Path B, $\beta_1 = 0$ null (off-window floor) & 94.8\% & 85.7\% & -- & -- & $-39.8$ \\
\addlinespace
Path A (spec v4) & 121.5\% & 112.4\% & 127.5\% & 118.4\% & $+164.1$ \\
Path B, concave-gap & 105.7\% & 96.6\% & -- & -- & $+43.7$ \\
\bottomrule
\end{tabular}
\end{threeparttable}
\end{table}

\noindent{\footnotesize \emph{Notes to Table~\ref{tab:bases}.} The book-WAC and composed columns are the committed mixed-basis measurements (note rates bucketed against pass-through coupons); converted to a single pass-through convention they read 107.2\% and 98.1\% for Path B central, with the identified marginal bit-invariant (run \texttt{coupon\_convention\_reweight}). \quad A prior fact frames the whole table. The engine renormalises the simulated pool to realized Fed holdings every month, so a leg's dollar roll-off is a \emph{rate} applied to an exogenous level and its balance path cannot drift from the realized one. No recovery percentage here is a balance-path fit, and none can compound error across months: each is close to a restatement of that leg's window-mean CPR against the empirical 5.14\%. That is why the sampling intervals are as tight as they are, and it is a limit on what any of these levels can corroborate. The final column is each leg's terminal cumulative runoff error, realized minus simulated roll-off (equivalently, simulated minus realized trapped liquidity), and it is the sign-transparent statement of the same information: recovery is an affine transform of it, $\text{recovery} = 1 + \text{error}/764.7$, so a large additive constant sits inside every percentage in the first four columns and makes all of them look better. Two things the percentages hide. Path~B's in-sample 107.0\% is an \emph{over}-prediction of trapped liquidity of \$53.8 billion, which is an 8.2\% \emph{under}-prediction of the \$652.8 billion of runoff actually realized---the referent flips even though the map does not. And the headline off-window central leg, at 91.3\% shared, carries the smallest error of any leg here on the standalone scorer: $+\$2.8$ billion, 0.4\% of realized runoff. Errors as a share of realized runoff: $+8.2\%$, $-2.5\%$, $+0.4\%$, $-6.1\%$, $+25.1\%$, $+6.7\%$ down the rows. The lock-in marginal is unaffected by any of this, being a difference of two trapped aggregates from which the constant cancels; that is the same fact as its basis-invariance. Standalone and shared are related by the flat 9.1-percentage-point curtailment wedge of Section~\ref{sec:method-benchmark} (\$69.56 billion over the window), derived in Section~\ref{sec:robustness-hybrid}. Path B's central elasticity band is 105.9--108.2\% standalone (96.8--99.1\% shared); the covariate-prior band is 92.2\%--99.2\% shared (approximately 94.3\%--101.3\% composed). The lock-in marginal is $+9.2$ points ($+\$70.3$ billion) on the U.S. leg at the in-sample calibration point, and $+5.6$ points ($+\$42.6$ billion) at the off-window floor (Section~\ref{sec:robustness-floor}); both are basis-invariant because the netting is common to the central and null legs. The off-window rows carry the same flat netting: the layer is a constant curtailment flow over a constant cap benchmark, identical across all five committed estimators spanning 29.6 points of standalone recovery, so it is path-invariant and therefore floor-invariant (\texttt{calibration\_reconciliation}). Full-book and composed columns were run only at the in-sample calibration and are not extrapolated to the off-window floor. Path A, the concave-gap variant, the elasticity band, and the covariate-prior band are all in-sample-calibration quantities. The Danish leg's counterfactual image, run at the production floor and so comparable to the in-sample point, is $+\$61.2$ billion, or 8.0 points, the difference reflecting the Danish leg's dynamic-balance compounding (Section~\ref{sec:abm-danish}). The ABM rows of Table~\ref{tab:estimators} run through the accounting layer natively (Figure~\ref{fig:recovery}) and are not restated here.\par}

The two hazard paths, despite sharing no household-level decision machinery with the ABM, together land closer to the empirical benchmark (97.9\% and 112.4\% of it on the benchmark-consistent shared accounting basis of Section~\ref{sec:robustness-hybrid}, 107.0\% and 121.5\% under their standalone scorers, both at the in-sample calibration point; Path B is 91.3\% shared at the off-window floor the paper headlines, and Path A was not re-run there) than any synthetic-population ABM specification tested in this paper, including the falsified burnout attempt. Overshooting less, however, is not validation: Path A overshoots the benchmark by \$164.1 billion---its terminal cumulative runoff error, and a 25.1\% under-prediction of the \$652.8 billion of runoff actually realized---with a monthly path fit no better than the ABM's, and Path B's recovery, reported across its full elasticity sensitivity band (\$809.9--\$827.7 billion, 105.9\%--108.2\%; Section~\ref{sec:pathb}), reflects an elasticity imported unmodified from \citet{liebersohn2024} rather than one fit to this paper's own benchmark, an import whose verification content Section~\ref{sec:pathb} limits to the marginal's bounded interval. The comparative claim I defend is therefore narrower than ``the hazard framework fits'', and narrower again than ``lands within 2.1\% of the benchmark'': that 2.1-point miss belongs to the in-sample calibration point, and quoting it beside a marginal measured at the off-window floor mixes two calibrations. At the calibration this paper headlines Path B (a loan-level survival structure with no household-utility content) recovers 91.3\% of the aggregate benchmark on the benchmark-consistent accounting basis, a miss of 8.7 points, and 88.2--93.7\% across the off-window floor's defensible range, a miss of 6.3 to 11.8 points; the 2.1-point miss (97.9\%; the composed companion is basis-dependent --- 100.0\% as committed on the mixed coupon basis, 98.1\% converted to a single pass-through convention, Section~\ref{sec:robustness-hybrid}) is the in-sample figure, retained as the model-fit diagnostic it is, with the ex-ante covariate priors spanning 92.2\%--99.2\% and approximately 94.3--101.3\% composed at that same in-sample calibration and not re-run off-window. What survives the reconciliation is an order-of-magnitude comparative statement rather than a precision one. Path A, whose forward simulation carries no useful propagated precision and whose own rate-gap coefficient is not statistically distinguishable from zero under any bias-respecting construction (Section~\ref{sec:patha}), lands at 112.4\% under a specification that shares data but no estimated parameters. I read that coincidence as order-of-magnitude agreement and nothing stronger: the same resampling pass puts Path A's mean at 78.3\% of the benchmark with a percentile interval spanning zero, so it cannot corroborate a sign. No synthetic-population household-choice specification I tested lands within 45\% of it. I attach no timing credential to this: Path B's raw $+0.190$ lag-0 correlation is carried by the window's shared trend (in first differences and after detrending the co-movement is indistinguishable from zero; Section~\ref{sec:pathb}), and the moving-block intervals of Table~\ref{tab:estimators} caution against weighing any single timing coefficient, including Path B's lag $-3$ peak, whose interval spans zero, and the benchmark side carries its own disqualification: its four exact-zero months are clip-induced boundary artifacts of the reporting series (Appendix~\ref{sec:robustness-seasonalfloor}, run \texttt{h1\_zero\_months\_diagnosis}). Timing alone, however, does not identify the lock-in channel: the $\beta_1 = 0$ null also peaks at lag $-3$ while still recovering 88.7\% of the benchmark on the shared basis at the in-sample calibration point and 85.7\% at the off-window floor (Table~\ref{tab:bases}), both under the production floor form; under the additive form the null recovers 35.6\% at that floor (Section~\ref{sec:robustness-floor}). The elasticity's identified contribution is its marginal recovery over that null, defined and bounded once in Section~\ref{sec:identification}; the design identifies the marginal, not the lag structure. At the in-sample calibration point, across 10 paired reruns with different draw seeds (central and null sharing each seed), the marginal is \$70.3 billion and varies by less than \$0.01 billion, because the production engine drains prepayments fractionally and the draw seed touches only the small delinquency channel; the uncertainty that matters for the marginal is the calibration band and the structural assumptions, not simulation noise. I restrict this claim to synthetic-population specifications: Section~\ref{sec:robustness-crossdesign} reports a real-covariate cross-design variant of the ABM that does land within 45\% of the benchmark (59.3\%, and 76.3\% once that covariate population is reweighted to the book's own coupon $\times$ vintage composition, which is within 24\%), a result that complicates rather than confirms the paradigm interpretation this paragraph would otherwise support---and the reweighted reading makes the complication larger, not smaller.

Because the hazard framework's only behavioral content is loan-level covariates, a cohort-level burnout term, and literature-calibrated elasticities (with no analogue to mobility desire, loss aversion, or utility maximization), its comparative success might be read as evidence for the loan-level and institutional mechanisms (heterogeneity in FICO, LTV, and geography; servicer forbearance and pipeline delay; and cohort-level burnout) that the synthetic-population ABM of Section~\ref{sec:abm} does not represent. The paper's own measurements cap that reading: the named mechanisms ablate to a point or less each (burnout 0.86 points at the in-sample floor and 0.45 at the off-window floor, FICO/LTV 1.3 points, the delinquency pipeline \$0.39 billion, with no servicer channel beyond that pipeline modeled), the synthetic companion of Section~\ref{sec:robustness-synthesis} recovers a comparable aggregate level with no real loan-level structure at all, and the data-source asymmetry of Section~\ref{sec:robustness-scale} remains unresolved---so what the level supports is the mechanical decomposition, with the named mechanisms bounded small and remaining candidates only for the residual. The scale-sensitivity replication in Section~\ref{sec:robustness-scale} rules out the ABM/hazard sample-size difference as the driver of this gap; the data-source difference is addressed only partially by the cross-design test in Section~\ref{sec:robustness-crossdesign}, which finds that endowing the ABM with real structural covariates changes its explanatory share substantially: enough, under one calibration choice, to raise the question of whether ``cannot represent'' is the right description at all. The rate-gap coefficient's sign stability across both hazard specifications does not indicate that the underlying par-payoff penalty is doing real work, and I restate it. Only Path A estimates that coefficient; Path B imports it from \citet{liebersohn2024} as a signed elasticity band, so it cannot come out the other way, and an agreement between one estimate and one import that has no freedom to disagree is an orientation check on the two specifications' unit and sign conventions, not corroboration. It is worth running---it is what would catch a transposed sign between the paths, and it passes---but the evidence that the penalty has a real effect is the external literature; the estimated Path A coefficient is directionally consistent without being statistically distinguishable from zero under bias-respecting constructions (Section~\ref{sec:patha}). Household rationality is not irrelevant, but the aggregation structure of loan-level survival is what a model of this shortfall must get right at minimum, and Section~\ref{sec:robustness-crossdesign} leaves open how much more it must get right. This sharpens the trade-off named in this paper's title, with the qualification Section~\ref{sec:robustness-crossdesign} adds: the price of the TBA market's homogeneity appears to be paid at least in part through loan-level and institutional frictions that the synthetic-population household model does not represent and that a market-value repurchase design, as Section~\ref{sec:abm-danish}'s counterfactual suggests, would substantially relieve.

\section{Discussion}\label{sec:discussion}

Sections~\ref{sec:abm} and~\ref{sec:hazard} presented the paper's central contrast: household simulation, rational or behaviorally extended, does not close the gap to the empirical benchmark, while a structurally distinct hazard framework (discarding household utility altogether, on a substantially larger and non-synthetic sample) recovers 91.3\% of it on the benchmark-consistent accounting basis in its microsimulation form at the off-window floor this paper headlines (88.2--93.7\% across that floor's range, a miss of 6.3 to 11.8 points), and 97.9\% at the in-sample calibration point, with the caveat that no estimator's monthly path co-movement, Path B's included, survives removal of the window's shared trend (Section~\ref{sec:pathb}). The contrast is easy to over-read. The $\beta_1 = 0$ null recovers 85.7\% on the same basis at the headline floor (88.7\% in-sample), so the components common to any prepayment model (scheduled amortization, the seasoning baseline, and the turnover floor) deliver the large majority of the level; the caps sat far above what any plausible prepayment environment would have produced. The named loan-level mechanisms are measured small wherever I can ablate them: switching burnout off moves Path B by \$6.58 billion ($0.86$ points) at the in-sample floor and \$3.48 billion ($0.45$ points) at the off-window floor this paragraph's other figures use, zeroing the FICO/LTV covariates moves it $1.3$ points, and the delinquency transition matrix is bounded at \$0.39 billion (Appendix~\ref{app:params}); no servicer channel beyond that pipeline is modeled at all. Three mechanisms therefore remain \emph{candidates} for the residual---modest at the in-sample calibration (2.1 points) but 8.7 points, roughly 1.6$\times$ the headline marginal itself, at the off-window calibration the paper headlines---and for the timing failure every estimator shares, not demonstrated drivers of the level; servicer behavior beyond the delinquency pipeline, updated-equity dynamics, and home-equity extraction carry no bound of any kind here, and each is plausibly rate-cycle-correlated. First, loan-level heterogeneity in credit score, geography, and life circumstance means that mobility decisions are not a smooth function of coupon and market rate alone; the hazard framework's 295 stratum fixed effects are one mechanism that lets Path A track pool-level heterogeneity the ABM's coupon-cohort buckets cannot represent. Second, servicer behavior (loss mitigation, forbearance, and loan modification practices that vary by servicer and were unusually active during the higher-rate, higher-delinquency-risk period of the QT window) operates on prepayment and default outcomes through channels outside a borrower-side decision model; Path B's explicit competing-risks structure and Markov delinquency pipeline are a direct, if partial, representation of this channel. Third, prepayment burnout is representable at the cohort level (as in both hazard paths) but not within a closed, per-agent household population, as Section~\ref{sec:abm-results}'s falsified survivor-selection experiment showed.

\subsection{The Dual Systemic Fallout}\label{sec:discussion-fallout}

While the macroeconomic literature frequently focuses on the housing mobility implications of the lock-in effect, its transmission into the fixed-income sector creates a two-pronged systemic risk.

The first channel runs through the central bank itself: the extension of MBS durations deprives the SOMA of its scheduled cash flows, breaking the predictability required for effective Quantitative Tightening; Appendix~\ref{sec:robustness-wal} expresses this extension in approximate weighted-average-life units. The \$764.7 billion shortfall-versus-cap is, per the decomposition above, mostly mechanical: a gap to a cap schedule that no plausible prepayment environment would have met. The duration extension itself (a 9.4-year portfolio WAL under the empirical path against 3.4 years at 2021 realized speeds, Appendix~\ref{sec:robustness-wal}, where the 3.4-year comparator is a refinancing-boom baseline) is the rate-driven object, and this extension, not the cap arithmetic, feeds the fragility channels below.

The second channel runs through private banks, which, unlike the Fed, must mark their portfolios to market. The structural duration extension of these low-coupon MBS assets generated hundreds of billions in unrecognized losses across the banking sector. \citet{jiang2023} show that duration extension alone was not sufficient to precipitate failure; rather, it was the interaction of mark-to-market losses with a high share of uninsured deposits that made specific regional institutions vulnerable to self-fulfilling runs, as uninsured depositors withdrew funds upon recognizing that reported asset values no longer covered liabilities. The duration extension that is a necessary precondition for this fragility channel is generated by the below-market coupon stock meeting rate normalization under any payoff rule---the rule-only counterfactual of Section~\ref{sec:abm-danish} shortens the portfolio by only 0.6 of the roughly six extension-years (Appendix~\ref{sec:robustness-wal})---so the elasticity channel this paper identifies is a second-order contributor to that precondition, and neither it nor the extension is, on its own, a sufficient condition for bank failure. What this paper prices on the institutional side is therefore cash-flow timing alone---principal arriving later than the cap schedule assumed---and not a market-value loss. No mark-to-market or DV01 image of that 0.6 of an extension-year is put on the portfolio, because Appendix~\ref{sec:robustness-wal} declines option-adjusted duration and convexity as out of scope, so the bank-side losses \citet{jiang2023} price have no counterpart here and the fragility channel is invoked as context rather than as something this design measures.

Beyond its transmission into fixed-income markets, the lock-in effect imposes a real economic cost by preventing households from relocating toward higher-productivity labor markets. \citet{hsieh2019} estimate that housing supply constraints in high-productivity U.S. cities impose substantial aggregate output losses through this kind of spatial misallocation of labor; their headline magnitudes are cumulative level effects over 1964--2009 and were revised in a published correction, so I cite the mechanism rather than a point estimate. The per-move surplus that would turn this leg into a magnitude is not quantified anywhere in this design: no per-household or per-move welfare estimate enters the paper, so the labor-reallocation leg of the mobility argument is directional only. While their analysis concerns supply-side constraints on housing availability, the demand-side mobility constraint I quantify here operates through the same channel: households who cannot economically justify relocating cannot respond to geographic differences in labor demand, regardless of housing supply elsewhere. I do not translate the \$764.7 billion shortfall into welfare terms: it is an accounting gap against a cap schedule, mostly mechanical by the decomposition above, and neither a lower nor an upper bound on the welfare cost of the par-payoff rule. The welfare-relevant objects are the constrained moves themselves, whose fiscal image is the lock-in marginal ($+\$42.6$ billion of trapped roll-off at the headline off-window floor, $+\$70.3$ billion at the in-sample calibration point), and the forgone labor-reallocation gains this literature identifies, which my accounting framework does not measure.

\subsection{Redemption-Cap Design: The Ex-Ante Arithmetic}\label{sec:discussion-capdesign}

That the caps would sit above achievable runoff is not my observation: New York Fed staff projected \$20--\$30 billion in monthly principal payments in the month the \$35 billion ceiling took effect \citep{nyfed2022}. What this paper adds is the quantification and its ex-ante construction. The New York Fed's May 2022 baseline path supplies the test: it implied agency MBS runoff of \$740.6 billion over the window against the caps' \$1{,}417.5 billion, so the \$35 billion cap was set roughly 1.7 to 1.9 times as high as the Federal Reserve's own contemporaneous projection of achievable runoff, the multiple depending on the same disclosed intra-2022 allocation choice as the anticipated share (Section~\ref{sec:method-benchmark}). The decomposition measures the same distance from the other side: both components of the $\beta_1 = 0$ null (scheduled amortization from the book's coupon, age, and term composition, and baseline turnover at any assumed floor) recover 88.7\% of the realized shortfall at the in-sample calibration point and 85.7\% at the off-window floor under the production floor form (35.6\% under the additive form), which is how far the fixed \$35 billion cap sat above the book's mechanical principal path. Table~\ref{tab:wal} prices that path in maturity units: under the empirical CPR the book's approximate weighted-average life is 9.4 years against 3.4 years at 2021 realized speeds, and the 3.4-year comparator is a refinancing-boom baseline rather than a turnover norm, so the speed that would have filled the cap is the one the window's rate configuration had already shut off.

\begin{table}[H]
\centering
\small
\begin{threeparttable}
\caption{Approximate SOMA portfolio WAL by prepayment scenario (years).}
\label{tab:wal}
\begin{tabular}{@{}lrr@{}}
\toprule
Scenario (mean CPR) & WAL, June 2022 & WAL, Nov.\ 2025 \\
\midrule
Scheduled amortization only (0\%) & 14.7 & 12.6 \\
Empirical path (5.14\%) & 9.4 & 8.5 \\
Production ABM (11.68\%) & 6.0 & 5.6 \\
Hazard Path A (3.34\%) & 10.9 & 9.7 \\
Hazard Path B (4.76\%) & 9.7 & 8.7 \\
Danish rule-only leg, production anchor (5.61\%) & 9.1 & 8.2 \\
Danish-level bracketing leg (3.39\%) & 10.8 & 9.6 \\
Baseline turnover floor, off-window headline read (4.99\%) & 9.5 & 8.6 \\
No-shock counterfactual, 2021 realized speeds (22.81\%) & 3.4 & 3.3 \\
\bottomrule
\end{tabular}
\begin{tablenotes}[flushleft]\footnotesize
\item Single-pool-per-cell approximation: vintage face shares from the SOMA CUSIP tabulation (2021: 43.9\%; 2020: 16.3\%; 2017--19: 6.0\%; 2022: 23.1\%; pre-2017: 10.6\%), a 90.7\%/9.1\% 30-year/15-year term split, coupon at the book's 2.49\% WAC, and each scenario's window-mean CPR held constant. Level-payment amortization with monthly SMM; November 2025 values age the same cells 42 months. Cell WAL formula: for remaining term $n$ months, $\mathrm{WAL} = \sum_{t=1}^{n} t\, p_t / (12 \sum_{t} p_t)$, with $p_t$ month-$t$ principal (scheduled level-payment principal plus $\mathrm{SMM} = 1-(1-\mathrm{CPR})^{1/12}$ of the surviving balance). Worked anchors for the scheduled-only row: a 2.49\% 30-year pool has origination WAL 16.9 years (the 2021-vintage cell, aged 12 months at June 2022, is at 16.3); the 9.1\% 15-year sleeve (origination WAL 8.0 years) and the older vintages (32.9\% of face is 2020 or earlier) pull the book blend to 14.7. Representative cell ages ($\{2022{:}\,2;\ 2021{:}\,12;\ 2020{:}\,32;\ 2017\text{--}19{:}\,60;\ \text{pre-}2017{:}\,120\}$ months at June 2022) are single-pool stand-ins recovered by matching the printed table (a reconstruction of the original computation's assumptions, disclosed as such), and anchored externally: ages derived from the SOMA CUSIP tabulation's own back-derived origination dates (maturity minus term, face-weighted within cells, computed with no reference to this table) run $\{2021{:}\,12;\ 2020{:}\,28;\ 2017\text{--}19{:}\,50\}$ months and give a scheduled-only blend of 14.9 years; naive vintage midpoints likewise give 14.9. The printed 14.7 sits within 0.3 years of both alternatives. Coverage reconciliation: the 90.7\%/9.1\% split (which sums to the 99.8\% covered) is the audit parse of the \$1{,}940.9 billion book, with 0.2\% of face in other terms never modeled; the 90.6\% of Appendix~\ref{sec:robustness-pipeline} is the production run's own as-of-date 30-year-only coverage share; the two parses differ by as-of date and their components round independently.
\item Turnover benchmark: the final row is a refinancing-boom baseline rather than a normal-turnover one --- 22.81\% is the 2021 realized speed, most of which is refinancing, and refinancing is what the QT window's rate configuration shut off. The turnover-floor row is the normal-turnover counterpart. The paper's own out-of-window turnover measurement is the floor read (Table~\ref{tab:oosfloor}): defensible 2018 reads of 4.70\%, 4.99\% and 5.33\% by depth of out-of-the-moneyness, against refi-contaminated pooled 2017--2019 and 2019 anchors at 6.07\% and 6.91\%. The turnover-floor row prices the middle read --- the paper's headline floor --- through the identical calculator (run \texttt{wal\_normal\_turnover}): it falls between Path B's 4.76\% and the empirical 5.14\% rows, and monotonicity in CPR puts the band ends on either side of it.
\end{tablenotes}
\end{threeparttable}
\end{table}

The design question that follows is speculative, and I flag it as design space rather than as a finding these estimates support. Whether this episode generalizes is a question this design cannot answer but can pose precisely. \citet{du2024}'s cross-country record covers seven central banks running active and passive runoff, and the mechanism identified here---a fixed nominal cap set above the mechanical principal path of a book whose coupons sit below market---is not specific to the Federal Reserve. What would settle it is narrow and public: each bank's announced cap or redemption schedule set against its realized redemptions, month by month, over its own tightening window. I do not assert that the other six were non-binding; I record that the comparison is constructible from published data and is not made here. Pre-committing a substitution instrument---a standing policy of substituting Treasury runoff or outright MBS sales when passive runoff undershoots---would convert extension risk from a silent shortfall into an explicit, priced policy choice, and it is the option consistent with the objective just stated, because it leaves the cap free to communicate the path instead of forcing it. Indexing the cap to the observable state of the outstanding stock (for example, to the share of the book sitting deeply below market coupon) is the weaker alternative, since it still works through the instrument the shortfall defeats. What the decomposition supplies such a design is computable in advance of the cycle, but as a band rather than as a point: scheduled amortization is fixed by the book's composition, while the turnover floor is measured, and its defensible out-of-window reads run from 4.70\% to 5.33\% depending on the depth at which out-of-the-moneyness is cut, carrying the central-elasticity marginal from $+6.8$ points at the deepest cut to $+4.3$ at the shallowest (Table~\ref{tab:oosfloor}); the floor read's own sampling error widens the identified margin further, and expressed in the units a cap designer would have to plug in, the mid-cut read's wild-cluster interval runs from 4.177\% to 5.800\% (Section~\ref{sec:robustness-floor}). What the cap is for bears on whether binding is desirable at all. A redemption cap paces and communicates the portfolio's composition path, and Section~\ref{sec:method-benchmark} concedes that these caps never bound over this window; choosing a level that would bind therefore trades the communication objective for a speed the schedule was not built to deliver, rather than improving on it. Read on its own terms, a cap intended to bind would be set at or below the lower edge of the book's projected scheduled-plus-turnover principal path, with the identified elasticity margin ($+2.1$ to $+13.2$ points of the shortfall across the calibration box) as the behavioral band around it --- which is the form the Federal Reserve's own contemporaneous projection took, a range rather than a point. Nothing in my estimates identifies the welfare or market effects of such designs; active sales in particular would realize the mark-to-market losses that passive runoff defers. I record the option to mark where the design space lies, not to recommend a point in it.

\subsection{Evaluating Alternative Securitization Policies}\label{sec:discussion-policy}

In response to the mobility collapse, U.S. policymakers, including the Federal Housing Finance Agency (FHFA), proposed expanding assumable and portable mortgages. Fixed-income analysis suggests these proposals address only part of the problem: while portability eases consumer mobility constraints, it explicitly prevents the loan from prepaying upon relocation. The same holds for assumability, with one added wrinkle: an assumed loan also fails to prepay (the below-market coupon transfers to the buyer and the extension stays on the investor's book), and the buyer must separately finance the gap between the purchase price and the assumed balance, typically at prevailing rates, which limits take-up. Both instruments therefore relocate rather than eliminate duration extension risk, solving the microeconomic mobility constraint while cementing the macroeconomic extension risk onto the investor's balance sheet.

\subsection{The Danish Market-Value-Repurchase Alternative}\label{sec:abm-danish}

A third alternative acts on the payoff rule itself rather than on the mortgage contract's portability or the redemption cap. Section~\ref{sec:abm-results} shows that no synthetic-population household-choice specification tested here explains the majority of the trapped liquidity---a finding Section~\ref{sec:robustness-crossdesign} confines to the synthetic population, its recalibrated real-covariate variant having crossed the pre-committed threshold past which the paradigm reading cannot be cleanly maintained---so the question is how much of that trap is a consequence of the par-payoff rule itself, rather than of household behavior under any institutional regime. The Danish market-value repurchase system, which severs the link between market rates and payoff cost, gives me a direct way to ask that question quantitatively. The operative specification calibrates each framework's Danish response against \citet{berger2026}'s separately estimated moving and refinance-in-place channels---a January 2026 SSRN working draft, and the single unrefereed source this channel rests on---and transplants them onto the U.S.\ loan population so that only the payoff rule changes; it yields the paper's finding on this question, an institutional gap of $+\$61.2$ billion (Table~\ref{tab:danish}). The rest of this subsection develops that specification and its sensitivity.\footnote{\label{fn:manifest}CPR means below are aggregates over the 42 active QT months (June 2022--November 2025), and the institutional wedge is the mean over those months of the cohort-weighted Danish-minus-U.S.\ simulated CPR difference. This operative specification supersedes an initial mechanism-substitution that had implied a much larger gap; that superseded figure, and the frozen run manifests underlying every ABM headline, are catalogued in Appendix~\ref{app:ledger}.}

The recalibration rests on two behaviorally distinct channels that \citet{berger2026} (my source for the Danish institutional mechanism itself) estimate separately, and that a single extrapolated mechanism would conflate. The first is a moving channel: does removing the par-payoff penalty change how often households relocate? The second is a refinance-in-place channel: does the option to buy back debt at a market discount induce discount-buying behavior independent of moving? Danish administrative data show the two channels behave very differently under a U.S. tax transplant. The moving channel is well-identified and close to flat in the coupon gap: Danish households move at an unconditional rate of 3.2\% per year regardless of how far their coupon sits below market \citep[\S4.6, structural parameters in tab.~3]{berger2026}, versus a U.S. moving-hazard slope of 0.57--1.20 percentage points per year per 100 bp of coupon gap \citep{fonseca2024}. The refinance-in-place channel, by contrast, is substantially a Danish tax artifact: discount buybacks are attractive there because mortgage interest is deductible at a flat 33\% rate and realized capital gains on the buyback are tax-exempt. Under U.S. tax treatment (a 22\% mortgage-interest deduction and a 15\% capital-gains tax on the discount, per Berger et al.'s own counterfactual --- a characterization I read differently, though I do not adjudicate it: on the ordinary reading a discount extinguishment is cancellation-of-debt income at ordinary rates, IRC \S 61(a)(11) and Treas.\ Reg.\ \S 1.61-12(c), while \S 108(a)(1)(B) and \S 108(a)(1)(E) exclude it for insolvent borrowers and for qualified principal-residence indebtedness respectively, so the treatment is borrower-state-dependent rather than uniform in either direction; post-TCJA the assumed interest deduction also overstates most households' marginal benefit. Read as a directional bracket rather than a correction, both considerations push the U.S.-transplant channel \emph{below} Berger et al.'s already-small estimate, and neither touches the realized Danish evidence, which is measured under Danish tax), this channel loses its shield, and their general-equilibrium counterfactual \citep[\S4.9.1]{berger2026} finds it shifts the equilibrium mortgage rate by only about one basis point, economically negligible.

Recalibrating both frameworks against these two estimated channels, rather than against a single extrapolated mechanism, first requires a choice: what to do with the moving channel's \emph{level}. Berger et al.'s identified fact about the moving margin is one of slope: Danish moving is flat in the coupon gap under market-value payoff; the 3.2\% per year is Denmark's descriptive unconditional moving rate, not an estimate of how often U.S. households would move under a Danish payoff rule. Importing that level onto the U.S. loan population (the \emph{Danish-level} anchor) therefore bundles two changes, the payoff rule and Danish baseline mobility itself, and produces Danish legs whose mean CPR (3.39--3.40\%, Table~\ref{tab:danish}) sits \emph{below} the 4\% turnover floor both frameworks impose on the U.S. side (Appendix~\ref{sec:method-frictions} and Section~\ref{sec:pathb}). Death, divorce, and forced relocation do not become rarer because the payoff rule changed, so a rule-only counterfactual on the same U.S. households cannot prepay below that floor. The construction that isolates the rule is the \emph{U.S.-intercept} anchor: keep the identified flatness, but anchor the level at the production U.S. hazard evaluated at zero rate gap, with the turnover floor, seasoning baseline, burnout, and covariate terms all retained and only the rate-gap suppression switched off, plus the same U.S.-transplant refinance-in-place channel. I adopt the U.S.-intercept anchor as the production reading (it is the only one consistent with the floor's stated semantics) and retain the Danish-level anchor as a labeled bracketing case in which U.S. households are assumed to adopt Danish baseline mobility along with the Danish rule. Holding the floor fixed while importing only the Danish moving slope treats baseline turnover and rate-elastic moving as separable channels; that separability is motivated by the floor's semantics rather than identified in the data available here, which is why the Danish-level anchor, importing the two jointly, is retained as the bracketing case.

The anchor choice determines the sign of the institutional gap. Under the shared-accounting hybrid pipeline (specified in Section~\ref{sec:robustness-hybrid}: the ABM's macro-accounting layer with Path B's microsimulated CPR as the micro-foundation), the U.S.-intercept anchor yields a Danish mean CPR of 5.61\% against the locked-in U.S. leg's 4.76\%, and an institutional gap of $+\$61.2$ billion (8.0\% of the benchmark) at Berger et al.'s general-equilibrium best estimate of a negligible refinance-in-place contribution (the U.S. leg reproduces the committed production artifacts exactly). That estimate is a price statement --- an equilibrium-rate effect of about one basis point --- and reading it as a near-zero buyback \emph{volume} is an inference realized Danish data contradict: over 2022M07--2023M12, deep-discount ($\leq 2\%$ coupon) Danish callable bonds redeemed at 26.4\% per year against a 4--5\% baseline, an extraordinary component of at least 22\% per year after a generous scheduled-amortization allowance (run \texttt{danish\_redemption\_validation}; Danmarks Nationalbank securities statistics, exact stock--transaction closure on the coupon-blocked universe). The transplant imports the payoff rule alone: it does not import the Balance Principle's series-level match funding, the tap-issued joint-and-several bond series that carry Denmark's liquidity, the advisory distribution that prompts conversions, or the Danish product mix --- in particular the large interest-only share, which would cut scheduled amortization, the null's largest component, and so move the shortfall in the opposite direction from the payoff rule. Each omission is a held-fixed feature of the U.S. leg, not a measured invariance. I therefore report the rule-only gap as a band over the refinance-in-place contribution rather than at the zero point: swept in 0.25\%-CPR steps under the production anchor, the gap runs from $+\$61.2$ billion at 0\% to $+\$256.8$ billion at 3\% (run \texttt{danish\_refi\_finegrid}; slope $\approx \$65$ billion per CPR point). The realized-implied anchor sits above the swept range entirely, pricing at $+\$1{,}212$ billion beyond the rejected partial-equilibrium ceiling. Four confounds travel with it, and together they are why it is an upper anchor on the transplant rather than an estimate of it: the realized rate is gross, so conversions bundle in; it includes non-household collateral; it is environment- rather than book-matched; and it is a Danish-tax realization of an option whose U.S.-transplant value the tax attenuation is meant to shrink. This is a return to a construction the pipeline already contained: the U.S.-intercept hybrid coincides numerically, leg for leg, with the NPV-identity hybrid the replication package's technical documentation records from before the Berger recalibration, because the buyback-at-market-value identity zeroes the effective rate gap, which is computationally the same object as the zero-gap U.S. hazard; what that earlier construction lacked, and Berger et al.'s moving-flatness estimate now supplies, is empirical evidence that Danish-rule moving is flat in the gap. The sign is invariant under this anchor, and it could not have come out otherwise. The specification defines the Danish moving hazard as the production U.S. hazard evaluated at a zero rate gap; the U.S. book carries a negative gap in essentially every cell over essentially the whole window (99.53\% of exposure below the window-minimum market rate, Section~\ref{sec:identification}); and the hazard is monotone in the gap. The zero-gap leg therefore must prepay faster, the Danish trapped balance must come in lower (\$687.8 billion against \$749.0 billion), and the gap must be positive. Sweeping the refinance-in-place input from 0\% to the rejected partial-equilibrium ceiling ($\approx$18\% CPR) only adds Danish prepayment on top of that, so the sweep's sign-robustness---positive at every point, rising to $+\$1{,}038.9$ billion at the ceiling (Figure~\ref{fig:gapsweep})---is forced by construction as well. I retain both as wiring checks and read neither as evidence. That premise is an inequality about this episode, not a property of the design. It holds because the simulated book is closed at 2017--2021 originations and the window's market rate stood above essentially all of its coupons in essentially every month; where a book's rate gap is mixed in sign---another cycle, a portfolio open to new originations at the prevailing rate---the inequality fails and the sign becomes a composition question this design does not answer, since the elasticity is imported rather than estimated. Neither this sign nor the marginal transports without re-checking that inequality. What is not forced is the magnitude: \$61.2 billion, 8.0\% of the benchmark, is the quantity the recalibration actually delivers, and the Danish-level anchor below shows that a different and defensible transplant choice reverses the sign outright. This positive gap is the lock-in marginal of Section~\ref{sec:hazard-interp} seen from the counterfactual side: switching off the rate-gap suppression is what the $\beta_1 = 0$ null does, so the rule-only Danish counterfactual and the paper's lock-in contribution at the in-sample calibration point are one object measured on two accounting legs ($+\$61.2$ billion here, where the Danish leg compounds its own counterfactual balances; $+\$70.3$ billion as the U.S.-leg marginal). At the 0\% edge of the swept band the two nearly coincide; under the band presentation that near-equality is a property of that edge and of the in-sample calibration on which both legs were run, not of the counterfactual. Being one object on a second accounting leg also means the Danish magnitude inherits every conditionality the marginal carries---the baseline seasoning convention and the floor's age transport among them---so the transplant supplies no independent information about the payoff rule. It re-denominates a quantity this paper has already conditioned, and it is conditional in exactly the same ways. That near-equality is a statement about the demoted in-sample point and not about the paper's headline: against the $+\$42.6$ billion off-window marginal the Danish rule-only gap is about 44\% larger, not equal, and both legs were run at the production floor; re-run at the off-window floor itself (run \texttt{danish\_offwindow\_floor}; the U.S. leg reproduces the committed off-window central exactly), the rule-only gap is $+\$28.2$ billion, 3.7\% of the benchmark, 34\% below the off-window marginal, so the near-equality is specific to the in-sample calibration and no equality survives at matched floors. Under the Danish-level anchor, by contrast, the gap is $-\$99.9$ billion (U.S. 4.76\%, Danish 3.39\% mean CPR) and crosses zero at a refinance-in-place contribution of 1.4\% CPR; the fragility reflects the bundled mobility-culture change. The ABM's own recalibrated run under the Danish-level anchor shows $-\$728.4$ billion (crossing at 12.6\%), but its U.S.-side leg is over-predicted relative to the empirical benchmark (11.76\% simulated mean CPR under the Berger run's own calibration-stage freeze, versus 5.14\% empirical; cf.\ the production tag's 11.68\%, Section~\ref{sec:abm-results}), and under the U.S.-intercept anchor the ABM's direction follows a fortiori without a run: the ABM's lock-in penalty only ever suppresses moving, so its zero-gap intercept exceeds its locked-in U.S. leg month by month and its rule-only gap is positive by construction. Table~\ref{tab:danish} collects all four specifications.

\begin{figure}[H]
\centering
\includegraphics[width=\textwidth]{fig4_institutional_gap_sensitivity}
\caption{Sensitivity of the institutional gap to the U.S.-transplant refinance-in-place contribution, swept from 0\% to the rejected partial-equilibrium ceiling ($\approx$18\% CPR), under both transplant anchors for the hybrid pipeline. Under the production U.S.-intercept anchor (rule-only counterfactual) the gap is $+\$61.2$ billion at the Berger et al.\ general-equilibrium best estimate ($\approx$0, marked) --- one anchor among the swept range, not a bound: the fine grid of run \texttt{danish\_refi\_finegrid} spans $+\$61.2$ to $+\$256.8$ billion over 0--3\% refinance-in-place CPR, and the realized-implied Danish anchor (22.4\% per year) lies above the plotted ceiling --- and positive at every sweep point; that positivity is forced by the zero-gap anchor and cannot be overturned by this sweep, since additional refinance-in-place only raises Danish prepayment further. Under the bracketing Danish-level anchor (rule plus Danish baseline mobility) the gap is $-\$99.9$ billion at the best estimate and crosses zero at a contribution of 1.4\% CPR. The anchor, not the refinance channel, decides the sign.}
\label{fig:gapsweep}
{\footnotesize\noindent Notes: The order-of-magnitude reduction against the superseded \$925.5 billion extrapolation is a property of the Berger et al.\ best estimate, not of the sweep: both anchors sit an order of magnitude below that figure at the best estimate, but the gap rises steeply in the refinance contribution and both anchors exceed it at the rejected partial-equilibrium ceiling ($+\$1{,}038.9$ billion under the U.S.-intercept anchor, $+\$972.0$ billion under the Danish-level anchor). That ceiling is rejected on partial-equilibrium grounds and is not reached by any calibration this paper adopts. The upper half of the sweep should not be read as a quantitative statement at all: at refinance contributions of 12\% and above the Danish leg's trapped balance turns negative ($-\$14.4$, $-\$158.0$ and $-\$290.0$ billion at 12, 15 and 18\%), which is not an economically meaningful quantity. Those points are retained only to show the direction of the sweep.}
\end{figure}

\begin{table}[H]
\centering
\small
\begin{threeparttable}
\caption{The Danish market-value-repurchase counterfactual under four specifications.}
\label{tab:danish}
\begin{tabular}{@{}p{3.4cm}rrr>{\raggedright\arraybackslash}p{4.0cm}@{}}
\toprule
Specification & U.S. leg & Danish leg & Gap (US$-$DK) & Mean CPR US/DK \\
\midrule
Mechanism-extrapolated ABM\tnote{a} & $+\$91.0$B & $-\$834.5$B & $+\$925.5$B & 11.68\% / 47.10\% \\
Berger ABM, Danish-level\tnote{b} & $+\$84.5$B & $+\$812.9$B & $-\$728.4$B & 11.76\% / 3.40\% \\
Berger hybrid, Danish-level\tnote{c} & $+\$749.0$B & $+\$848.9$B & $-\$99.9$B & 4.76\% / 3.39\% \\
Berger hybrid, U.S.-intercept\tnote{d} & $+\$749.0$B & $+\$687.8$B & $+\$61.2$B & 4.76\% / 5.61\% (band $+\$61.2$--$+\$256.8$B over 0--3\% refinance-in-place; run \texttt{danish\_refi\_finegrid}) \\
\bottomrule
\end{tabular}
\end{threeparttable}
\end{table}

\noindent{\footnotesize \emph{Notes to Table~\ref{tab:danish}.} Sign convention: each leg is signed trapped liquidity (the net shortfall of simulated roll-off relative to the phased redemption cap over the QT window), so positive values mean runoff fell short of the cap (liquidity trapped) and negative values mean simulated runoff exceeded the phased cap path. Gap $=$ U.S. leg $-$ Danish leg; a negative gap means the Danish regime traps more than the U.S. regime. \quad a Superseded: applies the U.S.-calibrated mobility function to a Danish-style rate gap it was never estimated to describe (see text). \quad b The Berger recalibration's U.S. leg is re-simulated under that run's own calibration-stage freeze (native 15-year gate), not the \$91.0 billion production headline, so both legs and the gap share one frozen calibration. The Danish leg prepays more slowly under the tax-adjusted elasticities, trapping more than the U.S. leg. Danish-level anchor: imports Denmark's descriptive 3.2\%/yr moving level, a rule-plus-country bundle whose Danish mean CPR sits below the 4\% turnover floor (see text). \quad c Bracketing case under the shared-accounting hybrid pipeline (Section~\ref{sec:robustness-hybrid}); both legs are printed from the shared-layer scoring artifact at the in-sample calibration point, the floor the Danish legs were run at (U.S. trapped \$749.0 billion = 97.9\% of benchmark; Danish trapped \$848.9 billion); neither leg was re-run at the off-window floor. Its sign is not robust: the sweep of the refinance-in-place input crosses zero at a 1.4\% CPR contribution. \quad d The paper's production estimate (rule-only counterfactual): Berger et al.'s moving-flatness fact anchored at the production U.S. hazard's zero-gap intercept, turnover floor retained, refinance-in-place at the general-equilibrium best estimate ($\approx$0). The U.S. leg reproduces the committed production artifacts exactly. The gap is positive at every point of the refinance-in-place sweep (Figure~\ref{fig:gapsweep}), which under this anchor is forced by construction rather than found: see Section~\ref{sec:abm-danish}. This row is face-denominated and ex-post on the U.S. coupon stack: under the cash-haircut reading its gap reverses sign, and the wedge that reverses it is a transfer from the SOMA portfolio to moving households rather than a resource cost (Section~\ref{sec:pathb}); and no ex-ante repricing of the repurchase option into the coupon is netted from either leg, an omission signed against the reported relief (see text).\par}

The initial mechanism-substitution that this specification corrects had implied a far larger gap, \$925.5 billion, by feeding a U.S.-calibrated mobility function a Danish-style rate gap it was never estimated to describe; it is an extrapolation rather than a measurement, catalogued in Appendix~\ref{app:ledger}. The case for the Danish alternative is modest and anchor-signed: roughly $+\$61$ billion, 8\% of the benchmark, at the in-sample calibration point ($+\$28$ billion, 3.7\%, at the off-window headline floor), the counterfactual image of the lock-in marginal. Its positivity throughout the refinance-in-place sweep is a property of the zero-gap anchor rather than a finding, so the magnitude is the informative part. The earlier ``sign we cannot identify'' reading survives only as a statement about the bracketing Danish-level anchor, which asks a different question (what if U.S. households also moved like Danes) and answers it with a Danish leg that violates the floor both frameworks impose on the U.S. side. This does not alter Sections~\ref{sec:abm} and~\ref{sec:hazard}'s central finding (that loan-level and institutional mechanics, not household choice, explain the U.S. empirical shortfall) since that comparison never depended on the Danish counterfactual. What it does revise is the payoff-rule question: the par-payoff rule's institutional cash-flow cost is real but an order of magnitude smaller than the original extrapolation implied, and of the same origin as the lock-in marginal the paper identifies directly---and of the same size as that marginal at the in-sample calibration point where both legs were run, though about 44\% larger than the off-window marginal the paper headlines. Both legs credit retired balance at par; under cash accounting the market-value buyback credit reverses the gap's sign, so the case as stated here stands on the benchmark's own face denomination, with the incidence bracket priced in Section~\ref{sec:pathb}. The $1-P$ wedge the U.S. rule collects from a moving household accrues to the bondholder, and here the bondholder is the SOMA portfolio, whose earnings reach the Treasury as remittances: the cash-haircut reading hands that wedge to the households who move and takes it from general taxpayers. So the reversal prices a transfer changing direction, not a resource cost, and its welfare sign turns on the marginal value of public funds against that of household liquidity---neither of which this framework measures. One further omission is a price rather than a held-fixed feature of the U.S. leg, and its direction runs against the reported relief. Every figure here is ex-post, on a coupon stack originated under the U.S. payoff rule; under a market-value payoff rule the repurchase option is written into the bond series and priced into the coupon ex ante, so the same households would have borrowed at a higher rate for the life of the loan \citep{campbell2013,berg2018}. That premium is netted nowhere, and the only ex-ante rate term the transplant carries---Berger et al.'s roughly one-basis-point equilibrium shift---prices the refinance-in-place channel alone. So $+\$61.2$ billion is gross institutional relief under face accounting, not a net social gain.

\section{Robustness Checks}\label{sec:robustness}

This section collects the methodological checks referenced throughout the paper. Two of them changed reported results rather than confirming them. The recalibrated cross-design variant of Section~\ref{sec:robustness-crossdesign} exceeds my ex-ante threshold and qualifies the household-choice reading at the aggregate level; and the off-window floor re-measurement of Section~\ref{sec:robustness-floor} is where the paper's headline marginal and level are actually derived ($+9.2 \to +5.6$ points; 97.9\% $\to$ 91.3\% shared), demoting the in-window calibration to an in-sample point. On the shared basis the paper quotes throughout, the hazard framework undershoots the benchmark at the headline calibration (91.3\%); it overshoots only on the standalone basis (107.0--121.5\%). The remaining checks confirm rather than change the central decomposition.

Four of this section's checks are mechanics rather than results and are collected in the online appendix: the behavioral-extension mechanics (Appendix~\ref{sec:robustness-extensions}), the weighted-average-life quantification (Appendix~\ref{sec:robustness-wal}), the fifteen-year fold-in and the rate-gap-units correction (Appendix~\ref{sec:robustness-pipeline}), and the seasonal-floor program (Appendix~\ref{sec:robustness-seasonalfloor}), whose verdict the calibration box already carries: what a seasonally varying floor moves is floor dispersion under the hard maximum, not cyclicality, and monthly timing earns no credential there---the frozen timing rule is cleared by a wrong-month rotation and by the $\beta_1 = 0$ null itself.

\subsection{Scale Sensitivity}\label{sec:robustness-scale}

The ABM's 10,000-household population and the hazard framework's 75,000-loan sample differ in both scale and data source. To test whether population size alone drives the gap between the two frameworks, I re-ran the ABM's core household mechanic (exponential mobility-desire heterogeneity---the distribution the production code draws, on a dollar scale---mapped through a rate-gap-driven CPR function) at N = 10,000 and N = 75,000 across 50 seeds each. The mean simulated CPR is 15.616\% at both population sizes (difference 0.0003 percentage points), while the 95\% confidence interval narrows by a factor of 2.4$\times$ at the larger $N$, of the same order as, though about 12\% short of, the 2.74$\times$ an idealized $\sqrt{N}$ law would predict ($\sqrt{7.5} \approx 2.74$), consistent with a non-i.i.d.\ component of seed noise; I do not lean on the exact ratio. Population size changes the precision of that mechanic's estimate, not its central tendency. Two limits on what the check shows. Both are large, so I state the conclusion conditionally. First, it exercised the household cost-benefit mechanic \emph{in isolation}: the production specification adds a 43\% front-end DTI wall, a discrete twenty-percent wait-and-see freeze trait, income-scaled curtailment, eleven-cohort multi-vintage weighting with a fifteen- and thirty-year term split, and a settlement-lag kernel, none of which are present here. The isolated mechanic's 15.616\% mean CPR accordingly runs 3.9 points---about a third---above the production ABM's 11.761\%, so the object whose scale-invariance is demonstrated is not the estimator whose recovery is at issue. Second, that check had no committed script, artifact or gate behind it; the figure survived from an earlier draft. Both limits are now addressed by a pre-committed rerun on the full production specification (run \texttt{production\_scale\_test}), fifty seeds at each population size, carrying the DTI wall, the freeze trait, curtailment, the eleven-cohort weighting with its term split and the settlement-lag kernel. It recovers 12.66\% of the benchmark at $N = 10{,}000$ against 12.18\% at $N = 75{,}000$---a difference of $0.48$ points against a fifty-seed standard deviation of $3.25$ points, so well inside one standard deviation---at mean simulated CPRs of 11.60\% and 11.64\%, close to the production run's 11.76\% rather than to the isolated mechanic's 15.6\%. Seed precision improves at the larger population, the standard deviation falling from $3.25$ to $1.05$ points (a ratio of 3.1 against the $\sqrt{7.5} \approx 2.74$ an idealized law would give), and again I do not lean on the exact ratio. Population size therefore does not drive the gap between this estimator and the hazard framework, and that statement now rests on the production specification rather than on a fragment of it. It does not follow, however, that the gap is attributable to modeling paradigm alone: that distinction cannot be closed by resampling, and Section~\ref{sec:robustness-crossdesign} takes it up directly with the completed cross-design pair.

\subsection{Benchmark Sensitivity to Data Source and Cap-Schedule Assumptions}\label{sec:robustness-benchmark}

The \$764.7 billion empirical benchmark (Section~\ref{sec:method-benchmark}) is constructed from the New York Fed's SOMA current-face-value series under the actual phased redemption cap schedule, netting realized roll-off against the cap within the active QT window. I treat net accumulation, rather than one-sided accumulation of shortfall months only, as the correct construction because it is the only version comparable to the Danish counterfactual of Section~\ref{sec:abm-danish}, where simulated roll-off routinely exceeds the cap; a one-sided construction would overstate the empirical benchmark relative to that counterfactual and bias the recovered percentages reported for both estimators. Researchers extending this work with alternative SOMA vintages or a different cap-schedule assumption should expect the benchmark level, and therefore the recovered percentages in Sections~\ref{sec:abm} and~\ref{sec:hazard}, to shift; the qualitative ranking of the three estimators reported in Table~\ref{tab:estimators} is not sensitive to this choice, since the ABM's shortfall against any plausible benchmark construction remains an order of magnitude smaller than either hazard path's. Section~\ref{sec:method-benchmark} separately addresses the conceptual status of the cap-relative construction itself. One such alternative is now measured rather than anticipated. The benchmark nets cash arriving in month $t$ against the cap in force in month $t$, which is what the cap constrains: the Committee set a ceiling on principal not reinvested in a given month, and that is a statement about cash. Agency remittance, however, means the cash arriving in month $t$ was generated by prepayment decisions in months $t$, $t-1$ and $t-2$, on the settlement kernel this paper already applies to simulated roll-off (0.10, 0.60, 0.30 at lags zero through two). Agency remittance is itself a deterministic lag of roughly one month, so that profile is not a description of the remittance cycle: it is an empirical smearing kernel over the months in which a given month's prepayment decisions arrive as cash, and the alternate kernels below price the choice. Convolving the cap schedule forward with that same kernel and re-netting the unchanged realized series gives a settlement-aligned benchmark of \$722.7 billion against the calendar \$764.7 billion, a shift of $-5.5$\% (run \texttt{settlement\_months\_benchmark}). The shift is mechanical, and its size is exactly accounted for: the QT window is half-open, so cap convolved past its final month is dropped---October 2025 loses its two-month tap and November 2025 its one- and two-month taps, \$42.0 billion of the \$1{,}417.5 billion schedule---and because the realized series is untouched and both legs sum over the same months, the benchmark moves by exactly that \$42.0 billion. Two alternate kernels bracketing that profile put the shift at $-\$28.0$ and $-\$52.5$ billion. I retain the calendar convention, for the reason Section~\ref{sec:method-benchmark} gives for retaining the cap at all: it is the only policy-committed quantity available, and it is committed in cash-arrival months. What the variant establishes is that re-basing moves no dollar quantity in this paper---the identified marginal is \$42.6 billion on either denominator---and moves the percentages only by the common factor $1.058$: the central Path~B level reads 96.6\% rather than 91.3\%, the $\beta_1 = 0$ null 90.7\% rather than 85.7\%, and the marginal $+5.9$ rather than $+5.6$ points. The decomposition is unchanged in substance, and the null still recovers the large majority of the benchmark under either convention. The variant also bears on the path-level timing question every estimator here leaves open. If the monthly series each estimator's path is scored against is a cash-arrival object---reported roll-off smeared over the months in which a month's prepayment decisions arrive as cash and, in its CPR form, clipped by a back-out whose non-negativity rule pins four months to exactly zero and posts the displaced mass into the following month's reading (Section~\ref{sec:method-benchmark})---then a monthly path is not a behavioral object, and the timing question may be ill-posed rather than a failure every estimator shares. I state that as a candidate account, not a finding: reporting mechanics and behavior arrive in the same series, and nothing here separates them. A related window-construction pitfall merits a reproducibility note: accumulating from the QT start date without an upper bound admits eight post-QT months (December 2025 onward). Where the dollar aggregation already applies the bounded active-QT mask, as it does throughout this paper, only the path statistics drift---the headline CPR means by roughly 0.3--0.5 percentage points; where it does not, the benchmark level itself moves. That asymmetry (path statistics silently sensitive to a window filter the masked dollar aggregates are robust to) is why every headline figure in this paper cites a single tagged run manifest (footnote~\ref{fn:manifest}).

\subsection{Cross-Design Test: Real Covariates in the ABM}\label{sec:robustness-crossdesign}

Section~\ref{sec:method-estimators} flags an unresolved confound between this paper's two estimators: they differ in modeling paradigm (household choice versus loan-level survival) and in data source (synthetic population versus real Freddie Mac loan records) simultaneously, so a low ABM share and a high hazard-framework share are equally consistent with either explanation. To make partial progress on separating them, I re-ran the ABM's decision rule with structural covariates (coupon, loan age, and original LTV) drawn directly from the same 75,000-loan Freddie Mac sample used throughout Section~\ref{sec:hazard}, in place of the log-normal, FRED-median-calibrated draws used elsewhere in this paper. Behavioral draws (mobility desire, loss aversion, patience, transaction cost) remain synthetic throughout, since the Freddie Mac data contains no analogue for these constructs; this test therefore varies real structural heterogeneity---and, until the post-stratification below, book composition jointly---not a fully real population.

Because the ABM's mobility-scale parameter is calibrated to hold a real-world baseline turnover floor fixed (Appendix~\ref{sec:method-abm}), swapping in a population with different LTV and coupon moments raises a calibration question that a naive covariate swap does not resolve on its own: should the turnover-floor anchor be re-derived against the new population's moments, or held fixed at its original value? I report both variants. The recalibrated variant applies the recalibration rule of Appendix~\ref{sec:method-frictions}, re-anchoring the mobility-scale parameter to reproduce the same 4--5\% turnover floor under the real population's moments; the frozen variant suspends that rule and holds the original synthetic-population calibration unchanged.

Before running either variant, I specified ex ante the following interpretive thresholds: a recovery above 50\% of the empirical benchmark would undercut the paradigm interpretation advanced in Sections~\ref{sec:abm}--\ref{sec:conclusion}; a recovery below 35\% would corroborate it, consistent with the synthetic-population ABM's own specifications; and a recovery between 35\% and 50\% would be ambiguous.

The recalibrated variant recovers \$453.5 billion, or 59.3\%, of the benchmark, above my ex-ante 50\% threshold. The frozen variant recovers \$159.5 billion, or 20.9\%, within the corroborating band. Both are single draws, so a pre-committed fifty-seed sweep varying both randomization layers (the behavioral draws and the covariate subsample; run \texttt{cross\_design\_seeds}) tests whether the threshold crossing is a property of the estimator or of the seed. It is the estimator: the recalibrated variant averages 60.2\% with a standard deviation of 3.3 points and a 95\% band of 55.1--66.0\%, the committed draw sits at the 48th percentile of its own distribution, and all fifty seeds classify as undercutting. The frozen variant averages 22.0\%. The band is wide enough to report---10.9 points, against the 7.7-point disclosure threshold fixed before the run---but not wide enough to reach the calibration swing, so the operative uncertainty here remains the calibration choice rather than seed noise, now measured rather than assumed. The same sweep records the deep-gap floor each variant implies, and the two variants separate cleanly on it. The recalibrated variant's CPR at an 8\% market rate averages 4.59\% and stays inside the observed 4--5\% turnover band on all fifty seeds. The frozen variant averages 7.05\% and lies outside that band on all fifty (range 6.31--7.63\%). This is the same admissibility test I apply to the externally anchored variant earlier in this subsection, and applied consistently it counts against the frozen leg, not the recalibrated one: the leg the paper reports as falling within the corroborating band is the leg that misses the turnover moment. I therefore read the 20.9\% as the less admissible of the two calibrations rather than as corroboration, which strengthens the undercutting reading rather than qualifying it. One further caveat on the word ``frozen'': re-deriving that variant's mobility-scale anchor per seed gives a mean of 43{,}493 with a standard deviation of 1{,}181 and a minimum of 39{,}984, against the committed 43{,}883, which sits at the very top of its own range. The anchor is frozen by convention, not because the data pin it there. Table~\ref{tab:crossdesign} summarizes both against the synthetic-population control, and Figure~\ref{fig:crossdesign} plots the paired comparison.

\begin{figure}[H]
\centering
\includegraphics[width=\textwidth]{fig5_cross_design}
\caption{The cross-design test as paired panels. Panel (a): dollar recovery for the two real-covariate variants against the pre-committed thresholds (above 50\% undercuts the paradigm interpretation; below 35\% corroborates it) and the synthetic-population control band; the recalibrated variant's 59.3\% crosses the undercutting threshold while the frozen variant's 20.9\% stays within the corroborating band. Panel (b): lag-0 path correlations, all negative, so neither variant reproduces the empirical path's timing regardless of calibration choice, and after Section~\ref{sec:pathb}'s detrending result no estimator in this paper does.}
\label{fig:crossdesign}
\end{figure}

\begin{table}[H]
\centering
\small
\begin{threeparttable}
\caption{Cross-design test results: real Freddie Mac structural covariates substituted into the ABM under two calibration choices, against the synthetic-population production ABM.}
\label{tab:crossdesign}
\begin{tabular}{@{}p{4.0cm}lll@{}}
\toprule
 & Recalibrated (primary) & Frozen (robustness) & Synthetic control \\
\midrule
Mobility-scale parameter & 36,086 & 43,883 & 43,883 \\
Trapped liquidity & \$453.5B (59.3\%; fifty-seed mean 60.2\%) & \$159.5B (20.9\%; 22.0\%) & \$91.0B (11.9\%) \\
Mean CPR & 8.14\% & 11.63\% & 11.68\% \\
CPR path corr. (lag 0) & $-0.336$ & $-0.311$ & $-0.318$ \\
Peak lag & 0 (no positive alignment) & 0 (no positive alignment) & 0 \\
\bottomrule
\end{tabular}
\begin{tablenotes}[flushleft]\footnotesize
\item Peak lag follows the max-$|r|$ convention of Table~\ref{tab:estimators}. Under the alternative most-positive-$r$ convention, the synthetic control's best alignment is $+0.195$ at lag $-3$, the ``$+0.19$ at the three-month lead'' of Section~\ref{sec:abm-results}; the two conventions summarize the same cross-correlation profile, and the substantive conclusion (no positive contemporaneous alignment, which, after Section~\ref{sec:pathb}'s detrending result, no estimator including Path B shows) holds under either.
\end{tablenotes}
\end{threeparttable}
\end{table}

Two results stand out. First, by my own ex-ante criterion, the primary recalibrated variant's 59.3\% recovery, 60.2\% averaged over fifty seeds, moves against the paradigm interpretation: with real structural covariates and a like-for-like turnover anchor, the household-choice ABM explains far more of the benchmark than the synthetic-population specifications throughout Sections~\ref{sec:abm} and~\ref{sec:robustness-scale} suggested. Second, the \$294 billion gap between the recalibrated and frozen variants (produced solely by how the turnover-floor anchor is set, with the underlying covariates held fixed) indicates that the ABM's aggregate share is highly sensitive to a calibration choice that is orthogonal to data source. The present test cannot, on its own, cleanly attribute Sections~\ref{sec:abm} and~\ref{sec:hazard}'s contrast to modeling paradigm rather than to data source: recalibration choice is doing comparable work to whichever of paradigm or data source the test was designed to isolate.

A pre-committed post-stratification separates composition from heterogeneity (run \texttt{cross\_design\_reweight}). Reweighting the real-covariate population to the SOMA coupon-by-vintage-group composition---the book the benchmark is measured on---with the mobility scale re-derived on the reweighted population under the same floor rule moves the recalibrated variant from 59.3\% to 76.3\% (importance weights; effective sample size 1{,}332 of 10{,}000, disclosed; a seed-42 re-draw within SOMA-weighted cells gives 70.6\%) and the frozen variant from 20.9\% to 12.6\% (re-draw 36.0\%). The composition mismatch therefore understated rather than manufactured the recalibrated recovery: at book composition, real-covariate household choice with a like-for-like floor anchor explains roughly seventy to seventy-six percent of the benchmark, the ex-ante undercutting verdict survives composition adjustment, and the calibration-choice spread this subsection flags widens to roughly \$488 billion. The reweight is a level exercise; the path-level failure of Table~\ref{tab:crossdesign} is not revisited by it.

A third variant constrains the calibration question from the other side: it drops the turnover-floor anchor entirely and parameterizes the two author-calibrated behavioral elements from the same external source Path B's elasticity imports, pinning the mobility-desire scale to \citepos{liebersohn2024} zero-gap moving level (a 1.5\% quarterly ZIP-code moving hazard when the average rate gap was essentially zero, 5.9\% annualized) and the wait-and-see freeze share to the magnitude their 6.5\%-per-quarter-per-100 bp mobility decline implies at the production 150 bp trigger (18.3\% against the author-calibrated 20\%; the DTI wall and the loss-aversion coefficient are already externally sourced). Recovery does rise without the floor recalibration (to 150.6\% of benchmark, overshooting it by half), but it does so by predicting essentially zero prepayment at QT-window rate gaps: the variant's CPR at an 8\% market rate is 0.0\%, against the observed 4--5\% discount-cohort turnover and against the roughly 1.44--1.45\% quarterly moving hazard \citet{liebersohn2024} themselves observe at 2022--2023 gaps. The externally anchored gate structure can match the zero-gap moving level or the observed deep-gap turnover, but not both: the exponential mobility-desire distribution that fits the former (scale 4{,}472 against the recalibrated 36{,}086) shuts moving down completely once the par-payoff penalty binds. I read this as clarifying the floor anchor's role: it is what keeps the ABM on the observed turnover data, not an artifact that suppresses the model's rate response. The anchored dimension itself is never claimed as fit: matching the floor is tautological by construction (Appendix~\ref{sec:method-frictions}), and what this variant measures is what that tautology purchases---empirical admissibility at deep gaps---one more instance of the deflate-as-well-as-inflate indeterminacy this subsection already carries (same harness as the two variants above). Scored against the pre-committed bands, as this variant's own specification required, its 150.6\% classifies as undercutting the paradigm reading, and the frozen artifact records that classification. I do not read it as informative about the paradigm question, and I label that reading post hoc rather than pre-committed, because this is the kind of distinction the paper's pre-commitment design is meant to surface. The pre-run specification committed only to reporting the deep-gap floor diagnostic ``alongside whatever the recovery number is, whichever direction it moves''; it did not commit to treating a failed diagnostic as grounds for removing the recovery from the band. The admissibility precondition---that a variant must reproduce an empirically admissible deep-gap turnover level before its recovery is interpretable---was formulated after the result was seen, and it is reported here as a labelled post-run qualification, not as a pre-registered exclusion. A reader who declines the qualification should score the variant at 150.6\% against the 50\% band. That strengthens the undercutting reading this subsection already concedes rather than weakening it, and no headline quantity turns on which way it is read.

To test whether a minimal simulated-minimum-distance calibration could reconcile the two external moments, I ran one: a joint two-parameter fit, the mobility-desire scale together with a forced-move point mass added to the structural family, with the moment set, admissibility bands, distance, and verdict vocabulary all fixed before execution and a parity gate replaying the committed 59.303\% first. The verdict is joint infeasibility, now with two degrees of freedom rather than one: at the distance-minimizing point (scale 1{,}000, forced-move mass 3.9\%) the zero-gap moving level overshoots its band just as the deep-gap turnover reaches its own (0.0644 against $[0.055, 0.063]$; 0.0394 against $[0.040, 0.050]$), because the forced mass lifts both moments together once the scale is at its floor. The near-fit, scored through the same harness, recovers 107.1\% of benchmark, far above the pre-committed cross-design bands, so external anchoring pushes this structural family past the benchmark, not toward it. I deliberately do not fit QT-window CPR moments: the recovery statistic is a functional of that path, and matching it would convert the test into a fitted quantity, the overfitting the pre-commitment architecture exists to prevent. The hazard-side headline quantities are not functions of ABM calibration and cannot move under any outcome of this run.

Both variants retain the ABM's central path-level failure: the monthly CPR path correlation stays negative at lag 0 ($-0.336$ and $-0.311$), as all estimators' are once the window's shared trend is removed (Section~\ref{sec:pathb}); Path B's distinction is level accuracy, not timing. The timing critique of Section~\ref{sec:abm} therefore survives regardless of which calibration variant is preferred, even as the dollar-recovery critique does not straightforwardly survive under the recalibrated variant.

\subsection{Cross-Design Synthesis: The Symmetric Companion Test}\label{sec:robustness-synthesis}

The symmetric companion test (estimating the hazard framework's survival structure on a synthetic population, to check whether that structure alone recovers the benchmark independently of real data) materially updates the cross-design picture: Path B recovers 106.0\% of the benchmark on a fully synthetic population with no real Freddie Mac data at all (96.9\% on the shared accounting basis, whose common netting applies equally here), close to its 107.0\% real-data recovery at the same in-sample calibration point, though the raw levels correlation drops from $+0.190$ to $-0.06$ at lag 0 (a levels comparison; neither value survives detrending, Section~\ref{sec:pathb}). That 106.0\% is a recovery level, however, and this paper demotes recovery levels everywhere else for a reason. Running the mechanism-removed leg on the same synthetic population --- $\beta_1 = 0$, nothing else changed --- recovers 95.6\%, so what the companion test identifies is a differential of 10.4 points, or \$79.3 billion, and not the 106.0\% level (\texttt{synthetic\_companion\_null}). That differential is of the same order as the 9.2-point real-book gap between the central leg's 107.0\% and the $\beta_1 = 0$ null's 97.8\%, which Section~\ref{sec:identification} already treats as too small to license a level-based reading; a superseded cross-basis figure put that gap at 18.3 points, which was an error of basis rather than of arithmetic --- it differenced a standalone-scorer central leg against a shared-basis null --- and on any single basis the real-book gap is 9.2 points at the in-sample calibration point and 5.6 points at the off-window floor, so the synthetic differential slightly exceeds its real-book counterpart rather than falling short of it; the synthetic recovery is floor-dominated in the same way the real one is. This decomposes the tension rather than resolving it uniformly: at the aggregate level, the household-choice paradigm required real covariates to reach 59.3\%, while the loan-level paradigm reaches a comparable aggregate recovery with none. At the path level, both paradigms still depend on real structure for correct timing. I treat the aggregate-level paradigm claim of Sections~\ref{sec:hazard} and~\ref{sec:conclusion} as corroborated by this result on its differential rather than substantially strengthened by its level, and the path-level claim as still open.

\subsection{Cross-Foundation Checks: Hybrid Pipeline and Full-Book Weighting}\label{sec:robustness-hybrid}

The hybrid pipeline referenced in Section~\ref{sec:abm-danish} holds the paper's accounting fixed while swapping the micro-foundation for loan behavior. Concretely, it runs the ABM's macro-accounting layer unchanged (SOMA balance tracking, phased-cap netting, income-scaled curtailment, term-aware scheduled amortization, and the Danish dynamic-balance loop), but replaces the ABM's CPR surface with Path B's microsimulated CPR path; no settlement kernel applies, since Path B settles prepayments in the month they occur. Holding the accounting fixed, the U.S.-side recovery is 97.9\% (the shared layer adds the income-scaled curtailment flow that Path B's standalone scorer omits: \$69.56 billion over the window), and re-running the Danish leg under the Berger-recalibrated regime yields, under the production U.S.-intercept anchor, the U.S. 4.76\% / Danish 5.61\% mean-CPR pair and the $+\$61.2$ billion institutional gap that Section~\ref{sec:abm-danish} adopts as its production estimate (the Danish-level bracketing anchor gives 3.39\% and $-\$99.9$ billion). The construction removes accounting differences between frameworks as a confound of the institutional comparison: both regimes flow through one accounting layer with one empirically grounded micro-foundation.

The same layer now scores every estimator, which is what makes the recovery figures commensurable with the benchmark: the actual SOMA roll-off underlying the empirical shortfall contains borrowers' curtailment dollars, a flow no hazard simulation models, so the shared layer adds it to each simulated roll-off. The netting is identical across estimators (\$69.56 billion for every run), and identical \emph{by construction}, not by coincidence, with the conversion architecture stated per leg: each estimator's hazard dynamics evolve on its own internal balances (Path A drains cohort balances, Path B drains loan balances, and the Danish legs deliberately compound a dynamic counterfactual balance, Appendix~\ref{sec:method-abm}), but every \emph{U.S.-leg} dollar conversion is common (simulated roll-off is rescaled month-by-month to the actual holdings path (WSHOMCB) by the book-to-sample ratio), so the curtailment flow, computed on those same actual holdings, is a common flow rather than a per-run quantity. It is a flat 9.1-percentage-point wedge between the two bases; the invariance argument covers U.S. legs only. Danish legs compute curtailment on their own counterfactual balances inside the Danish loop; the slower Danish drain leaves higher balances and hence more curtailment dollars, a differential I quantify and print: Danish-leg curtailment is \$70.33 billion against the U.S. legs' common \$69.56 billion for the central hybrid (+\$0.77 billion), already embedded in every printed gap and immaterial to its sign or size.

Table~\ref{tab:bases} collects every hazard leg's recovery on all four accounting bases (the lock-in margin is basis-invariant at either floor because the netting is bit-identical across legs and cancels exactly, an invariance that holds at every cell of Section~\ref{sec:robustness-floor}'s calibration box). The fully corrected figure comes from a joint run (the book-coupon-reweighted simulation scored through the shared layer in one pass), which recovers \$765.1 billion (unrounded: \$765.077 billion): 100.0\% of the \$764.7 billion benchmark at one decimal, and exactly 100.04\% as the ratio of the two unrounded figures, $765.077/764.748$. That is a mixed-basis measurement: the reweight matches borrower note rates to security pass-through buckets, and converting the sample side to the pass-through basis first (note rate less the vintage guarantee fee and base servicing, $\Delta = 0.80$ points, with legs at $0.60$--$0.90$) puts the like-for-like composed figure at 98.1\% (97.6--99.1 across the legs; run \texttt{coupon\_convention\_reweight}) --- so 100.0\% is the committed mixed-basis reading, not the corrected number. Path A's joint run reads \$905.1 billion, 118.4\%. The composed figure is conditional on the ex-ante covariate priors (the covariate band spans 92.2\%--99.2\% on the shared basis, approximately 94.3\%--101.3\% composed (Section~\ref{sec:pathb})), so its precision is a property of that calibration, not a tuned result. Because the netting is common, the joint results equal the one-at-a-time deltas composed ($107.0 + 2.1 - 9.1$) exactly; I report the identity as a cross-check on the joint run, not as its source. The abstract and Sections~\ref{sec:hazard-interp} and~\ref{sec:conclusion} state recoveries on this basis.

The curtailment layer invites one further stress, since its assumed rate is an external input rather than a fitted quantity. The monthly curtailment flow is the product of an income-driven fraction, which is independent of the assumed curtailment rate, and that rate, applied to the actual holdings path (WSHOMCB), so the netted total is exactly proportional to the assumed rate: it does not interact with any estimator's dynamics. A $\pm 25\%$ stress on the assumed rate (whose realized window average is approximately 0.84\% annualized, Section~\ref{sec:method-benchmark}) therefore scales the \$69.56 billion netting to \$52.17 and \$86.95 billion and the 9.1-point wedge to 6.8 and 11.4 points, a uniform shift in every U.S. leg's shared-basis recovery that leaves the lock-in marginal and every cross-estimator gap unchanged to the decimal, because both legs of each comparison net the identical flow. The same argument holds month by month for any time-varying curtailment process: the netting is the month-by-month product of a shared curtailment rate and the shared actual-holdings path, and the trapped-liquidity metric is a linear signed sum against a fixed cap-target path, so a seasonal profile, or an income elasticity that differs across 2022--23 and 2024--25, subtracts out of the marginal term by term whatever its shape; only the common level wedge and the Danish leg's sub-\$1 billion differential can move. A demonstration run makes the identity concrete: a $\pm 25\%$-amplitude seasonal profile and a 2022--23 versus 2024--25 regime split move the netted total to \$70.06 and \$68.09 billion and the wedge to 9.16 and 8.90 points, while the lock-in marginal is unchanged to below $10^{-13}$ points on both legs and the simulated CPR paths are bit-identical. A fourth profile makes the servicer case explicit: a deterministic two-servicer mix (fixed 0.60/0.40 shares, front-loaded and flat group factors aggregating to one monthly rate) moves the wedge to 9.30 points (\$71.12 billion netted) with the marginal again unchanged, because any servicer-heterogeneous pattern that aggregates to a monthly rate is a special case of the same identity. The Danish legs, which compute curtailment on their own counterfactual balances, scale only to first order, and an exact rerun of the accounting layer across the same range (the simulated CPR paths are bit-identical across scales, since curtailment enters only the accounting layer) puts their differential at $+\$0.76$, $+\$0.77$, and $+\$0.66$ billion at the $0.75\times$, $1.0\times$, and $1.25\times$ points respectively. Those figures are measured on the \emph{Danish-level bracketing} leg, because the committed microsim cache the scaling run consumes was built under that anchor, and they do not carry over to the production leg. Re-scoring the same range on the U.S.-intercept cache (run \texttt{curtailment\_danish\_scaling} with the production anchor) reverses the sign and crosses the materiality threshold: the differential is $-\$1.68$ billion at the unit point and $-\$2.39$ billion at $1.25\times$. The mechanism is the one this paragraph already names, read the other way. Curtailment accrues on retained balance, so the sign follows which leg drains faster: the bracketing Danish leg is \emph{slower} than the U.S.\ leg (3.39\% against 4.76\% mean CPR) and so retains more balance and accrues more curtailment, while the production Danish leg is \emph{faster} (5.61\%) and accrues less. I therefore withdraw the claim that the curtailment differential is immaterial to the Danish comparison. On the production leg it is not: at $-\$1.68$ billion it is 2.7\% of the $+\$61.2$ billion in-sample gap and 6.0\% of the $+\$28.2$ billion off-window gap, small enough not to threaten either sign but large enough that it belongs in the figure rather than in a dismissal.

The second check reweights the sample to the book. The Path A and Path B headline recoveries are balance-weighted estimates on the composition of the Freddie object each path uses: Path A the full 2017--2021 origination universe, Path B the 75,000-loan draw, whose unweighted mean note rate of 3.86\% sits above the book's 2.49\% face-weighted pass-through coupon. The origin of the Path B reweight is that draw; the unindented bucket shares Table~\ref{tab:composition} prints beside the book are the universe's, exposure-weighted at the window open, and the two are different populations. Reweighting the simulated hazards to the book's actual coupon mix shifts Path B from 107.0\% to 109.1\% of the benchmark and Path A from 121.5\% to 127.5\%, a composition effect of roughly two percentage points for Path B that bounds the sample-to-book extrapolation error along the coupon dimension --- on the committed mixed coupon basis; on a single pass-through convention the same reweight moves Path B only from 107.0\% to 107.2\%, so the like-for-like composition effect is $+0.2$ points and the bound tightens by an order of magnitude (run \texttt{coupon\_convention\_reweight}) (the full sample-versus-book composition comparison is Appendix~\ref{app:composition}). The agency dimension (Ginnie Mae collateral) has no internal bound; it is externally bounded and now scored directly by a published-series overlay in Section~\ref{sec:pathb}, whose Ginnie-specific contribution ($+\$38.5$ billion) lands inside the static bound.

\subsection{Floor and Elasticity-Band Sensitivity: The Calibration Box}\label{sec:robustness-floor}

Path B's dominant level-setter is the baseline turnover floor, and Section~\ref{sec:pathb} concedes that its 4\% calibration is measured inside the evaluation window. To quantify what rides on that choice, I swept the floor across seven values from 2\% to 6\% annual CPR jointly with the full Liebersohn--Rothstein elasticity band (5.5\%, 6.5\%, 7.7\%), rerunning the production central and $\beta_1 = 0$ null legs at every floor under the production convention: same committed 75,000-loan sample, same seed and regime tuple, one shared macro frame. Both sweeps were pre-committed: grid, parity gates, and interpretive thresholds entered the repository before any run executed (commit ordering recorded in Appendix~\ref{app:ledger}); the fresh runs at the production point reproduce the committed artifacts exactly (central \$818.530 billion, null \$748.185 billion, marginal $+\$70.345$ billion), and the instrumented floor-bind share reproduces the 36\% of 1,683,124 evaluated loan-months quoted in Section~\ref{sec:pathb}.

The sweep mechanics and the contamination separation behind the off-window anchors are collected in Appendix~\ref{app:floormech}; the exhibits they produce --- the calibration box (Table~\ref{tab:floorband}), the marginal as a curve in the elasticity (Table~\ref{tab:lowband}), and the out-of-sample anchors (Table~\ref{tab:oosfloor}) --- are in Section~\ref{sec:hazard}, and what follows are the readings the headline interval rests on.

The box sweeps the floor's level but never its functional form---the choice that censors the elasticity. A form test specified ex ante (spec and interpretive threshold fixed in the script's header before its runs executed, the ex-ante convention Appendix~\ref{app:ledger} records, with the script and its artifact entering the repository in one revision; the max-form pair at the production floor reproduces the committed artifacts exactly, as its parity gate) reran the central and null legs with the hard maximum of \eqref{eq:pathB} replaced by a competing-risks combination on the survival scale, $h = 1 - (1 - \underline{h})(1 - h^{\mathrm{vol}})$, under which involuntary and voluntary hazards add and the elasticity is never inert. The result changes the reading in two ways. The marginal moves to $+11.25$ points at the production floor, a $+2.05$-point shift that exceeds the $\pm 1$-point immateriality threshold fixed ex ante, so the form-dependence is real. But under the additive form the marginal is nearly floor-invariant, spanning only $+11.29$ to $+11.21$ points across the 3--5\% floors whose max-form span ($+10.8$ to $+5.5$) failed the stability threshold of Appendix~\ref{app:floormech}: the failed gate was measuring the censoring mechanics of the max form, not instability in the elasticity's contribution. The aggregate level, by contrast, is strongly form-dependent (the additive central leg recovers 67.2\% at the production floor, versus 107.0\% under the max form, because adding a 4\% involuntary hazard to moderate voluntary hazards raises book-wide prepayment), and the max form remains production because its floor semantics, a minimum total turnover matching the observed 4--5\% on deep-discount cohorts where the voluntary hazard is near zero, is the reading the calibration anchor supports; at that anchor the two forms nearly coincide. A pre-committed completion of this design (run \texttt{floor\_form\_offwindow}) runs the additive form at the off-window anchors themselves: the additive marginal is $+11.22$/$+11.21$/$+11.19$ points at the 4.695/4.991/5.334\% floors, within 0.06 points of its production-floor value, so the entire in-window-to-off-window demotion of the max-form marginal operates through the censoring channel this paragraph diagnoses, and the near-coincidence defense of the max form at the deep-discount anchor does not extend to the shallow-gap off-window anchors. The form dimension is now a measured curve rather than a two-point hull (run \texttt{floor\_form\_mixture}; both endpoints reproduce the committed forms bit-exactly). Let $s$ be the share of the measured floor entering as an additive competing involuntary cause, with $1-s$ remaining a floor on total turnover: $s = 0$ is the production hard maximum and $s = 1$ the additive form. The curve rises steeply and plateaus --- at the off-window anchors the marginal is $+7.5$ points by $s = 0.1$, $+9.7$ at $s = 0.25$, $+10.7$ at $s = 0.4$, and flat near $+11.2$ from $s = 0.6$ up (peak $+11.24$ at $s = 0.8$, a 0.03-point plateau dip to the additive endpoint) --- while the central leg's standalone-scorer recovery falls from 100.4\% to 55.9\% along the way, the level cost the additive end pays. Section~\ref{sec:pathb}'s reading that the strictly-involuntary share is plausibly well under half therefore does \emph{not} pull the marginal toward the max end: even modest additive shares close most of the gap to the additive form, because an additive component acts wherever the max form censors. What keeps the headline at the $s = 0$ endpoint is the max form's semantics, now displayed as the conservative endpoint of a measured curve rather than defended by fit. The level cost of that form-robustness carries over to the off-window anchors too, and it is large enough to report beside the hull rather than only at the production floor: at the 4.991\% headline anchor the additive central leg recovers 55.9\% of the benchmark against a 44.7\% null, and at the 4.695\% and 5.334\% anchors 59.3\% and 52.0\% at the central band point, so the additive form undershoots the max form's aggregate level by 43.5 to 45.3 points at every off-window anchor. The hull's upper endpoint is produced by the least accurate of these cells: $+13.10$ points comes from the additive form at the 4.695\% floor and the band's 7.7\% high edge, whose central leg recovers 61.2\% of the benchmark it is scored against. The form-conditional range is therefore a range over specifications that differ sharply in aggregate fit, and its upper end should not be read as an equally credentialed member. The max form remains production for its semantics --- the floor as a minimum on \emph{total} turnover, the reading the deep-discount calibration anchors --- and because it is the measured form curve's conservative endpoint (below); the aggregate-level comparison is context for the level, not the selection rule for the marginal's form, since the design does not identify levels. One further run measures the cell the hull leaves implicit. The form and transform dimensions were reported side by side without being crossed, which presumes they add. A pre-committed run of the joint cell (\texttt{concave\_additive\_marginal}) shows they do not. The concave transform under the additive form puts the marginal at $+9.25$ points at the production floor and $+9.22$ at the 4.991\% anchor---inside the hull at both, so the range itself stands---but the interaction is $-1.47$ points at the headline anchor, past the same $\pm 1$-point convention used throughout. The reason is the censoring this paragraph describes: the concave transform costs only $-0.52$ points under the max form, which absorbs most of it into the floor, against $-1.99$ under the additive form, which never censors and so passes the whole transform effect through. Two consequences. A successor run closes that grid (run \texttt{concave\_hull\_band\_ends}): the concave transform's remaining sixteen off-window cells put the form $\times$ transform product at thirty-six of thirty-six, and the hull does not survive intact. Seventeen of the eighteen concave cells fall inside it, but concave under the max form at the highest defensible floor and the band's low edge returns $+3.53$ points, below the $+3.89$ lower edge the log-linear cells set, so the form-conditional hull widens at the bottom to $+3.5$ to $+13.1$ points. Its upper edge is unchanged, because the concave transform is signed downward and its highest cell reaches only $+10.83$. Any reading that composes a form adjustment with a transform adjustment additively is unsupported at the calibration this paper headlines. The form-robust statement of the paper's identified quantity is a lock-in marginal of roughly $+9$ to $+11$ points at the production floor, positive under every level, band, and form combination tested.

Re-running the central and $\beta_1 = 0$ legs with the floor set to each defensible off-window value, under the production convention in every other respect (the fresh runs reproduce the committed production artifacts exactly at the 4\% floor: central \$818.530 billion, null \$748.185 billion, marginal $+\$70.345$ billion, floor-bind share 36\% of 1,683,124 loan-months), gives Table~\ref{tab:oosfloor}. Across the 2018 leg's defensible band the marginal runs $+6.8$ to $+4.3$ points ($+\$51.8$ to $+\$32.6$ billion), centered at $+5.6$ points ($+\$42.6$ billion) at that leg's mid-grid 4.99\% reading, and $+3.9$ to $+7.5$ points once the full Liebersohn--Rothstein band is swept jointly with it. The contaminated anchors are what produce the small numbers: the pooled 6.07\% read (evaluated at the box's adjacent 6.0\% row) gives $+2.3$ points, and the 2019 leg's 6.91\% read gives $+0.7$. The direction of the in-window concession stands---an off-window floor is higher than the in-window 4\%, and because the marginal decreases in the floor, a higher anchor shrinks the identified contribution rather than reversing it---but the conclusion that the plausible contribution therefore sits at the calibration box's lower edge does not. That lower edge was an artifact of the refinancing-contaminated pooled read; the defensible off-window floor places the contribution in the box's interior, at roughly $+5.6$ points, not at $+2.1$ to $+2.4$. Every read in that band is taken on cohort-months aged twelve to twenty-four, the only ages the 2018 leg contains, so the leg's mature-seasoning test is not computable rather than passed (Appendix~\ref{app:floormech}); where the mature contrast is measurable it runs one way in every cell (Section~\ref{sec:identification}), which puts the band below the mature book's baseline turnover and the marginal it implies above what that book would give.

The move from the in-window anchor to the off-window one changes two things at once, and until now the paper charged both to the window. The in-window anchor is measured at gap $\leq -2$ points and reads 3.972\%; the off-window anchor is measured at gap $\leq -0.25$ points and reads 4.991\%. Mapping both through the committed floor-to-marginal grid (PCHIP on the frozen sweep, not re-derived, and reproducing the committed grid to within \$0.69 billion at every read inside it), that is a move of $-\$28.312$ billion, or $-3.702$ points. Decomposed at matched depth, the move splits: the window component, measured with depth held at gap $\leq -0.25$ points on both sides, is $-\$20.874$ billion ($-2.729$ points, 73.7\% of the move); the depth component, measured within the in-window panel by moving from gap $\leq -2$ points to gap $\leq -0.25$ points, is $-\$7.438$ billion ($-0.973$ points, 26.3\%). Re-anchoring the same decomposition at the deepest depth both panels support (gap $\leq -1$ point) gives 66.9\% window and 33.1\% depth. The window dominates either way, so the demotion is not primarily a depth artifact. But depth owns a real and previously unattributed \$7.4 billion, and the consequence is a correction to how the demotion should be described: at matched depth the in-sample comparison point is $+\$63.5$ billion, not the $+\$70.3$ billion this paper narrates the demotion against. Roughly \$20.9 billion of the \$27.7 billion demotion is the window; the remainder is a change in the discount depth at which the anchor is defined, and has nothing to do with out-of-sampleness. The reverse ordering---window first at matched deep depth, then depth---is not computable: it would have to read the off-window leg at gap $\leq -2$ points, where that leg retains 0.02\% of its exposure, so no average of the two orderings is formed and the depth-first path is reported alone. Extending the 2018 depth ladder from three well-supported depths to five leaves the committed clean band exactly [4.695\%, 5.334\%], because the two additional reads (4.722\% at gap $\leq -0.75$ points, 4.869\% at gap $\leq -1$ point) fall strictly inside it, and the marginal band therefore stays \$33.31 to \$51.75 billion around the \$42.58 billion point (the \$33.31 billion grid-read endpoint and the \$32.6 billion direct-run endpoint quoted earlier differ by the PCHIP grid-read wedge, whose committed maximum is \$0.69 billion). Two further pre-committed measurements discipline this band (run \texttt{floor\_uncertainty}). First, the floor reads themselves carry sampling error the depth cuts do not express: a 1000-replicate stratum-cluster bootstrap puts the clean-band reads' standard errors at 0.39--0.46 points of CPR (31 clusters carry the 2018 leg), and propagating the point anchor's replicates through the committed floor-to-marginal grid gives a percentile 95\% interval of $+3.0$ to $+8.0$ points on the off-window marginal---double the depth-cut band's width, so the binding uncertainty on the headline marginal is the floor's sampling error, not the depth convention (the in-window reads, on 226 clusters, carry standard errors of 0.16--0.19 points). The percentile read under-covers at that cluster count and is demoted; the pre-committed corrections (run \texttt{floor\_inference\_correction}, with the percentile layer reproduced bit-exactly as its parity gate, and its Webb-weight successor \texttt{floor\_inference\_correction\_v2}) widen it mostly at the top, and the corrected ladder, with its degrees of freedom, is Table~\ref{tab:ladder}: the Webb wild-$t$ interval of $+2.9$ to $+8.7$ points is the binding layer, with the lower edge above zero on every corrected read. Second, an age-standardization of the off-window floor---reweighting the 2018 leg's age-bucket reads to the in-window deep-discount population's age mix, with the 2018 leg's empty mature buckets imputed from the in-window age gradient (84.0\% of weight imputed, and labeled as such)---reads 5.51\%, above the clean band's top, implying a marginal near $+3.8$ points: an indicative bound rather than a measurement, but signed one way, so the clean band is open below $+4.3$ and that endpoint is not a conservative floor. This read holds the conventional accounting fixed, and the Ginnie overlay of Section~\ref{sec:pathb} holds the floor at 4.991\%: both correct the same $+5.6$-point cell, both cut the same way, and the joint cell is now run (run \texttt{b5\_joint\_cell}). They do not look like the same adjustment counted twice---one moves the floor's level, the other confines the marginal to the conventional face share, and the overlay's share scaling is identical across the three speed series it was tested on---so carrying that $0.797$ scaling up to the higher floor, which the overlay run does not itself establish, would put the marginal near $+3.0$ points, at the bottom of that sampling interval and below the form-conditional hull's $+3.5$ lower end. That composition is now measured rather than projected: the joint cell returns $+2.9$ points, an interaction of $+0.0$ against the proportional projection --- where the first joint cell this paper ran for two adjustments to the marginal---the concave transform under the additive form, above---returned an interaction of $-1.47$ points at the headline anchor, which is why a composed point here would not be credible.
The same matched-depth machinery undercuts the argument by which the 2019 leg was rejected in Appendix~\ref{app:floormech}, though not the rejection itself. The rejection compares the 2019 leg's gap-$\leq$-0 read of 6.910\% against the retained 2018 leg's gap-$\leq$-0.25-point read of 4.991\%: a 1.919-point spread taken across two different depths, which is not a valid comparison as written. At matched depth the two legs do not converge either. Over the four depths at which both legs clear the 10\% support rule, the 2018-versus-2019 spread runs 1.576, 0.491, 2.112, and 0.857 points---non-monotone in depth, and at half a point the matched-depth spread exceeds the mismatched 1.919 points the earlier argument rested on. So neither ``the legs disagree by 1.9 points'' nor ``the legs agree once depth is matched'' is a defensible summary; the second would be a cherry-pick of one depth in four. The pooled 2017--2019 read is not a third independent leg, since it contains both 2018 and 2019 (and 2017), so any appeal to three legs agreeing double-counts; there are two. And the 2019 leg's own exposure collapses from \$1{,}085 billion at gap $\leq$ 0 to \$199 billion by half a point and \$0.29 billion by one and a half, so its deep reads are noise on the same grounds that disqualify the deep off-window cell. The honest verdict is therefore that the refinancing-contamination story is unadjudicable below gap $\leq$ 0 rather than refuted there: the only depth carrying solid exposure for both legs is gap $\leq$ 0, where 2019 sits 1.576 points above 2018, directionally consistent with refinancing contamination and consistent with the seasoning diagnostic that independently disqualifies the 2019 leg. The conclusion stands; the argument as earlier stated does not, and I state the weaker one.

A second instrument addresses the months dimension directly rather than the population dimension, and it is a floor-stability check rather than an outcome holdout: no realized prepayment from the held-out months enters either simulated leg, so the construction can falsify only the floor calibration's temporal stability, never the marginal against realized data. Calibrating the floor on deeply out-of-the-money turnover confined to the window's first nineteen months (June 2022 through December 2023) returns 3.97\%, and the central-minus-null marginal can then be summed over the remaining twenty-three months (January 2024 through November 2025) alone, which enter the floor's calibration nowhere. That held-out marginal is $+\$45.1$ billion, against $+\$44.8$ billion at the production floor over the identical months: very nearly the magnitude, and with it the sign, survives evaluation months that the calibration never touches. The decomposition is exact rather than approximate because trapped liquidity is a monthly sum against the redemption target, so the target cancels from the central-minus-null difference and the monthly marginals add back to the aggregate. Across every floor and band combination run here---twenty-seven cells spanning floors from 3.97\% to 6.91\%---the marginal is positive without exception, a consistency check the window's rate configuration cannot let fail (Section~\ref{sec:identification}) rather than a second identifying result. The identified content of the lock-in channel is therefore a magnitude that survives both instruments as a bounded interval: $+4.3$ to $+6.8$ points at the defensible off-window floor under the production max form ($+3.9$ to $+7.5$ jointly with the elasticity band), $+9.6$ to $+13.1$ under the additive form at the same anchors, a form-conditional hull of $+3.5$ to $+13.1$ points, against the $+9.2$ points the in-window production calibration returns, which this paper now reports as its in-sample calibration point rather than its headline.


\section{Conclusion}\label{sec:conclusion}

The shortfall of the Federal Reserve's mortgage runoff beneath its Quantitative Tightening caps was not solely a product of high interest rates, but a consequence of the structural design of the U.S. mortgage contract interacting with a prepayment environment mostly mechanical relative to the caps: under the production floor form the $\beta_1 = 0$ null recovers 85.7\% of the benchmark (35.6\% under the additive form) at the headline calibration under any rate response above the baseline turnover it embeds, and how much of the remainder sits outside household-level decision-making is a question the cross-design test of Section~\ref{sec:robustness-crossdesign} leaves open. This paper's two central empirical findings, established in Sections~\ref{sec:abm} and~\ref{sec:hazard} under the synthetic household population specified in Section~\ref{sec:method}, are complementary and neither is derived from the other, though they are not independent in calibration: both are anchored to a baseline turnover floor read off discount-cohort turnover, and at the in-sample calibration point to the same 2023--24 read (Section~\ref{sec:pathb}).

The ABM (rational, behaviorally extended, or subjected to a direct burnout test) cannot explain the majority of the Federal Reserve's trapped liquidity, while a structurally distinct hazard framework built entirely from loan-level survival dynamics, with no household utility function at all, recovers 91.3\% of the full empirical benchmark on the benchmark-consistent accounting basis in its microsimulation form at the off-window floor this paper headlines, 88.2--93.7\% across that floor's defensible range, and 97.9\% at the in-sample calibration point, so the miss is 8.7 points on the headline calibration and 2.1 points on the in-sample one (a level-accuracy distinction, not a timing one: no estimator's monthly co-movement survives detrending, Section~\ref{sec:pathb}; the estimated cohort-hazard form lands at a comparable aggregate level but does not corroborate it---its own coefficient's sign is not statistically distinguishable from zero---carries no valid production-spec uncertainty, and is excluded from headline figures, Section~\ref{sec:patha}).

The composition of that recovery, however, disciplines what the contrast can mean: under the production floor form the $\beta_1 = 0$ null recovers 85.7\% on the same basis at the headline floor (35.6\% under the additive form) and 88.7\% in-sample, and the paper's own ablations bound the named loan-level mechanisms (burnout $0.45$ points at the headline floor and $0.86$ in-sample, FICO/LTV covariates $1.3$ points, the delinquency pipeline \$0.39 billion; no servicer channel beyond that pipeline is modeled at all). What delivers the level is scheduled amortization, the seasoning baseline, and the turnover floor: the shortfall against the phased caps is mostly mechanical, because the caps sat far above what any plausible prepayment environment would have delivered---though against the Federal Reserve's own ex-ante projection rather than the non-binding cap, the lock-in channel accounts for roughly half the genuine surprise under the central allocation (23--49\% at the off-window floor and 38--80\% at the in-sample calibration point across the two disclosed intra-2022 allocations; Section~\ref{sec:method-benchmark})---and the lock-in elasticity's contribution is a bounded margin over that mechanical base rather than a pinned one: $+2.9$ to $+8.7$ points once the turnover floor is measured off the locked-in population it is meant to bound and that floor's own sampling error is propagated at the central elasticity with a wild-cluster correction (the percentile read is $+3.0$ to $+8.0$; the interval prices the floor read alone, and the elasticity dimension is Table~\ref{tab:lowband}), with $+5.6$ points the value at the mid-grid anchor inside it and every measurable correction to the floor or the accounting basis moving that value down within the interval, and $+9.2$ points at the in-sample calibration point (Section~\ref{sec:identification}; $+11.2$ under the competing-risks floor form). A separate question (whether the estimator contrast reflects the household-choice paradigm itself or the synthetic population it was tested on) receives an asymmetric answer in Section~\ref{sec:robustness-crossdesign}: the hazard framework's aggregate recovery survives a fully synthetic population, while the ABM's share moves with calibration choice and data source, its recalibrated real-covariate variant reaching 59.3\% (76.3\% at the book's composition), above the pre-committed threshold at which Section~\ref{sec:robustness-crossdesign} concludes the contrast cannot, on that test alone, be cleanly attributed to modeling paradigm; and the path-level timing question remains open for every estimator.

Section~\ref{sec:abm-danish}'s Danish counterfactual retires a superseded figure. The mechanism-extrapolated specification implied a \$925.5 billion institutional gap; recalibrating against \citepos{berger2026} separately-estimated Danish elasticities collapses it by an order of magnitude, and anchoring the transplant so that it changes only the payoff rule (the anchor consistent with the baseline turnover floor's semantics) signs it: $+\$61.2$ billion at the band's zero-refinance edge, 8\% of the benchmark, equal in origin to the lock-in marginal the paper identifies directly, and in size to that marginal at the in-sample calibration point on which both legs were run---though about 44\% larger than the off-window marginal the paper headlines; at the off-window floor itself the rule-only gap is $+\$28.2$ billion, 34\% below that marginal, and subject to the same payoff-credit incidence as the in-sample figure, since that leg too credits retired balance at par (Section~\ref{sec:pathb}). The sign is anchored rather than estimated: pinning the Danish leg at the U.S. zero-gap hazard guarantees it against a book that is out of the money almost everywhere, and the refinance-in-place sweep can only reinforce it, while the bracketing anchor that also imports Danish baseline mobility reverses it to $-\$99.9$ billion---on a Danish leg whose 3.39\% mean prepayment sits below the turnover floor the U.S. legs impose, which is the inconsistency that demoted it to a bracket. This does not disturb Sections~\ref{sec:abm} and~\ref{sec:hazard}'s central contrast, which never depended on the Danish comparison, but it leaves the payoff-rule question at a modest magnitude and a stated anchor rather than an unidentifiable sign: a Danish-style payoff rule would have relieved a real but small share of the Federal Reserve's cash-flow problem, not the \$925.5 billion an extrapolated mechanism once implied. One qualification travels with that magnitude, and it is now priced as a bracket rather than left open. The Danish leg credits retired balance at par, so the cash a market-value buyback would \emph{not} return to the Federal Reserve is absent from it; on the paper's own pricing proxy that channel exceeds the gap and moves it toward zero (Section~\ref{sec:pathb}). Because the pipeline cannot distinguish a cash haircut from a balance adjustment, the pre-committed bracket of Section~\ref{sec:pathb} (run \texttt{buyback\_credit\_bracket}) prices both: the gap stands at $+\$61.2$ billion under the balance-adjustment reading and reverses to $-\$89.4$ to $-\$117.7$ billion under the cash-haircut reading at the paper's own proxy discounts, so the honest reading is that the relief stands on the benchmark's own face denomination, while as reserve drain the rule costs the Federal Reserve exactly the par windfall it would otherwise have extracted from moving households---legs (2) and (3) of the trade-off are, on that reading, the same dollars seen from two sides. The mobility relief such a rule would provide remains the better-supported and larger effect.

A plausible reason the U.S. does not simply adopt the Danish payoff rule is the remaining leg of the trade-off, secondary-market liquidity: the Danish market-value repurchase system is difficult to reconcile with the U.S. fixed-income market as presently structured. The Danish system relies on strict one-to-one match funding (the Balance Principle); the United States relies on the To-Be-Announced (TBA) forward market, which requires highly homogenized, pooled mortgages to generate secondary-market liquidity. \citet{vickery2013} show that this homogeneity produces a measurable liquidity premium: TBA-eligible mortgage pools trade at significantly lower cost than otherwise comparable specified pools ineligible for TBA delivery. \citet{gao2017} measure the same differential directly in transaction data, and \citet{campbell2013} frames the U.S.--Danish contrast as a mortgage-design problem in exactly these terms. A market-value repurchase feature of the kind used in Denmark prices each payoff at the market value of the funding bond series --- a series-level price common to every borrower funded by that series on a given day, not a borrower-specific value --- so what it strains is not homogeneity itself but its \emph{locus}: Denmark standardizes at the level of large, tap-issued bond series under one-to-one match funding, where the TBA market standardizes at the level of the deliverable pool. The Danish callable market's own depth shows that market-value redemption and liquid secondary trading coexist under the series convention \citep{berg2018}; the binding cost for the United States is migrating its issuance architecture from pool-level to series-level standardization, not a liquidity impossibility. Adopting the Danish payoff rule without that broader redesign of agency securitization rules would therefore risk sacrificing a documented source of secondary-market liquidity in the transition. A hedging-side literature sharpens what the two payoff rules allocate. In the U.S. structure, prepayment risk is borne and priced by specialized MBS investors rather than by diversified households \citep{gabaix2007}, and the aggregate hedging of the duration swings that prepayment creates feeds back into the level and volatility of long-term rates \citep{malkhozov2016}; the behavioral component of prepayment is the part that cannot be hedged away, since behavioral uncertainty renders the market for the embedded option incomplete and even robust replication strategies can only bound the resulting exposure \citep{perotti2024}. The Danish market-value rule does not remove that behavioral residual (the borrower retains the par call on the rate-decline side), but it removes the rate-rise extension leg on which this cycle's exposure accrued, so the choice between payoff rules is also a choice about which side of the convexity profile the imperfectly hedgeable behavioral risk sits on, and at what feedback cost to the yield curve.

Ultimately, U.S. policymakers weigh three considerations, and this paper's evidence, under its production floor form (the additive form roughly doubles the margin at issue), is that they do not form the strict trilemma a three-way ``trade-off'' would imply: (1) deep, homogenized secondary market liquidity (the TBA market); (2) continuous household mobility (avoiding the lock-in effect); and (3) convexity-neutral duration profiles (predictable institutional cash flows). The evidence of Sections~\ref{sec:abm} and~\ref{sec:hazard} is that the U.S. has prioritized (1) at real cost to (2): household mobility is genuinely constrained, while the cash-flow shortfall against the caps is mostly mechanical --- a composition the Federal Reserve staff's own record and the ex-ante projection corroborate as composition rather than as numerical agreement, and not independently of it, since the projection may already have embedded part of the channel (Section~\ref{sec:lit}) --- with the lock-in channel contributing an identified margin of $+2.9$ to $+8.7$ points at the off-window floor, $+5.6$ at the mid-grid anchor inside it ($+9.2$ at the in-sample calibration point). What I no longer find support for is the assumption that (2) and (3) trade off sharply under face accounting, in which fixing mobility would cost face-value cash flow: under the rule-only transplant, Section~\ref{sec:abm-danish}'s recalibrated Danish counterfactual finds a market-value repurchase rule would have relieved (2), while its effect on (3) is incidence-conditional: an $\approx\$61$ billion improvement under face accounting (the figure that images the lock-in marginal at the production floor on which both were run), a $-\$89.4$-to-$-\$117.7$ billion cost under cash accounting at the committed proxy discounts, that cost being the foregone household-to-bondholder transfer (run \texttt{buyback\_credit\_bracket})---so this leg-(2)/(3) reading inherits the marginal's floor-and-band conditionality (Section~\ref{sec:identification}): its magnitude ranges with the $+2.1$-to-$+13.2$-point calibration sweep, it is stated at the in-sample calibration point rather than the off-window floor the headline marginal now uses, and its (3) sign is robust to the floor but not to the payoff credit's incidence, nor to bundling Danish baseline mobility into the transplant (Section~\ref{sec:abm-danish}). What this dissolution rests on is thinner than the framing suggests, so I state it plainly. The rule-only Danish gap is the lock-in marginal viewed from the counterfactual side, so ``(2) and (3) do not trade off sharply'' reduces to ``the marginal is small''---and under the production floor form the marginal is bounded small by construction; under the disclosed additive form, which never censors the elasticity, the same counterfactual is about $+11$ points and nearly floor-invariant (Section~\ref{sec:robustness-floor}). Because \eqref{eq:pathB} combines the turnover floor with the voluntary hazard by a hard maximum, switching the elasticity off can raise the hazard by no more than the excess of the $\beta_1 = 0$ hazard over the floor, and by nothing at all where the floor binds on the $\beta_1 = 0$ leg itself, 14.3\% of loan-months at the in-sample calibration point (Section~\ref{sec:identification}). The dissolution is therefore a statement about a calibrated ceiling, not an independent finding, and it is form-conditional. What survives across forms is that the trade-off is modest, not that it is absent. The binding constraint against adopting such a rule is therefore, on either form, not a mobility--cash-flow trade-off but the TBA liquidity premium described above. That last step is an argument from the cited literature rather than a result of this design: nothing here prices a series-level call, so the investor-side cost of the transplant is named and left unpriced, and the constraint is asserted on others' evidence, not on mine. This is the two-sided exchange the paper opened with: (1) purchased in the currency of (2), with a small side payment in (3) whose sign depends on the buyback's incidence. As long as the U.S. maintains the par-payoff MBS structure within the TBA framework, severe duration extension should be expected to recur in future tightening cycles that begin, as this one did, with a large outstanding stock of deeply in-the-money, below-market-coupon mortgages (the state dependence documented by \citet{berger2021} and \citet{eichenbaum2022} makes the extension conditional on that stock, not on tightening per se; the recurrence claim is itself an extrapolation beyond what this design tests, which is one episode). The duration extension itself, however, is mostly the mechanical arithmetic of scheduled amortization and baseline turnover against caps set above any plausible prepayment path; an alternative payoff rule would have trimmed it by the lock-in margin, not eliminated it, and would carry the liquidity cost above.

\subsection{Limitations and Avenues for Future Research}\label{sec:limitations}

Several limitations remain on both sides of the framework.

\textbf{Housing supply as a competing channel}: The mobility margin this paper imports is estimated in an environment of historically tight for-sale inventory, and housing supply enters here only as a friction scaler on the moving decision, never as a channel in its own right. A household that does not move because nothing is listed is observationally identical, in these data, to one that does not move because its coupon is below market, and this design does not separate them. The direction is argued rather than measured: a binding supply constraint would absorb some of what is attributed to lock-in, which is why the scaled-null exercise that removes the non-rate share of the moving decision is reported alongside the headline rather than as a robustness footnote (Section~\ref{sec:identification}; \citealp{aladangady2024,hsieh2019}). No inventory series is used to price it.

\textbf{Distributional silence}: This paper measures a book-wide institutional cost and never asks who bears it. Ginnie Mae borrowers---20.4\% of SOMA face, and the population whose structurally faster speeds Section~\ref{sec:pathb} documents---sit outside the Freddie Mac sample entirely, so nothing here establishes whether the mobility cost falls differently on FHA and VA households than on the conventional borrowers the sample contains. I record that as a distributional limitation of the design, not merely as a coverage bound on the estimate.

\textbf{Sample scale and comparability}: The ABM's 10,000-household synthetic population and the hazard framework's 75,000-loan real Freddie Mac sample differ in both size and data source. Section~\ref{sec:robustness-scale} settles the size question on the production specification (population size changes precision, not central tendency), and the completed cross-design pair of Section~\ref{sec:robustness-crossdesign} addresses the data-source question asymmetrically: the hazard framework's aggregate recovery is largely data-source-independent, while the ABM's is not.

\textbf{Data:} Freddie Mac 2017--2021 loan-level data (20 origination quarters), used for both hazard paths, is integrated via a dedicated ingest pipeline; where real data files are unavailable, results are validated against a synthetic fixture before real-data estimation. Fifteen-year MBS, approximately 9\% of SOMA face value, remain excluded from cohort modeling on the hazard side only; the ABM side folds them in (Appendix~\ref{sec:robustness-pipeline}).

\textbf{Model:} On the ABM side, the monthly CPR path fit remains weak (raw $R^2$ of $-6.984$) despite a directionally sensible level response to rates, and the production Danish counterfactual is a rule-only transplant onto U.S. households (Section~\ref{sec:abm-danish}): it does not estimate how Danish institutional frictions would themselves migrate, and the Danish-level bracketing anchor shows the gap's sign is not robust to bundling baseline-mobility differences into the transplant. The debt-to-income check remains front-end only, without a measure of total household debt. On the hazard side, Path A's Poisson pseudo-maximum-likelihood estimator requires a mild ridge penalty to remain well-conditioned at 295 stratum fixed effects (a stabilization device: Section~\ref{sec:patha} shows the prior-spec (v3) coefficients equal the unpenalized fit on the estimable strata), and its holdout RMSE (spec v3: approximately 38 percentage points unweighted, 2.5 points exposure-weighted; the frozen spec v4 fit re-scored on the same holdout: 37.9 and 3.0) together with a lag-0 CPR correlation of $-0.378$, more negative even than the ABM's $-0.318$, confines its evidential contribution to the aggregate dollar level, as Section~\ref{sec:patha} details; Path B's literature-calibrated hazard magnitudes, while directionally and quantitatively closer to the empirical benchmark than any ABM specification, are simulated on the 75,000-loan real Freddie Mac sample with elasticities calibrated from secondary literature rather than estimated jointly with the empirical benchmark itself, and its recovery percentage is reported across the calibration band (105.9\%--108.2\%) and now also with loan-level sampling uncertainty from an ex-ante-specified stratified bootstrap over the 75,000-loan draw (Section~\ref{sec:pathb}). An external second-agency replication (Fannie Mae loan-level data, same pipeline end-to-end) reproduces the lock-in marginal at $+8.68$ points against Freddie's $+9.20$, inside its pre-committed envelope: a sign-and-envelope replication, not a hazard-level match, and one whose envelope gate is 11.06 points wide, so the 0.52-point agreement, rather than the gate, is what should be read from it (Section~\ref{sec:pathb}). One global fact about the design, each of whose components is disclosed only locally: no outcome-holdout months exist anywhere in this paper. The nearest construction is the marginal's temporal holdout (Section~\ref{sec:robustness-floor}), in which the turnover floor is calibrated on the window's first nineteen months and the central-minus-null difference is then summed over the remaining twenty-three alone, months that enter that calibration nowhere; it returns $+\$45.1$ billion. That construction holds for the marginal, not for the level, and it is an input-stability check rather than an outcome validation: realized prepayment from the held-out months enters neither simulated leg, so it cannot falsify the marginal against data, and the statement above stands without exception. The production baseline turnover floor is still calibrated on 2023--24 turnover inside the evaluation window (the pre-committed off-window re-measurement of Section~\ref{sec:robustness-floor} finds the defensible off-window level higher than 4\%, so the direction of any anchoring error shrinks rather than inflates the identified marginal, and the headline marginal is now reported at that off-window level rather than the in-window one); Path A trains on 19 of the 42 QT months and its ridge-selection holdout doubles as its out-of-sample diagnostic; and Path B's \emph{level}, as distinct from its marginal, is still set by a floor measured in the window on which it is evaluated. Every recovery percentage in the paper is accordingly an in-sample or calibrated quantity in the months dimension, disciplined by ex-ante-specified sweeps, nulls, parity gates, and bootstraps rather than by clean holdout months. All headline figures are frozen in tagged run manifests, with reproduction details recorded in the run ledger (Appendix~\ref{app:ledger}). The marginal itself, at either floor, is a model-based decomposition conditional on the disclosed hazard structure (Section~\ref{sec:identification}): causal identification in the econometric sense would require exogenous variation in rate gaps across otherwise-comparable loans, which this design does not exploit.

\textbf{Status:} The Danish-regime recalibration (Section~\ref{sec:abm-danish}) calibrates both frameworks' Danish CPR response against \citepos{berger2026} separately-estimated moving and refinancing elasticities rather than an extrapolated U.S. mechanism, and anchors the rule-only transplant at the U.S. zero-gap intercept, collapsing the institutional gap by an order of magnitude to a signed $+\$61.2$ billion. Path B's full 5.5\%--7.7\% elasticity band is propagated through the microsimulation (Section~\ref{sec:pathb}), an exercise that also surfaced and corrected a units error in the rate-gap coefficient's application (Appendix~\ref{sec:robustness-pipeline}). The symmetric companion cross-design test is complete (Section~\ref{sec:robustness-synthesis}): the hazard framework's survival structure recovers 106.0\% of the benchmark on a fully synthetic population, closely matching its real-data recovery, while its path-level timing does not transfer to synthetic data. Four items are bound to the pre-submission freeze as a hard gate: (i) executed: the calendar-month specification of Section~\ref{sec:patha} is production Path A as of this freeze, with Table~\ref{tab:estimators} and the WAL table restated under it, while the coefficient-uncertainty table and the refit-and-resimulate interval remain labeled under the prior spec v3 because bootstrap replications under the seasonal design are branch-unstable (Section~\ref{sec:patha}; on the production panel no start reaches a better penalized objective, run \texttt{pathA\_cold\_starts}); (ii) executed for the LaTeX master: it carries no hardcoded cross-references (machine-checked by the replication package's liveness gates), builds clean with zero unresolved references or citations, and the build log ships with the submission bundle; the Word REF-field conversion applies when a Word submission master is cut from this source; (iii) executed: a claims-versus-code liveness audit ran at the freeze (77 production-mechanism claims traced adversarially to their executing code paths; the nine that did not trace as written are corrected in this draft or scoped to their frozen specifications, and the audit artifact ships in the replication package)---the defect class earlier caught in the Danish counterfactual's description; and (iv) the timing-sweep generator's golden-fixture test must pass at the freeze, three matcher generations having each produced a distinct silent defect. What remains open beyond that gate: adding a seasonal multiplier to Path B's baseline hazard (Section~\ref{sec:pathb}); a total-debt measure, or threshold sensitivity, for the ABM's front-end DTI gate (Appendix~\ref{sec:method-frictions}); and the monthly-frequency question the timing diagnosis leaves (the empirical path's variation at that frequency is seasonal, not rate-driven, and no estimator here models it). The second-agency replication is complete (Section~\ref{sec:pathb}; the envelope caveat is in the Model paragraph above). The cyclical-floor variant is likewise complete, with a correction to how its output should be read: a floor co-varying within the run with the window-standardized mortgage rate, nesting the production specification at $\kappa = 0$, moves the marginal between $+7.1$ and $+9.4$ points across $\kappa \in [-0.5, +0.5]$, every cell inside the constant-floor sweep of Section~\ref{sec:robustness-floor}; but 85.0\% of that span survives scrambling the floor path's time order and 98.0\% of it vanishes under the additive combination rule, so the span measures floor dispersion under the hard maximum rather than cyclicality (Section~\ref{sec:pathb} and Appendix~\ref{sec:robustness-seasonalfloor}). The co-movement component that survives is real, monotone in $|\kappa|$, and at most \$5.29 billion. The Ginnie composition bound has been upgraded from two published snapshots to a time-varying overlay built from the published monthly CPR series, which leaves the conventional-share marginal at $+7.3$ points, positive under every variant, with the Ginnie-specific level contribution inside the static bound (Section~\ref{sec:pathb}). Replication alone could not retire the vintage residual: the 2022 and pre-2017 vintages (23.1\% and 10.6\% of book face) lie outside the 2017--2021 universe of both agencies' samples used here. The observed-speed bound of Section~\ref{sec:pathb} now disciplines that extrapolation directly (\$11.7 billion, 1.5\% of benchmark, signed toward overstating trapped liquidity), so what remains open on this dimension is no longer the absence of a bound but the bound's own proxy assumption: observed Fannie segment speeds standing in for the SOMA book's, the same single-agency posture the production calibration itself carries. A modeled overlay in place of that proxy is not constructible under the committed sampler: the staged replication sample enforces the 2017--2021 vintage universe, the pre-QT snapshot rule leaves second-half-2022 originations with no admissible observation month, and the acquisition window closes at 2022Q4, so 2022 originations acquired in 2023 were never staged; the 2022 leg of the observed-speed bound (\$3.3 billion of the \$11.7 billion) therefore stands as that vintage's discipline. Appendix~\ref{app:composition} collects the formal composition comparison behind these bounds: sample versus book by agency, vintage, coupon, and term, with a cross-agency stability check on the dimensions the book's CUSIP-level disclosure cannot supply.

Sections~\ref{sec:abm} and~\ref{sec:hazard} establish that contrast under the synthetic household population this paper specifies. Section~\ref{sec:robustness}'s cross-design pair leaves the attribution asymmetric and inherits Section~\ref{sec:robustness-crossdesign}'s own hedge: the hazard framework's recovery survives a fully synthetic population (106.0\% versus 107.0\%), evidence that its survival structure, not its data source, carries the aggregate level, but the recalibrated real-covariate ABM's 59.3\%, and its 60.2\% fifty-seed mean, cross the pre-committed threshold above which the contrast cannot, on that test alone, be cleanly attributed to modeling paradigm rather than to calibration and data source; the path-level timing question remains open for every estimator. That result cuts both ways: recovering 106\% on a fully synthetic population shows the calibration constants, not real loan-level structure, deliver the level, which is consistent with the mechanical decomposition above rather than with a loan-level-mechanisms reading.

\newpage
\bibliographystyle{plainnat}
\bibliography{references}

\clearpage
\appendix

\noindent\textbf{Online Appendix.} The appendices that follow, through Appendix~\ref{app:verdicts}, are for online publication.

\section{Superseded Figures and Run Ledger}\label{app:ledger}

\paragraph{Data and code availability.} The replication package is the public repository at \url{https://github.com/eugeneoCMU/Lock-in-Effect}. It ships all estimation, simulation, and figure code; the frozen 75{,}000-loan Freddie sample and the cohort-month panel that every production run consumes; every committed results artifact this paper cites (the results JSONs, the production coefficient artifact, the exposure-weighted delinquency-transition matrix); the tagged run manifests with SHA-256 pins on their artifact files; and the machine-checked liveness gates and the claims-versus-code audit artifact. Hazards are estimated on the Freddie Mac Single-Family Loan-Level Dataset, which Freddie Mac makes publicly available at no cost after registration; the raw origination and performance files are not redistributed, and the committed code rebuilds every derived object from them under fixed seeds (engine seed 42 with per-regime offsets; every post-freeze run's specification is fixed ex ante in its script header, the convention this ledger records). Macro and portfolio inputs are the public FRED and New York Fed SOMA interfaces; a user-supplied FRED key is the only credential involved. Fannie Mae inputs are code-only, as Section~\ref{sec:pathb-fannie} states: no Fannie data and no derived cells ship, and the committed replication artifact carries headline statistics alone.

Headline figures in this paper changed across drafts as specification errors were found and corrected. Because some superseded values appear in earlier internal drafts and in the accompanying technical documentation, I catalogue them here rather than in scattered footnotes, so that a reader consulting both always knows which figure superseded which, and why.

\begingroup\sloppy
\begin{itemize}
\item \textbf{U.S. CPR 12.3\%, empirical CPR 5.45\%, wedge 35.1pp}: produced by a window filter that inadvertently included eight post-QT months; superseded by the bounded active-QT window (Section~\ref{sec:robustness-benchmark}).
\item \textbf{17.6pp wedge; 21--27\% Danish CPR}: a pre-multi-cohort, rational-only calibration; superseded by the multi-vintage production specification.
\item \textbf{\$101.2 billion / 13.2\% ABM estimate} (\texttt{run-2026-07-04}): the pre-fold-in, 30-year-only tag, raw path $R^2 = -6.443$; superseded by the 15-year fold-in production tag \texttt{run-2026-07-04-15yr-foldin} (\$91.0 billion / 11.9\%, $R^2 = -6.984$; Appendix~\ref{sec:robustness-pipeline}).
\item \textbf{Danish CPR 47.1\%; wedge $\approx$35.5pp; institutional gap \$925.5 billion}: the mechanism-extrapolated Danish specification; the 35.5pp wedge is the difference between its 47.1\% Danish and 11.68\% U.S. mean CPR. Superseded by the Berger-recalibrated \texttt{run-2026-07-05-berger} (Danish CPR $\approx$3.4\%; institutional gap $-\$728.4$ billion against that run's own \$84.5 billion U.S. leg; Section~\ref{sec:abm-danish}).
\item \textbf{``recovers 97.9\%, within 2.1\% of the benchmark'' as an unqualified headline}: the recovery level and the lock-in marginal were quoted at different turnover-floor calibrations after the marginal was demoted to the off-window floor. Both are retained, each labelled with its floor: 91.3\% recovery and an 8.7-point miss at the headline off-window floor (88.2--93.7\% across its range), 97.9\% and 2.1 points at the in-sample calibration point (\texttt{calibration\_reconciliation}). The related 0.2-point agreement between the null's 88.7\% and the Fed's 88.5\% ex-ante projection share is retracted as a corroborating coincidence: it requires both the in-sample floor and the uniform-spread allocation.

\item \textbf{18.3-point real-book central-minus-null gap}: differenced a standalone-scorer central leg against a shared-basis null. The single-basis gap is 9.2 points in-sample and 5.6 points at the off-window floor.

\item \textbf{Path B \$747 billion / 97.7\%}: the pre-unit-correction microsimulation, in which the calibrated lock-in elasticity was numerically inert; superseded by \$818.5 billion / 107.0\% (Appendix~\ref{sec:robustness-pipeline}). Its corrected analogue, the deliberate $\beta_1 = 0$ null of Section~\ref{sec:hazard-interp}, recovers 97.8\%.
\item \textbf{Path A \$915.0 billion / 119.7\%} (spec v3; $r = -0.444$ at lag 0, peak $-3$): the prior production specification, which omitted calendar-month seasonality; superseded by the calendar-month spec v4 adopted at the pre-submission freeze on its timing behaviour, the peak moving to $-2$; the recovery difference between the two specifications is not calendar-attributable and no part of the adoption case rests on it (\$928.9 billion / 121.5\%; \texttt{run-2026-07-14-pathA-seasonal}; Section~\ref{sec:patha}).
\item \textbf{Hybrid Danish gap $-\$99.9$ billion, ``sign not robustly identified''} (Danish-level anchor; Danish mean CPR 3.39\%): superseded as the production reading by the U.S.-intercept transplant (\texttt{danish\_us\_intercept}: gap $+\$61.2$ billion, Danish mean CPR 5.61\%, positive at every refinance-in-place sweep point, that positivity being forced by the zero-gap anchor; the U.S. leg reproduces the committed central artifacts exactly). The Danish-level rows are retained in Table~\ref{tab:danish} as a labeled rule-plus-country bracketing case; their Danish mean CPRs sit below the 4\% turnover floor, the inconsistency that motivated the restatement (Section~\ref{sec:abm-danish}).
\item \textbf{Post-freeze robustness artifacts}: the floor functional-form test (\texttt{hazard/floor\_form\_results.json}: additive-form marginal $+11.25$pp at the production floor, $+11.29$ to $+11.21$ across 3--5\% floors; max-form parity pair reproduces the committed artifacts exactly; Section~\ref{sec:robustness-floor}) and the Path B stratified loan-level bootstrap (\texttt{hazard/bootstrap\_pathb\_results.json}; Section~\ref{sec:pathb}). Neither changes a production headline; both are specified ex ante in their script headers, written before the runs executed. The seasonal-floor run follows the same convention (\texttt{hazard/seasonal\_floor\_timing\_results.json} and the placebo scoring in \texttt{hazard/data/b3\_timing\_scores.json}; Appendix~\ref{sec:robustness-seasonalfloor}), and carries one correction to its own provenance: an earlier module docstring described the twelve-cell calendar profile as estimated in-window, when it reproduces on the full cohort-month panel. Both bases are recorded in the artifact, the full-panel basis drives every run, and no figure elsewhere in this paper consumed the profile.
\item \textbf{``log-normal mobility-desire heterogeneity'' (scale-sensitivity description)}: the production ABM draws mobility desire from an exponential distribution on a dollar scale; the log-normal description was an error of description, not of code, caught by an adversarial claims-versus-code check and corrected in place (Section~\ref{sec:robustness-scale}).
\end{itemize}
\endgroup

Every current headline figure cites a frozen, tagged run manifest; two freezes ran on uncommitted working trees, so manifest git hashes alone do not name an exact reproduction checkout, and the repository's run ledger records the exact reproduction commits. The one that carries a headline quantity should be named here rather than left to the ledger: the fold-in Monte Carlo behind the ABM's seed distribution (\texttt{run-2026-07-04-15yr-foldin}) was frozen about two minutes before its own code was committed, so its manifest's commit field points at the \emph{previous} state of the tree. Its reproduction checkout is recorded in the ledger, and on that run, specifications must be identified by cohort-bucket count---seven for the thirty-year-only specification, eleven for the fold-in---rather than by the manifest commit, which is unreliable for exactly this reason. Anyone reproducing the ABM's 13.6\% seed mean or its 11.9\% frozen draw needs that rule, and it is the reason the two figures were once mistakenly attributed to the native-fifteen-year specification.

\paragraph{Reproducibility conventions.} Provenance details referenced from the main text are collected here. The ABM headline run's manifest lives at \path{abm/data/runs/run-2026-07-04-15yr-foldin/manifest.json} and is reproduced via \texttt{freeze\_run.py -{}-tag run-2026-07-04-15yr-foldin} against the fold-in specification's code; a manifest's recorded \texttt{git\_commit} is the pre-run parent of the freezing commit (the convention the replication package's technical documentation records). Post-freeze runs follow the spec-before-run convention: the grid, parity gates, and interpretive thresholds are fixed in the run script's header and committed before execution (for the floor sweeps, spec commits \texttt{ad52db6} and \texttt{651d1a0} precede result commits \texttt{d0ef130} and \texttt{57181b3}); where script and artifact enter the repository in one revision (the form test \texttt{hazard/floor\_form\_test.py}, the bootstrap), the header is the ex-ante record. Referenced artifact files not named in the main text: \texttt{wal\_anchor\_results.json} (WAL anchor), \texttt{floor\_sweep\_results.json} and \texttt{floor\_band\_cross\_results.json} (the calibration box), \texttt{theil\_data.json} (dynamic-fit decomposition), \texttt{vintage\_residual\_bound\_results.json} (vintage bound), and \texttt{ml\_comparator\_holdout\_results.json} (flexible-learner comparator).

\begin{table}[H]
\centering
\footnotesize
\begin{threeparttable}
\caption{Basis-and-inputs crosswalk for headline quantities. Bases are defined in Section~\ref{sec:method-benchmark} and collected numerically in Table~\ref{tab:bases}; every model quantity descends from a frozen, tagged run manifest (reproducibility conventions above).}
\label{tab:crosswalk}
\begin{tabular}{@{}p{4.9cm}p{3.4cm}p{2.2cm}p{3.9cm}@{}}
\toprule
Quantity & Accounting basis & Home exhibit & Source run / input \\
\midrule
Cap-relative benchmark \$764.7B & accounting identity on the actual-holdings path & Section~\ref{sec:method-benchmark} & FRED WSHOMCB + cap schedule \\
Expectations complement \$87.8B (11.5\%); ex-ante share 75.6--88.5\% & expectations-based complement & Section~\ref{sec:method-benchmark} & \texttt{expectation\_benchmark} \\
$\beta_1 = 0$ null 88.7\% (97.8\% standalone), in-sample calibration; 85.7\% (94.8\% standalone) at the off-window floor & shared (standalone in parentheses) & Table~\ref{tab:bases} & Path B null leg; off-window legs \texttt{calibration\_reconciliation} \\
Lock-in marginal, off-window floor, $+5.6$pp $= +\$42.6$B (headline) & basis-invariant & Eq.~\eqref{eq:marginal}; Table~\ref{tab:oosfloor} & central $-$ null paired legs at the off-window floor; \texttt{oos\_identification} \\
Binding layer on the headline marginal, $[+2.9, +8.7]$pp & basis-invariant & Table~\ref{tab:uncertainty}; Section~\ref{sec:robustness-floor} & \texttt{floor\_inference\_correction} (wild-cluster bootstrap-$t$ over the \texttt{floor\_uncertainty} replicate scheme) \\
Lock-in marginal $+9.2$pp $= +\$70.3$B (in-sample point) & basis-invariant & Eq.~\eqref{eq:marginal} & central $-$ null paired legs; \texttt{marginal\_decomposition}, \texttt{marginal\_monthly\_decomposition} \\
Path B central 97.9\% (107.0\% standalone; 100.0\% composed, mixed basis --- 98.1\% converted), in-sample calibration; 91.3\% (100.4\% standalone) at the off-window floor & shared / standalone / composed & Table~\ref{tab:bases} & production Path B leg; off-window legs \texttt{calibration\_reconciliation} \\
Path A 112.4\% (121.5\% standalone) & shared / standalone & Table~\ref{tab:bases} & spec v4 (non-headline corroboration) \\
Production ABM 13.6\% seed mean (11.9\% frozen draw) & shared layer, native & Section~\ref{sec:abm-results} & manifest \texttt{run-2026-07-04-15yr-foldin} \\
Danish rule-only gap $+\$61.2$B & rule-only counterfactual gap (U.S.-intercept anchor; face denomination) & Table~\ref{tab:danish} & \texttt{danish\_us\_intercept}; incidence bracket \texttt{buyback\_credit\_bracket} \\
Ginnie overlay marginal $+7.3$pp; differential $+\$38.5$B & shared & Section~\ref{sec:pathb} & \texttt{ginnie\_cpr\_overlay} \\
Vintage bound \$11.7B (1.5\%) & observed-speed bound & Section~\ref{sec:pathb} & \texttt{vintage\_residual\_bound} \\
\bottomrule
\end{tabular}
\end{threeparttable}
\end{table}

\paragraph{Robustness and diagnostic runs.} Headline-quantity runs are in Table~\ref{tab:crosswalk}; the diagnostic and robustness runs are indexed in Table~\ref{tab:runindex}.

{\footnotesize
\begin{longtable}{@{}p{5.4cm}p{8.8cm}@{}}
\caption{Robustness and diagnostic run index. Each row names a committed run cited in the main text by its result; tags are collected here rather than in the narrative.}
\label{tab:runindex}\\
\toprule
Run tag & Produces \\
\midrule
\endfirsthead
\toprule
Run tag & Produces \\
\midrule
\endhead
\texttt{wal\_normal\_turnover} & Portfolio WAL at the off-window turnover-floor read: the normal-turnover row of Table~\ref{tab:wal} (Appendix~\ref{sec:robustness-wal}) \\
\texttt{interp\_spot\_check} & CPR-surface interpolation spot-check (ABM specification, Appendix~\ref{sec:method-abm}) \\
\texttt{expectation\_spread\_variants} & Intra-2022 allocation ramp bounds on the ex-ante share (Section~\ref{sec:method-benchmark}) \\
\texttt{danish\_discount\_bound} & Danish buyback discount-rate sensitivity diagnostic (Section~\ref{sec:abm-danish}) \\
\texttt{floor\_inference\_correction} & Wild-cluster/CR small-sample correction of the floor-read sampling layer; sets the binding interval (Section~\ref{sec:robustness-floor}) \\
\texttt{band\_low\_extension} & Marginal-versus-elasticity curve below the adopted band at both anchor floors (Section~\ref{sec:robustness-floor}) \\
\texttt{moving\_share\_bracket} & Gap response confined to a share of the voluntary hazard; construct-mismatch bracket (Section~\ref{sec:pathb}) \\ \texttt{moving\_share\_bracket\_offwindow} & The same bracket re-run at the 4.991\% headline floor under a pre-fixed landing classifier (Section~\ref{sec:pathb}) \\ \texttt{psa\_level\_sweep} & Baseline seasoning ramp swept in level, 75--150 PSA, paired legs at both floors (Section~\ref{sec:identification}) \\ \texttt{covariate\_ablation\_offwindow} & Covariate endpoints re-run at the off-window floor, paired legs (Section~\ref{sec:pathb}) \\ \texttt{danish\_refi\_finegrid} & Rule-only gap swept over refinance-in-place 0--3\% CPR in 0.25\% steps under the production anchor; band presentation per the redemption-validation rule (Section~\ref{sec:abm-danish}) \\ \texttt{floor\_form\_mixture} & The floor's functional form as a measured curve in its strictly-involuntary share, endpoints the committed max and additive forms (Section~\ref{sec:robustness-floor}) \\ \texttt{bootstrap\_pathb\_cluster} (floor 4.0\%) & The stratum-cluster loan-sampling scheme re-run at the in-sample calibration; supersedes the within-stratum read there (Section~\ref{sec:pathb}) \\ \texttt{attenuation\_sensitivity} & Survival-selection attenuation of the imported elasticity, frailty-identity mapping at the measured pre-window survival share (Section~\ref{sec:pathb}) \\
\texttt{scaled\_null\_housing\_activity} & Housing-activity-scaled seasoning baseline calibrated to the Aladangady 44\% share; the assembly's upward-bias entry (Section~\ref{sec:identification}) \\
\texttt{h1\_zero\_months\_diagnosis} & Mechanism trace of the benchmark's four exact-zero months: clip-induced boundary artifacts (Appendix~\ref{sec:robustness-seasonalfloor}) \\
\texttt{layer\_convolution} & The floor-read and loan/stratum sampling layers convolved as independent draws; labeled second line, not a replacement (Section~\ref{sec:identification}) \\
\texttt{pathA\_ridge\_v4} & The ridge-device verification battery recomputed under spec v4; all three facts survive (Appendix~\ref{app:ridge}) \\
\texttt{pathA\_cold\_starts} & 150-start dispersed battery on the production panel: no better penalized optimum exists; the production basin is small (Section~\ref{sec:patha}) \\
\texttt{coupon\_convention\_reweight} & The full-book reweight and composed cell on a single pass-through convention; the marginal bit-invariant, the composed cell basis-dependent (Section~\ref{sec:robustness-hybrid}) \\
\texttt{coupon\_convention\_amortization} & The CPR back-out at the note-rate WAC: $+0.27$-point deterministic level shift, timing unchanged, dollars bit-identical (Section~\ref{sec:pathb}) \\
\texttt{pathA\_resimulate\_joint} & The refit-and-resimulate distribution rerun with full coefficient vectors at both specifications; floor-censored, reported beside the committed interval (Section~\ref{sec:patha}) \\
\texttt{vintage\_overlay} & The Ginnie-overlay template on the vintage axis: the marginal confined to the in-window-vintage share, $+3.7$ points at $0.66\times$ (Section~\ref{sec:pathb}) \\
\texttt{episode\_confrontation\_within} & The episode gradient composition-standardized and exactly decomposed; the confound's composition limb measured at $+0.77$ of $+4.20$ (Section~\ref{sec:identification}) \\
\texttt{buyback\_credit\_bracket} & Market-value buyback credit priced under both incidence readings (Section~\ref{sec:pathb}) \\
\texttt{ml\_comparator\_holdout} & Gradient-boosted Poisson comparator on the Path A holdout (Section~\ref{sec:patha}) \\
\texttt{landmark\_isotonic\_holdout} & Isotonic-recalibration comparator on the Path A holdout (Section~\ref{sec:patha}) \\
\texttt{floor\_cyclical} & Within-run cyclical turnover-floor variant (Section~\ref{sec:pathb}) \\
\texttt{vintage\_1516\_subleg} & 2015--16 subleg of the out-of-sample vintage bound (Section~\ref{sec:pathb}) \\
\texttt{fannie\_replication} & Fannie Mae external second-agency replication (Section~\ref{sec:pathb-fannie}) \\
\texttt{fonseca\_band\_anchor} & Fonseca elasticity-band corroboration (Section~\ref{sec:pathb}) \\
\texttt{cohort\_timing\_diagnostic} & Cohort-level timing falsification (Section~\ref{sec:pathb}) \\
\texttt{make\_theil\_data} & Theil dynamic-fit decomposition (Section~\ref{sec:pathb}) \\
\texttt{regime\_split\_marginal} & Baseline-break sensitivity of the lock-in marginal (Section~\ref{sec:identification}) \\
\texttt{subgroup\_marginals} & FICO/LTV/region ablation of the marginal (Section~\ref{sec:identification}) \\
\texttt{freeze\_sensitivity} & Wait-and-see freeze isolation sensitivity (Section~\ref{sec:robustness}) \\
\texttt{dti\_threshold\_sweep} & DTI-gate threshold sweep (Section~\ref{sec:robustness}) \\
\texttt{abm\_external\_gates} & ABM external CPR gate check (Section~\ref{sec:robustness}) \\
\texttt{smd\_two\_moment} & Simulated-minimum-distance two-moment fit (Section~\ref{sec:robustness}) \\
\texttt{curtailment\_profile\_demo} & Curtailment accounting-identity demonstration (Section~\ref{sec:robustness}) \\
\texttt{curtailment\_danish\_scaling} & Curtailment Danish-scaling invariance (Section~\ref{sec:robustness}) \\
\texttt{out\_of\_window\_floor} & Out-of-window turnover-floor anchor check (Section~\ref{sec:robustness-floor}) \\
\texttt{oos\_identification} & Out-of-sample identification of the lock-in marginal: contamination-split off-window floor, marginal re-run across the floor $\times$ band grid, and the early-window/held-out-months temporal holdout (Section~\ref{sec:robustness-floor}; Table~\ref{tab:oosfloor}) \\
\texttt{matched\_depth\_reconciliation} & Window-versus-depth decomposition of the floor demotion at matched discount depth, with the ex-ante 10\% exposure-support rule, the extended 2018 depth ladder, and the matched-depth leg comparison (Section~\ref{sec:robustness-floor}) \\
\texttt{seasonal\_floor\_timing} & Calendar profile on the turnover floor: four floor legs, the permutation and rotation placebo families, the convexity-matched flat control, and the frozen timing rule's placebo scoring (Appendix~\ref{sec:robustness-seasonalfloor}; Table~\ref{tab:seasonalfloor}) \\
\texttt{floor\_form\_offwindow} & Additive competing-risks floor form at the three off-window anchors; source of the form-conditional hull (Section~\ref{sec:robustness-floor}) \\
\texttt{concave\_marginal} & Concave gap-transform rerun of the paired legs at both floors (Sections~\ref{sec:pathb} and~\ref{sec:robustness-floor}) \\
\texttt{concave\_additive\_marginal} & Joint concave $\times$ additive cell: hull membership and the form/transform interaction (Section~\ref{sec:robustness-floor}; Table~\ref{tab:uncertainty}) \\
\texttt{floor\_uncertainty} & Cluster bootstrap of the off-window floor reads, propagated to the marginal; age-standardized floor variant (Section~\ref{sec:robustness-floor}) \\
\texttt{bootstrap\_pathb\_cluster} & Stratum-cluster bootstrap of the Path~B legs at the headline floor, against the within-stratum scheme (Section~\ref{sec:pathb}) \\
\texttt{burnout\_ablation} & $\beta_B = 0$ ablation of Path~B at both floors, with the marginal re-differenced (Section~\ref{sec:pathb}) \\
\texttt{danish\_offwindow\_floor} & Rule-only Danish gap re-run at the off-window floor (Section~\ref{sec:abm-danish}) \\
\texttt{ginnie\_overlay\_offwindow} & Ginnie-speed overlay composed with the off-window calibration (Section~\ref{sec:pathb}) \\
\texttt{patha\_sign\_test} & Bias-respecting test of the Path~A rate-gap coefficient's sign: BCa, temporal and permutation constructions (Section~\ref{sec:patha}) \\
\texttt{cross\_design\_reweight} & Cross-design population post-stratified to the SOMA coupon $\times$ vintage composition (Section~\ref{sec:robustness-crossdesign}) \\
\texttt{cross\_design\_seeds} & Fifty-seed sweep of both cross-design variants over both randomization layers, with per-seed floor reads (Section~\ref{sec:robustness-crossdesign}) \\
\texttt{episode\_confrontation} & Realized cumulative prepayment gradient by window-mean gap depth, against the model-implied gradient (Section~\ref{sec:robustness-floor}) \\
\texttt{production\_scale\_test} & Production-specification scale comparison at $N = 10{,}000$ and $N = 75{,}000$, fifty seeds each (Section~\ref{sec:robustness-scale}) \\
\bottomrule
\end{longtable}
}

\section{Agent-Based Model: Specification}\label{sec:method-abm}

I adopt an agent-based approach rather than a representative-agent framework because the phenomenon under study is heterogeneous: households differ in coupon, loan age, and exposure to search and sentiment frictions, and these differences interact non-linearly with the level of interest rates. \citet{geanakoplos2012} argue for, and demonstrate the feasibility of, agent-based models of mortgage prepayment in environments where behavioral and search frictions vary across the population and over time, a property directly relevant here since the frictions I model (housing search costs and consumer sentiment) are themselves time-varying. The broader methodological case for agent-based computational economics in settings with meaningful heterogeneity is set out in \citet{tesfatsion2006}.

The ABM simulates a population of N = 10,000 heterogeneous households. Initial agent endowments (income, home value) are drawn from log-normal distributions calibrated to Federal Reserve Economic Data (FRED) medians. In the production multi-vintage specification of Appendix~\ref{sec:method-frictions}, each household's mortgage coupon and origination date are drawn from the Federal Reserve's actual holdings composition: the live CUSIP-level SOMA parse of Section~\ref{sec:abm-results}, across seven coupon cohorts. A single-cohort variant, used for exposition and mechanism analysis, instead endows every agent with a 30-year fixed-rate mortgage at a 3.0\% coupon aged 60 months at the simulation start; where the text below discusses a single representative coupon, it refers to this pedagogical case, not to the production population.

Section~\ref{sec:robustness} reports a scale-sensitivity replication showing that the gap in explanatory power is not an artifact of sample size; whether it reflects modeling paradigm or data source is taken up in Section~\ref{sec:robustness-crossdesign}.

The decision to relocate turns on a comparison of the utility of moving against the aggregate financial penalty of doing so. The monthly payment is calculated via the standard annuity formula
\begin{equation}\label{eq:annuity}
M = P\,\frac{r\,(1+r)^{n}}{(1+r)^{n}-1},
\end{equation}
where $P$ is the outstanding principal, $r$ is the monthly interest rate, and $n$ is the remaining term in months. The mobility penalty evaluates \eqref{eq:annuity} twice: once at the borrower's locked-in note rate on the original principal and term, which gives the current payment (by the amortization identity, equal to the payment on the outstanding balance at the note rate over the remaining term), and once at the prevailing market rate on the outstanding balance over a fresh full original term---the replacement loan re-amortizes, a term-preference persistence the code implements deliberately---which gives the payment the borrower would take on by refinancing, so the penalty scales with the difference between the two; a native short-term variant of the penalty, active only for terms at or below 240 months, floors payment decreases at zero and is specified with the full decision rule in Appendix~\ref{app:abmspec}. Here $r$ is expressed in decimal monthly terms (the annual note rate divided by twelve; a 3.0\% coupon gives $r = 0.0025$), a rate \emph{level} distinct in both role and units from the annualized rate \emph{gaps} $g_{i,t}$ (in 100 bp units) and $\tilde g_{i,t} = g_{i,t}/100$ (decimal) that drive the hazard framework of Section~\ref{sec:pathb}; no rate symbol is reused across the two. Across a three-dimensional grid of market rate, dynamic friction, and rate velocity (and, in the multi-vintage specification, across seven coupon cohorts), the population's move/stay decisions are evaluated household-by-household at each grid point and aggregated to a pool-level annualized Conditional Prepayment Rate (CPR) surface (CPR is a pool-level convention; the household-level object is an annualized move intensity, consumed by the accounting layer at one-twelfth per month, Appendix~\ref{app:abmspec}), which is precomputed once and interpolated month-by-month against the live population of simulated market conditions, rather than re-solved via per-month household loops. This reduces the cost of the full grid by roughly two orders of magnitude and makes fifty-seed Monte Carlo estimation feasible. The surface stores each node's direct household-by-household evaluation, so interpolation reproduces the micro-simulation exactly at grid points; the operative approximation is the trilinear segment between 50 bp-spaced nodes, and a spot-check validates it directly: after an on-grid parity gate reproduced the committed surface values, direct micro-simulation and interpolation at worst-case off-grid midpoints differ by at most 0.23 points of CPR away from the freeze threshold, with the Danish column exact. The one systematic difference is the wait-and-see freeze itself: the engine applies it as a step at the 150 bp trigger while the surface renders it as a ramp across the adjacent velocity nodes, which binds in four realized months (August--November 2022; SOMA-weighted mean absolute difference 0.22 points of CPR over the window): surface semantics that the production figures share, and whose trapped-liquidity consequence the binary-step leg of Appendix~\ref{sec:robustness-extensions}'s freeze sweep prices directly. Figure~\ref{fig:cprsurface} renders one slice of that surface---the 3.0\% cohort of the expository variant---with the realized QT window located on it.

\begin{figure}[H]
\centering
\includegraphics[width=0.85\textwidth]{fig8_cpr_surface}
\caption{The precomputed CPR surface: the production surface's slice at the 3.0\% thirty-year SOMA cohort, at its parsed seasoning. Left: the U.S. par-payoff surface over the market-rate $\times$ dynamic-friction grid (no-freeze regime), with the 42 QT-window months overlaid at their realized (rate, friction) coordinates; the window sits entirely in the deep out-of-the-money corner, near the baseline turnover floor. Right: rate slices at the friction grid point nearest the window mean. Velocity enters the grid only through the wait-and-see freeze of Appendix~\ref{sec:method-frictions} (20\% of otherwise-mobile households held in place above a 150 bp trailing six-month rise), so the two regimes shown are its entire velocity structure; the Berger-recalibrated Danish counterfactual (Section~\ref{sec:abm-danish}) is near-flat, moving roughly a quarter point across the grid. The U.S. slice is identical across the frozen production surface and the current code head; the Danish curve is the Berger production specification.}
\label{fig:cprsurface}
\end{figure}

\begin{figure}[H]
\centering
\includegraphics[width=0.9\textwidth]{fig10_abm_cpr_paths}
\caption{Production ABM simulated CPR against the empirical SOMA back-out over the QT window (frozen fold-in run, seed 42). The simulated path runs at roughly twice the empirical level and does not track its month-to-month shape; the annotated means, lag-0 $r$, and $R^2$ are the manifest values quoted in the text. This is the level-and-shape mismatch behind the negative path statistics; the cross-correlation structure across lags is in Figure~\ref{fig:ccf}.}
\label{fig:abmpaths}
\end{figure}

\begin{figure}[H]
\centering
\includegraphics[width=0.82\textwidth]{monte_carlo_trapped_liquidity}
\caption{Monte Carlo population-draw distribution of ABM trapped liquidity (50 seeds, fold-in specification; CPR surface rebuilt per draw). The seed-level dispersion (SD \$24.5 billion) dwarfs the distance between the frozen seed-42 production draw (\$91.0 billion, 34th percentile) and the seed mean (\$103.7 billion); the empirical benchmark of \$764.7 billion lies far outside the distribution's support, which is the visual form of the ABM's 10--17\% recovery band.}
\label{fig:mc}
\end{figure}

The simulation contrasts two institutional payoff rules. Under U.S. par-payoff, the cost to retire the mortgage is strictly equal to the outstanding principal, the negatively convex case, in which the investor's upside is capped while extension risk is bounded only by the scheduled-amortization profile: a 14.7-year weighted-average life at zero CPR against the no-shock 3.4-year life (Table~\ref{tab:wal}), an order-of-magnitude extension ceiling rather than an unbounded one. Under Danish market-value repurchase, the payoff cost is the present value of the remaining scheduled payments discounted at the current 30-year primary mortgage rate (FRED MORTGAGE30US), capped at par. Two qualifications attach to this proxy. First, it is an option-free present value discounted at the primary rate: it ignores the embedded par call and any option-adjusted spread, and the primary rate embeds the primary--secondary spread (guarantee fee, servicing, originator margin), whereas the bond-market discounting a true buyback would use sits below it; both choices overstate the borrower's buyback discount as rates rise. These biases inflate the mechanism-substitution Danish CPR of Section~\ref{sec:abm-danish} (the since-superseded 47.1\%) but are exactly immaterial to the Berger-recalibrated production finding: under neither transplant anchor does that leg's prepayment hazard consume this pricing proxy (the Danish-level anchor imports the near-flat moving curve; the production U.S.-intercept anchor evaluates the U.S. hazard at zero rate gap), and a replication-package diagnostic that re-runs the Danish leg with the discount rate shifted down by 25--150 bp moves trapped liquidity by \$0.00 billion, while flipping 2,470--92,495 of the 1,683,082 loan-month present-value classifications\footnote{This Danish-leg classification universe of 1{,}683{,}082 loan-months is 42 fewer than the 1{,}683{,}124 evaluated in the U.S.\ production hazard panel (Section~\ref{sec:pathb}): the two figures count active loan-months in the two counterfactual legs, whose prepayment paths differ marginally even under the shared imported elasticity,} that no downstream quantity in that leg consumes. Second, what the Danish rule neutralizes is extension risk for the bondholder \citep{berger2026}: the borrower retains the par call, so the rate-decline side keeps its price cap, and I use ``convexity-neutral'' below strictly in this rate-rise sense. Under the Danish rule, as market rates rise above 3.0\%, the bond's present value falls, allowing the borrower to extinguish the debt at a discount and guaranteeing continuous cash flow generation. The Danish counterfactual balance is simulated forward on a dynamic, declining-balance basis, so that each month's elevated Danish roll-off is applied to the already-shrunk balance carried over from the prior month and the simulated payoff path compounds correctly over the multi-year QT window.

\section{Macroeconomic Friction and Structural Extensions}\label{sec:method-frictions}

To quantify the behavioral friction gap in standard models, I scale base transaction friction using two FRED series: (1) search friction, measured as the shortfall in Active Listing Count (ACTLISCOUUS) against a 2017--2019 baseline, the last stable pre-pandemic period before the 2020/21 rate-driven frenzy and the 2022 inventory collapse; and (2) psychological friction, indexed from University of Michigan Consumer Sentiment (UMCSENT) against a baseline of 85.0, with readings below that line scored as incremental friction at +5 bp per point short, capped at +150 bp during acute sentiment stress. This produces a dynamic friction parameter (a one-time moving transaction cost expressed as a share of the household's current home value) that scales between 8.23\% and 9.96\% of home value over the 42 active QT months.

Three behavioral extensions were added to the household decision function, none with any parameter tuned to match the empirical benchmark, though their grounding is uneven (the first two are parameterized from external sources, the third an author calibration): a debt-to-income (DTI) constraint that vetoes any move in which the new housing payment would exceed 43\% of gross monthly income, a front-end heuristic motivated by the pre-2021 General QM debt-to-income limit (12 C.F.R. \S{} 1026.43(e)(2)(vi)); the regulatory threshold was a back-end total-debt ratio and was superseded by the price-based General QM definition in 2021, so I use it as a behavioral anchor rather than a live regulatory constraint, and the model carries no total-debt data against which to test the threshold (Section~\ref{sec:limitations}); asymmetric loss aversion that scales perceived payment increases by 2.25$\times$ while leaving decreases unscaled (loss aversion: \citealp{kahneman1979}; the 2.25 coefficient: \citealp{tversky1992}; \citepos{andersen2022} housing-market estimate of roughly 2.5 is close, and because any coefficient in this range triggers the same turnover-floor recalibration described below, the artifact interpretation of Section~\ref{sec:abm-results} does not hinge on the choice); and a wait-and-see freeze that holds 20\% of otherwise-mobile households in place whenever the trailing six-month change in mortgage rates exceeds 150 bp, modeled as a fixed per-agent trait so the CPR surface remains deterministic; the 20\% share and 150 bp trigger are author calibrations with no external source, and the same is true of the sentiment-to-friction mapping and the base transaction-friction level of the preceding paragraph (the UMCSENT and ACTLISCOUUS inputs are data; the $+5$ bp-per-point mapping, its $+150$ bp cap, and the base friction level are mine). An externally parameterized variant of the mobility-desire scale and the freeze share, drawn from the same source as Path B's elasticity, is tested in Section~\ref{sec:robustness-crossdesign}.

Because loss aversion steepens the penalty function facing every household, and the model's mobility-desire parameter is independently calibrated to reproduce a real-world baseline turnover floor of 4--5\% CPR at 8\% market rates (the turnover speed observed on deeply out-of-the-money discount cohorts during 2023--2024, when refinancing incentives were absent and remaining prepayment reflected death, divorce, and forced relocation), I hold that empirical floor fixed and re-run the calibration under the steeper penalty function, which requires nearly tripling the mobility-desire scale parameter. This mechanical recalibration is the principal driver of the results in Section~\ref{sec:abm-results}.

Whenever a specification change alters the penalty or covariate structure facing every household, the mobility-desire scale, and only that parameter, is re-derived so that the simulation continues to reproduce the empirically observed 4--5\% baseline turnover floor. The same rule generates the recalibrated variant of the cross-design test (Section~\ref{sec:robustness-crossdesign}), whose frozen variant is defined by suspending it, and Appendix~\ref{sec:robustness-extensions} reports how much work this single choice does. Because the 4--5\% floor is imposed by construction, the ABM's agreement with observed discount-cohort turnover at deep rate gaps is an input, not a finding. The model earns no evidential credit on that dimension; its evidential exposure lies entirely on the dimensions the anchor leaves free, the dollar recovery and the timing of the CPR path. The anchor is also insulated from the behavioral gates by construction: the calibration point is evaluated at zero rate velocity, where the wait-and-see freeze cannot bind, so no freeze parameterization can move the floor (Appendix~\ref{sec:robustness-extensions} verifies the calibrated scale and floor CPR are bit-identical at every freeze share swept). What the anchor purchases, and what dropping it costs, is measured directly in Section~\ref{sec:robustness-crossdesign}: the externally parameterized variant that removes the floor overshoots the benchmark by half while predicting essentially zero turnover at QT-window rate gaps.

An income-scaled curtailment channel was added as an extension specified ex ante: a diagnostic correlation between the ABM's monthly CPR over-prediction and year-over-year growth in real disposable income (FRED series DSPIC96) of $r = -0.487$, concentrated in the worst over-prediction months (June--September 2022), led me to predict ex ante that curtailment, a strictly additive channel, would worsen the aggregate match. Curtailment was implemented as an income-scaled addition to simulated roll-off, mapping real income growth linearly from zero curtailment at $-2$\% growth to a full-scale annualized rate of 1.2\% CPR at +4\% growth, consistent with GSE daily prepayment report estimates.

Finally, the single representative 3.0\% coupon assumed elsewhere was replaced, as a diagnostic, with a multi-vintage composition: SOMA holdings bucketed into seven coupon cohorts yield a weighted-average coupon of approximately 2.55\% on the 30-year book, with the 2.0\% and 2.5\% buckets accounting for roughly 69\% of 30-year face value. Because a lower coupon implies a larger gap to prevailing market rates, my ex ante hypothesis was that this composition would increase predicted trapped liquidity toward the empirical benchmark. Fifteen-year MBS (approximately 9\% of SOMA face value) are excluded from this pass, and cohort origination dates are approximated from parsed maturity strings under a 360-month term assumption. Section~\ref{sec:abm-results} describes the production replacement of this diagnostic with a live, CUSIP-level pull of SOMA holdings.

\section{ABM Decision Rule and Parameter Provenance}\label{app:abmspec}
% Sources: abm/abm_lockin_simulation.py ("S:"), abm/fed_mbs_extension_risk.py ("F:"),
% abm/cross_design_test.py, abm/abm_external_gates.py. Line pins refer to the committed code.

This appendix specifies the ABM's move decision and every behavioral parameter at the level needed to reimplement the model without consulting the code; each equation and value is pinned to its source line in the replication package.

\paragraph{The move decision.}
Let $A(P,r,n) = P\,(r/12)\big/\big(1-(1+r/12)^{-n}\big)$ be the monthly payment on principal $P$ at annual rate $r$ over $n$ months. % S:130-136
Household $i$ carries gross income $y_i$ (\$/year), home value $H_i$ (\$), original principal $P_i = 0.8\,H_i$ (a fixed 80\% origination LTV in the synthetic population), note rate $c_i$, original term $n_i$ months, and outstanding balance $K_i$ (closed form at the cohort's seasoning; the U.S. par payoff). % S:141-155, S:314
The replacement-loan payment is $M^{\mathrm{new}}_{i,t} = A(K_i, m_t, n_i)$: the par payoff financed at the market rate $m_t$ over a fresh \emph{full original term}, not the remaining term (term-preference persistence). % S:222-225, S:366-381
With $\Delta_{i,t} = M^{\mathrm{new}}_{i,t} - A(P_i, c_i, n_i)$, whose second term is the current payment (by the amortization identity, equal to the payment on $K_i$ at $c_i$ over the remaining term), the loss-aversion-scaled penalty, in \$/month, is term-dependent:
\begin{equation}\label{eq:abmpenalty}
\pi_{i,t} =
\begin{cases}
\lambda\,[\Delta_{i,t}]_{+} + [\Delta_{i,t}]_{-} & n_i > 240 \text{ (30-year book)},\\
\lambda\,\big[M^{\mathrm{new}}_{i,t} - A(K_i, c_i, n_i)\big]_{+} & n_i \le 240 \text{ (native short-term gate)},
\end{cases}
\end{equation}
% S:383-409; the 240-month switch is S:408
with $\lambda = 2.25$, $[x]_{+} = \max(x,0)$, $[x]_{-} = \min(x,0)$: payment increases are scaled by $\lambda$; on the 30-year book, decreases pass through unscaled as gains, while the short-term gate floors them at zero to isolate the pure rate-lock cost. The frozen production headline draws 15-year voluntary CPR from the same-coupon 30-year surface, so the short-term branch is operative only in the Berger-run calibration stage (Section~\ref{sec:abm-danish}).
The effective transaction rate is $\tau_{i,t} = \mathrm{clip}\big(\phi_t + (\tau_i - 0.07),\,0.02,\,0.12\big)$: the agent's idiosyncratic draw enters as an offset around the 7\% base mean and is re-clipped after the macro friction $\phi_t$ shifts it. % S:241-243, S:459-461
Household $i$ moves in month $t$ iff
\begin{equation}\label{eq:abmmove}
\underbrace{\frac{M^{\mathrm{new}}_{i,t}}{y_i/12} \le 0.43}_{\text{DTI veto}}
\;\wedge\;
\underbrace{D_i > S\,\pi_{i,t} + \tau_{i,t} H_i}_{\text{cost--benefit}}
\;\wedge\;
\underbrace{\neg\big(v_t > 0.015 \,\wedge\, u_i < 0.20\big)}_{\text{wait-and-see freeze}},
\end{equation}
% Gates: DTI S:455-457, cost-benefit S:458-463, freeze S:464-467, conjunction S:468
where $D_i \sim \mathrm{Exponential}(\theta)$ is the mobility-desire draw in dollars, $S = 60$ months is the expected-stay horizon converting the monthly penalty into a dollar stock, $v_t$ is the trailing six-month mortgage-rate change (decimal), and $u_i \sim U[0,1]$ is a fixed per-agent trait (the code's ``patience'' draw, consumed only by the freeze gate; it is not a time-preference parameter). % S:48, S:192, S:302-309
All inequalities are as coded: the DTI gate passes at exactly 43\%, the cost--benefit gate requires strict exceedance, and the freeze binds only strictly above 150 bp. The share of a pool passing \eqref{eq:abmmove} at an evaluation point is read as that point's annualized CPR (the calibration anchor is an annual floor) and consumed by the macro accounting layer at $\mathrm{CPR}/12$ per month. % F:957

\paragraph{Parameter provenance.}
Table~\ref{tab:abmparams} collects every behavioral parameter. None is tuned to the \$764.7 billion benchmark; the mobility-desire scale $\theta$ is the one parameter re-derived under the floor rule, and only ever that one.

\begin{table}[H]
\centering
\footnotesize
\begin{threeparttable}
\caption{ABM behavioral parameters: value, units, distribution, and provenance.}
\label{tab:abmparams}
\begin{tabular}{@{}p{3.3cm}p{5.9cm}p{1.6cm}p{4.0cm}@{}}
\toprule
Parameter & Value / distribution & Units & Provenance \\
\midrule
Population & $N = 10{,}000$; engine seed 42 & households & author choice \\
Income $y_i$ & log-normal, $\sigma = 0.45$, median from FRED MEHOINUSA672N & \$/year & FRED-calibrated draw \\
Home value $H_i$ & log-normal, $\sigma = 0.35$, median from FRED MSPUS & \$ & FRED-calibrated draw \\
Origination LTV & 0.80, fixed & --- & author choice \\
Mobility desire $D_i$ & exponential, scale $\theta$ & \$ & floor recalibration (last row) \\
\quad $\theta$, production (frozen) & 43{,}883 & \$ & floor anchor, synthetic population \\
\quad $\theta$, cross-design recalibrated & 36{,}086 & \$ & floor anchor, real Freddie covariates (Section~\ref{sec:robustness-crossdesign}) \\
\quad $\theta$, external variant & 4{,}472 & \$ & \citepos{liebersohn2024} zero-gap 5.9\%/yr moving level, no floor (Section~\ref{sec:robustness-crossdesign}) \\
Expected-stay horizon $S$ & 60 & months & author calibration \\
Loss aversion $\lambda$ & 2.25 & --- & \citet{tversky1992} \\
DTI wall & 0.43, front-end & --- & 12 C.F.R. \S{} 1026.43(e)(2)(vi), behavioral anchor \\
Freeze share; trigger & 0.20 of agents; $v_t > 0.015$ over six months & ---; decimal & author calibration \\
Patience trait $u_i$ & uniform $[0,1]$, fixed per agent & --- & implementation device (freeze membership) \\
Transaction-cost draw $\tau_i$ & normal$(0.07,\,0.015)$, clipped $[0.02, 0.12]$ & share of $H_i$ & author calibration \\
Dynamic friction $\phi_t$ & $0.07$ base $+$ search penalty (listing shortfall vs.\ 2017--19 baseline, proportional, cap $+0.02$) $+$ sentiment penalty (5 bp per UMCSENT point below 85, cap $+0.015$); realized 8.23--9.96\% over the 42 QT months & share of $H_i$ & FRED ACTLISCOUUS, UMCSENT; mappings and base level are author calibrations \\
Curtailment & linear in real-income YoY growth: 0 at $-2\%$ to 1.2\% CPR at $+4\%$; $/12$ monthly, additive, macro layer only & annualized CPR & GSE daily-prepayment report estimates; FRED DSPIC96 (Appendix~\ref{sec:method-frictions}) \\
Turnover floor & binary search $\theta \in [10^{3}, 5\times10^{5}]$, at most 30 iterations, until CPR at an 8\% market rate under the U.S. rule, base friction, zero velocity lands in $[4\%, 5\%]$, evaluated on the dominant-weight 30-year SOMA cohort & annual CPR & 2023--24 discount-cohort turnover (Appendix~\ref{sec:method-frictions}) \\
\bottomrule
\end{tabular}
\begin{tablenotes}
\footnotesize
\item Committed artifacts carry the unrounded binary-search scales (43{,}882.8125; 36{,}085.9375; 4{,}472.05). The module default $\theta = 12{,}500$ is a placeholder overwritten by calibration at every run. Code pins: draws \texttt{abm\_lockin\_simulation.py} 296--309; constants 38--69; floor search 649--685; friction \texttt{fed\_mbs\_extension\_risk.py} 66--70, 294--334; curtailment 85--87, 141--155; variant scales \texttt{cross\_design\_test.py} 81--111 and \texttt{abm\_external\_gates.py} 88--89.
\end{tablenotes}
\end{threeparttable}
\end{table}

\paragraph{Order of operations.}
Within a simulated month the three gates of \eqref{eq:abmmove} are evaluated on the surviving population and combined conjunctively; they are side-effect-free, so their order cannot change the result, and the code evaluates them in vectorized form in the order DTI, cost--benefit, freeze. % S:449-468
In cohort-path (burnout) mode, movers are removed permanently before the next month's evaluation, so the turnover tail sets a natural floor on the depleted pool. % S:470-518
Two named components are not in-month gates. The 4--5\% floor enters once, before simulation, through the desire-scale search of Table~\ref{tab:abmparams}, re-run, for that parameter only, whenever a specification change alters the penalty or covariate structure facing every household; because the anchor is evaluated at zero velocity, no freeze parameterization can move it (Appendix~\ref{sec:method-frictions}). % S:649-685, S:520-524
Curtailment never touches the household decision: it is an additive term in the macro layer's simulated roll-off. In production the decision rule is also not re-solved monthly: the CPR surface is precomputed on the rate $\times$ friction $\times$ velocity grid with all gates applied at each node and interpolated month-by-month, with the freeze's step-versus-ramp semantics disclosed in Appendix~\ref{sec:method-abm}. % F:873-895, S:526-548

The correction of this paper's earlier ``log-normal'' description of the mobility-desire distribution---it is exponential, on a dollar scale---is recorded in the run ledger (Appendix~\ref{app:ledger}).

\section{Path B Parameterization and Delinquency Transitions}\label{app:params}

\begin{table}[H]
\centering
\footnotesize
\begin{threeparttable}
\caption{Objects and notation for the two hazard specifications. Path A's macro coefficients are reported in Section~\ref{sec:patha}, and every Path B parameter with its provenance is collected in Table~\ref{tab:params}.}
\label{tab:hazard-notation}
\begin{tabular}{@{}lllc@{}}
\toprule
Symbol & Object & Units & Path \\
\midrule
\multicolumn{4}{@{}l}{\emph{Sample and exposure}}\\
$s$ & stratum (296 four-way cells) & -- & A \\
$(s,t)$ & unit of observation (stratum, month) & -- & A \\
$i$ & loan (75,000 sampled; 40,234 active) & -- & B \\
$(i,t)$ & unit of observation (loan, month) & -- & B \\
$y_{s,t}$ & prepaid unpaid principal balance & dollars & A \\
$n_{s,t}$ & UPB exposure (log offset) & dollars & A \\
\midrule
\multicolumn{4}{@{}l}{\emph{Link and baseline}}\\
$\alpha_{s}$ & stratum fixed effect (295 enter) & log-rate & A \\
$f(a_{s,t})$ & loan-age seasoning spline & log-rate & A \\
$h_{0}(a_{i,t})$ & PSA baseline hazard (Section~\ref{sec:pathb}) & monthly & B \\
$\underline{h}$ & baseline turnover floor & monthly SMM & B \\
\midrule
\multicolumn{4}{@{}l}{\emph{Covariates and coefficients}}\\
$g_{s,t}$ & rate gap & 100 bp units & A \\
$b_{s,t}$ & stratum burnout state & demeaned share & A \\
$\phi_t$ & dynamic friction index & index & A \\
$\beta_g,\ \beta_b,\ \beta_\phi$ & rate-gap, burnout, friction coeff. & log-rate & A \\
$g_{i,t}$ & rate gap (negative when locked in) & 100 bp units & B \\
$\tilde g_{i,t}=g_{i,t}/100$ & rate gap in decimal units & decimal & B \\
$\sigma_{i,t}=\max\{0,-\tilde g_{i,t}\}$ & rate stress & decimal & B \\
$z^{\mathrm{FICO}}_{i},\ z^{\mathrm{LTV}}_{i}$ & standardized static covariates & SD & B \\
$B_{s(i),t}\in[0,1]$ & cohort-level burnout state & cumulative share & B \\
$\beta_1,\ \beta_F,\ \beta_L,\ \beta_B$ & gap, FICO, LTV, burnout coeff. & log-hazard & B \\
$\gamma_r,\ \gamma_F,\ \gamma_L$ & default-hazard coefficients & log-hazard & B \\
\bottomrule
\end{tabular}
\begin{tablenotes}[flushleft]\scriptsize
\item Case and subscript separate the two paths throughout, and no coefficient symbol is reused across them. Path A's burnout coefficient $\beta_b$ (on $b_{s,t}$, estimated on the panel) and Path B's $\beta_B$ (on the cohort state $B_{s(i),t}$, a calibrated prior) are distinct objects carrying distinct symbols, as are Path A's friction coefficient $\beta_\phi$ (on the friction index $\phi_t$) and Path B's FICO coefficient $\beta_F$. The rate-gap elasticities $\beta_g$ (Path A, log-rate, standardized) and $\beta_1$ (Path B, log-hazard, per 100 bp) are on different scales; the paper compares them on sign, not magnitude (Section~\ref{sec:hazard-interp}).
\end{tablenotes}
\end{threeparttable}
\end{table}

Table~\ref{tab:params} collects every Path B parameter. All values are fixed a priori in the configuration; none is estimated on, or tuned to, the \$764.7 billion benchmark. One parameter is nonetheless measured inside the evaluation window, and benchmark-independence should not be read as window-independence: the 4\% baseline turnover floor is anchored to turnover observed on discount cohorts during 2023--2024 (Appendix~\ref{sec:method-frictions} and Section~\ref{sec:pathb}), so the floor (which sets much of the aggregate level, binding in 36\% of loan-months) is in-sample with respect to the QT window even though it never touches the benchmark dollar figure. Section~\ref{sec:robustness-floor} sweeps the floor jointly with the elasticity band: the aggregate level and the lock-in marginal both move materially with the floor, within the bounded interval that section reports; the marginal's sign does not move, but under this window's rate configuration it could not (Section~\ref{sec:identification}). The covariate and default-hazard coefficients are round calibration priors. The two prepayment covariates' signs run opposite the canonical credit-risk direction (higher FICO $\Rightarrow$ better refinancing access $\Rightarrow$ faster prepayment; higher LTV $\Rightarrow$ less equity $\Rightarrow$ slower prepayment): in the deeply out-of-the-money 2022--2025 window the refinancing channel is shut and prepayment is turnover-dominated, so credit covariates proxy mobility rather than refinancing access; high-LTV loans are disproportionately recent purchases that churn. A Path-A-style estimation on the same stratum-month panel (Poisson PML with vintage-by-coupon fixed effects, so the credit buckets are identified from cross-bucket variation; bucket-mean covariates z-scored to the microsim pool's own moments) yields point estimates consistent with the priors' direction (and, taken at face value, roughly double their magnitudes), but with intervals that reject neither zero nor the opposite signs: $\hat{\beta}_F = -0.39$ (stratum-cluster bootstrap 95\% CI $[-1.51, +0.65]$) and $\hat{\beta}_L = +0.21$ ($[-0.59, +1.00]$). The estimation does not settle the sign question; it fails to contradict the priors while quantifying their sensitivity. The sensitivity brackets are correspondingly modest: a $\beta_F = \beta_L = 0$ rerun moves the recovery from 107.0\% to 108.3\%, and the panel-estimated pair moves it to 101.3\% --- in-window figures; at the off-window floor the same endpoints span 8.2 points and `modest' does not carry over to the level, while the marginal's covariate sensitivity shrinks to $+0.1$ points (run \texttt{covariate\_ablation\_offwindow}); benchmark-independence is unaffected, since the estimation touches only the Freddie panel. The $\beta_B = 0$ rerun of Section~\ref{sec:pathb} quantifies the burnout channel the same way.


Delinquent loans transition monthly through the matrix of Table~\ref{tab:delinq}, estimated from consecutive stratum-month \emph{modal} servicer states in the Freddie cohort panel, exposure-UPB weighted (committed as an artifact in the replication package). The division of labor with the hazard draws: the default-hazard draw moves loans \emph{into} the pipeline (Current $\to$ D30) and terminates already-delinquent loans it hits, while the matrix governs month-to-month movement \emph{within} the pipeline; the matrix runs after the draw each month, so a loan entering D30 via the draw can also take a matrix transition in its entry month, an overlap bounded along with everything else by the \$0.39 billion extreme-matrix test below. Only the delinquency rows are used in simulation, terminal states are absorbing by construction, and the forbearance state is unpopulated in the estimation panel and never entered in production runs.
I publish the full crosstab. D30 row (21 transitions, \$58.3M row exposure): 13 cures carrying 95.6\% of the exposure weight, two D30 continuations (0.8\%), five D60 progressions (3.1\%), and one Defaulted transition (0.4\%); the printed 0.004 is the exposure share of that single event, whose count share would be $1/21 = 0.048$. D60 row (9 transitions, \$3.1M): six cures carrying 56.8\% and three D90+ progressions carrying 43.2\%; the count shares are $6/9$ and $3/9$; the divergence between count and exposure shares means that a few larger-balance stratum-months drive the matrix, which is what an exposure-weighted estimator is designed to do. D90+ row (38 transitions, \$76.2M): seven cures (30.3\%), one D30 (0.8\%), thirty self-continuations (68.9\%), and zero Defaulted events; the ``--'' cells are observed zeros in a thin panel, not structural zeros, which is why serious delinquency shows no direct default exit while D30 shows one. Probabilities are printed to three decimals from these exposure shares and rows may not sum to exactly one. The matrix is accordingly reported for completeness, not precision: with cells this thin it describes the estimation panel rather than estimating a transition law, and the paper's estimates rely on the materiality bound below, not on the matrix's point values.

The thinness is immaterial for the paper's estimates: replacing the entire matrix with either extreme (all delinquent loans cure next month; all default next month) moves trapped liquidity by \$0.39 billion (0.05\% of the estimate). The 482 loans delinquent at the window start are routed by the matrix from the first month (mostly cured, given the cure-heavy rows); the month-end delinquency stock thereafter never exceeds 154 of the 40,234 active loans, and delinquent UPB never exceeds 1.1\% of the pool. Both extreme-matrix runs include the initial 482, so the bound covers them.

\begin{table}[H]
\centering
\small
\begin{threeparttable}
\caption{Exposure-weighted delinquency transition matrix (modal stratum-month states).}
\label{tab:delinq}
\begin{tabular}{@{}lrrrrr@{}}
\toprule
From $\backslash$ To & Current & D30 & D60 & D90+ & Defaulted \\
\midrule
D30 & 0.956 & 0.008 & 0.031 & -- & 0.004 \\
D60 & 0.568 & -- & -- & 0.432 & -- \\
D90+ & 0.303 & 0.008 & -- & 0.689 & -- \\
\bottomrule
\end{tabular}
\begin{tablenotes}[flushleft]\footnotesize
\item Probabilities are exposure-UPB shares of thin modal-state transition counts (D30 row: 21 transitions; D60: 9; D90+: 38; full crosstab in the text and replication artifact). ``--'' cells are observed zeros in a thin panel, not structural zeros; rows may not sum to exactly one at printed precision. Reported for completeness, not precision: the paper's estimates rely on the \$0.39 billion extreme-matrix materiality bound in the text, not on these point values.
\end{tablenotes}
\end{threeparttable}
\end{table}
Sampling design: the draw is not a probability sample of the SOMA book, and the allocation rule is why. A pool built at four times the target size takes an equal allocation from each quarterly Freddie origination file; the pool is then drawn down to 75,000 loans proportionally within stratum, and every retained loan carries weight one (\texttt{hazard/loan\_sample.py}). Nothing in that chain reweights toward the book. Each origination year from 2017 through 2021 therefore takes 20.0\% of the draw by count, so 2017--19 is 60.0\% of it against 6.0\% of book face. That count basis is not the weighting the aggregation applies, which is the surviving balance at the window open, and on that basis the picture is materially different: 2017--19 is 27.0\%, 2020 30.7\% and 2021 42.3\%, against book face shares of 6.0\%, 16.3\% and 43.9\%. So the 2021 cohort is close to matched, and the 2017--19 overweight is a factor of about four and a half rather than the factor of ten the count basis suggests; origination balance, which compares originations to current face and is the wrong basis for this comparison, gives 55.5\%, 22.2\% and 22.3\%. What no weighting repairs is out-of-support mass. The draw spans 2017 through 2021 only, so pre-2017 (10.6\% of book face) and the 2022 cohort (23.1\%) carry no support at all: 33.7\% of face, which must be netted against the Ginnie share on the book's own joint cells rather than added to it. The draw's 3.86\% unweighted mean note rate also sits above the book's 2.49\% face-weighted pass-through coupon. Both tilts move the rate-gap distribution the floor takes its maximum against under the production floor form, and this paper signs neither: Section~\ref{sec:robustness-hybrid}'s coupon reweight rescales the aggregate level and leaves the marginal unchanged, so it does not test a composition effect running through the floor's bind share.

Attrition accounting: of the 75,000 sampled loans, 34,734 prepaid and 32 defaulted before the QT window opens, so 40,234 are active at the June 2022 start (39,752 current, 482 delinquent). The production U.S. run therefore evaluates 1,683,124 loan-months over the 42-month window, a mean of 40,074 active loans per month. The baseline turnover floor $\underline{h}$ binds in 36\% of those loan-months (Section~\ref{sec:pathb}).

The $\delta$-to-$\beta_1$ conversion of \eqref{eq:beta1} is insensitive to the turnover proxy behind it: in the small-$P_q$ limit the transform reduces to $\beta_1 = -\ln(1-\delta) = 0.0672$ at the central $\delta$, and over $P_q \in (0, 0.12]$ (from arbitrarily low turnover to double the adopted proxy) the central $\beta_1$ ranges only from $0.0672$ to $0.0701$, so even a large error in $P_q$ moves the conversion by a few percent.

\section{Path A Estimation Detail and Timing Diagnostics}\label{app:patha-detail}

This appendix collects the estimation detail of Path A (Section~\ref{sec:patha}), the direct input-timing diagnosis, the Danish discount-rate diagnostic, and the monthly-timing exhibits of Section~\ref{sec:identification}.

The second hazard path aggregates Freddie Mac loan-level performance data to the stratum-month level: strata $s$ are the 296 observed four-way cells of vintage $\times$ coupon $\times$ FICO bucket $\times$ LTV bucket, and the unit of observation is the stratum-month pair $(s,t)$. The bucket edges are the estimation configuration's: FICO below 680, 680--740, and 740 and above; LTV at or below versus above 80; coupons in 0.5-percentage-point steps; vintages 2017--2021. Geography is not a stratum axis, and no geographic term enters \eqref{eq:pathA}. The model is a discrete-time hazard with a log-exposure offset,
\begin{equation}\label{eq:pathA}
\log \mathbb{E}\!\left[y_{s,t} \mid x_{s,t}\right] = \log n_{s,t} + \alpha_{s} + f\!\left(a_{s,t}\right) + \beta_g\, g_{s,t} + \beta_b\, b_{s,t} + \beta_\phi\, \phi_t ,
\end{equation}
where $y_{s,t}$ is stratum $s$'s prepaid unpaid principal balance (UPB) in month $t$, in dollars, $n_{s,t}$ its UPB exposure, $x_{s,t}$ the right-hand-side covariates, $\alpha_{s}$ the stratum fixed effect (295 enter, one absorbed as reference), $f(\cdot)$ a loan-age seasoning spline in mean loan age $a_{s,t}$ (truncated quadratic power basis: a linear term plus $(a - k)_+^2$ at knots $k \in \{12, 24, 36, 60, 84, 120\}$ months, seven basis columns), $g_{s,t}$ the rate gap in 100 bp units, $b_{s,t}$ the stratum burnout state, and $\phi_t$ the dynamic friction index. The outcome is dollars, not a count: with the log-exposure offset the modeled object is the prepayment rate $y_{s,t}/n_{s,t}$, a nonnegative fraction, so the model is estimated by Poisson pseudo-maximum likelihood, whose score equations remain consistent for the conditional mean without any distributional assumption on $y_{s,t}$; no dispersion statistic is meaningful for this pseudo-likelihood. The 295 stratum fixed effects raise no incidental-parameters problem for the coefficients of interest: because the model is Poisson, concentrating the fixed effects out of the log-likelihood returns exactly the conditional-likelihood score for the slope vector, so the unconditional dummy-variable fit computed here yields the same slope estimates as the conditional estimator and stays consistent as the number of strata grows with the 36 training months held fixed, unlike the fixed-effect logit or probit, whose slope estimates are inconsistent at fixed panel length \citep{wooldridge1999}. That consistency requires only the multiplicative conditional mean $\mathbb{E}[y_{s,t}\mid x_{s,t},\alpha_{s}] = n_{s,t}\exp(\alpha_{s} + x_{s,t}^{\top}\beta)$ with no restriction on the distribution, dispersion, or serial dependence of $y_{s,t}$. All uncertainty statements below come from the two bootstrap schemes of Table~\ref{tab:bootstrap}: a stratum-level cluster bootstrap, robust to arbitrary within-stratum dependence and overdispersion but structurally unable to measure uncertainty in the friction coefficient ($\phi_t$ varies only over time, so resampling strata reproduces the identical friction series in every replication), and a temporal moving-block bootstrap that resamples the 36 training months in blocks of six with each month's full cross-section intact, which carries that burden (and, partially, burnout's). Table~\ref{tab:panel} reports the estimation panel's accounting; both it and Table~\ref{tab:bootstrap} are collected in Appendix~\ref{app:ridge}, because Path A is excluded from every headline range by pre-committed specification and neither table enters a headline quantity. The displayed form \eqref{eq:pathA} is the no-seasonal spec v3 design; the production specification (spec v4) adds eleven calendar-month dummies and is the design under which the 121.5\% recovery reported below is computed, though that level is not calendar-attributable (Section~\ref{sec:patha}), while every Path A inferential object stays on spec v3. Table~\ref{tab:hazard-notation} collects the notation shared by the two hazard paths.

The 295 stratum fixed effects absorb each cohort's baseline prepayment speed so that the burnout term captures within-pool adverse selection over time rather than cross-sectional differences in pool quality. The burnout regressor is a within-stratum demeaned measure of the surviving loan stock, orthogonalized against the loan-age spline so that it is not mechanically collinear with seasoning. This design is the estimation counterpart of \citepos{lesniewski2026} selection identity: the exposure-offset outcome is itself a pool-level, survival-weighted (here UPB-exposure-weighted) hazard, the 295 strata cut the cross-sectional dispersion that drives the identity's negative selection term down to within-cell variation, and the demeaned surviving-stock regressor proxies the residual composition drift that the time-invariant fixed effects cannot absorb. An earlier specification using four vintage-year dummies in place of the 295 stratum fixed effects returned a burnout coefficient of the wrong sign (+0.76); replacing the coarse vintage fixed effects with the full 295-stratum specification and demeaning burnout within stratum removes that wrong-signed artifact. The corrected estimate ($-0.13$) carries the theoretically expected sign (a more heavily burned-out pool should prepay more slowly, holding age and rate gap fixed), but the sign statement is a point-estimate statement only: the temporal bootstrap scheme of Table~\ref{tab:bootstrap} centers this coefficient on the opposite sign (mean $+1.23$, median $+0.52$), and under both schemes it is imprecise and not statistically distinguishable from zero, so I treat it as a control rather than a finding. Because 295 fixed effects make the model ill-conditioned under plain iteratively reweighted least squares (20 strata record no training prepayment event at all, so the likelihood is monotone in their fixed effects and the maximum-likelihood estimates diverge to $-\infty$), the fit carries a nominal ridge penalty $\lambda_{\mathrm{ridge}}$ on the full coefficient vector, with $\lambda_{\mathrm{ridge}}$ chosen over the grid $\{10^{-5}, 10^{-4}\}$ on a temporally blocked holdout comprising all stratum-months from January 2024 onward, so that no within-stratum information crosses the train/holdout split. The penalty is not what stabilizes the fit: at these magnitudes it is a numerical no-op (the two grid points return coefficient vectors identical to machine precision from the IRLS warm start, so the production coefficients are the IRLS solution and the grid selection is inconsequential), the 20 zero-event strata are held at the finite value the warm-started optimizer leaves them at (the penalty shrinks them only negligibly); and because the split is temporal over a fixed set of strata, that same fixed effect scores each stratum's post-January-2024 holdout months, so no fixed effect is estimated on, or transferred to, a stratum absent from the training window. The penalty's operative role is as the fallback for resampled refits in the bootstrap machinery. Nor is the production point an artifact of the device: dropping the 20 zero-event strata and fitting unpenalized Poisson PML reproduces the spec v3 macro coefficients to within $0.002$ in standardized units (rate gap $+0.6727$ against the spec v3 $+0.6734$; under spec v4 the same check holds to within $0.0022$, Appendix~\ref{app:ridge}). The verification details, including a reference-stratum sensitivity of penalized cold-start refits that motivates the warm-started bootstrap below, are recorded in Appendix~\ref{app:ridge}, now under both specifications: the battery has been recomputed under spec v4 and the device is unchanged (run \texttt{pathA\_ridge\_v4}). Two caveats still attach: the same holdout later serves as the out-of-sample diagnostic, so the RMSE reported below is at least weakly favored by selection, and months never touched by any selection step do not exist in this design, one reason Path A is held to a robustness role.
Estimated on this panel, the production specification (spec v4, the calendar-month seasonal design adopted at the freeze) yields a rate-gap coefficient of $+0.65$ (standardized; the training standard deviation of the rate gap is 221 bp, so the equivalent per-100 bp coefficient is $+0.29$), a burnout coefficient of $-0.17$, and a friction coefficient of $-0.007$, each with the theoretically expected sign. The friction coefficient's shrinkage relative to the prior spec ($-0.037$) is expected under the seasonal design, whose eleven month dummies absorb most within-year variation in the time-only friction index. Point estimates alone overstate what is identified here, so Table~\ref{tab:bootstrap} reports cluster-bootstrap uncertainty for all three (estimated under the prior spec v3; see the table note), including the bias diagnostics: the rate-gap point estimate sits near the 3rd percentile of its own bootstrap distribution (bootstrap mean $+1.49$, median $+0.76$, estimated bias $+0.82$), a signature of resampling-induced upward skew under which percentile and basic intervals disagree (the percentile interval excludes zero while the bias-reflecting basic interval does not). I therefore rest no claim on interval-based significance for any Path A coefficient. The defensible statements are about direction and imprecision: the rate-gap coefficient's point estimate is positive, and positive in 197 of 198 converged replications (99.5\%), but that share is a property of a right-skewed resampling distribution rather than a test, and the pre-committed bias-respecting test of $H_0\!: \beta_g \le 0$ does not reject under any construction (BCa $p = 0.093$ stratum / $0.412$ temporal; permutation $p = 0.241$; run \texttt{patha\_sign\_test}), so the coefficient is directionally consistent with the imported elasticity without being statistically distinguishable from zero; its magnitude is identified only up to a wide, construction-dependent interval, and burnout and friction are not distinguishable from zero on any construction. All three enter the production simulation as fitted inputs, none as findings, and no result in this paper depends on their point values. Path A's aggregate recovery now carries the interval that pass produces: resimulating the QT window under each of the 198 converged replications (the three macro coefficients swapped in per draw; intercept, age spline, and fixed effects held at the production fit) yields a median of \$922.6 billion and an interquartile range of \$900.1--\$1{,}130.6 billion, but a 95\% percentile interval of $[-\$4{,}459.8, +\$1{,}190.1]$ billion. A joint rerun that swaps in each replication's full coefficient vector --- intercept, age spline, and at spec v4 the month effects, drawn together with the macros --- returns $[+\$689.0, +\$1{,}217.6]$ billion, far narrower, and it is not adopted: 106 of its 198 draws, among them all 79 rogue-mode draws (drawn rate-gap coefficient nearer the reference-swap attractor $+1.96$ than the production $+0.67$), sit at exactly the zero-voluntary-prepayment floor of \$1{,}217.6 billion --- the trapped figure a pool returns when voluntary prepayment is switched off and only scheduled amortization rolls off, measured by a probe through the same coefficient handoff rather than by formula (run \texttt{pathA\_resimulate\_joint}). The committed interval remains the quoted one, because its $-\$4{,}459.8$ billion tail and the joint rerun's floor pile are the same rogue-mode replicates rendered without and with their compensating intercepts: one optimizer pathology of penalized refits on resampled panels, not two sampling statements. Measured at spec v4, the same joint rerun fails 1 of 200 replications against spec v3's 2 and its draws avoid the floor entirely (0 of 199), so the seasonal design's instability manifests as rogue-mode attraction (73 of 199 draws), not as fit failure. Fifteen replications with extreme burnout or friction draws (the two coefficients the bootstrap classifies as statistically zero) destabilize the 42-month forward compounding entirely, and 83.8\% of draws land above the benchmark. Those two facts pull in opposite directions, and the unfavorable one belongs in the text as much as the favorable one: the same 198 draws have a \emph{mean} of \$598.9 billion, or 78.3\% of the benchmark---below the mechanical $\beta_1 = 0$ null's 85.7\% and below Path~B's headline 91.3\%---against a median of 120.6\%, and on the share basis the percentile interval is $[-583.2\%, +155.6\%]$ with a resampling standard error of 201.9 points. I report the mean alongside the median for the same reason the coefficient table does: these distributions are severely right-skewed, so no single summary is sufficient, and quoting only the median and the fraction above the benchmark would flatter the estimator. Propagated coefficient uncertainty therefore leaves the forward simulation with no useful precision, an additional reason the aggregate dollar figure is held to the aggregate-level role assigned below and is not carried in the abstract. Every Path A inferential object (Table~\ref{tab:bootstrap}, this interval, and the train/holdout decomposition below) remains estimated under the prior spec v3, because bootstrap replications under the spec v4 seasonal design are branch-unstable (the resampled penalized likelihood has a second optimum that attracts replications away from the production branch); on the production panel itself 1 of 50 dispersed starts lands on the production branch and no start reaches a better penalized objective --- the objective evaluated at the frozen coefficient vector is the best value any of 150 starts attains, and every non-production terminus sits at least seven orders of magnitude above it, an optimizer stall of a strictly convex objective rather than a rival optimum (run \texttt{pathA\_cold\_starts}) --- and until the resampled instability is stabilized, the production specification carries no valid uncertainty of its own, and Path A is excluded from the paper's headline figures for this reason, retained as an aggregate-context, non-corroborating exhibit. Simulating forward (draining each cohort's balance deterministically each month by balance times hazard, with fractional prepayment settling directly to simulated SOMA roll-off and no settlement-lag convolution), this model recovers \$928.9 billion, or 121.5\%, of the \$764.7 billion empirical benchmark under the production calendar-month specification (spec v4), adopted at the pre-submission freeze under pre-committed parity gates that reproduce the committed seasonal robustness artifact exactly (the prior specification is preserved as a spec-v3 exhibit).

Three diagnostics temper this aggregate result because they bear directly on how much weight the \$928.9 billion figure can carry. First, the holdout fit (evaluated under spec v3; the frozen spec v4 fit, re-scored on the same January-2024-onward holdout without refitting, gives 37.9 unweighted and 3.0 exposure-weighted percentage points against spec v3's 37.8 and 2.5, so the seasonal terms do not improve the out-of-sample fit) is weak in both weightings. The units matter: the unweighted RMSE of approximately 38 percentage points of annualized stratum-month CPR is dominated by small-exposure strata with unstable realized rates, while the exposure-weighted RMSE (the weighting the fit itself uses) is 2.5 percentage points against an empirical mean CPR of 5.14\%, half the mean. Neither weighting supports out-of-sample predictive power at the stratum-month level (and the ridge grid's two candidate fits are numerically identical on this holdout, so no selection hinges on the weighting). A flexible-learner ceiling test bounds how much of this weakness is functional form: a gradient-boosted Poisson learner given the same covariate information on the identical pre-committed holdout improves the exposure-weighted RMSE materially, to 0.9 points against the frozen spec v4's 3.0, while leaving the unweighted RMSE essentially unmoved (37.6 against 37.9) and out-of-sample $R^2$ at zero ($+0.003$); a penalized Poisson GLM on the exact v4 design does no better than the production fit (2.6 and 38.1 points). Read against the run's pre-committed thresholds, flexibility buys recalibration of the high-exposure cells but no cross-cell predictive power in the 2024-onward regime. The landmarking-with-recalibration strand receives the same test: an exposure-weighted isotonic map fit on the 2023 inner-validation year only, applied to the frozen spec v4 holdout predictions, improves the exposure-weighted RMSE to 2.53 points against the frozen 3.04 (short of the pre-committed 0.80$\times$ threshold) while leaving the unweighted RMSE at 37.8 and out-of-sample $R^2$ at zero ($-0.007$); the pre-committed verdict is no material change: a monotone level recalibration fit in the training regime does not transfer predictive power across the 2024 boundary. Second, the recovery is concentrated where the model is unsupervised (decomposition measured under spec v3; it was not recomputed at adoption): of the spec v3 fit's \$915.0 billion, \$331.0 billion accrues in the 19 QT months overlapping the training window (June 2022--December 2023), against \$337.5 billion empirical (98\%), while \$584.1 billion accrues in the 23 post-training simulation months from January 2024 onward (21 of which have performance data; the empirical comparison uses those), against \$427.3 billion empirical (137\%). Components are independently rounded, and the overshoot lives almost entirely out of sample. Third, the raw CPR correlation against the empirical path is $r = -0.444$ at zero lag (more negative even than the ABM's $-0.318$), and \eqref{eq:pathA} omits calendar-month seasonality, one of the four factors of \citepos{richard1989} canonical prepayment function and a natural first suspect for a monthly path fit this poor. Re-estimating with eleven calendar-month dummies confirms strong, correctly signed seasonality (spring and early-fall log-effects up to $+0.43$ against a January base) and improves the simulated path's contemporaneous correlation from $-0.444$ to $-0.378$, with the peak moving from lag $-3$ to $-2$. That timing improvement is calendar-attributable, and I test it rather than assert it. Permuting the assignment of the twelve estimated month log-effects to calendar months preserves their multiset exactly --- and therefore every level channel they drive --- while destroying only their alignment with calendar time; scored on the conjunctive criterion the claim actually states (a lag-0 correlation better than $-0.378$ \emph{and} a peak at lag $-2$), scrambled assignments win 1 of 64 draws, a one-sided $p$ of 0.031, with an independent sixteen-draw family at a different seed also returning a single win (\texttt{pathA\_seasonal\_permutation}). The two components scored separately, 11 and 8 of 64, are diagnostics only: each necessarily exceeds the joint count, so reading either alone understates a conjunctive claim. The recovery figure is a different matter, and I no longer credit it to the calendar terms. Across those same 64 scrambles the recovery spans 121.14--122.29\%, a range of 1.15 points that contains the adopted 121.5\%, so the level is calendar-invariant: permutation holds the multiset of month effects fixed, and it is the multiset, not its alignment, that sets the level. Nor is the spec-v4 total decomposable as production plus seasonality in the first place --- the seasonal refit with all twelve month multipliers zeroed moves the recovery by $+\$80.27$ billion on its own, and restoring the multipliers claws back $-\$66.44$ billion of that, against a net reported seasonal effect of $+\$13.83$ billion. The timing improvement is real but far from sufficient: the named first suspect accounts for only a modest share of the negative path fit. A model that recovers 121.5\% of the dollar benchmark while mismatching the monthly path this badly may be recovering the aggregate partly through offsetting errors: over-predicting prepayment in some cohort-months and under-predicting in others in a way that nets out near the benchmark without the underlying mechanism being correctly timed. I therefore treat Path A's recovery as evidence about the aggregate level only. No timing claim rests on Path B either: its apparent levels advantage does not survive detrending (Section~\ref{sec:pathb}), and every claim this paper makes for Path B is a level-accuracy claim. Table~\ref{tab:pathadiag} collects Path A's diagnostics in one place.

\begin{table}[H]
\centering
\caption{Path A diagnostic summary. An aggregate-context, non-corroborating exhibit: Path A is excluded from the paper's headline figures by pre-committed specification, and no cell of this table enters a headline quantity. Values repeat committed figures derived in this section and Appendix~\ref{app:theil}.}
\label{tab:pathadiag}
\footnotesize
\begin{threeparttable}
\begin{tabular}{@{}p{4.6cm}p{9.6cm}@{}}
\toprule
Diagnostic & Value \\
\midrule
Aggregate recovery \emph{(spec v4)} & \$928.9 billion $=$ 121.5\% standalone; 112.4\% shared basis. Order of magnitude coincides with Path B; the sign does not corroborate it, the coefficient being statistically indistinguishable from zero (row below and Section~\ref{sec:patha}) \\
Refit-and-resimulate distribution \emph{(spec v3)} & point \$915.1 billion; median \$922.6 billion; \emph{mean \$598.9 billion $=$ 78.3\% of benchmark}; IQR \$900.1--\$1{,}130.6 billion; 95\% percentile interval $[-\$4{,}459.8, +\$1{,}190.1]$ billion ($[-583.2\%, +155.6\%]$ on the share basis); 83.8\% of draws above benchmark \\
Train vs post-training split \emph{(spec v3)} & \$331.0 billion (98\% of empirical) vs \$584.1 billion (137\%) \\
Fit diagnostics \emph{(spec v4)} & exposure-weighted holdout RMSE 3.0 points (2.5 under spec v3); Theil $U_1$ 0.431, $U_2$ 1.010; terminal cumulative error $+\$164.1$ billion \\
\bottomrule
\end{tabular}
\begin{tablenotes}
\footnotesize
\item \emph{Notes:} Each row carries the specification it was measured under, because the table mixes two: the aggregate level and the fit diagnostics are the production spec v4, while the train/post-training split remains spec v3; the resampling distribution has since been rerun jointly at both specifications (Section~\ref{sec:patha}), and the committed spec v3 interval remains the quoted one --- the joint rerun is narrower for the mechanical reason stated there. The spec v3 point (\$915.1 billion) is shown so the median is not read as centred on the spec v4 level printed above it. The 95\% interval spanning zero is the stated reason for the pre-committed headline exclusion; the flexible-learner and isotonic comparators of this section bound what refitting the same information could recover.
\end{tablenotes}
\end{threeparttable}
\end{table}

\begin{figure}[H]
\centering
\includegraphics[width=0.9\textwidth]{fig2_ccf}
\caption{Cross-correlation functions between each estimator's simulated CPR path and the empirical SOMA-derived path, lags $-6$ to $+6$. Under the implemented convention (Section~\ref{sec:pathb}), a peak at negative lag means the empirical series leads and the simulated path trails. Path B's $-3$/$-4$ peak and Path A's and the ABM's negative contemporaneous correlations summarize the same fact from Table~\ref{tab:estimators}: no estimator shows positive contemporaneous alignment, and once the window's shared trend is removed no estimator's co-movement is distinguishable from zero under the moving-block bands.}
\label{fig:ccf}
\end{figure}

A moving-block bootstrap (block length 6, 4,000 replications) puts the 95\% interval on the peak-lag correlation at $[-0.21, +0.49]$ and on the lag-0 correlation at $[-0.13, +0.44]$; with seven lags searched, individual band-point differences in peak $r$ carry no evidential weight. These intervals are not sensitive to the block-length choice in any verdict-relevant way: at block lengths 4 and 8 the peak-lag endpoints move to $[-0.27, +0.39]$ and $[-0.22, +0.51]$, and the lag-0 endpoints to $[-0.11, +0.45]$ and $[-0.17, +0.42]$, every interval still containing zero.

A Theil decomposition formalizes where each estimator's monthly error lives; it is a fit diagnostic, not a timing credential. On levels, the ABM's mean-squared error is 68\% bias share (the level failure already visible in its aggregate), while Path B's is 11.3\% bias and 82.5\% variance share on the corrected amortization basis (its residual monthly error is under-dispersion against a seasonal empirical path, not level bias; run \texttt{coupon\_convention\_amortization}), with Path A intermediate (26.5\% bias, 63.0\% covariance share); after linear detrending the variance share still dominates for Path B (93.7\%). On first differences, Theil's $U_2$ is 1.23 for the ABM, 1.01 for Path A, and 1.00 for Path B: no estimator beats a naive no-change forecast of month-over-month CPR changes, which restates the timing concession above in forecast-accuracy units. The full decomposition by transform, with cumulative dollar-error profiles whose terminal values reproduce Table~\ref{tab:estimators}'s aggregates by construction, is collected in Table~\ref{tab:theil}.

A direct input-timing diagnosis rules out the dating of the rate input and narrows, but does not isolate, the mechanism. In monthly \emph{changes}, the simulated CPR responds to the contemporaneous rate change with the correct sign and near-mechanical strength ($\Delta$-correlation $-0.94$ at lag 0), while the empirical CPR's month-to-month changes show no detectable rate response ($+0.25$, inside the $\pm 0.31$ zero band at $n = 41$): at monthly frequency the empirical path's variation is dominated by non-rate factors the simulation does not model, calendar seasonality above all. One construction asymmetry also belongs in this reading: the empirical CPR back-out retains the income-scaled curtailment component (Section~\ref{sec:method-benchmark}), while the simulated CPR paths exclude it (curtailment enters only as a common dollar flow in the accounting layer), so the two series being correlated differ by a small, slowly-varying component as well as by seasonality.

\begin{table}[H]
\centering
\caption{Danish-leg sensitivity to the buyback discount rate. The diagnostic executed on the Danish-level Berger-recalibrated leg (mean Danish CPR 3.39\%), which is the bracketing anchor of Section~\ref{sec:abm-danish} rather than the production U.S.-intercept leg. Neither leg consumes the present-value classifications the discount rate enters, so the zero deltas hold on the standalone, shared, and printed bases alike; under the production anchor that independence is structural rather than run-verified. The baseline CPR path reproduces the frozen parquet for the leg that ran exactly. The zeros bound the hazard channel only: the leg's dollar conversion credits retired balance at par, an un-priced channel sized in Section~\ref{sec:pathb}.}
\label{tab:discount}
\begin{tabular}{lccc}
\toprule
Discount-rate shift & PV classifications flipped & $\Delta$ CPR path (pp) & $\Delta$ trapped (\$B) \\
\midrule
0 bp (production) & -- & 0.00 & 0.00 \\
$-25$ bp & 2{,}470 of 1{,}683{,}082 & 0.00 & 0.00 \\
$-50$ bp & 7{,}805 of 1{,}683{,}082 & 0.00 & 0.00 \\
$-75$ bp & 17{,}493 of 1{,}683{,}082 & 0.00 & 0.00 \\
$-100$ bp & 32{,}488 of 1{,}683{,}082 & 0.00 & 0.00 \\
$-125$ bp & 55{,}210 of 1{,}683{,}082 & 0.00 & 0.00 \\
$-150$ bp & 92{,}495 of 1{,}683{,}082 & 0.00 & 0.00 \\
\bottomrule
\end{tabular}
\end{table}

The decomposition also carries its time axis. The monthly marginal path (the central-minus-null difference of the two committed simulated legs, whose sum reproduces the $+\$70.3$ billion bit-exactly) is positive in all 42 window months, splits 21.3\%/41.9\%/36.8\% across the window's thirds, and peaks at \$2.29 billion in November 2023, the window's eighteenth month. By coupon bucket the terminal marginal divides 38.0\%/43.7\%/18.3\% across the $<3.0\%$/3.0--4.0\%/$\geq 4.0\%$ groups on 40.3\%/40.4\%/19.3\% of sample balance; Figure~\ref{fig:marginaltiming} draws the stack.

\begin{figure}[H]
\centering
\includegraphics[width=0.94\textwidth]{fig12_marginal_timing.png}
\caption{Monthly lock-in marginal by coupon bucket. Stacked areas are the per-bucket monthly marginals from the month-by-cell capture of the group-ablation engine; the dark line is the aggregate central-minus-null path, whose sum reproduces the $+\$70.3$ billion in-sample marginal bit-exactly. Dotted rules mark the window's thirds; the marker is the November 2023 peak. The stack differs from the aggregate by the decomposition's 1.03\% terminal additivity residual, asserted in the generator.}
\label{fig:marginaltiming}
\end{figure}

\section{Path A Estimation Panel, Coefficient Uncertainty, and Ridge Device}\label{app:ridge}

Three facts about the ridge device of Section~\ref{sec:patha} are verified in the replication package, under both specifications: the spec v3 records stand and each fact has been recomputed under spec v4 (run \texttt{pathA\_ridge\_v4}); the device is unchanged. First, at production magnitudes the penalty is a numerical no-op on the full panel: the two grid points ($\lambda_{\mathrm{ridge}} \in \{10^{-5}, 10^{-4}\}$) return coefficient vectors identical to machine precision from the IRLS warm start (maximum parameter difference $4.5\times10^{-17}$ in the replication artifact; $8.4\times10^{-17}$ under spec v4), so the grid selection is inconsequential and the production coefficients are the IRLS solution to machine precision; ``mildly shrunk'' would overstate the penalty's effect, which is nil. Second, ridge on the full coefficient vector with one stratum absorbed as reference is not invariant to the reference choice: in cold-start refits on resampled panels, swapping the reference stratum can move the standardized rate-gap coefficient by more than its own magnitude (an optimizer-path artifact of penalized refits with the 20 inestimable fixed effects), which is why the bootstrap machinery warm-starts every replication and discards non-converged draws; the warm-started spec v4 swap refit is reference-invariant to $7.3\times10^{-14}$, consistent with the warm-start policy that concession motivated rather than evidence the cold-start sensitivity is gone. Third, the production point itself is not an artifact of these choices: a refit that drops the 20 zero-event strata (423 of 10,176 cells) and fits unpenalized Poisson PML converges, is exactly invariant to the reference choice, and reproduces the spec v3 macro coefficients to within $0.002$ in standardized units (rate gap $+0.6727$ against the spec v3 $+0.6734$ and friction $-0.0373$ against $-0.0371$, third-decimal agreement; burnout $-0.1317$ against $-0.1301$, the larger burnout gap reflecting that dropping the 423 cells changes the burnout regressor's within-stratum demeaning and age-orthogonalization sample). The same refit at spec v4 --- the identical 20 strata and 423 cells drop out, since the month dummies do not change the zero-event set --- is again exactly reference-invariant ($1.1\times10^{-14}$) and reproduces the spec v4 macro coefficients to within $0.0022$: rate gap $+0.6463$ against $+0.6469$, friction $-0.00715$ against $-0.00711$, burnout $-0.1734$ against $-0.1713$, against a tolerance pre-committed at $0.005$, loosened from v3's $0.002$ because the month dummies absorb time variation the friction coefficient formerly carried (run \texttt{pathA\_ridge\_v4}).

Table~\ref{tab:panel} reports the estimation panel this fit is built on, and Table~\ref{tab:bootstrap} the coefficient uncertainty the two bootstrap schemes of Section~\ref{sec:patha} return. Both sit here rather than in the body because Path A is an aggregate-context, non-corroborating exhibit excluded from every headline range by pre-committed specification, and no cell of either table enters a headline quantity.

\begin{table}[H]
\centering
\small
\begin{threeparttable}
\caption{Path A estimation panel.}
\label{tab:panel}
\begin{tabular}{@{}lrr@{}}
\toprule
 & Training & Holdout \\
\midrule
Performance months & Jan.\ 2021--Dec.\ 2023 & Jan.\ 2024--Sep.\ 2025 \\
Stratum-months & 10,176 & 6,077 \\
Prepaid UPB (PML outcome) & \$650.0B & \$131.3B \\
UPB exposure (month-sums) & \$72.5T & \$37.4T \\
\bottomrule
\end{tabular}
\begin{tablenotes}[flushleft]\footnotesize
\item The panel aggregates the full Freddie Mac 2017--2021 origination universe (8.84 million active loans at the window's peak month, not the 75,000-loan sample that Section~\ref{sec:pathb} simulates) to 296 strata, and carries no sampling weights: exposure and the outcome are raw dollars, so the pseudo-maximum-likelihood fit is UPB-weighted by construction (a dollar of prepayment counts the same in any stratum-month). The temporal split at January 2024 is the ridge-selection holdout of the text; the performance window is January 2021 through September 2025.
\end{tablenotes}
\end{threeparttable}
\end{table}

\begin{table}[H]
\centering
\footnotesize
\setlength{\tabcolsep}{4pt}
\begin{threeparttable}
\caption{Path A coefficient uncertainty under two bootstrap schemes (estimated under the prior specification, spec v3; see note).}
\label{tab:bootstrap}
\begin{tabular}{@{}lrrrrccr@{}}
\toprule
Coefficient & Point & Mean & Median & SE & Percentile CI & Basic CI & Draws $\le 0$ \\
\midrule
\multicolumn{8}{@{}l}{\textit{Stratum-cluster scheme (198 of 200 replications converged)}} \\
Rate gap (standardized) & $+0.673$ & $+1.491$ & $+0.763$ & $1.152$ & $[+0.60,\,+4.11]$ & $[-2.76,\,+0.74]$ & 0.5\% \\
Burnout (orthogonalized) & $-0.130$ & $-0.056$ & $-0.130$ & $1.363$ & $[-1.76,\,+1.43]$ & $[-1.69,\,+1.50]$ & 75.3\% \\
Friction & $-0.037$ & $+0.352$ & $+0.015$ & $0.506$ & $[-0.25,\,+1.38]$ & $[-1.46,\,+0.18]$ & 48.0\% \\
\midrule
\multicolumn{8}{@{}l}{\textit{Temporal moving-block scheme, block length 6 months (196 of 200 converged)}} \\
Rate gap (standardized) & $+0.673$ & $+3.313$ & $+2.556$ & $3.845$ & $[+0.44,\,+12.18]$ & $[-10.84,\,+0.91]$ & 0.5\% \\
Burnout (orthogonalized) & $-0.130$ & $+1.229$ & $+0.523$ & $2.586$ & $[-3.16,\,+6.16]$ & $[-6.42,\,+2.90]$ & 40.8\% \\
Friction & $-0.037$ & $+1.050$ & $+0.567$ & $2.802$ & $[-3.71,\,+7.69]$ & $[-7.76,\,+3.64]$ & 37.2\% \\
\bottomrule
\end{tabular}
\begin{tablenotes}[flushleft]\footnotesize
\item Stratum-cluster scheme: the observed vintage--coupon--FICO--LTV strata are resampled with replacement, each resampled stratum receives a fresh fixed-effect label so duplicated strata contribute independent fixed effects, and per-replication standardization (including the burnout orthogonalization, which is refit inside every replication) is rescaled to production standardized units. This scheme is structurally silent about the friction coefficient: the friction index varies only over time, so every replication sees the identical series, and the stratum-scheme friction row understates that coefficient's uncertainty by construction. Temporal scheme: the 36 training months are resampled in moving blocks of six with each month's cross-section kept intact, same refit and rescaling; it is the informative scheme for friction and, partially, burnout; its friction SE (2.80) is roughly $5.5\times$ the stratum scheme's (0.51). The ridge penalty is held at the production $\lambda_{\mathrm{ridge}} = 10^{-4}$ rather than re-selected per replication (re-selection would bootstrap a different, coefficient-plus-model-selection estimator). Replications returning non-finite or implausible ($|\beta| > 20$) coefficients are counted as failures, not draws; the point refit reproduces the production coefficients to four decimals on the same 10,176 training cells. Coefficients are in production standardized units (training SD of the rate gap: 221 bp). The stratum-scheme burnout median coincides with the point estimate at reported precision only (unrounded: median $-0.1297$, point $-0.1301$). Mean, median, and the implied bias (mean $-$ point) are reported because both schemes' rate-gap distributions are heavily right-skewed; under that skew the percentile and basic 95\% intervals disagree about zero, so interval-based significance is construction-dependent and the text relies on sign stability only; but the bootstrap sign-share is itself percentile logic, and the pre-committed bias-respecting test of $H_0\!: \beta_g \le 0$ (run \texttt{patha\_sign\_test}; BCa with the 296-refit jackknife acceleration this note previously deferred, plus a month-permutation leg) does not reject under any construction ($p = 0.093$/$0.412$/$0.241$), so the 99.5\% positive share is a description of a right-skewed resampling distribution, not a sign test, and no sign claim for this coefficient survives.
\end{tablenotes}
\end{threeparttable}
\end{table}

\section{Dynamic-Fit Decomposition}\label{app:theil}

Table~\ref{tab:theil} collects the Theil fit diagnostics summarized in Section~\ref{sec:pathb}: $U$ statistics, the bias/variance/covariance decomposition of each estimator's mean-squared error on three transforms of the monthly CPR paths, and the terminal cumulative dollar errors. The $\beta_1 = 0$ null's cumulative-error profile follows from the committed legs by arithmetic: its terminal cumulative error is $-\$16.6$ billion against the empirical path, versus Path B's $+\$53.8$ billion; the null under-runs the empirical roll-off by approximately the lock-in marginal, by construction, and the full monthly profile is committed with the marginal-timing data of Section~\ref{sec:identification}.

\begin{table}[H]
\centering
\caption{Theil dynamic-fit decomposition for the three estimators' monthly CPR paths over the 42-month QT window. Fit diagnostics, not timing credentials: no estimator's co-movement survives detrending (Section~\ref{sec:pathb}), and $U_2 > 1$ means worse than a naive no-change forecast of month-over-month changes.}
\label{tab:theil}
\footnotesize
\begin{threeparttable}
\begin{tabular}{@{}lrrr@{}}
\toprule
 & ABM & Path A & Path B \\
\midrule
Theil $U_1$, levels & 0.412 & 0.438 & 0.271 \\
Theil $U_1$, detrended & 0.807 & 0.839 & 0.980 \\
Theil $U_2$, first differences & 1.226 & 1.009 & 1.003 \\
\midrule
\multicolumn{4}{@{}l}{\textit{MSE shares, levels (\%)}} \\
\quad Bias & 68.1 & 30.6 & 11.3 \\
\quad Variance & 0.4 & 10.5 & 82.5 \\
\quad Covariance & 31.5 & 58.8 & 6.3 \\
\multicolumn{4}{@{}l}{\textit{MSE shares, first differences (\%)}} \\
\quad Bias & 0.1 & 0.3 & 0.1 \\
\quad Variance & 14.0 & 71.4 & 96.8 \\
\quad Covariance & 85.8 & 28.3 & 3.1 \\
\multicolumn{4}{@{}l}{\textit{MSE shares, linearly detrended (bias $\equiv 0$, \%)}} \\
\quad Variance & 0.5 & 43.6 & 93.8 \\
\quad Covariance & 99.5 & 56.4 & 6.2 \\
\midrule
Terminal cumulative error (\$B) & $-673.8$ & $+164.1$ & $+53.8$ \\
\bottomrule
\end{tabular}
\begin{tablenotes}[flushleft]\footnotesize
\item $U_1 \in [0,1]$, where zero is a perfect fit; $U_2$ is relative to the naive no-change forecast. Path A and Path B columns are computed against the hazard-side empirical CPR back-out at the pool's note-rate WAC (the corrected amortization basis; run \texttt{coupon\_convention\_amortization}, committed 2.5\%-convention values in \texttt{theil\_data.json}); the levels-bias shares are the cells the correction moves ($+4.6$pp for Path B), the timing rows move at or below a thousandth. The ABM column retains its committed basis --- the converted ABM reconstruction failed its own replay gate and is not quoted (the ABM's cumulative-error leg uses its own committed back-out). Terminal cumulative errors equal each estimator's trapped-liquidity aggregate minus the empirical \$764.7 billion on its own accounting basis, reproducing Table~\ref{tab:estimators} by construction. Artifact \texttt{theil\_data.json}; inputs and correlations are gated to the committed run manifests.
\end{tablenotes}
\end{threeparttable}
\end{table}

Grouped calibration localizes the same verdict (run \texttt{grouped\_calibration}). With the frozen Path A cell predictions binned into five exposure-weighted prediction-CPR quintiles (edges fit on the training window only) crossed with calendar half-years, the training halves calibrate at a pooled predicted-to-realized ratio of 1.000 (exact by construction of the exposure-offset fit), with per-half ratios spanning 0.80--1.17, while the post-2024 holdout drifts into monotonically worsening under-prediction: 0.64, 0.57, 0.44, 0.43 across 2024H1--2025H2, pooled predicted 2.24 against realized 4.21 points. The departure is temporal drift rather than a reliability-slope failure: pooled across periods the four upper prediction bins calibrate at ratios 0.94--1.11, and the shortfall concentrates in the lowest-prediction bin (43.5\% of exposure), where the 2024--25 cells accumulate: the finer-resolution form of the $U_2 > 1$ diagnosis above. This is a cell-panel CPR object, distinct from the cumulative dollar-runoff split of Section~\ref{sec:patha}, and like every Path A diagnostic it moves no headline figure.

\section{Composition Shift: Estimation Sample versus SOMA Book}\label{app:composition}

Table~\ref{tab:composition} reports a formal covariate-shift analysis between the Freddie Mac data behind the hazard paths and the SOMA book the benchmark spans, on the dimensions both sides carry (run \texttt{composition\_shift}). The book side is the New York Fed's CUSIP-level SOMA holdings parsed by the same production code that builds the ABM cohorts, anchor-gated to the committed marginals (book WAC 2.49\%, Ginnie share 20.4\%, term split, vintage-group shares); the API also served the historical June 29, 2022 as-of, whose composition is consistent with those anchors and is reported in the artifact as a non-gated secondary block. The shifts on coupon, vintage, and agency are large (the expected result, disclosed in prose throughout the paper), and each carries a committed, signed bound as its consequence, listed in the table's final column. The joint agency$\times$vintage cells settle the one netting question the marginals leave open: of the 2022 vintage's 23.1\% of book face, 20.7\% is conventional and 2.4\% is the Ginnie intersection, so vintage and agency marginals cannot be added without double-counting that cell. Geography, loan size, FICO, and LTV do not exist at the CUSIP level of the Fed's published holdings (par value and security description only); panel (b) supplies the deliverable substitute (the Freddie and Fannie loan-level samples agree to negligible shift on exactly those dimensions), and the hazard specification's exposure to the loan-level covariates is bracketed directly by the printed covariate-block endpoints (99.2\%/92.2\% shared-basis, Section~\ref{sec:pathb}).

\begin{table}[H]
\centering
\caption{Composition comparison. Panel (a): the two Freddie objects the hazard paths use --- the estimation universe Path A is fitted on (exposure-weighted at the June 2022 window open) and, indented beneath it, the 75,000-loan draw Path B simulates (by count and on its surviving balance at the window open) --- against the SOMA agency-MBS book (CUSIP-level holdings, anchor-gated). The three are distinct populations and none is mixed within a row; only the unindented rows carry a PSI against the book. Panel (b): cross-agency stability of the loan-level samples on the dimensions the book cannot supply. PSI is the population stability index $\sum_i (p_i - q_i)\ln(p_i/q_i)$; values above 0.25 are conventionally read as large shift, below 0.10 as negligible.}
\label{tab:composition}
\footnotesize
\begin{threeparttable}
\begin{tabular}{@{}p{2.1cm}p{4.7cm}p{4.7cm}p{3.6cm}@{}}
\toprule
Dimension & Freddie side (\%) & Book (\%) & PSI; committed consequence \\
\midrule
\multicolumn{4}{@{}l}{\emph{(a) Freddie side vs SOMA book: estimation-universe rows, with Path B's draw indented beneath}} \\
Agency & UMBS 100.0; Ginnie 0.0 & UMBS 79.6; Ginnie 20.4 & 2.54; static bound \$20--47B, overlay differential $+\$38.5$B \\
Vintage & pre-2017: 0.0; 2017--19: 12.5; 2020: 37.3; 2021: 50.2; 2022: 0.0 & 10.6; 6.0; 16.3; 43.9; 23.1 & 4.33; vintage residual bound \$11.7 billion \\
\quad Path B draw & \emph{count}: 0.0; 60.0; 20.0; 20.0; 0.0 \newline \emph{surviving balance}: 0.0; 27.0; 30.7; 42.3; 0.0 & & no committed draw-vs-book vintage PSI; equal allocation per origination quarter, never reweighted (Appendix~\ref{app:params}) \\
Coupon & $<3.0\%$: 18.0; 3.0--4.0\%: 68.2; $\geq 4.0\%$: 13.8 & 73.9; 18.7; 7.3 & 2.29; full-book reweighting 107.0\% $\to$ 109.1\% \\
\quad Path B draw & \emph{count}: 7.1; 34.7; 58.2 \newline \emph{surviving balance}: 14.0; 58.8; 27.2 & & 3.38 on the draw's origination-balance shares; draw mean note rate 3.86\%; the reweight moves the level, not the marginal \\
Term & 30-year collateral (15-year face excluded from hazard cohorts, disclosed in Section~\ref{sec:limitations}) & 30yr 90.7; 15yr 9.1; other 0.2 & book-only dimension \\
\midrule
\multicolumn{4}{@{}l}{\emph{(b) Freddie vs Fannie loan-level samples (75{,}000 loans each, origination-balance weighted)}} \\
State & \multicolumn{2}{l}{PSI 0.018} & negligible \\
Loan size & \multicolumn{2}{l}{PSI 0.000} & negligible \\
FICO & \multicolumn{2}{l}{PSI 0.005} & negligible \\
LTV & \multicolumn{2}{l}{PSI 0.000} & negligible \\
\bottomrule
\end{tabular}

\end{threeparttable}
\end{table}

\noindent{\footnotesize \emph{Notes to Table~\ref{tab:composition}.} Three populations appear in panel (a) and none is mixed within a row: the unindented rows are the hazard estimation universe, exposure-weighted at the window open; the indented rows are the 75,000-loan Path B draw, by count and on the surviving balance the aggregation applies; the Book column is SOMA face. All three sides use rounded half-point coupon buckets, but the two Freddie sides bucket borrower note rates while the Book column buckets security pass-through coupons --- a wedge of $+0.80$ points --- so the buckets are aligned in width, not in object. Within the draw the $\geq 4.0\%$ share moves with both conventions and they must be read separately: on count weights, rounding to half points raises it from 47.0\% to 58.2\%; holding the rounding fixed, the weighting lowers it from 58.2\% by count to 53.7\% on origination balance and 27.2\% on surviving balance. The WAC pair is mixed-basis, as Section~\ref{sec:pathb} records: the draw's 3.86\% is the unweighted mean of borrower note rates, the book's 2.49\% the face-weighted mean of security pass-through coupons. Every shift is signed in the same conservative direction: the out-of-sample segments prepaid faster, so trapped liquidity is if anything overstated (Section~\ref{sec:pathb}). That signing covers observed speeds in the segments outside the draw's support; it does not cover what the draw's composition within 2017--2021 does to the floor's bind share, which no committed run scores (Appendix~\ref{app:params}). Panel (b) uses the committed Fannie replication sample under its aggregates-only license posture. Artifact \texttt{composition\_shift\_results.json}.\par}

\nocite{cfpb_atr,frbny_soma,fred_actliscou,fred_wshomcb,fred_dspic96,fred_umcsent}

\section{Floor-Sweep and Contamination Mechanics}\label{app:floormech}

This appendix collects the mechanics behind the calibration box of Section~\ref{sec:robustness-floor}: the floor-level sweep, the marginal as a curve in the imported elasticity, the separation of the two refinance contaminations behind the off-window floor anchors, and the out-of-sample anchor construction. The three exhibits themselves are read where the headline is set, in Section~\ref{sec:hazard} (Tables~\ref{tab:floorband}, \ref{tab:lowband}, and~\ref{tab:oosfloor}).

Table~\ref{tab:floorband} reports the lock-in marginal at every cell of the resulting calibration box. Three facts shape the reading. First, the pre-committed stability threshold ($\pm 2$ points across the empirically plausible 3--5\% floor range) fails: at the central elasticity the marginal spans $+10.8$ to $+5.5$ points across that range. The mechanism is structural rather than statistical: the hazard specification takes a maximum against the floor, so wherever the floor binds the elasticity is inert, and raising the floor mechanically crowds out the lock-in channel. Second, both pre-committed monotonicities hold, the marginal decreasing in the floor within every band row and increasing in the band at every floor, and the box maximum lands at the predicted corner: $+13.2$ points ($+\$100.7$ billion) at a 2\% floor on the 7.7\% band edge. Third, every cell is positive, with a minimum of $+2.1$ points at a 6\% floor on the 5.5\% band edge---a consistency check on the specification, not an identifying result, since the window's rate configuration fixes the sign before any floor or band is chosen (Section~\ref{sec:identification})---and every marginal in the box is basis-invariant because the shared-accounting netting is common to the central and null legs of each cell (Section~\ref{sec:robustness-hybrid}). The verification reading of Section~\ref{sec:pathb} accordingly weakens from a point to a bounded interval: the lock-in elasticity's identified contribution is $+9.2$ points at the production calibration, $+5.0$ to $+12.6$ points across the empirically plausible floor range, and $+2.1$ to $+13.2$ points across the full box, with the off-window re-measurement below placing the headline reading at $+4.3$ to $+6.8$ points, while the aggregate level remains floor-dominated throughout.

The box sweeps the elasticity only over its published band, and the band's lower edge is a literature boundary rather than an evidential one: the in-sample estimates cannot reject zero for the elasticity (Section~\ref{sec:pathb}), so a reader who weighs that non-rejection over the imported band needs the marginal at values the box never visits. A pre-committed extension (run \texttt{band\_low\_extension}; grid, parity gates against the committed production and off-window cells, and a monotonicity check fixed in the script header before any run) re-runs the paired legs at quarterly-decline values $\delta \in \{0, 1.0, 2.0, 3.0, 3.25, 4.0, 5.0\}\%$ below the band at both anchor floors, and Table~\ref{tab:lowband} reports the resulting curve beside the band's own cells. The curve is monotone at both floors, passes through zero at $\delta = 0$ by construction, and is concave at the off-window floor through the censoring mechanism of Section~\ref{sec:identification}: raising the response crowds the hazard into the floor (the bind share climbs from 35.8\% to 73.5\% along that curve), so each increment of elasticity buys less marginal. At half the adopted central, $\delta = 3.25\%$, the marginal is $+5.1$ points ($+\$39.1$ billion) at the in-sample floor and $+3.5$ points ($+\$26.5$ billion) at the headline off-window floor. This curve prices the dependence on the imported elasticity rather than identifying it: how small the marginal could be, if the moving response is weaker than the literature band, is read off the curve, and the identified content remains the bounded range stated in Section~\ref{sec:identification}.


Those two checks vary the floor's level and its functional form; a third pre-committed check interrogates the level's provenance---whether the 4\% figure is an artifact of being measured in-window on the locked-in cohorts the marginal later decomposes (runs \texttt{out\_of\_window\_floor} and \texttt{oos\_identification}; spec, anchor definitions, acceptance thresholds, and parity gate fixed in each script header before its runs, the ex-ante convention Appendix~\ref{app:ledger} records). Run in-window on the deeply out-of-the-money 2023--24 discount cohorts, the identical exposure-weighted-CPR procedure reproduces the production floor at 3.8--3.9\% annual CPR, flat across seasoning cuts and validating the measurement where it is used. Applied out-of-window to pooled 2017--2019 discount-cohort turnover, before both the QT window and the 2020--21 refinancing wave, it recovers a materially higher 6.1\%. That pooled reading, however, does not survive its own diagnostics, and I do not take it as the operative off-window anchor.

Two contaminations are separable, and separating them is what changes the reading. The first is seasoning. A floor meant to capture strictly involuntary turnover---death, divorce, relocation---should be approximately flat in loan age, but the pooled 2017--2019 out-of-the-money turnover \emph{rises} monotonically with it: 3.68\% below twelve months, 5.78\% at twelve to twenty-four, and 10.99\% at twenty-four to thirty-six. The second is rate regime. Splitting the pooled window at the 2018/2019 turn separates a rising-rate leg (the 30-year rate moved 4.03\% to 4.64\% over 2018, suppressing refinancing incentive) from a falling-rate leg (4.46\% to 3.72\% over 2019, a refinancing wave). The legs read very differently: the 2018 leg returns 4.70--5.33\% across the out-of-the-money gap grid at the age-$\geq$-12 cut, while the 2019 leg returns 6.91\% and carries the entire 10.99\% age-$\geq$-24 cell. The pooled 6.1\% is therefore a refinancing-inflated composite, not a baseline turnover read, and the 2019 leg is disqualified on both diagnostics at once. The rate-regime half of that disqualification, as stated, compares the two legs at different discount depths; the matched-depth reconciliation below repairs the comparison and reports what survives it, which is less than the sentence just written implies.

What remains is the 2018 rising-rate leg, and its status must be stated precisely because the paper's off-window anchor now rests on it. Its mature-age seasoning test is \emph{not computable}: the leg contains no out-of-the-money cohort-months at all past the twenty-four-month age cut, because a loan that is out of the money while rates are rising is necessarily a recent, low-coupon origination and therefore young. The 2018 leg is thus clean by rate-regime construction rather than clean by test, and its one computable age contrast (3.78\% to 5.33\% across the twelve-month cut) is confounded by the PSA ramp, on which young loans genuinely turn over less, so it is not evidence about the contrast the leg is missing. The out-of-agency read is: in every cell where a mature contrast exists at all, mature turnover is the higher one, and Section~\ref{sec:identification} reports by how much. A second limitation concerns depth rather than age. The gap-$\leq$-two-point cell at which the in-window anchor is measured does exist off-window: the 2018 leg's gap runs to $-2.87$ points, and it holds 35 cohort-months reading 3.440\% annual CPR. Grids that stop at half a point do so because their thresholds were fixed at zero, a quarter point, and half a point in the script header, not because the rate history forbade anything deeper. The true limitation is exposure collapse, and a matched-depth run (\texttt{matched\_depth\_reconciliation}) measures it: down the 2018 leg's depth ladder the share of its gap-$\leq$-0 exposure that survives runs 100\%, 83.7\%, 72.9\%, 35.9\%, 11.3\%, 3.9\%, and 0.02\%, so the two-point cell carries \$0.282 billion of \$1{,}264.5 billion. That cell is present and worthless, and an ex-ante support rule---a depth is admissible for a leg only if the leg retains at least 10\% of its gap-$\leq$-0 exposure there---excludes it before any reading is taken. These are the reasons the off-window floor resolves to a \emph{range} rather than a point, and I report it as one.



\section{Behavioral Extension Mechanics}\label{sec:robustness-extensions}

This appendix reports the mechanism behind each of the three behavioral extensions summarized in Section~\ref{sec:abm-results}. As Appendix~\ref{sec:method-frictions} describes, loss aversion steepens the penalty every household faces, which on its own would suppress simulated mobility below the historically observed 4--5\% baseline turnover floor; to preserve that floor, I re-calibrate the mobility-desire parameter upward, nearly tripling it. That recalibration mechanically inflates the model's overall responsiveness to rate gaps, so the jump to 54.9\% is an artifact of holding the turnover floor fixed, not a validation of loss aversion as a mechanism. Because that floor is anchored at maximum lock-in, the sign of the recalibration's effect on rate-sensitivity is not fixed: a specification change that eases turnover at the anchor lowers the recalibrated mobility scale (in the real-covariate cross-design of Section~\ref{sec:robustness-crossdesign} it falls to 36,086 from the synthetic 43,883) and deflates rather than inflates rate-sensitivity, so the headline recovery, itself resting on a floor-anchored calibration, inherits that indeterminacy, which is why the ABM's negative verdict is read only alongside the cross-design hedge and never on its own. The second extension, income-scaled curtailment, pulls the estimate down to 45.1\%; because curtailment is a strictly additive drag on prepayment, this decline is what I predicted ex ante from a diagnostic correlation between the ABM's over-prediction and disposable-income growth (Appendix~\ref{sec:method-frictions}). The third extension, multi-vintage coupon composition, pulls the estimate down further, to 33.7\%. A lower weighted-average coupon widens the gap to prevailing market rates, and my hypothesis was that this should increase predicted trapped liquidity, not reduce it. The inversion suggests the ABM's CPR mechanism is not translating a wider coupon gap into slower prepayment the way the household-choice framework predicts it should. I have not diagnosed this inversion, and its implication cuts in both directions: an aggregate response to the rate gap running backward is equally consistent with ``household choice explains little'' and with ``this specific ABM's gap response is mis-specified,'' which is why the ABM's recovery share is never quoted in this paper without the cross-design qualifier of Section~\ref{sec:robustness-crossdesign}, why the paradigm attribution inherits that section's hedge, and why the two-estimator contrast, not the ABM's point estimate, is what the paper asks the reader to carry. I cannot rule out that a differently specified household model would perform better, but that possibility does not bear on the results reported in Section~\ref{sec:abm}.

The wait-and-see freeze now carries a direct sensitivity run, scored on the frozen Danish-anchor production accounting (trapped \$84.5 billion at the production 150 bp trigger and 20\% share; parity gate $+\$0.02$ billion). Removing the freeze entirely drops trapped liquidity to \$44.6 billion (an isolated contribution of $+\$39.9$ billion, $+5.2$ points of benchmark), and moving the six-month trigger to 100 or 200 bp spans \$105.9 to \$70.5 billion, so the trigger's placement matters at the few-points scale and its author-calibrated status remains a disclosed limitation. Sweeping the freeze share with the floor anchor retained gives \$65.1 billion at a 10\% share, \$81.5 billion at the externally implied 18.3\%, and \$103.3 billion at 30\%: re-anchoring the share to the Liebersohn--Rothstein-implied value moves recovery by only $-0.4$ points of benchmark, and the calibrated mobility scale and floor CPR are bit-identical at every share, confirming that the freeze cannot compromise the baseline turnover floor. The front-end DTI gate receives the parallel treatment: moving the 43\% wall to 36\% or 50\% on the same frozen surface moves ABM trapped liquidity to \$153.3 and \$119.7 billion (from \$84.5 billion at the production 43\%), with the turnover floor re-anchored into the empirical 4--5\% band at every threshold, so the gate is material to the ABM's own estimate while never breaching the floor. That sweep is non-monotone in the wall, and the production 43\% sits at its minimum, which invites the reading that the paper's own parameter choice is the one minimizing its own estimate. It does not survive decomposition (run \texttt{dti\_nonmonotonicity\_decomposition}). The floor-retention rule re-calibrates the mobility-desire scale only where the fixed-scale floor would leave the 4--5\% band, and across these three thresholds that happens exactly once: at 50\%, where loosening the wall lifts the fixed-scale floor to 5.09\% and the search cuts the scale from 43{,}883 to 41{,}934. Holding the scale at its production value instead, the sweep is monotone---\$153.4, \$84.6 and \$46.4 billion at 36, 43 and 50\%---and the 50\% leg's entire rebound is the re-calibration, worth $-\$73.4$ billion. The frozen and re-calibrated legs are bit-identical at 36 and 43\%, where the rule is a no-op, which is what identifies the channel. The wall is therefore monotone in the direction the mechanism implies, and the published non-monotonicity is the free parameter absorbing the floor constraint rather than a property of the DTI gate. Nothing here moves a headline: the ABM is a stress-test device and the hazard marginal carries no DTI gate.

\section{Duration Quantification: Approximate Portfolio Weighted-Average Life}\label{sec:robustness-wal}

The mechanism the paper quantifies is duration extension, so I express it in maturity units. Table~\ref{tab:wal}, printed with the cap-design arithmetic in Section~\ref{sec:discussion-capdesign}, reports the SOMA book's approximate weighted-average life (WAL) under scheduled amortization only, under the empirical CPR path, and under each estimator's implied path, at the QT window's start and end; this appendix documents its construction and reads it against each estimator.

Read against the empirical column, each estimator's CPR error maps into effective-maturity error: the ABM's over-predicted prepayment implies a book roughly 3.4 years too short at window start, Path A's under-prediction a book 1.5 years too long, and Path B lands within 0.3 years. Relative to the no-shock world of the table's final row (prepayments continuing at the 22.81\% mean 2021 realized CPR, implying a 3.4-year WAL) the empirical 9.4-year WAL is a 6.0-year extension: that duration gap is what the paper's dollar benchmark prices month by month. The two Danish rows read the counterfactual in the same maturity units: feeding the Danish legs' mean CPRs through the identical calculator gives a rule-only WAL of 9.1 years at the window's start against the U.S. leg's 9.7 (a 0.6-year rule-only shortening, with the Danish-level bracketing case at 10.8), and the rule-only leg's extension is additionally rate-invariant on the rate-rise side by the identified moving-flatness (Section~\ref{sec:abm-danish}). I deliberately price extension rather than convexity: option-adjusted duration and convexity require a term-structure and prepayment-option model outside this paper's scope, and ``convexity-neutral'' is used strictly in this rate-rise sense.

\section{Fifteen-Year MBS Fold-In and Path B Rate-Gap Units}\label{sec:robustness-pipeline}

Two features of the production pipeline set headline figures cited throughout this paper: the treatment of fifteen-year MBS, and the units in which Path B applies its rate-gap coefficient. Appendix~\ref{app:ledger} catalogues the figures each supersedes.

The first is the fifteen-year MBS fold-in. The SOMA cohort parser buckets 15-year holdings separately by coupon, never merging 15-year and 30-year buckets, and 15-year scheduled amortization enters both the empirical CPR back-out and the ABM's simulated roll-off. Restricting to 30-year MBS covers 90.6\% of face value, leaving the 15-year book (approximately 9.1\% of holdings) to contribute to total portfolio size but not to cohort-correct amortization or CPR weighting; folding it in brings coverage to 99.8\% of SOMA face value at a SOMA weighted-average coupon of 2.49\% (Table~\ref{tab:wal}'s note reconciles these coverage components with the audit parse's 90.7\%/9.1\% term split).

I tested, but did not adopt, a fully native treatment in which 15-year loans also draw voluntary CPR from a 15-year-specific surface. This produced out-of-bounds predicted values: the ABM's payment-delta gate (the change in \eqref{eq:annuity}'s monthly payment upon refinancing at the prevailing rate, which drives loss-aversion and mobility-desire responses elsewhere in this paper) breaks down for seasoned 15-year borrowers, whose same-term replacement payment is nearly flat relative to their current payment. Native 15-year surfaces predicted 32\%--61\% CPR against an empirical range of roughly 5\%--8\%, enough to flip the ABM's aggregate trapped-liquidity estimate to a large negative value, far outside the empirically observed range. Production therefore uses a structural-only fold-in: 15-year cohorts contribute real weights and real scheduled amortization, while voluntary CPR for those cohorts is drawn, in the frozen headline specification, from the same-coupon 30-year surface. This is a scope limitation for the production headline rather than for the framework: a 15-year-specific behavioral gate (a pure rate-lock penalty for short-amortization loans) is implemented and is the basis of the Berger run's calibration-stage freeze (Table~\ref{tab:danish}, note b); its adoption for the production headline is deferred to the pre-submission freeze.

The second is the units in which Path B applies its rate-gap coefficient $\beta_1$, calibrated from \citet{liebersohn2024} via the survival-function transform described in Section~\ref{sec:pathb} and expressed as a log-hazard-ratio per 100 bp increase in the rate gap. Applied to the gap in decimal rather than per-100 bp units the coefficient is numerically inert, yielding a multiplier of approximately $\times$1.0026 at a representative QT-era state (a 3.0\% coupon against a 6.8\% market rate) against the intended suppression of approximately $\times$0.77. Production applies it as $\exp(100\,\beta_1\,\tilde g_{i,t})$, with $\tilde g_{i,t}$ the decimal form of \eqref{eq:pathB}'s rate gap (equivalently $\exp(\beta_1 g_{i,t})$ in \eqref{eq:pathB}'s 100 bp units; the gap is negative when the borrower is locked in, so \eqref{eq:beta1}'s positive $\beta_1$ suppresses prepayment), consistent with Path A's convention of a positive coefficient applied to the gap in comparable units, and expresses the Danish counterfactual's rate-gap function in rate-equivalent units so that the buyback discount properly offsets the locked-in spread (live in the mechanism-substitution variant; the production Danish leg's prepay response, under either Berger-recalibrated anchor of Section~\ref{sec:abm-danish}, does not consume that function's classification, which is nonetheless computed, as Section~\ref{sec:pathb} discloses). Section~\ref{sec:pathb} reports the results under this convention; the superseded \$747 billion (97.7\%) microsimulation is catalogued in Appendix~\ref{app:ledger}.

Both are structural or mechanical rather than changes in modeling philosophy, and both are additive: disabling either (30-year-only filtering; the decimal-unit $\beta_1$ convention) exactly reproduces the corresponding superseded figures catalogued in Appendix~\ref{app:ledger}, which I treat as evidence that each is isolated to its intended code paths.

\section{Seasonalizing the Turnover Floor: Floor Dispersion, and Why No Timing Claim Follows}\label{sec:robustness-seasonalfloor}

The three checks above vary the floor's level, its functional form, and its provenance, holding it constant within the window in each case. Two runs in this paper let the floor vary within the window instead: the cyclical variant of Section~\ref{sec:pathb}, which indexes it to the rate path, and the pre-committed seasonal check taken up here, which indexes it to the calendar. Both move the marginal, and this check is the one built to isolate why. It was designed to answer two questions at once: whether letting the turnover floor vary by calendar month moves the lock-in marginal, and whether the resulting monthly path tracks the empirical benchmark well enough to license a timing claim (run \texttt{seasonal\_floor\_timing}; spec, leg definitions, the placebo families, and a frozen timing rule fixed in the script header before any run executed). The answers are asymmetric, and the second is the more consequential: the marginal does move, but not for the reason the design was built to test, and the timing rule is met without being informative.

The calendar profile is a twelve-cell exposure-weighted annualized turnover CPR measured on deeply out-of-the-money cohort-months. It is estimated on the full cohort-month panel (7{,}756 cohort-months, \$74.2 trillion of exposure, pooled annual CPR 4.187\%), not on the QT window; the window-restricted variant differs in four calendar cells and drives nothing here, and I state the basis because an earlier module docstring described the profile as in-window and was wrong. The profile peaks in June with a peak-to-trough ratio of 2.13. Four legs are run, each as a paired central and $\beta_1 = 0$ null: the production constant 4.0\% floor, a constant 4.5\% floor, a \emph{shape-only} leg carrying the calendar profile with its exposure-weighted mean pinned at 4.0\%, and the pre-committed 4.5\%-mean seasonal leg. The shape-only leg is the one that isolates the calendar, because the pre-committed leg confounds shape with the level move the committed sweep already prices. Every leg re-runs the production microsimulation with only \texttt{INVOLUNTARY\_CPR\_ANNUAL} varied by calendar month; the constant-4.0\% control reproduces the committed production artifacts bitwise (central \$818.530 billion, null \$748.185 billion, marginal $+\$70.345$ billion, $+9.20$ points).

Giving the floor a calendar profile at a held-constant 4.0\% exposure-weighted mean moves the marginal from $+\$70.345$ billion to $+\$65.448$ billion, $+9.20$ to $+8.56$ points: a shape effect of $-\$4.897$ billion. A permutation family disposes of the natural reading of that number. Assigning the \emph{same} twelve floor values to scrambled calendar months, so that the multiset of values is preserved exactly and only the month labels change, retains on average 89.2\% of the shape effect across fourteen draws (69.9\% to 102.8\%), with permuted marginals of \$65.311 to \$66.924 billion against the true \$65.448 billion. Calendar order is very nearly irrelevant to the movement, so the movement is not a calendar effect.

What it is instead is a Jensen effect operating through the floor's combination rule. Production combines the floor and the voluntary hazard by a hard maximum (\texttt{FLOOR\_MODE = 'max'}, the semantics \eqref{eq:pathB} and the form test above discuss), and at the production floor the floor binds in 36.3\% of evaluated loan-months. Under a maximum, dispersion in the floor is one-sided: in months where the voluntary hazard already exceeds the floor, the low floor cells are truncated away and contribute nothing, while the high cells bind and lift the effective hazard. Any dispersion in the floor therefore raises book-wide prepayment regardless of which month carries which value, which is exactly what the permutation family measures. A convexity-matched control separates the small remaining piece: because the annual-CPR-to-monthly-hazard transform is itself convex, the twelve-cell profile is equivalent in level to a flat 4.0090\% floor rather than a flat 4.0000\% one, and running that flat leg gives $+\$70.156$ billion, a $-\$0.189$ billion move. Level leakage is thus 3.9\% of the shape effect and floor dispersion under the maximum rule is the remaining 96.1\%. I therefore report this as a floor-dispersion and functional-form result. It is not evidence that seasonality moves the marginal.

Level and shape do not separate cleanly, and I do not present an additive split as though they did. Measured against the constant-4.0\% control, the level move contributes $-\$12.704$ billion ($-1.66$ points), the shape move $-\$4.897$ billion ($-0.64$ points), and the two together $-\$16.040$ billion ($-2.10$ points), leaving an interaction residual of $+\$1.560$ billion ($+0.20$ points). On the simulated monthly paths the residual is larger relative to what it qualifies: mean absolute interaction residual 0.164 points of monthly CPR against mean absolute total movement 0.533 points, so the interaction is 30.9\% of the total movement. Any level-plus-shape decomposition of this floor must carry that residual; quoting either component alone misstates both.

What the experiment does buy is a narrow bound. Holding the turnover floor flat within the window is a modeling convenience, and this run bounds how much of the constant-floor marginal that convenience was responsible for: at most $\$4.897$ billion of the $+\$70.345$ billion in-sample marginal, 7.0\% of it, which is 11.5\% of the $+\$42.6$ billion headline marginal the same absolute bound is measured against, and 96.1\% of even that is the maximum rule's response to dispersion rather than anything about calendars. The identified lock-in marginal is not an artifact of a flat floor. That is the whole of what the seasonal design establishes, and it establishes it about levels only.

\paragraph{The timing concession.} The design's second purpose was to test whether a seasonalized floor earns a monthly-timing claim, under a rule frozen before the runs: Theil $U_2$ on first differences below one \emph{and} a positive detrended lag-0 correlation against the empirical benchmark. That rule is met as written by both seasonal legs (shape-only $U_2 = 0.995$, $r = +0.016$; pre-committed leg $U_2 = 0.990$, $r = +0.070$), and I report it as met, because it was frozen ex ante and I will not reopen a rule after seeing which way it fell. I nonetheless claim nothing from the pass, because five diagnostics computed afterwards show the rule cannot discriminate. First, the benchmark carries it. The SOMA back-out has four exact-zero months (June 2022, February 2023, April 2024, September 2025), three of them immediately followed by a spike of up to 14.01\% annualized. The zeros are now traced to their mechanism: each is clip-induced, a month whose reported rolloff falls below the scheduled-amortization floor so that the back-out's non-negativity clip pins it at exactly zero, with the displaced mass posting into the following month's reading (unclipped values $-1.59$/$-0.49$/$-0.84$/$-0.84$ points; run \texttt{h1\_zero\_months\_diagnosis}) --- a month-boundary allocation artifact of the reporting series, not a behavioral observation. The June 2022 zero opens the window, is followed by an ordinary month, and so contributes one first difference where the others contribute two. Those pairs touch 7 of 41 first differences in total but carry 0.813 of the empirical differenced signal and 0.796 to 0.811 of every leg's $U_2$ numerator. Repairing them by reallocating each spike evenly across its pair, conserving the pair total, flips both seasonal legs from pass to concede and turns the shape-only leg's detrended correlation from $+0.016$ to $-0.134$. Second, a calendar-rotation placebo moves the true profile to the wrong months while holding its amplitude and dispersion exactly; the rule is cleared by 1 of 11 wrong-month rotations for each seasonal leg, and the five-month rotation \emph{beats} the true profile on $U_2$ (0.9934 against 0.9947). Third, 3 of the 14 scrambled calendar orders above clear the rule, and the best of them attains $U_2 = 0.975$, better than either real seasonal leg. Fourth, and decisively, the $\beta_1 = 0$ null clears the rule on both constant-floor legs, the 4.0\% null doing so with $U_2 = 0.9997$ and a detrended correlation of $+0.115$, stronger than either seasonal central leg's. A criterion satisfied by a specification with the lock-in elasticity removed cannot be evidence for the lock-in mechanism. Fifth, the correlations are not distinguishable from zero on 42 months in any case ($p = 0.921$ and $p = 0.661$), and the leg that isolates the mechanism is the fragile one: the shape-only leg's verdict flips on 5 of 42 single-month deletions while the pre-committed leg's flips on none.

Two further diagnostics point the same way without being needed for the verdict. The seasonal legs concentrate 79.7\% (shape-only) and 89.4\% (pre-committed) of their simulated detrended variance in the twelve-month calendar harmonic, while the empirical benchmark carries 0.8\% of its own detrended variance there: the model acquires calendar structure the benchmark can neither confirm nor refute. And on the committed cross-correlation metric the seasonal legs peak at lag $+6$ with a negative correlation, on the boundary of a 42-month window, which I mark non-interpretable rather than read.

The right reading of this is a finding about the criterion, not a failure of the experiment. The frozen rule was a reasonable ex-ante test and it turns out to have no power against the alternatives that matter, because a lumpy benchmark, wrong-month profiles, scrambled calendars, and an elasticity-free null all satisfy it. I therefore restrict every claim in this paper to levels: no estimator here earns a monthly-timing claim, which extends to the seasonalized floor the levels-only concession the production microsimulation's monthly path already carries, where the raw lag-0 correlation does not survive detrending (Section~\ref{sec:pathb}). This is a restriction on what the design identifies, not a retraction of any level result, and no aggregate, marginal, or bound reported elsewhere in this paper depends on timing.

\paragraph{Carried caveats.} Four qualifications travel with the seasonal floor wherever it is quoted, and none of them is repaired by the runs above. First, the profile's identification is circular in the same way the in-window floor's is: the deep-out-of-the-money filter that defines it captures 98.2\% of QT exposure, so it is close to the identity map on the panel and the profile is not measured off a lock-in-free population. Second, the clean off-window leg that rescues the floor's \emph{level} cannot rescue its \emph{shape}: the 2017--2019 out-of-the-money cells number 335 cohort-months on \$86.2 billion of exposure and record zero prepaid balance in February and April, so no twelve-cell profile is estimable there at all. Third, the seasonal share is not stable across windows, the calendar terms explaining $R^2 = 0.011$ of the full-panel variation against 0.755 in the QT window, and the profile inherits that instability. Fourth, this floor profile peaks in June while the committed \emph{voluntary} seasonality of the estimated cohort-hazard form peaks in March and October. They are different objects and need not agree, but the tension is real and I state it rather than let the two profiles sit unreconciled in the same paper.

\begin{table}[H]
\centering
\footnotesize
\begin{threeparttable}
\caption{Seasonalizing the turnover floor. Each leg re-runs the production microsimulation as a paired central and $\beta_1 = 0$ null with only the turnover floor's calendar profile varied; the marginal is the central-minus-null difference and is basis-invariant. $U_2$ and $r$ are scored on the central leg against the SOMA back-out over the 42-month window. The constant-4.0\% row reproduces the committed production artifacts bitwise. Pre-committed run \texttt{seasonal\_floor\_timing}; artifacts catalogued in Appendix~\ref{app:ledger}.}
\label{tab:seasonalfloor}
\begin{tabular}{@{}lrrrrrr@{}}
\toprule
Floor leg & Null (\$B) & Central (\$B) & Marginal (\$B) & Marginal (pp) & $U_2$ (diffs) & $r$ (detr.) \\
\midrule
Constant 4.0\% (production) & 748.185 & 818.530 & $+70.345$ & $+9.20$ & 1.0029 & $-0.203$ \\
Constant 4.5\% & 738.869 & 796.510 & $+57.642$ & $+7.54$ & 1.0023 & $-0.255$ \\
Shape only, mean pinned 4.0\% & 742.625 & 808.073 & $+65.448$ & $+8.56$ & 0.9947 & $+0.016$ \\
Pre-committed seasonal, mean 4.5\% & 729.298 & 783.602 & $+54.305$ & $+7.10$ & 0.9896 & $+0.070$ \\
\midrule
\multicolumn{7}{@{}l}{\textit{Placebo families, against the shape-only leg}} \\
Permutation, 14 scrambled orders\tnote{a} & --- & --- & $+65.311$--$66.924$ & --- & 0.975 best & --- \\
Rotation, 11 wrong-month shifts\tnote{b} & --- & --- & --- & --- & 0.9934 best & --- \\
Convexity-matched flat 4.0090\%\tnote{c} & --- & --- & $+70.156$ & $+9.17$ & 1.0029 & $-0.204$ \\
\bottomrule
\end{tabular}
\begin{tablenotes}[flushleft]\footnotesize
\item[a] The same twelve floor values assigned to scrambled months, multiset preserved exactly. The family retains 89.2\% of the shape effect on average (69.9\%--102.8\%), so the movement is floor dispersion under the maximum combination rule, not calendar timing. Three of the fourteen clear the frozen timing rule, and the best beats both seasonal legs on $U_2$.
\item[b] The true profile moved to the wrong months, amplitude and dispersion held. One of eleven clears the frozen rule for each seasonal leg, and the five-month shift beats the true profile's $U_2 = 0.9947$.
\item[c] The flat floor whose monthly hazard matches the calendar-frequency-weighted mean of the twelve seasonal hazards, reproducing the transform's convexity with zero dispersion. Its $-\$0.189$ billion move is pure level leakage, 3.9\% of the shape effect; the remaining 96.1\% is floor dispersion.
\end{tablenotes}
\end{threeparttable}
\end{table}

\section{Pre-Committed Rules and Their Post-Run Adjudications}\label{app:verdicts}

This paper's checks are pre-committed, but pre-commitment fixes what will be measured, not what a measurement means once it exists: adjudicating an outcome against its rule involves post-run judgment, and that judgment has run in both directions in this paper---failures have been diagnosed rather than merely reported, and passes have been weakened rather than banked. A green gate suite is therefore not self-certifying, and this table exists so the judgment itself can be audited. Its domain is narrower still, and worth stating plainly, because a pass count invites a confidence it cannot support. The suite checks three things: that every printed literal matches the committed artifact it is drawn from, that artifacts agree with one another where they overlap, and that named spans are present and, where order carries meaning, ordered. It checks no premise. Nothing in it establishes that a run's specification was the right one, that an imported elasticity transports to this population, or that any identification claim holds; those are argued in the body, and they are what a referee should read the body for. Nor did it see the page: a suite of 108 green gates once coexisted with roughly 919 text items set off the physical sheet across nine pages---among them a full page of this appendix's own ledger and a row of Table~\ref{tab:headline}---which is why a separate check now runs against the built PDFs rather than the source. I report that here because it is the sharpest evidence I have for the claim this appendix rests on: these checks constrain what the paper may say about its own runs, and nothing more. Its inclusion rule: every pre-committed interpretive rule or threshold whose post-run treatment was anything other than mechanical enforcement, in either direction, plus the enforced rules that bound a headline quantity. Each row states the rule as committed, the outcome, and the adjudication with its direction relative to the paper's headline claims.

{\scriptsize
\begin{longtable}{@{}p{4.4cm}p{4.2cm}p{5.6cm}@{}}
\caption{Pre-committed rules and their post-run adjudications. ``Enforced'' means the committed rule was applied exactly as written. ``Pass upheld, claim withheld'' means the rule was met as written and the pass is reported, but the claim was then discounted on diagnostics computed after the run. Every other entry is a post-run judgment, with its direction stated. Rows reference the section that reports them.}
\label{tab:verdicts}\\
\toprule
Rule as committed & Outcome & Adjudication (direction) \\
\midrule
\endfirsthead
\toprule
Rule as committed & Outcome & Adjudication (direction) \\
\midrule
\endhead
Floor-sweep stability, $\pm 2$pp across the 3--5\% floor range (Appendix~\ref{app:ledger}) & FAILED: $+10.8$ to $+5.5$ points & Reinterpreted with a diagnosis: the instability is the max form's censoring mechanics, established by the additive-form rerun (Section~\ref{sec:robustness-floor}). Softens a failure (for the headline) \\
Twenty-seven-cell positivity across all floor/band anchors (Section~\ref{sec:identification}) & PASS: positive in every cell & Pass upheld, claim withheld: demoted to a specification check because the window's rate configuration forces the sign (against the paper's own pass) \\
Fannie replication acceptance envelope (Section~\ref{sec:pathb-fannie}) & PASS: $+8.68$ vs $+9.20$ & Pass upheld, claim withheld: the 11.06-point envelope judged too wide to carry information; the 0.52-point agreement quoted instead (against the paper's own gate design) \\
Percentile cluster-bootstrap layer on the floor reads (Section~\ref{sec:robustness-floor}) & Reported: $[+3.0, +8.0]$ & Superseded by rule: the pre-committed 0.5pp comparison fired MOVES and the wild-cluster interval became the binding layer (against the headline) \\
Production floor calibrated in-window, disclosed as circular (Section~\ref{sec:pathb}) & 4.0\% on locked-in cohorts & Reinterpreted with a diagnosis: re-measured off-window; the headline demoted from $+9.2$ to $+5.6$ points (against the headline) \\
Off-window pooled 2017--2019 floor read (Section~\ref{sec:robustness-floor}) & 6.1\%, implying $+2.3$ points & Disqualified post hoc on two diagnostics (seasoning, rate regime), which raised the marginal; the disqualification argument itself later re-adjudicated as weaker than first stated, conclusion retained (both directions) \\
Ex-ante stratified loan bootstrap, strata held fixed (Section~\ref{sec:pathb}) & CI $[+9.17, +9.23]$ & Superseded: the cluster scheme showed a $\sim$37$\times$ understatement of loan-sampling uncertainty; the committed claim retracted (against the headline)  --- now measured at both floors: $30.8\times$ in-window at the 4.0\% calibration (round-28 rerun)\\
Committed cyclical-floor $\kappa$ grid (Section~\ref{sec:pathb}) & $+7.1$ to $+9.4$ points & Reinterpreted with a diagnosis: the span is floor dispersion under the hard maximum, not cyclicality, by scramble and additive decompositions (for the headline) \\
Frozen seasonal timing rule (Appendix~\ref{sec:robustness-seasonalfloor}) & Cleared by the $\beta_1 = 0$ null and by a wrong-month rotation & Pass upheld, claim withheld: rule met as written, levels only (against the headline) \\
Joint concave$\times$additive cell, $\pm 1$pp interaction convention (Section~\ref{sec:robustness-floor}) & Interaction $-1.47$ points & Enforced, then applied to run \texttt{b5\_joint\_cell} under its pre-committed landing rule: the interaction is within the convention and the composed point ($+2.9$) is now reported (neutral)\\
Hull closure over the form$\times$transform grid (Section~\ref{sec:robustness-floor}) & $+3.53$ below the $+3.89$ edge & Enforced: the hull widened at the bottom to $+3.5$ (against the headline) \\
Cross-design 50\% undercutting threshold, fixed before the run (Section~\ref{sec:robustness-crossdesign}) & Crossed on all fifty seeds & Enforced: the paradigm contrast cannot be cleanly attributed (against the headline) \\
ABM-null feasibility probe (Section~\ref{sec:abm-interp}) & Run before any spec existed & Scoped post hoc: admitted only as a labeled non-existence claim, stamped \texttt{pre\_committed: false} (disclosure) \\
Scale-test parity tier, \$0.05B (Section~\ref{sec:robustness-scale}) & Halted on a FRED revision & Halted without widening the tolerance; re-run under a second pre-commitment (neutral) \\
Form-test $\pm 1$pp immateriality convention (Section~\ref{sec:robustness-floor}; Section~\ref{sec:pathb}) & Exceeded: $+2.05$ points & Enforced for the form test; the same absolute convention strained by the transform test's relative reading, which the text flags itself (mixed) \\
Buyback-credit incidence rule (Section~\ref{sec:pathb}) & REVERSES at every proxy discount & Enforced: the signed leg-(3) claim replaced by the incidence bracket (against the headline) \\ Moving-share bracket, in-sample and off-window (Section~\ref{sec:pathb}) & Descriptive bracket; ordering check PASS at both floors & Enforced: reported as committed; the off-window rerun landed under a pre-fixed A/B/C classifier, Branch A (neutral) \\ Sub-band elasticity extension (Table~\ref{tab:lowband}) & Curve reported at both floors & Enforced: descriptive exhibit; no verdict rule attached ex ante, so none claimed post hoc (neutral) \\ PSA-sweep expectation bands, fixed ex ante (Section~\ref{sec:identification}) & In-sample cells inside; off-window cells outside (steeper under censoring) & Enforced: verdict from widths as committed; the band misses reported, not absorbed (neutral) \\ Covariate compression band $[0.45, 0.60]$, fixed ex ante (Section~\ref{sec:pathb}) & Realized 1.17: the level swing EXPANDS off-window & Enforced: prediction miss reported as measured; the strengthened sensitivity statement lands with its counter-fact (against the headline's level, neutral to the marginal) \\ Danish redemption validation, thresholds fixed pre-fetch (Section~\ref{sec:abm-danish}) & RULE-A$'$: realized deep-discount redemptions 26.4\%/yr, more than twice the threshold & Enforced: the $\approx$0 pin's contradiction stated and the rule-only gap re-presented as a band (against the headline's Danish reading) \\ Within-stratum feasibility gate, disposition fixed ex ante (Section~\ref{sec:identification}) & The literal within-stratum gradient is not computable (bucket and stratum share the coupon axis); FG resolved to full-axis standardization (18 of 22 cells common, 0.3\% imputed) BEFORE any outcome was read & Enforced: branch (a) --- the composition defence falls as stated and the tested sentence replaces the asserted one; the entry is directional, not a marginal re-estimate (upward, on the assembly's counterweight side) \\ Vintage-overlay expectations, fixed ex ante (Section~\ref{sec:pathb}) & Scale realized $0.6645\times$ inside the pre-committed $[0.655, 0.672]$; placebo spread $0.0002$ against a $0.01$ alarm; vintage-specific component $+0.001$ points & Enforced: landed as a change of estimand --- the Ginnie row is relabelled the same way, and neither overlay is a correction to the headline (neutral) \\ Joint-resimulation floor, measured twice (Section~\ref{sec:patha}) & The first NARROWS print was provisionally rejected (labeled A7); the first floor probe then measured the wrong path and its own audit guard said so --- the pile is the clip-everywhere path, matched exactly by the corrected probe (A7-C1b); point-head parity $0.0000$ at both specs & Enforced: (ii-b) --- the joint interval is reported beside, never in place of, the committed one; the floor disclosure is mandatory wherever it appears (neutral) \\ Amortization-wedge additivity, as drafted (Section~\ref{sec:pathb}) & The drafted `additive constant' condition was impossible twice over: four clip-bound months force a zero wedge (the H1 mechanism), and two amortization curves diverge with age, so the truth is the elementwise schedule-difference identity, which holds at $10^{-10}$ (two labeled amendments); the ABM-basis companion failed its own replay gate and is not quoted & Enforced: branch (i) adjudicated on the corrected forms with both adjudications recorded; Theil table refreshed, timing unchanged, dollars bit-invariant (neutral) \\ Coupon-conversion projection bracket (Section~\ref{sec:robustness-hybrid}) & Realized composed 98.1\% fell below the pre-committed projection band (98.7 projected, band edge 98.6): the one-bucket linear projection under-predicts the $\Delta = 0.80$ two-bucket split into the book's 2.0\%-coupon mass & Reported as a prediction miss; the landing branch does not depend on the bracket (downward, on the composed diagnostic only --- the marginal is bit-invariant) \\ Cold-start battery, verdict precedence fixed ex ante (Section~\ref{sec:patha}) & Realized (iii-b): 1/50 on the production branch --- the unfavorable count --- with production still the best objective of 150 starts; the anchor is an evaluation of the frozen vector, not a refit (labeled amendment), and non-finite or exploded starts are excluded by the committed pipeline's own criterion & Enforced: the exclusion is strengthened, not weakened; the honest count lands in both branch-unstable sentences (downward, on Path A's credibility only --- excluded from every headline) \\ Ridge-battery blocking forms (Appendix~\ref{app:ridge}) & The spec-literal ddof-1 refit parity failed at exactly the predicted $3.2\times10^{-5}$ (the production fit standardizes with ddof 0); re-anchored pre-run to the production construction and the frozen-artifact replay (labeled amendments) & Enforced: SURVIVES on all three facts; the two recomputed-at-adoption scopes this run covers are made affirmative, the resampling-scope pair awaits the joint rerun (neutral) \\ Convolution parity re-anchor (Section~\ref{sec:identification}) & The drafted per-draw 1e-9 parity was mathematically unattainable (percentile interpolation does not commute with the curved floor map); split pre-run into a bit-exact quantile-route gate and a per-draw gate with a shared order-statistic bracket (labeled A7) & Enforced: realized $[+2.80, +8.99]$, $+0.37$ wider than the floor read, inside the pre-committed prediction; floor-invariance probe 1.11/1.12 against a $\pm 25\%$ tolerance (neutral) \\ Zero-months mechanism probe (Appendix~\ref{sec:robustness-seasonalfloor}) & First run stopped on its own P1: the spec targeted raw-diff zeros where the manuscript's zeros are back-out-clip zeros; amended post-run (labeled H1-A1), re-run OK & Enforced: outcome ARTIFACT lands wording only --- the concession can only strengthen, no re-score, no committed number moves (neutral) \\ Scaled-null direction, derived ex ante (Section~\ref{sec:identification}) & The review's `signed upward' shorthand reconciled before the run: the mechanics derive the member DOWN; realized monotone down, roots inside the expected bracket, $\phi = 1$ parity bit-exact & Enforced: an upward bias documented; the corrected member lies below the headline (downward) \\ Attenuation prediction bands, fixed ex ante (Section~\ref{sec:pathb}) & All three off-window cells inside their bands; elasticity ordering held & Enforced: reported as committed; listed in the assembly as a bracket in $\theta$ rather than as a correction (a transport sensitivity on an imported coefficient, and $\theta$ is not estimable) (downward) \\ Mixture-curve probe thresholds (Section~\ref{sec:robustness-floor}) & First run stopped on two threshold misses & Amended post-run, labeled: continuity tolerance corrected for float noise ($10^{-15} \to 10^{-13}$; identity probes passed at $10^{-16}$), monotonicity re-stated as monotone-to-plateau (measured 0.03-point dip, 0.1-point band); engine parity gates passed unamended (neutral) \\
\bottomrule
\end{longtable}
}

\end{document}
